\subsection{Restriction of the theta index to a paired row}
\label{subsec:paired-row-indices}
\hypertarget{source-paired-row-indices}{}

The geometric index estimate must be applied away from the distinguished prime of its run.

The auxiliary construction is made twice, with distinguished primes
\(p_*\ne p_*'\). Assign the two runs to the places by
\[
 \operatorname{run}(v)=\begin{cases}1,&v\nmid p_*,\\2,&v\mid p_*.\end{cases}
\]
The assigned run avoids its distinguished residue characteristic, and conjugate places use the same run.

\paragraph{The common form coordinates.}
\label{par:selected-form-coordinates}

The geometric comparison begins over the construction field, whereas division takes place over the row field. The same form coordinates must be used on both sides.

At the selected point let \(K_0\) be the coefficient field of the
common moduli datum and \(K\) the field of its row matrices. The
two runs extend this datum by short-theta covers at their
respective distinguished primes. The cusp space \(C_j\) and its
integral lattice \(M_0\) belong to the common datum. The core
level is \(g_0g_1\) with the specified half-level ideals.
A basis \(F_i\) of the cusp space over \(K_0\) maps to a \(K\)-basis \(b_i\) of the ambient row space under the identification in form coordinates
of Section~\ref{subsec:selected-cusp-reading}. Define
\[
       z_j:K^{R_j}\longrightarrow K\otimes_{K_0}C_j
\]
by \(z_j(b_i)=1\otimes F_i\). The cusp rules and the form map
are defined on the common datum, so the same \(z_j\) works
for both covers at every place.

\begin{definition}[Lattices in common row coordinates]\label{def:common-row-lattices}
For a module \(A\) of forms write
\[
 \langle A\rangle_j=\operatorname{span}_{\mathcal O_K}
            \{z_j^{-1}(1\otimes F):F\in A\}.
\]
Put \(M=\langle M_0\rangle_j\). Over the valuation ring at \(v\)
use its localization and set \(B_v=\langle A_v\rangle_j\) locally,
where \(A_v\) is the inverse image of the scaled product module
under the weight-\(\kappa_j\) inclusion on the cover used at \(v\).
Use the same rule for the unscaled short sections.
\end{definition}

\begin{lemma}[Compatibility of weights and actions on the two covers]\label{lem:selected-form-level-tie}
In the lattices of Definition~\ref{def:common-row-lattices}, the weights of the three factors satisfy
\[
 w_{F,j}+Tw_{L,j}=\kappa_j,\qquad n_P+Tn_W=n_V .
\]
The action on \(M\) pulls back to the level action on either cover.
Its scaled products belong both to \(M\) and to \(\sigma_v\) times
the unscaled short lattice.
\end{lemma}
\begin{proof}
The form frame has positive block given by the nonnegative large
entries and negative block given by the negatives of the small entries
in reverse order. Their block sizes are \(p_j=k+j-1,\ q_j=m-j\);
the central degree is \(n_V\). Removing the \(T\) light factors of
degree \(n_W\) therefore gives the indicated complementary weight.
Both covers have trivial extra level parameter. The core action
commutes with \(z_j\) on each basis vector \(b_i\), hence on their span.
The product law for short rules gives the unscaled inclusion at
weight \(\kappa_j\); the scaled inclusion in \(M\) is the defining
integrality condition in \(A_v\). The host and coefficient embeddings
agree on the half-level generators, the common level product and the
small and ramified prime powers. Their product thus gives the same
\(\sigma_v\) in both inclusions. \label{cor:paired-row-indices-1}
\end{proof}

\paragraph{The product module on the selected cover.}

In those coordinates the integral object to which the row is compared is a sum of modules generated by products of sections.

Fix a run and lift. Write \(\Theta\) for the module of
products of the scaled theta sections. More
precisely, let \(L\) be the scaled module at light weight,
\(A_F,A_0,A_\kappa\) the scaled short sections in the three
auxiliary weights. With \(s\) the dimension of the parameter
space and \(T\) the chosen product exponent, the module in the
algebra of sections is
\[
                L^T A_F+ L^{T-s} A_0+A_\kappa .
\]
The middle term uses the weight of \(A_F\) plus \(s\) copies of
the light weight. For \(M_0\) take the inverse image under its
weight inclusion when intersecting with this sum. 

\begin{auxlemma}[The branch and product lattices on the two covers]
\label{lem:parent-slice-lattices}
Fix a lift of a completed member of Definition~\ref{def:completed-member}.
On either of its two auxiliary covers use the common form coordinates
of Section~\ref{subsec:completed-division} and the scales of
Lemma~\ref{lem:common-slice-scales}. The light columns and short
sections realize Lemma~\ref{hyp:slice-chart-branch-lattices}
and~\ref{hyp:slice-product-lattices} on the selected slice.
Writing \(B\) for the slice algebra, the basis has the order of choice
\[
 (t'_i)_{i=1}^{\dim B}\text{ over }K(X)\ ;\qquad
 (t'_i\otimes1)_i\text{ over }K_\lambda(X)\quad(\lambda\text{ finite}).
\]
\end{auxlemma}

\begin{proof}
\proofstep{Choose sections before choosing a valuation.}
Write \(K\) for the coefficient field of the run and \(X\) for
the independent coefficients of the light rows \(L_0,\ldots,L_s\).
In the chart of \(L_0\) the equations \(L_i/L_0=0\)
give a finite reduced algebra \(B\) over \(K(X)\) by
Proposition~\ref{prop:selected-generic-cut}. Choose the sections \(g_0,t_2\) of complementary weights as in that proof, avoiding the finitely many componentwise vanishing subspaces. Over the infinite coefficient field this is possible. The function \(g_0t_2/L_0^{s+4}\) is a unit on the cut, so \(g_0\) trivializes the weight-\(w_F\) line there.

Writing \(d=\dim_{K(X)} B\), interpolation on the cut gives short
integral sections \(t'_i\) as in
Lemma~\ref{lem:slice-short-basis-before-localization}, with
\[
   B=\bigoplus_{i=1}^{d}K(X)\,(t'_i/g_0)|_B .
\]
The determinant giving this equality remains nonzero on scalar
extension. At a branch \(\zeta\), let \(e_\zeta\) be its order on
the uniformizer, and normalize the point valuation \(v_y\) so that
the uniformizer has value one. The integral frame comparison of
Propositions~\ref{prop:slice-chart-metric-readings}
and~\ref{prop:slice-completed-branch-comparison} gives
\[
  \operatorname{ord}_\zeta(t'_i/g_0)+c_\zeta
       =e_\zeta v_y(t'_i)\ge0,\qquad c_\zeta=e_\zeta v_y(g_0).
\]
Thus the same sections are integral in the shifted branch lattice
and span it rationally. For the scaled sections on the chosen cover their span has colength
\[
 I_\lambda=\operatorname{length}_{O'}(B_0/\textstyle\sum_i O't_i).
\]

\proofstep{Recover the products in the given row.}
Put \(a=T-s\) and let \(H_g\) be the kernel of the low Fourier
coefficient tests of Proposition~\ref{prop:shifted-low-mode-descent}.
At the definite lift take the whole space. The jet terms of
degrees less than \(a\) lie in the span of \(N_1\).
For the difference from this jet, Proposition~\ref{prop:lift-and-divide-at-slice}
gives terms \(L^\alpha h_\alpha\) with \(|\alpha|=a\) and
\(h_\alpha\) of weight \(w_F+s w_L\). At that weight the short
sections span by Proposition~\ref{prop:shifted-rational-cusp-span}.
Descending the containment by Proposition~\ref{prop:descent-of-slice-division}
therefore gives, over the generic coefficient field,
\[
              H_g\subseteq \operatorname{span}(N_1+Z).
\]

At weight \(\kappa\) take short sections whose images form a
basis modulo \(H_g\), as in Proposition~\ref{prop:slice-complement-determinant-cost}.
Their images are integral in the quotient of the same row. Its
rank and its index cost are zero at the definite lift. These
last sections and the two preceding sets of products are the
generators allowed in \(\Theta\). This verifies both sets of lattice
conditions before the choice of a test subspace or depth.
\end{proof}

\begin{proposition}[Uniform remainder for the product module]\label{prop:geometric-run-index}
Fix an allowed lift of a selected run. Use the light and short
section data of Lemma~\ref{hyp:slice-chart-branch-lattices},
the depth and integral lattice data of
Proposition~\ref{prop:local-product-colength}, and the common
scale of Proposition~\ref{prop:scaled-geometric-remainder}. There are
remainders \(r_\lambda\) for its finite places, chosen before the
subspace and cut. Write \(M_{0,\lambda}\) and \(\Theta_\lambda\) for
the completed lattices. For every finite place \(\lambda\) of \(K_0\),
every integer \(h\ge\max(0,T\nu_\lambda)\), and every subspace \(V\) of the
completed span of \(M_0\), the length over the completed valuation ring and the total
weighted remainder satisfy
\begin{equation}\label{eq:geometric-run-index}
 \begin{aligned}
 \operatorname{length}\left(\frac{V\cap M_{0,\lambda}}
          {V\cap M_{0,\lambda}\cap\Theta_\lambda}\right)
     &\le h\dim V+r_\lambda,\qquad r_\lambda\ge0,\\
 \sum_\lambda\frac{\log {\rm N}\lambda}{[K_0:\Q]}r_\lambda
     &\le R_j fmb\,\delta_2 .
 \end{aligned}
\end{equation}
Here \(R_j=\operatorname{rank}M_0\), \(0<\delta_2\le1\) as required by the scale in \eqref{eq:slice-remainder-scales},
and the \emph{same} \(r_\lambda\) works for all \(V,h\).
\end{proposition}

\begin{proof}
Lemma~\ref{lem:parent-slice-lattices} supplies the branch and product lattices for the completed member. For this run and lift take the four-term majorant from \eqref{eq:geometric-four-terms} for \(r_\lambda\). Its branch lengths, discriminant contribution and complementary determinant are computed before choosing \(V\). Proposition~\ref{prop:local-product-colength} bounds the local intersection for every permitted \(V,h\) by \(h\dim V+r_\lambda\). Proposition~\ref{prop:global-product-colength}, together with the choices of Proposition~\ref{prop:scaled-geometric-remainder}, gives the asserted sum with the same majorants. Thus the order of choices is
\[
 (\text{run},j)\ \longmapsto\ (r_\lambda)_\lambda,
 \qquad (\lambda,V,h)\ \longmapsto\ h\dim V+r_\lambda .
\]
This order is needed when the subspace is defined over the larger field.
\end{proof} \label{cor:paired-row-indices-2}

\paragraph{Finite evaluations for the geometric index.}
\label{par:slice-evaluation-net}

A colength in the product module can be estimated by restricting sections to an evaluation set. The norm used for that restriction is fixed on the same run.

Work on the run at weight \(\kappa\) with coefficient field
\(K_0\), and write \(V_\kappa\) for its rational cusp space.
The analytic evaluations below use the uniformization of this
space on the level of the run. In particular they can be
restricted to any of the cusp subspaces tested by the index.
Put \(Z_b=frL\), where \(2f=[E:\Q]\) for the CM field \(E\), \(r\)
is the rank and \(L=\log b\).

\begin{proposition}[Finite evaluations on the selected run]\label{prop:slice-evaluation-net}
Fix a run and lift in the form coordinates of
Definition~\ref{def:common-row-lattices}. Use its rational cusp space
\(V_\kappa\) of weight \(\kappa\) over \(K_0\), with \(Z_b=fr\log b\)
and the
archimedean and level sizes of \eqref{eq:slice-net-run-sizes}.
For every embedding \(\sigma:K_0\hookrightarrow\C\) there is a nonempty finite set
\(\mathcal N_\sigma\) of scalar evaluation functionals, linear over \(\sigma\).
If \(m_\sigma=|\mathcal N_\sigma|\), they satisfy
\begin{equation}\label{eq:slice-net-size}
 \frac14\|F\|_{\sup,\sigma}\le
    \max_{y\in\mathcal N_\sigma}|e_{\sigma,y}(F)|
      \le\|F\|_{\sup,\sigma},\qquad
 \log(16m_\sigma)\le1041Z_b^3 .
\end{equation}
Complex conjugation identifies the sets of functionals at conjugate embeddings, and the simultaneous evaluation map is injective after extending
scalars from \(K_0\) to \(\C\) at each embedding.

\end{proposition}

\begin{proof}
\proofstep{Reduction to a self-dual Hermitian lattice.}
First use the embeddings sending the marked imaginary
quadratic root to the chosen root. The components there are
framed Hermitian domains with the character, principal
congruence subgroup and scalar weight of the run. In each
component write \(\Lambda\) for the symplectically self-dual
period lattice with pairing
\(\operatorname{Tr}_{E/\Q}(\xi H)\), for \(H\) its Hermitian
form and \(\bar\xi=-\xi\). Divide \(\Lambda\) by the fixed
rank-one fractional ideal \(I\) of the principal datum:
\(L'=I^{-1}\Lambda\). The rank-one trace pairing is perfect,
in the sense that
\[
 \operatorname{Tr}_{E/\Q}(\xi a\bar b x)\in\mathbb Z
 \quad(a,b\in I)\qquad\Longleftrightarrow\qquad x\in\mathcal O_E.
\]
For every \(u\in E^r\), test the trace pairing of \(Iu\) with
\(I L'=\Lambda\). The rank-one identity gives
\[
 \{u:H(u,L')\subseteq\mathcal O_E\}
   =\{u:\operatorname{Tr}_{E/\Q}(\xi H(Iu,\Lambda))
                        \subseteq\mathbb Z\}
   = I^{-1}\Lambda = L'.
\]
The second equality uses the self-duality of \(\Lambda\).

Take positive integers
\(N_A,N_C,m_0\) with \(N_AI\subset\mathcal O_E\),
\(N_C\in I\), and \(2m_0\xi\in\mathcal O_E\).
Thus \(N_A\Lambda\subset L'\) and \(N_C L'\subset\Lambda\).
The principal level \(g\) is the product of the two prescribed levels. Write \(T_v(\mu)=H(\mu,v)v\) in the convention
linear in \(\mu\). For \(a\in\mathcal O_{E^+}\) the nilpotent
part of the transvection with parameter \(Ma\),
\(M=gN_AN_Cm_0\), takes values in \(g\Lambda\) because
\[
  N_A H(\Lambda,v)\subseteq\mathcal O_E,\qquad
  N_C v\in\Lambda,\qquad 2m_0\xi\in\mathcal O_E .
\]
Since \(v\) is isotropic, \(T_v^2=0\), so the inverse has the
negative of the same nilpotent part. The transvection preserves
\(\Lambda\) and is trivial modulo \(g\). It belongs to the
principal subgroup and the section character is trivial
there. The same construction applies with the opposite
choice of \(\xi\) on a component.

\proofstep{The size conditions on the selected run.}
Here are the bounds used on the selected run. Write \(d_E\)
for the absolute discriminant, \(k_0=2|\tau(\xi)|\) on the
chosen half, \(\lambda_{\tau,h}\) for the frame weight,
\(S=\sum_{\tau,h}|\lambda_{\tau,h}|\), and choose
\(W>0\) at least the sum of the spreads
\(\max_h\lambda_{\tau,h}-\min_h\lambda_{\tau,h}\).
The bounds on the selected parameters give
\begin{equation}\label{eq:slice-net-run-sizes}
 \log N_A\le4L,\quad \log N_C\le2L,\quad
 \log m_0\le4L,\quad \log g\le19L,\quad 0<k_0\le2 .
\end{equation}
Here \(N_C\) is the norm of the ideal used in the rank-one
datum. Thus \(\log M\le30L\).

At \(\kappa\) we have
\(f\le r\le n_\kappa\), \(f\ge16\),
\(\log f\le2L\), \(\log n_\kappa\le L\) and
\(\log d_E\le4fL\). The differential and central parts of the
weight give \(W=D_F+fn_\kappa+1\) and
\(S\le D_F+frn_\kappa/2\), where \(D_F\) is the sum of the
nonnegative flag degrees. The size inequalities for this weight
give \(\log(1+D_F)\le2L\). Since \(f\le r\le n_\kappa\le e^L\) and \(1+D_F\le e^{2L}\),
\[
 W\le2e^{2L}\le e^{3L},\qquad
 1+S\le e^{2L}+\tfrac12 e^{3L}\le e^{3L}
\]
in the selected range \(L\ge2\).
These bounds use the fixed run's prescribed central weight,
level and field. The two run orientations satisfy them
separately. \label{cor:slice-evaluation-net-1}

\proofstep{Bounding the section away from short isotropic directions.}
For a cusp section on one component, write \(P(x)\) for its
scalar in the frame \(x\), \(n(x)\) for the product of the
absolute flag-minor ratios to powers
\(-\lambda_{\tau,h}/2\), and \(U=\sup_x|P(x)|n(x)\).
This is a finite supremum by the cusp condition. Let \(Q_{x,\tau}\)
be the positive majorant of \(H_\tau\) obtained by reversing
the negative part in the orthogonal flag of \(x\). Set
\[
 \varepsilon_* = \min\left(1,
       \frac{\pi}{M(fW+2)k_0\log 2}\right).
\]
If \(U>0\), take a frame \(x\) with \(|P(x)|n(x)>U/2\).
Every nonzero \(v\in L'\) then has
\(Q_{x,\tau}(v)\ge\varepsilon_*\) at some \(\tau\).

To prove this, assume all the majorants are smaller.
The signed form is dominated by the positive majorant, so
\(|\tau H(v,v)|\le Q_{x,\tau}(v)<\varepsilon_*\le1\) for every
chosen embedding. Because \(H(v,v)\in\mathcal O_{E^+}\),
\[
      |\mathrm N_{E^+/\Q}H(v,v)|
          =\prod_\tau |\tau H(v,v)|<1 .
\]
This rational integer is zero. Thus the small vector is isotropic
and the rank-one transvection applies.
On the complex transvection through \(x\) in direction \(v\),
use the same parameter \(s\) at every \(\tau\).
The condition on the \(\tau\)-flag is
\(1+k_0Q_{x,\tau}(v)\Im s>0\).

On the line \(\Im s=-t\), the upper bound \(Q_{x,\tau}(v)\le\varepsilon_*\)
makes the margin explicit:
\[
 1-k_0t Q_{x,\tau}(v)\ \ge\ 1-k_0t\varepsilon_*=\tfrac12,
 \qquad t=(2k_0\varepsilon_*)^{-1}.
\]
This criterion and the norm comparison follow by taking the
leading minors of the rank-one change of the frame Gram
matrix. On \(\Im s=-t\), with \(t=1/(2k_0\varepsilon_*)\),
the ratio of the original metric factor to the new one is
at most \(2^{W/2}\), by telescoping the minor ratios with
the weight differences. The whole function \(P(x(s))\) is
holomorphic for \(\Im s>-2t\), and is \(M\)-periodic there.

\proofstep{The cusp expansion and the maximum principle.}
To see that it vanishes at the upper end, use the cusp
expansion in the integral isotropic direction \(N_Cv\).
The transvection is quadratic in its isotropic direction. Thus its parameter \(s'\) for \(N_Cv\) has to satisfy
\[
  N_C^2s'=s,\qquad s'=s/N_C^2.
\]
The Fourier frequencies are nonnegative in each real embedding
by boundedness in the normalized metric: the metric changes
at most polynomially in the positive imaginary directions,
so a negative frequency would grow exponentially there.
The zero frequency vanishes by cuspidality. The expansion
is absolutely convergent at the original frame. Along the
positive diagonal, each remaining exponential tends to zero
and is bounded by one in modulus. This proves the upper-end
vanishing. Put 
\[
  z=e^{2\pi i s/M}
\]
 and \(R=e^{2\pi t/M}\). Periodicity gives
a holomorphic function \(F(z)=P(x(s))\); the upper-end limit gives
\(F(0)=0\).

Apply the maximum principle to \(F(z)/z\) on \(|z|\le R\):
\[
 |F(1)|\le R^{-1}\max_{|z|=R}|F(z)|
        =e^{-2\pi t/M}\max_{\Im s=-t}|P(x(s))|.
\] The result is
\[
 |P(x)|n(x)\le e^{-2\pi t/M} U\,2^{W/2}.
\]
Since \(2\pi t/M\ge(fW+2)\log 2\), the right side is strictly
less than \(U/2\). Thus the set of frames whose squared first
minimum for \(\sum_\tau Q_{x,\tau}\) on \(L'\) is at least
\(\varepsilon_*\) contains a point at which each given nonzero component section nearly attains its supremum. This threshold is independent of
the section. \label{cor:slice-evaluation-net-2}

\proofstep{A finite grid of reduced frames.}
We now construct one finite set on this part with larger weights. Use
\[
 B=\max\left(1,\frac{(2fr)^4d_E^2}{4\varepsilon_*}\right),
 \qquad
 \eta=\frac{1}{1000r^2(1+S)}.
\]
The Euclidean metric \(\sum_\tau Q_{x,\tau}\) has real
dimension \(2fr\). Its dual lattice is obtained by applying
twice the isometric sign involution of the majorant to the
trace dual of \(L'\). Multiplication by \(d_E\) sends that
trace dual into \(L'\), by Hermitian self-duality and the
different. The Euclidean dual first minimum is therefore
at least \(2\sqrt{\varepsilon_*}/d_E\). The real transference bound gives
\[
 \lambda_{2fr}(L')
 \le\frac{2fr}{\lambda_1((L')^*)}
 \le\frac{(2fr)d_E}{2\sqrt{\varepsilon_*}} .
\]

The vectors attaining the real rank span the \(E\)-space, since
any proper \(E\)-subspace has smaller real dimension. Select \(r\)
independent vectors over \(E\); the square of the displayed bound
is at most \(B\). The larger factor \((2fr)^4\) also suffices in the transference estimate. \cite{banaszczyk}
In an orthonormal flag their Hermitian Gram entries
are bounded in absolute value by \(B\) at all embeddings.
Orthonormalization here is right translation by an upper
triangular frame change and preserves the normalized value of the section.

There are at most \((1+2B)^{2fr^2}\) such Gram matrices.
For each entry, Minkowski space packing in the maximum of the
complex moduli gives \((1+2B)^{2f}\), since distinct integers
of \(E\) differ by an element with some conjugate of modulus
at least one.
Using the entry bound independently at the \(r^2\) matrix positions gives
\[
  \big((1+2B)^{2f}\big)^{r^2}=(1+2B)^{2fr^2}.
\]
 Fix a Gram matrix \(A\).

To classify the
short bases modulo the principal congruence group, pair
\(N_A\Lambda\) with each basis, and read the pairing vectors
modulo the image of \(gN_AN_CA\) on \(\mathcal O_E^r\).
The values on \(r\) generators of the corresponding finite
quotient of \(\Lambda\) suffice. Such generators exist because
a finite quotient by a nonzero scalar of a projective module
of rank \(r\) is free of rank \(r\) over the quotient ring.
With \(C=gN_AN_CA\), the required scalar is \(\det C\). Indeed
\[
 (\det C)\,\mathcal O_E^r
 =C\,\operatorname{adj}(C)\mathcal O_E^r\subseteq C\mathcal O_E^r .
\]
If the residues agree, the isometry taking one basis to
the other sends every \(x\in\Lambda\) to a vector differing
from \(x\) by an element of \(gN_C\operatorname{span}_{\mathcal O_E}v\).
Apply the same argument to the inverse to obtain equality of
lattices. The isometry belongs to the principal subgroup.
The size of the pairing quotient is

\[
  |N_{E/\Q}\det(gN_AN_CA)|\le(rgN_AN_CB)^{2fr},
\]
by the determinant expansion. The number of basis types per
Gram matrix is at most its \(r\)-th power.

For each representative basis \(v\), consider all orthonormal
frames \(x\) in which its columns have squared majorant length
at most \(B\). At \(\tau\) put \(C_\tau(x)=x_\tau^{-1}v_\tau\).
Its entries have modulus at most \(\sqrt B\), and its
determinant has fixed modulus
\(d_\tau=|\tau(\det A)|^{1/2}\).
Since \(\prod_\tau d_\tau\ge1\), one can use mesh sizes
\[
 \frac{\eta d_\tau}{2r^{r+1}(\sqrt B)^r}
\]
on the real and imaginary coordinates at \(\tau\).

On an occupied cell the difference of the two coefficient
matrices is bounded entrywise by the mesh, up to the factor
\(\sqrt2\) from its real and imaginary coordinates. The identity
\[
 x_\tau^{-1}y_\tau-1
   =(C_\tau(x)-C_\tau(y))\,
        \frac{\operatorname{adj}C_\tau(y)}{\det C_\tau(y)}
\]
and the bound \((r-1)!(\sqrt B)^{r-1}\) on each cofactor
give an entrywise error at most \(\eta\). Choose one
actual frame per occupied cell. Write \(C=6r^2(rB)^r/\eta\). There are \(2r^2\) real coordinates at
each of the \(f\) embeddings, each with at most \(C/d_\tau\) cells.
Their product is bounded by
\[
 \prod_\tau (C/d_\tau)^{2r^2}
 = C^{2fr^2}\bigl(\prod_\tau d_\tau\bigr)^{-2r^2}
 \le C^{2fr^2}.
\]

\proofstep{Stability of section values on the grid.}
For an orthonormal frame, an entrywise perturbation of radius
\(10\eta\) stays in the domain. Indeed the principal
Gram minors differ from the signed unit minors by perturbations
of the identity. With entry radius \(u\le10\eta\), the Gram
perturbation has entries bounded by \(2u+ru^2\le3u\).
After multiplication by the diagonal signs on each principal
block its operator norm is at most \(3ru\). The singular
values lie between \(1-3ru\) and \(1+3ru\), so the log
of the absolute determinant is bounded by \(9r^2u\). The signed determinant is positive by continuity along the
segment from the orthonormal frame. If \(D_{\tau,h}\) denotes
the absolute leading minor and \(D_{\tau,0}=1\), the metric
has logarithm
\[
 \log n=-\tfrac12\sum_{\tau,h}\lambda_{\tau,h}
             (\log D_{\tau,h}-\log D_{\tau,h-1}).
\]

Applying the minor estimates to both terms gives
\(|\log n|\le9r^2uS\le0.09\). In particular \(n>0.91\)
on the ball, so \(|P|\le U/0.91\) there.
The Schwarz bound on each complex line through the origin,
applied to the difference from the value at the orthonormal
frame, changes the scalar by at most \(0.21U\) at radius
\(\eta\). Thus a representative of a frame of norm greater
than \(U/2\) has normalized norm at least
\(0.91(1/2-0.21)U>U/4\).
Both isometry and orthonormal flag change preserve the norm
before this perturbation. A component with an empty set of indices with larger weights can use one arbitrary frame, since then every cusp
section on it vanishes. \label{cor:slice-evaluation-net-3}

\proofstep{Counting the grid and the components.}
The resulting number \(N\) on a component satisfies
\[
 \log N\le
   2fr^2\log(3N_AN_Cgr B^2)
   +2fr^2\log(6r^2/\eta)+2fr^3\log(rB).
\]
The run bounds give \[
 -\log\varepsilon_*\le41L,\qquad \log B\le65fL,
\] and bound the three logarithms
on the right by \(223fL,20fL,66fL\), respectively.
After adding \(\log16\) the estimate is at most \(1000Z_b^3\).

For the number of components, first pass from the period
lattices to Hermitian self-dual lattices as above. They all
lie in the same rational Hermitian space. In fact, if the
determinant lattice in a chosen basis has fractional ideal
\(J\), self-duality gives 
\[
  (\det H)J\bar J=\mathcal O_E.
\]
The ratio of two Hermitian determinants is therefore a local
norm at every inert finite place, where its valuation is even
and unit norms are surjective.
At an inert unramified prime, for a uniformizer \(\pi\) fixed by conjugation and a unit \(u\), the determinant norm has
\[
  v\big(N(\pi^k u)\big)=2k .
\]
 Every element is a norm at a
split place, and the selected extension \(E/E^+\) is unramified
at every finite place. Equal signatures give equal determinant
signs at the real places. Landherr's classification by signature
and local determinant classes gives the asserted isometry.
\cite{landherr}\label{cor:slice-evaluation-net-4}

An integral self-dual lattice in this space has a majorant
with squared first minimum at least \(f\). To see this, choose
a complete flag in a maximal rational isotropic subspace
and append its annihilators in reverse order. The annihilator
grades pair perfectly with the opposite isotropic grades
on the lattice. Indeed all the intersections defining the
flag are saturated. A functional on one grade extends to
the lattice, vanishing on the preceding flag term, and
self-duality realizes it by a vector in the corresponding
annihilator.
For successive saturated isotropic stages \(I\subset J\), self-duality gives the perfect Hermitian pairing between the two grades
\[
  (I^\perp/J^\perp)\times(J/I)\longrightarrow{\cal O}_E,\qquad ([x],[y])\longmapsto h(x,y).
\]
 In a Witt basis the paired line ideals are
therefore \(J_i\) and \(\bar J_i^{-1}\).

At each chosen
embedding give these lines squared metric factors
\((\mathrm N J_i)^{-1/f}\) and \((\mathrm N J_i)^{1/f}\).
They are reciprocal, as required for a majorant on a
hyperbolic pair. Give the anisotropic middle space an
orthogonal majorant at each embedding. For a nonzero vector \(v\), take its first nonzero
graded image in the increasing flag. If that image is \(a\)
in a line ideal \(J\), the product of its squared lengths
at one embedding from each conjugate pair is
\[
 \prod_\tau\frac{|\tau(a)|^2}{(\mathrm N J)^{1/f}}
      =\frac{|\mathrm N_{E/\Q}a|}{\mathrm N J}\ge1 .
\]
This applies to the opposite grade \(\bar J_i^{-1}\) as well,
since its norm gives the reciprocal metric factor.

If the
image is in the anisotropic middle grade, its Hermitian
square is a nonzero element of \(\mathcal O_{E^+}\); the
product of the absolute values of its conjugates is at
least one. Each absolute value is at most the squared
majorant on this grade. In either case the full squared
majorants dominate these lengths, so
\(\prod_\tau Q_\tau(v)\ge1\) and AM--GM gives
\(\sum_\tau Q_\tau(v)\ge f\).
Transference at this majorant gives bounded bases with the
above value of \(B\) with \(\varepsilon_*\) replaced by \(f\);
call it \(B_0\). A Gram matrix \(A\) and the image of the
lattice under pairing with its basis determine its isometry
class. In the integral pairing coordinates an image \(M\) satisfies
\[
             \operatorname{im} A\subseteq M\subseteq\mathcal O_E^r .
\]
The finite-module argument counts \(M/\operatorname{im}A\)
by at most \(r\) generators.

This generator count gives at most \(((1+2B_0)rB_0)^{2fr^2}\) unmarked lattices. Evaluating the
level marking on \(r\) generators gives at most
\((g^\sharp)^{2fr^2}\) markings per lattice. Thus the number of
components is bounded by
\[
          ((1+2B_0)rB_0 g^\sharp)^{2fr^2}.
\] Here \(g^\sharp\) is the level for the full
marking and \(\log g^\sharp\le30L\).
Polarized complex torus uniformization identifies equal
marked classes. In particular the components number at most
\(\exp(41Z_b^2)\): use
\(\log B_0\le (12+8f)L\), \(1+2B_0\le e^L B_0\) and
\(2(56+16f)\le41fL\) for \(f\ge16, L\ge2\).
Taking the disjoint union over components gives
\[
 \log(16m_\sigma)\le1000Z_b^3+41Z_b^2\le1041Z_b^3.
\]
\label{cor:slice-evaluation-net-5}

\proofstep{Compatibility of scalar evaluation at the test frames.}
At a selected frame \(y\) put
\(e_{\sigma,y}(F)=l_\kappa n(y)P_F(y)\), where
\(l_\kappa>0\) is the fixed Lie-frame factor of the
run and \(P_F\) is the scalar representing \(F\) in these coordinates. The supremum defining
\(\|F\|_{\sup,\sigma}\) uses this same factor. This proves
\eqref{eq:slice-net-size}. These are complex frames. The
field of definition used in the index is \(K_0\) for the
sections, maps and integral lattices. Each evaluation functional is \(\sigma\)-semilinear on this \(K_0\)-space. The height and
local chart construction both use this original field.

\proofstep{Injectivity of the evaluation maps.}
The map to analytic sections is injective after extension to
\(\C\). Let \(A=K_0[t_1,\ldots,t_n]/I\) be a test
algebra, where \(I\) is radical, and let \(y\) be a test
frame over \(A\). For each \(K_0\)-algebra map

\[
  \phi:A\to\C
\]
 over \(\sigma\), base change of the
section rule and complex uniformization identify
\(\phi(F(y))\) with the analytic value at the induced complex frame times a scalar independent of \(F\).
Suppose \(F_1,\ldots,F_h\) are independent over \(K_0\)
and their analytic realizations satisfy \(\sum_i c_iP_{F_i}=0\).
For a fixed test frame \(y\), the tensor
\[
                    b_y=\sum_i c_i\otimes F_i(y)
                    \in\C\otimes_{K_0,\sigma} A
\]
evaluates to zero under every complex point, since the
uniformization scalar at that point is independent of \(i\). This
algebra is reduced because \(K_0\) has characteristic
zero and is perfect. It is of finite type over \(\C\), so
the Nullstellensatz says that the intersection of the
kernels of its complex points is zero. \cite{test-algebra-base-change}
Let \(V\) be the \(K_0\)-span of the coefficients \(c_i\) in
\(\C\), using \(\sigma\) for the scalar action.

The inclusion
\(V\hookrightarrow\C\) remains injective after tensoring with
\(A\) over the field \(K_0\), so the relation is already zero
in \(V\otimes_{K_0}A\). For every dual basis element \(\theta\in V^\vee\),
\[
          0=(\theta\otimes1)b_y=\sum_i\theta(c_i)F_i(y)
                                  \qquad\text{for every }y.
\]
The section \(\sum_i\theta(c_i)F_i\) is therefore zero,
and independence forces all its coefficients to vanish. The dual basis separates
points of \(V\), so all \(c_i\) vanish.
The net controls the supremum for complex linear combinations of these sections as well, since these sums are still holomorphic
cusp sections with the same transformation law. Thus simultaneous evaluation on the finite set is injective over \(\C\). At the other
embeddings use the conjugates of the evaluation functionals on the chosen half.
The two halves use the same supremum by the convention for
the marked imaginary root, and conjugation preserves both
the norm bounds and complex linear independence. \label{cor:slice-evaluation-net-6}

\end{proof}

Define the evaluation norm on the selected weight by
\[
 \|F\|_{2,\sigma}^2
    =\frac1{m_\sigma}\sum_{y\in\mathcal N_\sigma}|e_{\sigma,y}(F)|^2 .
\]
Put \(P_d=1041Z_b^3(1+\log|\kappa|)\).

\begin{proposition}[Evaluation norm and restricted rank]\label{prop:slice-evaluation-norm}
This is a conjugation-compatible Hermitian norm with
\[
 \|F\|_{2,\sigma}\le\|F\|_{\sup,\sigma}
       \le 4\sqrt{m_\sigma}\,\|F\|_{2,\sigma},\qquad
 \dim_{K_0} V_\kappa\le m_\sigma.
\]
When the restricted share \(d_s\) is positive and \(T>s\),
\(\log d_s\le1041Z_b^3\). On the selected scale
\(P_d\le4164Z_b^4\).
\end{proposition}

\begin{proof}
The absolute value of each evaluation is bounded by the supremum. Conversely, by
\eqref{eq:slice-net-size}, \(\|F\|_{\sup,\sigma}\) is at most
four times the largest absolute value on the finite set, which is at most
\(\sqrt{m_\sigma}\,\|F\|_{2,\sigma}\). The finite evaluation
map is injective, so the dimension bound follows from the
number of complex coordinates.
The run's weight size \(|\kappa|\ge1\) makes this one bound
for the log of the squared-norm multiplier and for the log
of the dimension.

When \(d_s>0\) sections are selected by
restriction to the cut, the share bound is
\[
 \binom{T-1}{s}d_s\le R_j\le\dim_{K_0}V_\kappa.
\]
The selected range has \(T>s\), so
\(\log d_s\le1041Z_b^3\) before even using \(P_d\).
For \(d_s=0\) its determinant height contribution is zero.
This supplies the count in the height bound for the same
restricted short sections. It also supplies the metric and
rank bounds for the boundary determinant on the full cusp
lattice. With \(1+\log|\kappa|\le4Z_b\) on the selected
scale, the numerical absorption uses \(P_d\le4164Z_b^4\).

\end{proof}

\paragraph{Use of the evaluation norm in the index.}
The distinguished-prime construction selects the light,
short and auxiliary sections and the indicated restricted
basis using the running scalars. The bounds on the local
charts and lifts then apply to this family at all finite
places. Replacing the auxiliary Hermitian norm by the
finite evaluation norm carries these very sections, points,
trivializations and branch orders through the comparison.
For each embedding \(\sigma\) used for evaluation, the cost of using the finite evaluations in place of the supremum is controlled by
\[
  \|F\|_{2,\sigma}\le\|F\|_{\sup,\sigma}\le4\sqrt{m_\sigma}\,\|F\|_{2,\sigma}.
\]
The degrees of the projective cut and auxiliary section
still use their specified metrics, and the heights of the
generators still use the supremum. The finite norm enters the
boundary determinant by the inequality just proved. This is
the norm and count input to the four-term remainder below,
on each orientation of the run separately.
\label{cor:slice-evaluation-net-7}

\paragraph{A generic linear section and the local remainder.}
We compute the remainder of Proposition~\ref{prop:geometric-run-index}
on the generic slice. The local argument constructs a cokernel
length at a fixed place. Its four terms will be averaged only
after the same slice and sections have been chosen over the
number field.
\hypertarget{source-run-index-length}{}
Fix one run, one weight 
\[
  \kappa=w_F+Tw_L
\]
 and one finite
place \(\lambda\). In this calculation let \(O\) be the completed
DVR, \(\pi\) its uniformizer and \(M\) the completed cusp lattice.
Write \(s=\dim Y\), \(T\ge\max(1,s)\), 
\[
  a=T-s.
\]

Take \(s+1\)
generic combinations \(L_i=\sum_k X_{ik}w_k\) of the integral
generators of the light sections. The coefficients are independent
indeterminates. Localize \(O[X]\) at \((\pi)\) to obtain a DVR
\(O'\) with the same uniformizer. The valuation of a nonzero
polynomial is the minimum of the valuations of its coefficients.
The residue field is the rational function field in the reductions
of the \(X_{ik}\). Flatness sends a simple factor \(O/(\pi)\) in a finite
length filtration to \(O'/(\pi)\), again a simple factor.
In particular for every finite-length \(O\)-module \(C\),
\[
 \operatorname{length}_{O'}(C\otimes_O O')
   =\operatorname{length}_O(C).
\]

\begin{established}[Integral sections on the branches of the generic cut]\label{hyp:slice-chart-branch-lattices}
On the chart \(L_0\ne0\), put \(u_i=L_i/L_0\) for \(1\le i\le s\).
The algebra \(A\) of \(u_1=\cdots=u_s=0\) over \(K'=\operatorname{Frac}O'\)
is finite and reduced. Choose a weight-\(w_F\) frame \(g_0\) on this cut.
For a prime \(\zeta\) of its normalization over \(O'\), write \(e_\zeta\)
and \(f_\zeta\) for its ramification index and residue degree and
\(c_\zeta=e_\zeta v_{y_\zeta}(g_0)\). In the \(g_0\)-frame the module
\[
 B_0=\{v:\operatorname{ord}_\zeta(v/g_0)+c_\zeta\ge0
                    \ \text{for every }\zeta\}
\]
is a full \(O'\)-lattice. The short integral sections \(t'_i\) can be
chosen before localization so that their restrictions \(t_i\) belong to
\(B_0\) and \(I_\lambda=\operatorname{length}_{O'}(B_0/\sum_i O't_i)<\infty\).
\end{established}
\begin{proof}
Proposition~\ref{prop:selected-generic-cut} gives the finite reduced
intersection and identifies the branch readings with point evaluations
in integral frames. Its algebra is a product of separable fields. A
generator of a line over this product is obtained by taking a nonzero
vector in every field factor. Each \(c_\zeta\) is the order of this
vector in the integral point frame, hence an integer. In each field
factor the normalization of the excellent DVR \(O'\) is finite and
semilocal Dedekind. Writing \(\mathfrak p_\zeta\) for the respective maximal ideals in
this normalization, the defining inequalities identify the module:
\[
               g_0^{-1} B_0=\prod_\zeta\mathfrak p_\zeta^{-c_\zeta}.
\]

Each fractional ideal is a full lattice.
Lemma~\ref{lem:slice-short-basis-before-localization} gives one basis
over the generic coefficient field that remains a basis over \(K'\).
At each branch an integral section satisfies
\(\operatorname{ord}_\zeta(t_i/g_0)+c_\zeta=e_\zeta v_{y_\zeta}(t'_i)\ge0\).
Its span is thus a sublattice of the same rank, as asserted.
\end{proof}

\begin{established}[The products and the complementary lattice]\label{hyp:slice-product-lattices}
Let \(N_1\) be the \(O'\)-span of the products on the left below.
Among those on the right one can select a rational basis for their
span. Denote its \(O'\)-span by \(Z\):
\[
 L_1^{\alpha_1}\cdots L_s^{\alpha_s}L_0^{T-|\alpha|}t'_i
 \quad(|\alpha|<a),\qquad
 L_1^{\alpha_1}\cdots L_s^{\alpha_s}t_k^0
 \quad(|\alpha|=a).
\]
Here \(t_k^0\) belongs to the integral short family of weight
\(w_F+s w_L\). The subspace \(H_g\) of forms with vanishing designated
low coefficients lies in the rational span of \(N_1+Z\) over \(K'\).
One can select integral products \(\theta_i\) whose images are a
rational basis modulo \(H_g\). In particular,
\[
 L_{Q,\lambda}
 =\operatorname{length}_{O}\left(
 M/\left((M\cap H_g)+\textstyle\sum_i O\theta_i\right)\right)<\infty .
\]
At the last lift this length is zero.
\end{established}
\begin{proof}
The shifted families of Proposition~\ref{prop:three-shifted-spanning-sets}
span the required spaces. Their integral products therefore span the
quotient used to choose the \(\theta_i\), by
Proposition~\ref{prop:shifted-low-mode-descent}. Extract independent columns
there, and independently extract columns at the boundary over the generic
coefficient field. Proposition~\ref{prop:descent-of-slice-division} puts
\(H_g\) in the resulting span by equality of the matrix ranks before and
after adjoining each of its basis vectors. Extending the coefficient
field preserves these ranks. Since each \(L_j\) is an \(O'\)-linear
combination of integral light columns, both displayed families belong
to \(\Theta\otimes O'\). The \(\theta_i\) are chosen over the original
coefficient field from the integral products. The module in the
denominator of the length has full rank. At the final lift the low-mode
quotient is zero and \(M\cap H_g=M\).
\end{proof}

\begin{lemma}[Saturation length of the boundary products]\label{lem:slice-boundary-saturation}
After extending to \(O'\), write \(Z^{\rm sat}\) for the
intersection of the rational span of \(Z\) with \(M_{O'}\).
For the selected independent boundary columns \(z_k(X)\),
let \(G_\lambda\) be the minimum Gauss valuation among nonzero
determinants \(\det(\phi_i(z_k(X)))\), where the \(\phi_i\)
are integral at \(\lambda\). Then
\[
          \operatorname{length}_{O'}(Z^{\rm sat}/Z)\le G_\lambda .
\]
For \(Z=0\), use the empty determinant and \(G_\lambda=0\).
\end{lemma}

\begin{proof}
The exterior product of the boundary columns lies in
\(\pi^b\bigwedge^{\operatorname{rk}Z} Z^{\rm sat}\), where
\(b=\operatorname{length}_{O'}(Z^{\rm sat}/Z)\) by Smith form.
An integral system of linear functionals evaluates that product
into \(\pi^b O'\), so every nonzero determinant in the
definition has Gauss valuation at least \(b\). Some determinant
is nonzero because the chosen columns are rationally independent
and integral functionals span the rational dual. Taking the
minimum gives the claimed bound. The argument is on the generic
DVR, where specialization of \(X\) at one tuple might raise the
valuation.

\end{proof}

\paragraph{Lengths in the thickened intersection.}

Simple restriction to the reduced intersection would discard its higher orders. The quotient by a power of the defining ideal keeps these contributions in the local length.

Write \(J=(u_1,\ldots,u_s)\), and restrict sections to the
quotient of the weight-\(w_F\) module by its product with \(J^a\).
The restriction map \(\rho\) first divides a weight-\(\kappa\)
section by \(L_0^T\). Let \(A_{\boldsymbol m}\) be the lattice
generated by classes admitting representatives whose coefficient
Hasse derivatives through degree \(a-1\), restricted to the
intersection, have orders at least \(-m_\zeta\) in the shifted
trivializations above. The Hasse derivatives are Taylor
coefficients for the substitution 
\[
  X\mapsto X+U
\]
 in the
generic sections. They include the derivative of order zero.

Write \(A=A_{\boldsymbol 0}\). For \(a=0\) the quotient is
zero and both lengths below vanish. For \(a>0\) we require
the local lifting hypothesis of the construction: every
\(b\in B_0\) is the restriction of a section with class in
\(A\). This condition involves a lift through order \(a-1\).

\begin{lemma}[Lengths from the \(J\)-adic filtration of a slice]\label{lem:slice-thickening-length}
Use the reduced cut and shifted branch lattices of
Lemma~\ref{hyp:slice-chart-branch-lattices}. For \(a>0\)
assume every \(b\in B_0\) lifts to a class in \(A=A_{\boldsymbol 0}\)
through order \(a-1\). Let \(m_\zeta\ge0\) be integers. In the thickened intersection,
with \(R=\binom{T-1}{s}\), one has
\[
 \operatorname{length}(A/\rho N_1)\le R I_\lambda,\qquad
 \operatorname{length}(A_{\boldsymbol m}/A)\le
                             R\sum_\zeta f_\zeta m_\zeta .
\]
The map \(\rho\) is injective on the rational span of \(N_1\),
kills \(Z\), and has kernel \(Z^{\rm sat}\) on the integral
lattice in the span of \(N_1+Z\). At \(a=0\) the thickening
and both displayed quotients are zero.
\end{lemma}

\begin{proof}
\proofstep{Which graded pieces contribute.}
Here is why only
\[
               R=\#\{\alpha\in\mathbb N^s:|\alpha|<a\}
                      =\binom{T-1}{s}
\]
copies of the branch lengths occur. Filter this quotient by
\(J^k\). A representative in filtration degree \(k\), modulo
degree \(k+1\), has the form 
\[
  \sum_{|\alpha|=k}u^\alpha b_\alpha
\]
by ideal multiplication; only the restrictions of \(b_\alpha\)
to the intersection matter. To read those restrictions, write
\(v_l=(w_l/L_0)|_{J=0}\).

The valuation identification at
\(\zeta\) reads its order as \(e_\zeta\) times the difference
between the value of \(w_l\) and the depth of the light lattice.
This is nonnegative since \(w_l\) belongs to that lattice.
Furthermore \(\sum_l X_{0l}v_l=1\). The \(v_l^k\) therefore
generate the unit ideal in the integral closure at every degree
\(k\): a maximal ideal containing all positive \(k\)-th powers
would contain all the \(v_l\). Choose coefficients \(d_{kl}\)
integral on all the branches with 
\[
  \sum_l d_{kl}v_l^k=1;
\]
for \(k=0\) any one column suffices.

For each column \(l\), vary only its coefficients in rows
\(1,\ldots,s\), leaving row zero fixed. Write \(U\) for these changes with \(L_0\) fixed. All \(k\)
factors of \(u^\beta\) must receive a derivative to survive
at \(J=0\). Thus, for \(|\alpha|=|\beta|=k\), the coefficient map is
\[
  [U^\alpha](u^\beta b_\beta)(X+U)\big|_{J=0}
        =\mathbf1_{\alpha=\beta}v_l^k b_\beta\big|_{J=0}.
\]
The same map kills \(J^{k+1}\). Summing the maps over \(l\)
with weights \(d_{kl}\), whose pairing with \((v_l^k)_l\)
is one, gives the \(O'\)-linear reading

\[
  \operatorname{gr}^k_J\to\prod_{|\alpha|=k} (B_0\otimes_{O'}\operatorname{Frac} O')
\]
with components \(b_\alpha|_{J=0}\).
If every output is zero, all the \(b_\alpha\) vanish on the
slice, and the class has filtration degree at least \(k+1\).
The resulting map into the direct sum of the branch spaces is
therefore injective.

\proofstep{The integral degree-zero lattice on each graded piece.}
Let \(B_{\boldsymbol m}\) denote the degree-zero lattice
with lower orders decreased by \(m_\zeta\). An element of
\(A_{\boldsymbol m}\) in filtration degree \(k\) has each
reading in \(B_{\boldsymbol m}\). Indeed, changing its
representative by order \(a\) does not affect a derivative
of order \(k<a\), and every coefficient of the Taylor
expansion of an \(O'\)-scalar is integral in \(O'\).
For the latter assertion use the geometric series for the
inverse of a polynomial with Gauss-unit constant term in its
Taylor expansion.
For an invertible constant term \(b\) and an increment \(h\) of positive order, the expansion used here is
\[
  (b+h)^{-1}=b^{-1}\sum_{j\ge0}(-h/b)^j .
\]
 Thus formation of the indicated spans
preserves the derivative conditions.

For \(A\), the products
in \(N_1\) of degree \(k\) realize a separate copy of
\(\sum_i O't_i\) in each slot. They lie in \(A\): after
division by \(L_0^T\) they are \(u^\alpha t'_i\). For all
coefficient variations the Taylor expansion of \(u_i\) on
the intersection is
\[
             \frac{\sum_l U_{il}v_l}
                       {1+\sum_l U_{0l}v_l}\quad(i>0).
\]
It has integral coefficients. Multiplying a class in \(A\)
by any \(u^\alpha\) preserves the zero-pole condition, and
\(t'_i\) is constant under changes of \(X\).
For the comparison of \(A_{\boldsymbol m}\) with \(A\), use
the lifting hypothesis: for any tuple \(b_\alpha\in B_0\)
choose lifts \(\widetilde b_\alpha\in A\). In degree \(k\)
the coefficient reading just constructed sends
\[
 \sum_{|\alpha|=k}u^\alpha\widetilde b_\alpha
 \ \longmapsto\ (b_\alpha)_{|\alpha|=k}.
\]
As the tuple varies, the image contains the entire product of the
integral branch lattices \(B_0\).
Thus one can subtract a lift of the integral part before
passing to the next degree.

\proofstep{Computing lengths in the normalized branches.}
The integral closure of a DVR in the finite reduced
characteristic-zero algebra is a product of Dedekind rings;
here the field factors are finite separable extensions.
The shifted lattices in its one-dimensional section module
are fractional ideals componentwise. Changing the bound by
one at \(\zeta\) adds a simple residue factor of length
\(f_\zeta\) over \(O'\). Adding the lengths of the graded
quotients gives
\begin{equation}\label{eq:thick-intersection-lengths}
  \operatorname{length}(A/\rho N_1)\le R I_\lambda,\qquad
  \operatorname{length}(A_{\boldsymbol m}/A)
                        \le R\sum_\zeta f_\zeta m_\zeta .
\end{equation}
For \(s>0\) the count for \(R\) follows by adjoining a slack
coordinate to compositions of \(a-1\).

For \(s=0\) there is
one empty multi-index, and \(T\ge1\). At \(T=s>0\) both
sides of the count are zero. The derivative reading also
proves that \(\rho\) is injective on the rational span of
the first family: take the least degree occurring in a
nontrivial relation and use linear independence of the \(t_i\)
in each slot. All of \(Z\) maps to zero by its degree \(a\).
Consequently the kernel on the integral lattice in the
span of \(N_1+Z\) is \(Z^{\rm sat}\).
\label{cor:run-index-length-1}

\end{proof}

\paragraph{Choice of the light generators.}
On each run use the scalar \(s_W\ne0\) assigned to its
distinguished prime and light weight, and put
\(M_0=g_0g_1\). If \(D_F\) denotes the polynomial degree
of the light weight, then \(D_F\le\operatorname{DF}(\kappa)\)
and \(2\operatorname{DF}(\kappa)+2<g_i\) for \(i=0,1\).

In particular 
\[
  M_0\ge4D_F+4.
\]
 The half-level theta
module at this weight is generated by harmonic polynomial
differential rules. Its short generating family includes
the translates of a single block determinant rule of
degree \(D_F\). Multiply the entire family by \(s_W\).
The scaled theta lattice is inside the integral cusp
lattice by the weight and level integrality comparison.
The count and upper norms of its generators are those of
the unscaled theta family with the norm of \(s_W\) included.
In the tables this gives
\(\log N\le n_B\), 
\[
  \log N\le P_{rw}(S_L-1),
\]
 and the
mean log bound \(G_c S_L\) for the maximum of their
norms over the generator list. In particular the number
of terms in the sum can be charged once using \(\log N\).
\label{cor:slice-light-1}

\begin{auxlemma}[Interpolation on the quotient by the short periods]
\label{lem:light-thick-interpolation}
Let \(Y=\mathbb C^h/\Gamma\) have a principal Riemann form
\(\operatorname{Im}H\) and its Gaussian theta metric. Use Haar
probability in the squared section norm. Let \(M\ge1\), \(c\ge0\),
and suppose
\[
 \sqrt{\pi M}\min_{0\ne\gamma\in\Gamma}\|\gamma\|\ge R_0,\qquad
 R_0^2\ge\max(270,90(h+c)).
\]
For \(h=0\) the lattice condition is empty. For each unit direction
\(y\), or for the constant polynomial when \(c=0\), there is a
level-\(M\) theta section \(\beta\) whose Taylor polynomial through
degree \(c\) is \(p(z)=H(z,y)^c\). It satisfies
\[
 \|\beta\|_2^2\le2M^{-h}(\pi M)^{-c}c!,\qquad
 \sup_z |\beta(z)|^2e^{-\pi M\|z\|^2}\le4(\pi M)^{-c}c!.
\]
The constant jet \(p=1\) can likewise be prescribed through any
order \(c\) in this range. Its bounds are
\(\|\beta\|_2^2\le2M^{-h}\) and Gaussian supremum squared at most four.
\end{auxlemma}
\begin{proof}
Put the Gaussian norm
\(\|Q\|_F^2=\int_{\mathbb C^h}|Q(z)|^2e^{-\pi M\|z\|^2}dz\)
on polynomials of degree at most \(c\). Let \(U_\gamma Q\) be
\(Q(z-\gamma)\) multiplied by the factor of automorphy
\(\epsilon_\gamma\exp(\pi M H(z,\gamma)-\pi M\|\gamma\|^2/2)\),
with the semicharacter of the line. This operator preserves the
Gaussian norm. Define \(SQ=\sum_{\gamma\in\Gamma}U_\gamma Q\).
Polynomial growth times Gaussian decay gives normal convergence on
compact sets, also after each fixed differentiation. The cocycle
identity makes \(SQ\) a section.

Let \(TQ\) be the Taylor polynomial of \(SQ\) through degree \(c\).
Rotation in each complex coordinate makes distinct monomials
orthogonal in the Gaussian integral. For any polynomial \(Q'\) of
degree at most \(c\) it follows that
\[
 \langle (T-1)Q,Q'\rangle_F
     =\sum_{\gamma\ne0}\langle U_\gamma Q,Q'\rangle_F .
\]
One can first integrate on balls and then increase the radius.
Normal convergence and Gaussian integrability justify the limit.
Write \(t=\|\gamma\|/2\). The two sets \(\|z\|\ge t\) and
\(\|z-\gamma\|\ge t\) cover the integration domain. For any
polynomial of degree at most \(c\), dilation and orthogonality of
its homogeneous terms give
\[
 \int_{\|z\|\ge t}|Q(z)|^2e^{-\pi M\|z\|^2}dz
 \le 2^{h+c}e^{-\pi M t^2/2}\|Q\|_F^2.
\]
Cauchy--Schwarz on the two sets therefore bounds the absolute
value of each inner product by
\(2^{1+(h+c)/2}e^{-\pi M\|\gamma\|^2/16}\|Q\|_F\|Q'\|_F\).

For the scaled lattice \(\sqrt{\pi M}\,\Gamma\), disjoint balls
of radius \(R_0/2\) give at most \((3u)^{2h}\) nonzero points
of norm at most \(R_0u\) for \(u\ge1\). Put \(a=R_0^2/16\).
Integrate this counting bound against \(2au e^{-au^2}\,du\).
Since \(2\log u\le u^2-1\) and \(a\ge2h\), this gives
\[
 \sum_{\gamma\ne0}e^{-\pi M\|\gamma\|^2/16}
 \le 3^{2h}\frac{a}{a-h}e^{-a}\le2\,9^he^{-a}.
\]

The product of this bound and \(2^{1+(h+c)/2}\) is less than
\(1/2\), by the two lower bounds for \(R_0^2\). In particular
\(T\) is invertible and \(\|T^{-1}\|\le2\). Set \(Q=T^{-1}p\),
\(\beta=SQ\). Unfolding the section norm against one copy of
the sum gives
\[
 \|\beta\|_2^2=\langle SQ,Q\rangle_F
        =\langle p,Q\rangle_F\le2\|p\|_F^2 .
\]

Here the Euclidean covolume is one because the Riemann form is
principal in the compatible unitary metric. A unitary change of
variables and a radial Gaussian integral give
\(\|p\|_F^2=M^{-h}(\pi M)^{-c}c!\).

For the point estimate, the evaluation kernel on the covering
space is \(M^h\exp(\pi M H(z,w))\). Gaussian integration of the
Taylor series after translation to \(w\) gives its reproducing
identity. Its periodization gives the evaluation kernel on \(Y\)
by unfolding. On the diagonal, in the Gaussian metric, its modulus
is at most \(M^h\sum_\gamma e^{-\pi M\|\gamma\|^2/2}\le2M^h\)
by the preceding counting bound. The reproducing identity and
Cauchy--Schwarz give the asserted supremum bound.
\end{proof}

\begin{auxlemma}[Lifting a jet in a direction orthogonal to the short periods]
\label{lem:light-thick-lift}
Use the parameters of Lemma~\ref{lem:joint-theta-parameters}.
Let \(I\) be a saturated isotropic period lattice on which the
semicharacter is trivial. Write \(W=I\otimes\mathbb R\),
\(V_B=(W+iW)^\perp\), and \(U\subset W\) for the projections onto
\(W\) of the imaginary coordinates of all periods along \(W+iW\).
Assume
\[
 \operatorname{rank}I\le d=2fm,\qquad
 \|u-u'\|\ge\sqrt B\quad(u\ne u',\ u,u'\in U),
\]
and that the induced
principally polarized quotient torus in \(V_B\) has scaled first
minimum at least \(R_0\) as in Lemma~\ref{lem:light-thick-interpolation}.
For each unit \(y\in V_B\) in one embedding space and
\(1\le c\le D\), there is a level-\(M\) section \(s\) on the original
torus with
\[
 \begin{gathered}
 \left|d^cs(0)[v_1,\ldots,v_c]-c!\prod_lH(v_l,y)\right|
         \le \varepsilon c!,\qquad \|s\|_2^2\le K_6(c!)^2,\\
 \varepsilon=7\,2^g(4.3)^d e^{-T_0},
 \qquad K_6=6(2d)^{d/2}
 \end{gathered}
\]
for unit directions \(v_l\) in this embedding. An unused column admits
a section with \(|s(0)|\ge1/2\) and \(\|s\|_2^2\le K_6\).
\end{auxlemma}
\begin{proof}
Take orthonormal real coordinates on \(W\), write \(z\) for their
complexification and \(w\) for the \(V_B\)-coordinate. Lift the
section \(\beta\) of Lemma~\ref{lem:light-thick-interpolation} by
\[
 F_0(z,w)=\exp\!\left(\frac{\pi M}{2}\sum_j z_j^2\right)\beta(w).
\]

The multiplier under a period with imaginary coordinate zero is
precisely the multiplier of the original line. In fact the real
translation on \(W\) contributes the difference of the two
quadratic exponents, and the remaining multiplier is the one
defining the quotient line. The trivial semicharacter on \(I\)
makes the quotient multiplier independent of its lift.

Sum the unitary translates of \(F_0\) over the remaining period
classes, indexed by \(u\in U\), to form \(s\).
A translated term, in the Gaussian metric, has absolute value
\[
 e^{-\pi M\|\operatorname{Im}z-u\|^2}
       |\beta(w-w_u)|e^{-\pi M\|w-w_u\|^2/2},
\]
where \(w_u\) is the perpendicular component of the representative.
The Gaussian and the lattice separation imply normal convergence
on compact sets; reindexing gives the full theta transformation law.

Integration over \(W/I\) kills cross terms because their characters
have different frequencies \(M u\). Average next over the quotient
torus in \(w\). Every translate of \(\beta\) has the same squared
norm. The remaining Gaussian sum is bounded uniformly in
\(t=\operatorname{Im} z\): choose a nearest \(u_0\in U\), so
\(\|t-u\|\ge\|u-u_0\|/2\), and apply the preceding packing integral
in real dimension at most \(d\). With \(T_0=MB\) this gives
\[
 \sum_{u\in U}e^{-2\pi M\|t-u\|^2}
 \le1+2\,3^d e^{-\pi T_0/2}\le 2,\qquad
 \|s\|_2^2\le 2\|\beta\|_2^2 .
\]

The last comparison uses Haar probability on every quotient.
In particular \(\|s\|_2^2\le4M^{-\dim_{\mathbb C}V_B}
(\pi M)^{-c}c!\le K_6(c!)^2\).

It remains to bound the derivatives of the terms \(u\ne0\).
In the holomorphic frame the Gaussian bound and
\(\|u-t\|^2\ge\|u\|^2/2-\|t\|^2\) yield
\[
 |s(v)-F_0(v)|\le
 2\sqrt{c!}(\pi M)^{-c/2}
 e^{3\pi M\|v\|^2/2}\,2\,3^d e^{-\pi T_0/2}.
\]

To differentiate in the \(c\) prescribed unit directions, restrict
to \(v=\sum_{l=1}^c z_lv_l\), with all \(|z_l|\le(M(c+1))^{-1/2}\).
Here \(M\|v\|^2\le c\), and the coefficient of
\(z_1\cdots z_c\) is the desired derivative. Cauchy's inequality,
followed by \(c!\ge(c/e)^c\), bounds the error divided by \(c!\) by
\[
 4\,3^d e^{-\pi T_0/2}
 e^{3\pi c/2}\frac{((c+1)/\pi)^{c/2}}{\sqrt{c!}}
 \le 4\,3^d e^{6c-\pi T_0/2}\le\varepsilon .
\]

Indeed \(T_0\ge28D\) by \eqref{eq:joint-theta-log-cost}.
The zero mode has the jet of \(H(w,y)^c\), since its lower
Taylor coefficients vanish and the exponential factor is one
at zero. This proves the derivative comparison. With \(c=0\)
the same construction starting from the constant jet gives
\(|s(0)-1|\le4\,3^de^{-\pi T_0/2}\le1/2\) and the asserted norm.
\end{proof}

\begin{auxlemma}[Separation of harmonic jets at the marked level]
\label{lem:ideal-lattice-jet}
Let \(\Lambda\) be the marked self-dual Hermitian lattice at a complex
frame and \(\mathfrak m\bar{\mathfrak m}=(M_0)\) its chosen half-level
ideal. Let \(D\) be a nonzero constant-coefficient differential operator
of order \(d\) in the \(n\) columns of the frame, with \(M_0\ge4d+4\).
Write \(\theta_\nu\) for the theta sections with the characters of the
marked quotient in those columns. In the theta trivialization at the
origin,
\[
             (D\theta_\nu(0))_\nu\ne0 .
\]
\end{auxlemma}
\begin{proof}
Let \(A'\to A\) be the isogeny of complex abelian varieties given by
the inclusion of period lattices \(\mathfrak m\Lambda\subset\Lambda\).
Take its product \(\pi:(A')^n\to A^n\) and denote by \(L\) the product
of the principal theta lines on \(A^n\). If \(B\) is the alternating
Riemann form on the original lattice, the Riemann form on the ideal
lattice is \(B/M_0\) for its principal polarization. Indeed the
\(B\)-dual of \(\mathfrak m\Lambda\) is
\(\bar{\mathfrak m}^{-1}\Lambda\), so its dual for \(B/M_0\) is
\[
        M_0\bar{\mathfrak m}^{-1}\Lambda=\mathfrak m\Lambda .
\]

It follows that \(\pi^*L\) is an \(M_0\)-th power of an ample line
up to a degree-zero twist. Absorb the twist in one ample factor.
Since \(M_0\ge d+2\), grouping factors gives a product of \(d+2\)
ample line bundles. Theorem~1 of Bauer and Szemberg gives the
surjection \cite{bauer-szemberg-jets}
\[
 H^0((A')^n,\pi^*L)\longrightarrow
 (\pi^*L)_0/\mathfrak m_0^{d+1}(\pi^*L)_0 .
\]

Here \(\mathfrak m_0\) denotes the maximal ideal at the origin.
The functional induced by \(D\) on this jet space is nonzero. In local
coordinates a monomial of degree \(d\) paired with a nonzero
top coefficient has nonzero image in characteristic zero.

Write \(H=\ker\pi\). Character decomposition of the finite cover
over \(\mathbb C\) and the projection formula give
\[
 H^0((A')^n,\pi^*L)
  =\bigoplus_{\nu\in\widehat H} H^0(A^n,L\otimes P_\nu),
\]
where \(P_\nu\) is the degree-zero line of character \(\nu\).
Each summand is one-dimensional. The polarization is principal,
so \(L\otimes P_\nu\) is a translate of \(L\). The half-level marking
indexes these translates by the characters of \(H\). Their pullbacks
in the common theta trivialization are precisely the \(\theta_\nu\),
up to nonzero constant factors. Thus \(D\) cannot vanish on all the
summands. Changing an orthonormal basis makes an invertible linear
change of variables on the jet and preserves this conclusion.
\end{proof}

\begin{auxlemma}[Selecting the column directions of a block]
\label{lem:light-block-directions}
In an embedding space of dimension \(N\le r\), take an orthonormal
basis \(w_1,\ldots,w_N\) and \(i\le N\) orthonormal row
directions \(e_a\). Suppose that for \(1\le j\le t\) the available
direction \(y_j\) satisfies
\[
 H(v,y_j)=H(v,w_j)+\sum_{l\le t}E_{jl}H(v,w_l),
 \qquad \sum_{l\le t}|E_{jl}|\le\frac1{30t^{3/2}},
\]
for all row directions \(v\). For \(j>t\) take \(y_j=w_j\).
One can choose \(i\) distinct indices \(j_b\) so that
\[
 \left|\det(H(e_a,y_{j_b}))_{a,b}\right|\ge 2^{-r/2}/2 .
\]
\end{auxlemma}
\begin{proof}
The \(i\)-by-\(N\) matrix \(C_{aj}=H(e_a,w_j)\) satisfies
\(CC^*=1\). Cauchy--Binet shows that the squared absolute values
of its \(i\)-rowed minors sum to one. The maximum modulus \(\Delta\)
is thus at least \(\binom Ni^{-1/2}\ge2^{-r/2}\). Choose its columns.
At most \(t\) of them change. Expand the determinant multilinearly
in these columns. Each minor of the original matrix appearing in
an error term has modulus at most \(\Delta\); repeated columns
give zero. For \(t>0\) the sum of the absolute errors is therefore
bounded by
\[
 \Delta\bigl((1+(30t^{3/2})^{-1})^t-1\bigr)
 \le\Delta(e^{1/30}-1)<\Delta/2 .
\]
For \(t=0\) there is no error.
\end{proof}

\begin{auxlemma}[The short periods and their dual lattice]
\label{lem:light-isotropic-periods}
Use the principally polarized CM torus \(V/L\) of
Lemma~\ref{lem:eta-period-decay} and the parameters of
Lemma~\ref{lem:joint-theta-parameters}. Write \(A=\operatorname{Im}H\)
and \(\lambda=\lambda_1(L)\). There is a saturated
\(\mathcal O_E\)-submodule \(I\subset L\) with isotropic real
span \(W\). Write \(U(x)\in W\) for the vector representing
\(w\mapsto A(x,w)\) in the Euclidean metric and \(X=U(L)\).
Then
\[
 I=W\cap L,\qquad X=I^*,\qquad \dim_{\mathbb R}W\le d=2fm,
 \qquad I=0\ \Longleftrightarrow\ \lambda\ge h .
\]
The induced lattice on the complex orthogonal complement of
\(W+iW\) has minimum at least \(h\). If \(I\ne0\), then
\begin{equation}\label{eq:light-dual-successive}
 \lambda_1(X)^2\ge B,\qquad
 \lambda_{\dim W}(X)\le d^2/\lambda .
\end{equation}
\end{auxlemma}
\begin{proof}
Start with \(W=0\), and keep \(W\) isotropic, rational with
respect to \(L\), and stable under \(E\otimes\mathbb R\).
Maintain the inequality \(\|U(y)\|>(\Omega h)^{-1}\) whenever
\(y\in L\) has nonzero restriction to \(W\). If some
\(x\in L\setminus W\) has distance less than \(h\) from \(W\),
write \(x=w+z\) with \(w\in W\), \(z\perp W\). Then
\[
 \|U(x)\|^2=A(z,U(x))\le\|z\|\,\|U(x)\|,
 \qquad \|z\|<h\le(\Omega h)^{-1}.
\]

Thus \(A(x,W)=0\). Take integral \(\omega_j\) spanning \(E\)
over \(\mathbb Q\) with embedding norms bounded by \(\Omega\), as in
Lemma~\ref{lem:eta-period-decay}. Integrality and the Rosati
identity give
\[
 |A(x,\omega_jx)|=|A(z,\omega_jz)|\le\Omega\|z\|^2<1.
\]

Hence \(A(x,Ex)=0\). Enlarge \(W\) by the real span of
\(\mathcal O_E x\) and intersect with \(L\). This preserves
isotropy. For a nonzero restriction of \(A(y,\cdot)\) to the
larger space, its norm dominates the previous norm if the
restriction was already nonzero. Otherwise it has nonzero integral
value at some \(\omega_jx\). Its representing vector \(u\) on the
new space is orthogonal to the old one, so
\[
       1\le |(u,\omega_j z)|<\|u\|\Omega h .
\]

This proves the maintained inequality. The dimension increases at
each enlargement. At termination each lattice point outside \(W\)
has distance at least \(h\) from \(W\). In particular for
\(y\in L\) with \(A(y,W)=0\), the component perpendicular
to \(W\) lies in \((W+iW)^\perp\) and has that length unless
\(y\in I\). These components form the quotient lattice. Its alternating
form is again unimodular: extend a basis of the primitive
isotropic sublattice \(I\) to an integral symplectic basis of \(L\).
This can be done one primitive vector at a time, choosing a dual
vector of pairing one and splitting off their hyperbolic plane.
The remaining symplectic pairs give a basis of \(I^{\perp_A}/I\).
In the quotient Gaussian coordinates its scaled minimum satisfies
\[
     \sqrt{\pi M}\,h=R_0\sqrt{\pi f\rho^{21}}\ge R_0 .
\]


The initial enlargement occurs exactly when \(\lambda<h\).
The rational span of \(I\) has an \(E\)-dimension \(a\).
Isotropy implies \(W\cap iW=0\), so \(W+iW\) realizes its
complexification inside \(V\). Each embedding occurs there with
dimension \(a\). Since some embedding block of \(V\) has
dimension at most \(m\), we obtain \(\dim_{\mathbb R}W=2fa\le d\).
The theta semicharacter is trivial on \(I\), since
\(A(k,uk)=0\) there for the half-trace element \(u\).
Every integral functional on \(I=L\cap W\) extends to \(L\).
Perfectness of \(A\) gives \(U(L)=I^*\).
Finally
\[
  B\Omega^2h^2=4f^4/\rho^5\le1,\qquad
  \lambda_{\dim W}(I^*)\lambda_1(I)\le\dim W\le d
\]
by Euclidean transference. With \(\lambda_1(I)\ge\lambda\) this
implies \eqref{eq:light-dual-successive}. \cite{banaszczyk}
\end{proof}

\begin{auxlemma}[Frequencies in one CM embedding]
\label{lem:light-triangular-frequencies}
In Lemma~\ref{lem:light-isotropic-periods}, let \(V_\sigma\) be
the space for one embedding and \(A_\sigma=(W+iW)\cap V_\sigma\)
have dimension \(a>0\). There exist an orthonormal basis
\(e_1,\ldots,e_a\) of \(A_\sigma\), frequencies \(u_s\in X\),
and real \(R_s>0\) such that, with \(p_\sigma\) the orthogonal
projection to \(V_\sigma\),
\begin{align}
 p_\sigma u_s/R_s &=e_s+\sum_{i<s}t_{si}e_i,
 &\sum_{i<s}|t_{si}|^2&\le (30a^2)^{-2},\label{eq:light-triangular-projection}\\
 1&\le2\pi R_s, & \|u_s\|&\le 3R_s/2,\nonumber\\
 |2(u_s,y)|&\le 90m^2d^2\rho\,\|y\|^2\quad(y\in X),
 &\|u_s\|^2&\le 8100m^4d^4\rho^7 P.\label{eq:light-frequency-cost}
\end{align}
\end{auxlemma}
\begin{proof}
Put \(K=d^2\rho\). Let \(A'\) be the span of the projections
already selected. Choose a shortest \(v\in X\) with
\(p_\sigma v\notin A'\), and put \(\mu=\|v\|\). The choice
exists by discreteness and \eqref{eq:light-dual-successive}, which
also gives
\[
            1\le\mu\le d^2/\lambda .
\]
Let \(Z\subset W\) be the real span of the \(\mathcal O_E\)-multiples
of frequencies of length at most \(\mu\). It is spanned by the
products of these frequencies with the \(\omega_j\), hence by
vectors of \(X\) of length at most \(\Omega\mu\). Here stability
of the dual follows from the Rosati identity. The spaces \(Z\)
are nested, since the successive shortest lengths cannot decrease.
The projections chosen at earlier steps lie in \(p_\sigma Z\).

On the \(\{\sigma,\bar\sigma\}\)-summand of the real CM
decomposition of \(Z\), the real CM algebra contains \(J=(i,-i)\).
Isotropy, applied to \(z,Jz\), gives
\[
 0=A(z,Jz)=\pm\big(\|p_\sigma z\|^2-\|p_{\bar\sigma}z\|^2\big),
 \qquad \|p_\sigma z\|^2=\tfrac12\|z\|^2 .
\]
Polarization of this equality gives the corresponding identity
for real inner products. Orthogonally to \(A'\) in \(p_\sigma Z\)
choose a vector of length \((30a^2+1)K\mu\), and take its inverse
image \(t\) in this pair summand. The choice makes \(t\) orthogonal
to every frequency of length less than \(\mu\), since its
projection already lies in \(A'\).

Choose a nearest \(u_s\in X\cap Z\) to \(t\). Take a real basis
of \(Z\) from the generators of length at most \(\Omega\mu\).
Rounding coordinates in its Gram--Schmidt directions, starting
with the last, gives
\[
   \|u_s-t\|\le\tfrac12\sqrt d\,\Omega\mu\le K\mu .
\]
Indeed at each step the residual orthogonal coordinate is at
most half the length of its Gram--Schmidt vector. The last
inequality follows from \(d\ge f\) and
\(d\Omega^2/4=df^3\rho^2\le d^4\rho^2\).
Define \(R_s=\|p_{(A')^\perp}p_\sigma u_s\|\) and \(e_s\)
to be that residual divided by \(R_s\). Then
\[
 R_s\ge30a^2K\mu,\qquad
 \|p_{A'}p_\sigma u_s\|\le K\mu,\qquad
 \|u_s\|\le\bigl(\sqrt2(30a^2+1)+1\bigr)K\mu\le45a^2K\mu .
\]
The first two estimates give \eqref{eq:light-triangular-projection};
they also prove that the induction increases the dimension.

For \(\|y\|<\mu\) compare \(u_s\) with \(u_s\pm y\in X\cap Z\).
Their distances from \(t\) and \((t,y)=0\) give
\(|2(u_s,y)|\le\|y\|^2\). For \(\|y\|\ge\mu\) the same
quantity is at most \(90a^2K\|y\|^2\). Finally use
\(P=2f/(\rho^2\lambda^2)\) in the thin range. Since
\(d^4=16f^4m^4\le16f^5\le8f\rho^3\), we get
\[
 (45a^2K\mu)^2\le
 \frac{2025}{2f}\,a^4d^8\rho^4P
 \le8100m^4d^4\rho^7P .
\]
\end{proof}

\begin{auxlemma}[One column with a prescribed short-period frequency]
\label{lem:light-frequency-column}
For \(u=u_s\) and \(R=R_s\) of
Lemma~\ref{lem:light-triangular-frequencies}, let \(N=2\pi R\)
and \(1\le c\le D\). In the theta frame at the origin there is a
level-\(M\) section \(s\) such that for unit vectors
\(v_j\in V_\sigma\),
\begin{align}
 \left|d^cs(0)[v_1,\ldots,v_c]
      -(-2\pi i)^c\prod_j H(v_j,u)\right|
      &\le 7\,2^g(4.3)^d e^{-T_0}N^c,
       \label{eq:light-frequency-jet}\\
 \|s\|_2^2 &\le K_6\exp(2\pi\|u\|^2/M)N^{2c},
 \qquad K_6=6(2d)^{d/2}.\nonumber
\end{align}
\end{auxlemma}
\begin{proof}
Choose \(x\in L\) with \(U(x)=u\) and put \(a=x/M\).
Use the lift of Lemma~\ref{lem:light-thick-lift} in coordinates
\((z,w)\in(W+iW)\oplus V_B\). Start with a quotient section
\(\beta_0\) whose Taylor polynomial through degree \(c\) is
\(1\). Use on the quotient the inverse unitary translation by
\(w_a\), the \(V_B\)-coordinate of \(a\), and choose its
semicharacter to give the required line after that translation.
The proof of Lemma~\ref{lem:light-thick-interpolation}
gives the uniform bounds
\[
 \|\beta_0\|_2^2\le 2M^{-\dim_{\mathbb C}V_B},\qquad
 \sup_w |\beta_0(w)|e^{-\pi M\|w\|^2/2}\le2 .
\]

Let \(z_a\) be the complex coordinates of the projection of
\(a\) to \(W+iW\). For the real orthonormal basis of \(W\) used
there, \(\operatorname{Im}z_a=u/M\). Apply the unitary translation
by \(a\) on the full torus after the lift, and divide by
\[
 C=\exp\!\left(\frac{\pi M}{2}
                 \left(\sum_j z_{a,j}^2-\|z_a\|^2\right)\right),
 \qquad |C|=e^{-\pi\|u\|^2/M}.
\]

The condition \(Ma\in L\) gives a translation of the same level
line. Multiplying its Gaussian factor by the translated zero mode
and completing the square gives
\[
   C^{-1}\,(\text{translated zero mode})(v)
    =e^{-2\pi i H(v,u)}
       \exp\!\left(\frac{\pi M}{2}\sum_j z_j(v)^2\right)
       \beta_0(w(v)).
\]

In fact the coefficient of the linear term in \(z_j(v)\) is
\(\pi M(\bar z_{a,j}-z_{a,j})\); the quotient translation has
cancelled the linear term in \(w(v)\). On \(V_\sigma\) the
quadratic form \(\sum_j z_j(v)^2\) is zero. It is the restriction
of the complexified Euclidean form on \(W\), invariant under the
real CM unit-circle action with characters \(e^{it},e^{-it}\)
on the conjugate pair. A square on the \(\sigma\)-eigenspace
has character \(e^{2it}\). The Taylor polynomial of \(\beta_0\)
now gives the main derivative in \eqref{eq:light-frequency-jet}.
The full translation is an isometry, so the norm comparison for the
lift and \(N\ge1\) give
\[
 \|s\|_2^2\le2|C|^{-2}\|\beta_0\|_2^2
             \le4e^{2\pi\|u\|^2/M}\le K_6e^{2\pi\|u\|^2/M}N^{2c}.
\]

It remains to estimate the nonzero modes after translation.
Put \(L_c=2\pi\|u\|+32\sqrt{M(c+1)}\) and \(b=c/L_c\).
For \(\|v\|\le b\) the term indexed by \(y\in X\setminus\{0\}\)
in the translated sum, divided by \(C\), has modulus at most
twice the exponential of
\[
 \tfrac12\pi Mb^2+2\pi b(\|u\|+M\|y\|)
               -\pi(M-M_1)\|y\|^2,\qquad
 M_1=90m^2d^2\rho .
\]
Indeed its Gaussian includes the difference
\(\pi M(\|u/M\|^2-\|\operatorname{Im}z(v)-u/M-y\|^2)\).
Expand the square, use \(\|\operatorname{Im}z(v)\|\le b\),
and apply \eqref{eq:light-frequency-cost} to the cross term
in \(u,y\). The remaining Gaussian at \(v\) costs at most
\(\pi Mb^2/2\). The factor two is the quotient supremum bound.

Using \(M_1\le M/2\), \(Mb^2\le c/1024\),
\(2\pi b\|u\|\le c\) and \(2b\|y\|\le16b^2+\|y\|^2/16\),
we bound the exponent by
\[
 c+\frac{33\pi c}{2048}-\frac{7\pi}{16}M\|y\|^2
       \le 1.06c-1.37M\|y\|^2 .
\]
Distinct frequencies are separated by \(\sqrt B\); disjoint
half-radius balls bound the count through radius \(t\sqrt B\)
by \((3t)^d\) for \(t\ge1\). Integrating this count against
\(2\alpha t e^{-\alpha t^2}\,dt\), where \(\alpha=1.37 T_0\),
and using \(2\log t\le t^2-1\) and \(T_0\ge d\), gives
\[
 \sum_{y\ne0}e^{-1.37M\|y\|^2}
   \le 3^d e^{-\alpha}\frac{\alpha}{\alpha-d/2}
   \le 2\,3^d e^{-1.37T_0}.
\]

For the mixed derivative use the polydisc
\(v=\sum_{j=1}^c t_jv_j\), \(|t_j|\le1/L_c\). The coefficient
of \(t_1\cdots t_c\) is the required derivative, so Cauchy's
formula costs exactly \(L_c^c\). Finally
\[
 L_c/N\le34\sqrt{M(c+1)},\qquad
 c\log\!\bigl(34e^{1.06}\sqrt{M(c+1)}\bigr)\le0.18T_0.
\]

The second inequality follows from \(c+1\le2c\) and
\(4c(7+\log(Mc))\le T_0\) in
\eqref{eq:joint-theta-log-cost}. The error is at most
\(4\,3^d e^{-1.19T_0}N^c\), which gives
\eqref{eq:light-frequency-jet}.
\end{proof}

\begin{auxlemma}[From estimates in one column to a determinant jet]
\label{lem:light-block-jet-product}
Let \(V\) have a Hermitian inner product \(H\), linear in the first
variable, and let \(\mathcal H\) be the Hilbert space of level-\(M\)
theta sections in one column with its Gaussian \(L^2\) norm. At the
origin use the common holomorphic theta frame. Partition some of the
\(n\) columns into disjoint ordered blocks \(B_j\) of size
\(1\le i_j\le r\). Fix orthonormal directions
\(e_{j1},\ldots,e_{ji_j}\) in the embedding space for each block.
For integers \(c_j\ge1\), write
\[
 \mathcal D=\prod_j
   \det\bigl(\partial_{e_{ja},k}\bigr)_{1\le a\le i_j,\,k\in B_j}^{c_j},
 \qquad D=\sum_j i_jc_j\ge1 .
\]

The second index of a derivative denotes its column. Let \(0<\omega\le1\)
and \(0\le\varepsilon\le\exp(-(D+1)(r\log r+3r+3))\).
For each \(k\in B_j\) suppose that there are
\(s_k\in\mathcal H\), \(\mu_k\ne0\), and \(y_k\in V\) such that
\begin{gather}
 \label{eq:light-single-column}
 \|y_k\|\le2,\qquad
 \omega\|s_k\|_2^2\le|\mu_k|^2,\qquad
 \left|\det(H(e_{ja},y_k))_{1\le a\le i_j,\,k\in B_j}\right|
       \ge\delta:=2^{-r/2}/2,\\
 \label{eq:light-single-column-derivative}
 \left|d^{c_j}s_k(0)[v_1,\ldots,v_{c_j}]
       -\mu_k\prod_l H(v_l,y_k)\right|\le\varepsilon|\mu_k|
\end{gather}
for unit directions in the indicated embedding space. In each unused
column suppose that \(s_k\ne0\) satisfies
\(\omega\|s_k\|_2^2/4\le|s_k(0)|^2\).
Then a unit vector \(F\in\mathcal H^{\otimes n}\) satisfies
\[
 |\mathcal D F(0)|^2\ge (\omega/4)^n 2^{-(r+2)D}.
\]
\end{auxlemma}
\begin{proof}
Take the product \(F_0=\bigotimes_k s_k\). In the determinant expansion
of \(\mathcal D F_0\), each term is a product of the derivatives of the
individual sections. The number of signed terms, before collecting
like terms, is
\[
 N=\prod_j(i_j!)^{c_j}\le r^D .
\]

Divide by \(s_k(0)\) in unused columns and by \(\mu_k\) in used ones.
The main term \(P_0\) and its lower bound are
\[
 P_0=\prod_j\det(H(e_{ja},y_k))^{c_j},\qquad
 |P_0|\ge\delta^{\sum_jc_j}\ge\delta^D .
\]
The bound on \(y_k\) shows that the normalized main factor in a column
of block \(j\) has modulus at most \(2^{c_j}\). Replacing it by the
actual jet changes it by at most \(\varepsilon\).
Let \(m_1=\sum_j i_j\le D\). In one signed product the error is at most
\[
 2^D\big((1+\varepsilon)^{m_1}-1\big)
       \le 2^D m_1\varepsilon(1+\varepsilon)^{m_1}
       \le D4^D\varepsilon .
\]

Multiplying by \(N\) still costs at most half the main term. Indeed
\[
 2D(4r)^D\delta^{-D}\varepsilon\le1.
\]

To check the last inequality, take logarithms and use
\(\log(2D)\le2D\), \(\log4<2\), and
\((r/2+1)\log2\le r/2+1\) in the assumed exponential bound
(or combine \(\log(2D)\le D\) and \(\log 4=2\log 2\)).
Thus
\[
 |\mathcal D F_0(0)|^2\ge
 \tfrac14\delta^{2D}
   \prod_{k\ {\rm used}}|\mu_k|^2
   \prod_{k\ {\rm unused}}|s_k(0)|^2 .
\]

There are \(m_1\ge1\) used columns, so the normalization gives
\[
 \tfrac14\omega^{m_1}(\omega/4)^{n-m_1}
       =4^{m_1-1}(\omega/4)^n\ge(\omega/4)^n .
\] Divide by
\(\|F_0\|_2^2=\prod_k\|s_k\|_2^2\) to obtain the claim.
\end{proof}

The norm on the tensor power also explains why this test bounds the
sum over the original marked translates. Let \(s\) generate the
principal theta line and let \(\Theta_\nu\) be the finite translates
of the pullback of \(s^{\boxtimes n}\) under the half-level isogeny.
The kernel acts unitarily. Its distinct characters give orthogonal
lines, and the translates run through them, by the character
decomposition in Lemma~\ref{lem:ideal-lattice-jet}. Their squared
norms are all \(h_s=\|s\|_2^{2n}\), because the Gaussian metric on the
pullback is the pullback metric and Haar probability pushes forward
to Haar probability. Expanding the unit vector of the lemma gives
\[
 F=\sum_\nu a_\nu\Theta_\nu,\qquad
 h_s\sum_\nu|a_\nu|^2=1,\qquad
 \sum_\nu|\mathcal D\Theta_\nu(0)|^2
       \ge h_s|\mathcal D F(0)|^2 .
\]
The last inequality is Cauchy--Schwarz. It is important here that the
test gives the size of a derivative of a \emph{unit} section.

To apply these estimates at the selected frame, set
\(d=2fm\), \(K_6=6(2d)^{d/2}\), and write \(P\) for the cusp
penalty of Lemma~\ref{lem:eta-period-decay}. Define
\[
 \omega=K_6^{-1}\exp\!\left(-\frac{2\pi}{M}
                 8100m^4d^4\operatorname{rd}(E)^7 P\right).
\]

There are enough distinct columns at this weight. With block ranks
\(p+q=r\), using a separate group of columns for every basic minor
requires at most
\[
 f\left\{\tbinom{p+1}2+\tbinom{q+1}2\right\}
 \le fr(r+1)/2\le n_W .
\]
Indeed the formula of Lemma~\ref{lem:rank-window} gives
\(n_W\ge10rf^2(k-1)^3L^2/m\), while \(m\le k\), \(r<2k\) and
\(k\ge5\). Repetitions of a basic minor are its exponent \(c_j\) in one
group of columns.

At each embedding take the triangular frequency directions of
Lemma~\ref{lem:light-triangular-frequencies} in the span of the short
periods, and complete by an orthonormal basis of its perpendicular
complement. Their one-column sections come respectively from
Lemmas~\ref{lem:light-frequency-column} and~\ref{lem:light-thick-lift},
on the same lattice of Lemma~\ref{lem:light-isotropic-periods}.
Cauchy--Schwarz gives the row-sum bound for the \(t\) short
directions in Lemma~\ref{lem:light-block-directions}:
\[
 \sum_{j=1}^t |E_{ij}|\le\sqrt{t}\,\|E_{i\bullet}\|_2
       \le \frac{\sqrt t}{30t^2}=\frac1{30t^{3/2}}\qquad(t>0).
\]
The frequency vectors used
after normalization have norm at most \(3/2\), and the perpendicular
directions have norm one. This verifies the determinant and
direction conditions in \eqref{eq:light-single-column}. For a thin column take \(\mu=(-2\pi iR)^c\), and for a perpendicular column take \(\mu=c!\). The two norm comparisons then give
\[
 \frac{\|s\|_2^2}{|\mu|^2}\le
 \begin{cases}
 K_6\exp(2\pi\|u\|^2/M),&\text{thin},\\
 K_6,&\text{perpendicular},
 \end{cases}
 \qquad
 \frac{\|s\|_2^2}{|\mu|^2}\le\omega^{-1}
\]
by \eqref{eq:light-frequency-cost}. For unused columns the
section constructed with the constant jet gives the stated
factor \(1/4\). Finally, \eqref{eq:joint-theta-rank-cost} gives
\[
 \log\varepsilon=\log7+g\log2+d\log(4.3)-T_0
 \le -(D+1)(r\log r+3r+3).
\]

Thus the marked sum has the lower bound of
Lemma~\ref{lem:light-block-jet-product} with this particular
\(\omega\).

\begin{lemma}[The scaled light generators]
\label{lem:scaled-light-metric}
The chosen scaled light generators span the integral light lattice.
On the framed generic fibre they have no common zero. For the generators \(w_l\) at a nondefinite lift,
\[
 \sum_l\|w_l\|^2\ge K_1^{-1}e^{-2cP}.
\]
In addition,
the scalar section of weight \(i_\eta g_0^*u_0\) in
\eqref{eq:selected-eta-product} is nonzero on every component and has squared norm at most \(K_2e^{-2i_\eta P}\) for the same
\(P\). At the definite lift the first-minimum penalty is zero.
\end{lemma}

\begin{proof}
\proofstep{Nonzero values for the scaled light family.}
Here is the nonvanishing argument in this normalization.
A triangular change to an orthonormal flag frame multiplies
every value at the same weight by its nonzero character.
Indeed, for the common weight character \(\chi\) and frame change \(b\),
\[
  (w_l(yb))_l=\chi(b)(w_l(y))_l,\qquad \chi(b)\ne0 .
\]
On the orthonormal frame apply Lemma~\ref{lem:ideal-lattice-jet}
with \(M_0\ge4D_F+4\) to the differential operator of a nonzero harmonic polynomial
of degree \(D_F\) in the prescribed flag type. The harmonic
polynomials used here exist by the central and flag weight
conditions. The frame-coordinate inclusion preserves
its leading tensor. The test gives a nonzero reading on
a theta section on the indicated sublattice. The finite
theta translation identity writes that section as a linear
combination of the allowed half-level translates.

Writing this translation identity as \(\Theta=\sum_\nu a_\nu\Theta_\nu\), the jet functional \(D\) from the harmonic polynomial therefore satisfies
\[
  0\ne D\Theta=\sum_\nu a_\nu D\Theta_\nu\quad\Longrightarrow\quad D\Theta_\nu\ne0\text{ for some }\nu.
\]
 Applying
the polynomial to this identity gives a nonzero reading
of some permitted theta rule. A representative of the
same level marking changes its reading by a nonzero
theta factor. This rule lies in the module spanned by
the unscaled family, so its multiple by \(s_W\) is a
linear combination of the chosen light columns.
Indeed, for a linear combination \(s_W H=\sum_j a_j w_j\), evaluation gives
\[
  (s_WH)(y)\ne0\quad\Longrightarrow\quad (w_j(y))_j\ne0 .
\]
 Some
column is nonzero. A common algebraic zero on a generic
framed chart would give a complex zero after the good
embedding, by faithful flatness and the Nullstellensatz.
The complex uniformization on the framed fibre excludes
it. This proves the required no-common-zero property on
the generic fibre.
\label{cor:slice-light-2}

\proofstep{Comparing the normalized light metric.}
For the metric comparison let \(P\) be the cusp penalty,
zero above the chart threshold for the first minimum and
equal below it to
\(2f/(\operatorname{rd}(E)^2\lambda_1^2)\).
Put \(M=g_0g_1\), \(D=D_F\ge1\), \(n=n(w_L)\), and
\(d_*=2fm\), where \(m\) bounds the smaller signature rank.
Define
\[
 B_*=24(2d_*)^{d_*/2},\qquad
 a=\frac{8100\pi n m^4 d_*^4\operatorname{rd}(E)^7}{M},\qquad
 J=B_*^{-n}2^{-(r+2)D}\exp(-2aP).
\]

Let \(b_\nu\) be the unscaled theta values for the translated block
determinant polynomial in the orthonormal frame. Use Lemma~\ref{lem:light-block-jet-product} at the parameters of
Lemma~\ref{lem:joint-theta-parameters}. Its one-column norm allowance is
\[
 \omega=\bigl[6(2d_*)^{d_*/2}\bigr]^{-1} e^{-2aP/n},
 \qquad (\omega/4)^n 2^{-(r+2)D}=J .
\]

The period comparison here uses the two elements \(\pi_i\) of
Proposition~\ref{prop:two-type-primes}. Write \(\alpha=\pi_0\pi_1\)
with the halves chosen by this orientation. The ideal computation in
Proposition~\ref{prop:scaled-integral-light-reading} gives
\(\mathfrak m=(\bar\alpha)\) and \(\alpha\bar\alpha=M\).
If \(H_A\) is the Hermitian form of the original principal polarization,
the principal form on the ideal lattice is \(H'=H_A/M\). Thus
\[
 (\Lambda,H_A)\longrightarrow(\mathfrak m\Lambda,H'),\qquad
 v\longmapsto\bar\alpha v,\qquad
 H'(\bar\alpha v,\bar\alpha v)=H_A(v,v).
\]

This map identifies the two first minima, the cutoff \(h\) and the
penalty \(P\). It also identifies the CM theta signs because it
commutes with the half-trace element. A marked unit letter has
squared \(H_A\)-length \(2\xi_R\) by definition of the polarization
scale; on the ideal torus its squared \(H'\)-length is \(2\xi_R/M\).
The unitary factors at each embedding preserve orthonormality.
Homogeneity of the order-\(D\) differential operator therefore
gives squared scale \((2\xi_R/M)^D\). Haar probability on the
ideal torus pushes forward under its isogeny to Haar probability
on the original torus, so the original theta section of unit
mean square contributes the factor one. The squared sum satisfies
\[
                    \sum_\nu |b_\nu|^2\ge (2\xi_R/M)^D J .
\]

The list of light generators includes these indexed translates after
multiplication by \(s_W\). Combining the bound
\(|\sigma(s_W)|^2\ge M^D\) with the squared differential normalization
\((2\pi\xi_R)^{-D}\) gives
\[
 S:=\sum_l\|w_l\|_{\rm rec}^2
 \ \ge\ (2\pi\xi_R)^{-D}M^D(2\xi_R/M)^D J=\pi^{-D}J.
\]

No minimum over the individual translates occurs here. Since \(J>0\),
taking logarithms gives
\[
 -\tfrac12\log S\le
 \tfrac12\big[D\log\pi+n\log B_*+(r+2)D\log2\big]+aP .
\]

The numerical allowances in the two runs are
\[
 A_5\ge\tfrac12\big[\log\pi+(r+2)\log2+n\log B_*\big],
 \qquad c\ge a .
\]

As \(D\ge1\), \(B_*>1\) and \(P\ge0\), this is at most \(A_5D+cP\).

If \(\ell_L\) converts the metric used in the height
calculation to \(\|\cdot\|_{\rm rec}\), the first-minimum
lower bound for the sum in that metric has
\(K_1=e^{2A_5D_F}\ell_L^2\). The scalar auxiliary form
is chosen independently of the later cut rows. Squaring the cusp estimate for this \(P\) gives
\[
 \ell_\eta^2\|\eta\|^2\le e^{-2i_\eta P}e^{2i_\eta A_2},
 \qquad
 \|\eta\|^2\le K_2e^{-2i_\eta P},\quad
 K_2=e^{2i_\eta A_2}\ell_\eta^{-2}.
\]
The two component bounds on \(\eta\) are its nonvanishing
as a section on each component and the displayed decay.

Strict dominance of \(i_\eta g_0^*w_L-w_\eta\)
comes from the weight gaps in the page. The inequality
\(2cg_0^*\le1\) ensures the strict cusp margin
\(c i_\eta g_0^*<i_\eta\).
The definite case uses \(P=0\), with no scalar auxiliary
form. These comparisons reduce the geometric light-section
inputs to the half-level jet and translated block estimates,
the auxiliary form and the specified page inequalities.
\label{cor:slice-light-3}
\end{proof}


\paragraph{The slice for the chosen sections.}
Write \(K=K_0\). A component \(Z\) of the moduli space over \(K\)
has a field \(\Omega_Z\) obtained by adjoining the coordinates of the
two open flag cells to its function field. If the signatures at a lift
are \((p,q)\), their dimensions and that of the underlying unitary
moduli space give
\[
 s=f\left(pq+\tbinom p2+\tbinom q2\right).
\]
Thus \(\Omega_Z/K\) is finitely generated of transcendence degree
\(s\). Evaluation in the generic flag frame sends a section to an
element of this field.

The components are finite in number and
open and closed. Indeed the generic fibre is smooth of finite type
over a field, and each of its local rings is regular and a domain.
Multiplying by the idempotent of one component extends a section
on it by zero.
If \(e_Z\) denotes this idempotent, its restrictions are
\[
  e_Z|_Z=1,\qquad e_Z|_{Z'}=0\quad(Z'\ne Z).
\]
 This preserves the weight and cuspidality, since the
indicator is locally constant on the analytic uniformization.
Vanishing on every generic flag cell detects the zero section by
density on the reduced charts.

Use the generators and the common first-minimum penalty from
the preceding calculation. The selected light generators span the scaled light lattice
and have no common zero on the framed generic fibre. Before the last
lift the first-minimum estimates on a good complex uniformization give
a function \(P\) and constants \(K_1,K_2>0\) such that
\[
 e^{-2cP}\le K_1 S,\qquad \|\eta\|^2\le K_2 e^{-2i_\eta P},
 \qquad S=\sum_l\|w_l\|^2.
\]
At sufficiently small first minimum \(\lambda_1\) we have
\(P\ge c_0/\lambda_1^2\), with \(c_0>0\).
The scalar section \(\eta\) is nonzero on each component.
With \(a=i_\eta g_0^*\), the run has \(ac<i_\eta\) and
\(a w_L-w_\eta\) strictly dominant in the flag and domain
directions. At the definite last lift use \(\eta=1\), \(a=1\)
and strict dominance of \(w_L\). We explain how these inputs
supply the generic dimension needed for the section calculation.
\label{cor:selected-slice-1}

\begin{proposition}[The reduced generic intersection and evaluation on its branches]
\label{prop:selected-generic-cut}
The generic cut defined by the chosen light rows is finite and reduced.
It is computed from the product of the component fields by the
equal-weight ratios of the light sections. At every branch where
the reference section is nonzero, the independent zeroth row attains
the minimum of the light-column orders. If \(v\) is the branch order,
\(e\) its relative ramification and \(c_\zeta=v(g_0(y_\zeta))\), then
for weight \(w_F+Tw_L\)
\[
 v\big((F/(g_0L_0^T))|_\zeta\big)+c_\zeta
      =e\big(v_{y_\zeta}(F)-T\nu_\zeta\big),
 \qquad \nu_\zeta=v(L_0(y_\zeta))/e.
\]
\end{proposition}

\begin{proof}
\proofstep{Finding a generic slice with the chosen sections.}
At a point with \(\|\eta\|^2/S^a\ge d>0\), substitute the bounds for
\(S^{-1}\) and \(\|\eta\|^2\):
\[
 d\le \|\eta\|^2S^{-a}\le K_2 e^{-2i_\eta P}K_1^ae^{2acP},
 \qquad
 P\le\frac{\log(K_1^aK_2/d)}{2(i_\eta-ac)} .
\]
Here the strict margin \(ac<i_\eta\) permits division.
There the first minimum is bounded below. Reduction theory gives
compact representatives up to the arithmetic group and changes
of flag frame. The definite space already has compact representatives.
The invariant ratio therefore attains a positive maximum.
Writing \(\Phi=\|\eta\|^2/S^a\), the maximizing point \(x_0\) thus satisfies
\[
  \Phi(x_0)=\max_x\Phi(x)>0,\qquad S(x_0)>0 .
\]
 Choose a
nonvanishing column there.

In holomorphic frames the Levi form of
the logarithm of the invariant ratio is the curvature of
\(a w_L-w_\eta\) minus \(a\) times the pullback of the
Fubini--Study form under \(x\mapsto[w_l(x)]\). This follows by
writing the logarithm as the sum of \(\log|\eta|^2\), the weight
potential, and \(-a\log\sum_l|w_l|^2\). The first summand is
pluriharmonic near the maximum. The invariant metric on a strictly
dominant weight has positive Levi form on these tangent directions
(the coefficients on the flag roots are the adjacent weight
differences).

Write \(\Phi\) for the invariant ratio. Along a nonzero vector \(v\) killed by the projective differential, the pulled-back projective form vanishes. The remaining strictly positive term would give
\[
  {\cal L}\log\Phi(v,\bar v)>0,
\]
contradicting negative semidefiniteness at the maximum. The projective differential is therefore injective. Some \(s\) ratios
to the nonvanishing column then have an open set of complex values.
A polynomial relation over \(K\) on \(Z\), homogenized using this
column, would vanish on the corresponding analytic component.
It follows that the field of the light ratios has transcendence
degree \(s\). The compact set and the maximizing point here may
depend on the component and weight. Only the dimension conclusion
is used.

\paragraph{Constant-coefficient ratios.}

Before eliminating the slice variables, we need to know the field generated by the ratios of the chosen light sections on each component.

Fix a column \(w_b\) nonzero at the generic frame of \(Z\).
The light ratios and the component field have the same transcendence degree. Additivity therefore gives
\[
  \operatorname{trdeg}_{K(w_l/w_b)}\Omega_Z=s-s=0.
\]
Since \(\Omega_Z\) is finitely generated over the field of ratios, this algebraic extension is finite. Let \(u\in\Omega_Z\)
have a monic equation of degree \(n\) over the ratio field. Each
coefficient \(a_i\) for \(0\le i<n\) is a quotient
\(A_i(w)/B_i(w)\) of homogeneous polynomials of the same degree.
Put 
\[
  P(w)=w_b\prod_{i<n} B_i(w),
\]
 of degree \(k>0\).
Then \(z=uP(w)\) satisfies
\[
 z^n+\sum_{i<n} C_i(w)z^i=0,\qquad
 C_i=P^{n-i}A_i/B_i,\qquad \deg C_i=(n-i)k .
\]
All \(C_i\) are polynomials.

In each flag chart the weight
transition of \(z\) is the \(k\)-th power of that of the light
weight. The transformed equation is monic with regular coefficients.
The chart ring is smooth and integrally closed in its fraction field.
Thus the transformed rational function is regular on the chart.
These sections agree on overlaps by agreement at the generic point.
Extend by zero to the other components.

\proofstep{Returning the ratio to a cusp section.}
The resulting weight-\(k w_L\) section \(g\) is cuspidal. For an
explicit check the homogeneous root bound in a normalized frame
gives \(\|g\|\le C(\max_l\|w_l\|)^k\): divide the equation by
the corresponding power of this maximum and apply the bound for
roots of a monic polynomial with bounded coefficients. Where all
coefficients vanish the root is zero. Multiplying the monic relation by \(g\) gives
\[
               g^{n+1}=-\sum_{i=0}^{n-1} C_i(w)g^{i+1}
\]
on all components, including the components where \(g=0\).
Each coefficient \(C_i(w)\) has positive degree \((n-i)k\)
in cusp forms, so every summand on the right is cuspidal.

On a rational boundary face, the allowable Fourier indices
are in the positive semidefinite cone. A sum of indices lies
in the face only if each index lies in it. Let \(\operatorname{CT}_{\mathfrak f}\) denote restriction to
the indices of this face \(\mathfrak f\). A product can contribute
there only through factors whose indices all lie on the face.
Consequently
\[
 \bigl(\operatorname{CT}_{\mathfrak f}g\bigr)^{n+1}
   =\operatorname{CT}_{\mathfrak f}(g^{n+1})
   =-\sum_{i=0}^{n-1}
     \operatorname{CT}_{\mathfrak f}(C_i(w))\,
           \bigl(\operatorname{CT}_{\mathfrak f}g\bigr)^{i+1}
   =0.
\]
Every \(C_i(w)\) is cuspidal. The boundary reading is over a reduced
complex chart, so \(\operatorname{CT}_{\mathfrak f}g=0\). We have obtained
\(u=g/P(w)\) at the generic frame. Both numerator and denominator
may be padded to any larger degree using \(w_b\); on a finite set
of components take the largest degree and extend by zero before
summing. No degree bound from this auxiliary ratio calculation
enters the index. \label{cor:selected-slice-2}

\paragraph{The algebra of the generic intersection.}

With these ratios fixed, the slice equations eliminate one coefficient in every row. This gives the algebra on which branch orders and trace discriminants will be computed.

Let \(F=K(X)\), with all \((s+1)N\) coefficients independent.
On a component put \(x_l=w_l/w_b\) and eliminate \(X_{ib}\) for
\(1\le i\le s\) using 
\[
  X_{ib}=-\sum_{l\ne b}X_{il}x_l.
\]
Denote the other variables by \(X'\), including every variable of
row zero. The substitution embeds \(F\) in \(L_Z=\Omega_Z(X')\).
Indeed the differentials of \(s\) of the \(x_l\) are independent
in characteristic zero.

Let \(l_1,\dots,l_s\) index such a set. Modulo \(dX'\) the
determinant test is
\[
 \left[\prod_{i=1}^s X_{i l_i}\right]
      \bigwedge_{i=1}^s\!\left(-\sum_{l\ne b}X_{il}\,dx_l\right)
       =(-1)^s dx_{l_1}\wedge\cdots\wedge dx_{l_s}\ne0 .
\]
The variables of different cut rows are independent. Hence the
full substituted coefficient family is algebraically independent.
The field \(L_Z\) is finitely generated, and its transcendence
degree over \(K\) is 
\[
  s+((s+1)N-s)=(s+1)N.
\]
 Thus \(L_Z/F\)
is finite, and separability follows from characteristic zero.
More precisely in \(\Omega_Z[X]\) the kernel of substitution is
the ideal of the \(s\) row equations, by division by their pivot
variables. After localizing by \(K[X]\setminus\{0\}\) the quotient
is a domain algebraic over \(F\), hence a field with fraction field
\(L_Z\). The independent zeroth row has nonzero value there.
After the coefficient extension to \(F'=K_{0,\lambda}(X)\) the
component algebra is 
\[
  F'\otimes_F L_Z,
\]
 finite and reduced.

\proofstep{Passing from fields to the actual cut.}
We verify why this calculation in function fields applies to the
weight-zero chart of the original forms. Evaluation embeds the chart
before cutting in the product of the component chart rings
\((F'\otimes_K\Omega_Z)[1/L_0]\); injectivity follows by putting
the fractions over a common power of \(L_0\), using independence
of scalar coefficients and detection of a homogeneous numerator at
the generic frames. For each component its indicator is in the chart:
replace each \(w_l\) in the zeroth row by its extension by zero from
that component and divide the result by \(L_0\).
If \(w_{l,Z}\) is the extension by zero of \(w_l|_Z\), the component indicator is
\[
  e_Z=\frac{\sum_lX_{0l}w_{l,Z}}{L_0},\qquad e_Z|_{Z'}=\begin{cases}1&Z'=Z,\\0&Z'\ne Z.\end{cases}
\]
 These orthogonal
idempotents split the chart.

On one summand let \(Q\) be any prime containing the cut equations,
and read the light columns as \(w_l/L_0\) in its residue field. Some
\(w_b/L_0\) is nonzero there by the zeroth-row identity. Invert it.
The algebra of constant homogeneous numerator ratios with pivot
denominator \(w_b^k\) lies in \(\Omega_Z\) and specializes at \(Q\).
Let \(R\subset\Omega_Z\) be that algebra of ratios and
\(\bar R\) its image at \(Q\). The original independent
coefficients are recovered from \(\bar R\) and the free
coefficients by the pivot equations. Therefore
\[
 (s+1)N
 \le \operatorname{trdeg}_K\operatorname{Frac}(\bar R)
                            +((s+1)N-s),
\]
while \(\operatorname{trdeg}_K\Omega_Z=s\). The two bounds
force the residue transcendence degree to be \(s\).
The specialization is then injective, as follows.

To see this, choose \(s\) ratios with
algebraically independent residues (the same dimension count applies
to the algebra generated just by the light ratios). Invert all
nonzero polynomials in them. The localized domain is algebraic over
the resulting rational function field and hence a field.
Indeed a nonzero element \(a\) algebraic over the rational function field has a monic minimal polynomial with \(c_0\ne0\), so
\[
  a^{-1}=-c_0^{-1}\big(a^{d-1}+c_{d-1}a^{d-2}+\cdots+c_1\big).
\]
 Its map to
the residue field has zero kernel. Thus every nonzero homogeneous
denominator of the constant-coefficient ratio construction is
invertible in the local ring at \(Q\), after dividing by a power of
\(L_0\).

\proofstep{Checking that the two localizations agree.}
Consequently the component chart, localized near \(Q\), already
contains all functions of \(\Omega_Z\) as ratios in weight zero.
It also recovers the inverse of the generic zeroth row in the
chosen generic frame: write \(1/w_b\) as a component function
and multiply by \(w_b/L_0\).
In the chosen frame the equality used here is
\[
  (1/w_b)(w_b/L_0)=1/L_0 .
\]
 Conversely evaluation embeds this
localization in the corresponding localization of the component
chart ring. More explicitly, localize at the elements of the
weight-zero chart which are nonzero at every prime of its slice.
The same denominator argument gives \(\Omega_Z\) there, using a
fixed generically nonzero \(w_b\), which stays nonzero at every
such prime by the specialization just proved. The constants, the
coefficients and \(1/L_0\) generate the component chart ring.
Upon cutting, all the inverted denominators were already units.

In the other direction these denominators stay units on the cut
of the component chart ring, since a prime there contracts to a
prime containing the cuts. This identifies the actual slice with
\[
                   \prod_Z (F'\otimes_F L_Z).
\]
In this description \(L_0\) in the component chart ring denotes
its generic-frame reading. The independent frame scalings disappear
in its weight-zero fractions.

\paragraph{Integral frames and the order at a branch.}

A local order is an order of a section in a chosen frame. The cusp form and its complementary section furnish such frames for the branches of the slice.

For the weight-\(w_F\) sheaf the run uses a cusp form \(g_0\)
nonzero on every component and a form \(t_2\) of complementary
weight \((s+4)w_L-w_F\), also generically nonzero on each.
For \(g_0\), nonzero values of the theta sections of weight \(w_F\) provide
a proper vanishing subspace on each of the finitely many analytic
components. Their union cannot cover a vector space over the
infinite coefficient field. The level test here uses the positive
last weights and 
\[
  3\operatorname{DF}(w_F)+2\le M
\]
 where \(M\)
is the full level. Multiplication by \(g_0\) is injective by generic
evaluation. The degree-zero element
\(b=g_0t_2/L_0^{s+4}\) is invertible on the slice just identified.
Indeed it is a unit on each incidence field and stays a unit after
extension of scalars.

Take a chart element \(b'\) with 
\[
  b'b=1\pmod J.
\]
 For
\(a\ge1\), an inverse on the thickened chart is
\[
    b^{-1}\equiv b'\sum_{k=0}^{a-1}(1-b'b)^k\pmod{J^a},
 \qquad
 \frac{h}{g_0}
  =b^{-1}\frac{h t_2}{L_0^{s+4}}\pmod{J^a}
       \quad(h\text{ of weight }w_F).
\]
Indeed \((1-b'b)^a\in J^a\), and the numerator in the second
formula has weight zero after division by the indicated power.
This gives the asserted trivialization by \(g_0\) on the cut
and its positive-order thickenings.

\proofstep{The branch order in the reference frame.}
At a branch \(\zeta\) take the completion of its normalization
DVR. Write \(U\) for the fraction field, \(v\) for its order
and \(e=v(\pi)\). The embeddings 
\[
  \Omega_Z\to L_Z\to U
\]
and \(F'\to U\) evaluate the original rows. Over \(U\) pull
the generic framed point on \(Z\) through the first embedding.
The integral model used for this operation is the proper
compactification with its semiabelian family and invertible
pushforward of the theta line. Its good-level condition at a
residue prime is that the prime-to-residue-characteristic part
of the level is at least three. The two different odd level
primes ensure it even in their own residue characteristics.
The valuation ring of the coefficient model maps into the
branch DVR since the branch order restricts to \(e v_\lambda\).

In particular nonnegative coefficient orders stay nonnegative at the branch:
\[
  v_\lambda(a)\ge0\quad\Longrightarrow\quad v(a)=e\,v_\lambda(a)\ge0.
\]
The valuative criterion extends the underlying point to that DVR.
In the pulled-back Lie lattices the intersections with the
successive generic flag subspaces are saturated. Their successive
quotients are free over the DVR, so choose compatible integral
bases. Choose also a generator of the pulled-back invertible
theta module. These choices reframe the original field point
by an upper triangular change in each flag and a scalar in the
theta line. If \(A\) denotes the block triangular and theta change,
a section \(g\) of weight \(u\) has reading

\[
  g(y_\zeta)=\chi_u(A)g|_\zeta.
\]
 The multiplier is the
same for every section of that weight. A second integral
basis for the same pulled-back flag and theta line changes
it by a unit, since its diagonal entries and theta scalar
are units. Hence for equal-weight \(g,h\) with nonzero
denominator,
\[
   v\big(g(y_\zeta)/h(y_\zeta)\big)
         =v\big((g/h)|_\zeta\big).
\]
This places the branch calculation in the extended set
of integral tests. It gives nonzero
evaluations of \(g_0\) and, before the final lift, \(\eta\).

\proofstep{The minimum from the independent zeroth row.}
The independent zeroth row also computes the minimum on the
light lattice. Choose a column \(b\) of smallest finite order
on \(Z\) at this branch. In the pivot elimination in the next
paragraph the free variables have Gauss valuation over the valued
component field, including all the variables of row zero.
It follows that
\[
 v\left(\sum_l X_{0l}w_l/w_b\right)=
 \min_l v(w_l/w_b)=0 .
\]

For clarity, this uses only the residue transcendence-degree
calculation below, which needs no minimum assertion about \(L_0\).
The light columns span the scaled light lattice over the integers.
The ultrametric inequality and attainment at \(b\) identify its
depth at \(y_\zeta\) with \(\nu_\zeta=v(L_0(y_\zeta))/e\).
Similarly \(c_\zeta=v(g_0(y_\zeta))\). Writing
\(v_y=v/e\), we have for weight \(w_F+Tw_L\)
\[
 v\big((F/(g_0 L_0^T))|_\zeta\big)+c_\zeta
                   = e\big(v_{y_\zeta}(F)-T\nu_\zeta\big).
\]
All entries on the left are integer orders or infinity. In particular
each integral source section has nonnegative unshifted point valuation
by the definition of the cusp lattice on integral frames. These are
the same points for the branch length, the depth test and the lifting
through \(J\). \label{cor:selected-slice-3}
\end{proof}

\begin{established}[Function fields and point valuations on the cut]
\label{hyp:slice-valued-point-lift}\label{cor:generic-slice-lifting-1}
Let \(D'\) be the normalization of \(O'\) in the finite cut algebra.
Each branch of \(D'\) belongs to a factor associated to a component
\(Z\) and its valued point field contains \(\Omega=K_0(Z)\), of
transcendence degree \(s\). Equal-weight ratios are evaluated through
this inclusion and rational functions in the cutting coefficients
through the original Gauss field. Every \(f\in\Omega\) can be expressed as
\[
       f=g/P(w),\qquad \operatorname{wt}(g)=k w_L,\quad \deg P=k,
\]
where \(P\) has constant coefficients and the numerator is a cusp
section, extended by zero to the other components.
\end{established}
\begin{proof}
Choose on \(Z\) a light column \(w_b\) nonzero at its generic point.
The incidence equation solves for one coefficient in every cutting
row. Write \(X'\) for all coefficients except the pivot coefficients
\(X_{ib}\) with \(i>0\). The field on the component is
\[
 X_{ib}=-\sum_{l\ne b} X_{il}\,w_l/w_b \quad (i>0),\qquad
 K_0(\mathrm{Inc}_Z)=\Omega(X').
\]

Normalizing and
passing to a branch retains the indicated embeddings. The point-frame
calculation of Proposition~\ref{prop:selected-generic-cut} identifies
the order of each ratio with its branch order. Its homogeneous
denominator argument clears a monic equation for \(f\) over the light
image, extends \(P(w)f\) in weight \(k w_L\), and proves that its boundary
coefficients vanish. Extension by zero is defined on the disjoint
components of the smooth variety. On \(L_0\ne0\) the cut is defined
by \(J=(L_1/L_0,\ldots,L_s/L_0)\), so its weight-zero algebra is the
chart quotient by \(J\), and its weight-\(w_F\) line is free in the
frame of Lemma~\ref{hyp:slice-chart-branch-lattices}.
\end{proof}

\begin{proposition}[Lifting the completed branch lattice]
\label{prop:generic-branch-jet-lift}
For the data of Lemma~\ref{hyp:slice-valued-point-lift}, every
\(b\in B_0\) in the integral branch lattice of
Lemma~\ref{hyp:slice-chart-branch-lattices} has a
weight-\(w_F\) chart lift with integral Hasse derivatives modulo \(J^{n_0}\)
for \(n_0=\max(1,a)\). If \(s_b\) denotes this lift and \(D^{[0]}\) its restriction to the cut (the zeroth Hasse derivative), then
\[
 D^{[0]}s_b=b,\qquad n_0=\max(1,a).
\]
 The approximating
numerator may be taken at whatever degree is needed for the clearing.
\end{proposition}

\begin{proof}
\proofstep{Normalized orders in the branch field.}
At a branch write \(v\) for its integer order and \(k\) for
the residue field of \(K_0\) at \(\lambda\).
Choose a pivot \(w_b\) with least finite order among the
light columns in a common frame. All \(r_l=w_l/w_b\) on
its component have nonnegative order. Let \(n\) be the number
of variables \(X_{il}\), counting row zero. On the intersection
we have
\[
          X_{ib}=-\sum_{l\ne b} X_{il}r_l
                  \qquad(1\le i\le s).
\]
Call the other \(n-s\) variables \(X'\). All \(n\) residues
of the \(X_{il}\) are algebraically independent over \(k\)
by restriction to the original Gauss DVR. The residue field
of \(\Omega\) has transcendence degree at most \(s\) over
\(k\). Indeed any residue-independent family lifts to an
algebraically independent family over \(K_0\), by normalizing
a proposed polynomial relation at a coefficient of minimum order.

The displayed relations put all \(n\) residues in the field
generated by this residue field and the residues of \(X'\).
If \(\bar\Omega\) is the residue field of \(\Omega\), the
two dimension bounds read
\[
 n\le\operatorname{trdeg}_k\bar\Omega(\bar X')
       \le \operatorname{trdeg}_k\bar\Omega+(n-s)\le n.
\]
Thus the \(n-s\) residues in \(\bar X'\) are algebraically
independent over \(\bar\Omega\). For \(P=\sum_\alpha
a_\alpha (X')^\alpha\ne0\), divide by a coefficient \(a_\beta\)
of minimum order. Its reduction has a nonzero coefficient and
does not vanish at \(\bar X'\). If \(Q=\sum_\alpha b_\alpha(X')^\alpha\ne0\) as well, the two
normalized reductions give exactly
\[
 v(P/Q)=\min_\alpha v(a_\alpha)-\min_\alpha v(b_\alpha).
\]
In particular its values already belong to the value group of
the restricted valuation on \(\Omega\).

Evaluations of slice functions at this branch lie in the
valuation closure of \(\Omega(X')\). A homogeneous fraction
of weight zero, written with denominator \(L_0^t\), can be
evaluated by dividing its numerator and denominator by
\(w_b^t\). All ratios of the global homogeneous forms involved
then lie in \(\Omega\), while the coefficients belong to
\(K_{0,\lambda}(X)\). Any such coefficient is approximated
arbitrarily closely at the Gauss valuation by elements
of \(K_0(X)\), on approximating its finitely many polynomial
coefficients and then taking the quotient. At a finite list
of branches this approximation can be made simultaneously.
Nonvanishing denominators remain nonzero under sufficiently
close approximation.
For a denominator \(Q\ne0\) one may choose its approximation \(Q'\) so that
\[
  v(Q'-Q)>v(Q)\quad\Longrightarrow\quad v(Q')=v(Q).
\]

\proofstep{Restricting the coordinate valuation.}
In particular the restriction of \(v\) to \(\Omega\) takes
the value one. Choose an element of the slice algebra of
order one at this branch, by localization in the finite
product of fields, and approximate its value strictly closer
than its own order. Write the approximant as \(P/Q\) with coefficients \(a_\alpha,b_\beta\)
and denote the order-one element by \(z\). The ultrametric inequality
and the minimum formula give
\[
 v(z-P/Q)>v(z)=1
 \ \Longrightarrow\ 
 1=v(P/Q)=\min_\alpha v(a_\alpha)-\min_\beta v(b_\beta)
 \in v(\Omega^\times).
\] Moreover distinct branches over the
same component give inequivalent valuations on \(\Omega\).

If their restrictions were equivalent, normalize to give
the base uniformizer value one. The valuations would agree
on \(\Omega\), allowing the same pivot, and agree on every
rational expression in \(X'\) over it by the minimum rule.
For each slice element use the same homogeneous expression
at both branches and the simultaneous coefficient approximation.
Its two valuations would agree, including infinite order by
arbitrarily close approximation. Distinct branches are
distinguished by an element of their different maximal
ideals in \(D'\). This proves the assertion.

Let \(A_{\rm loc}\) be the weight-zero chart with all
denominators nonvanishing on the generic slice inverted.
Use the simultaneous branch Taylor expansion
\[
 A_{\rm loc}\longrightarrow \prod_\zeta
                           K_\zeta[[\Delta X]],
\]
where \(\Delta X\) are the increments and \(K_\zeta\) the
completed branch fields. Pull back the subrings
\(O_\zeta[[\Delta X]]\) and take the image \(C\) of this
pullback under constant-term evaluation. The Taylor
product identity makes \(C\) a subring of \(D'\).
The inverse-series formula on the Gauss ring gives
\(O'\subset C\).
A tuple of functions of the components, integral at all
branches, belongs to \(C\). Indeed write each component
function as a ratio of homogeneous forms of equal degree,
using the birational generation.

The denominator is nonzero
as a function there, hence at each of its embedded generic
points. Pad the degrees by powers of a nonzero light section
on each component, extend by zero, and sum numerators and
denominators separately.
Once all forms have the same degree, extension by zero gives the desired componentwise ratio
\[
  \left.\frac{\sum_Z g_Z}{\sum_Z h_Z}\right|_{Z'}=\frac{g_{Z'}}{h_{Z'}} .
\]
 This represents the whole tuple
by one ratio with denominator nonzero on the slice. To put
the ratio in the localized chart one first divides both forms
by the same power of \(L_0\). If \(h\) denotes this ratio of component functions, the same cancellation
holds after replacing \(X\) by \(X+\Delta X\). Thus
\[
 h(X+\Delta X)=h(X),\qquad
 [(\Delta X)^\alpha]h(X+\Delta X)=0\quad (|\alpha|>0).
\]
Its constant coefficient is the prescribed integral tuple, which is
therefore in \(C\).

\proofstep{Identifying the completed lattices.}
We show \(C=D'\). The restricted component valuations just
proved are nontrivial and pairwise inequivalent on each
component. Weak approximation at this finite set realizes
arbitrary assigned integer orders, using an element of order
one locally. It can also approximate a prescribed value
at one branch while prescribing integrality at all the other
branches. Do this component by component and use the homogeneous
ratio just constructed. Thus \(C\) realizes all nonnegative
order profiles, and a locally integral component function has
the same residue at a given branch as some element of \(C\).
\cite{weak-approximation}

Fix a branch \(\zeta\). For an element of \(D'\) approximate
its value there by \(P/Q\in\Omega(X')\) with positive valuation
of the difference. After dividing both polynomials by one
minimum-order coefficient of \(Q\), the Gauss formula gives

\[
  v_\zeta(Q)=0,
\]
 \(v_\zeta(P)\ge0\). Thus their coefficients
are integral at this branch. Replace their residues by the
elements of \(C\) supplied by weak approximation; the variables
already belong to \(O'\). Write \(C_\zeta\) for the image of \(C\) in the branch residue
field. After coefficient replacement the desired residue is
\[
            \bar P/\bar Q\in\operatorname{Frac}(C_\zeta),
                      \qquad \bar P,\bar Q\in C_\zeta,\quad \bar Q\ne0 .
\]
The field of \(O'\)-residues lies in \(C_\zeta\). Since \(C_\zeta\)
is a domain finite-dimensional over that field, multiplication by
each nonzero element is surjective, so \(C_\zeta\) is a field.
Thus the reduction of \(C\) surjects onto the residue field at
each branch.

\proofstep{Separating the branch idempotents.}
For completeness write \(\mathfrak q_\zeta\) for the branch
prime in \(D'\). The ideals \(C\cap\mathfrak q_\zeta\) are
maximal by residue surjectivity and distinct by the order
profiles. At each branch take an element of \(C\) of order
one. Lift the residue digits successively with its powers
as place values. This surjects onto the local quotient modulo
\(\mathfrak q_\zeta^{e_\zeta}\). Put \(\mathfrak m_\zeta=C\cap\mathfrak q_\zeta\).
Distinctness of these maximal ideals permits simultaneous
choice of the lifted digits by the Chinese remainder theorem.
Since \(\mathfrak m_\zeta^{e_\zeta}\) maps to zero modulo
\(\mathfrak q_\zeta^{e_\zeta}\), the maps give a surjection
\[
 C/\!\prod_\zeta\mathfrak m_\zeta^{e_\zeta}
   \longrightarrow
    \prod_\zeta D'/\mathfrak q_\zeta^{e_\zeta}
       \simeq D'/\pi D'.
\]
The last equality uses the branch ramification exponents.
Thus the finite \(O'\)-module \(D'/C\) has zero reduction
modulo \(\pi\), and Nakayama gives \(C=D'\).

\proofstep{Clearing the common denominator.}
Finally choose a ratio of homogeneous forms, independent
of \(X\), of order \(-c_\zeta\) at each branch by weak
approximation. Multiplying \(g_0\) by this ratio gives a
frame of point order zero everywhere and with vanishing
positive Taylor coefficients. In this frame the shifted ball
\(B_0\) is just \(D'\). The equality \(C=D'\) supplies a
localized section lifting every \(b\in B_0\) with integral
Hasse readings.

To return to the chart before localization
take a common weight-zero denominator \(d\) invertible on
the finite reduced slice. Lift its inverse modulo \(J\) to
a chart function \(d_1\). Then
\[
                 d_1\sum_{i=0}^{n_0-1}(1-dd_1)^i
              \equiv d^{-1}\pmod {J^{n_0}},\qquad
                 n_0=\max(1,a).
\]
After clearing the possible equal-weight denominator of the
frame in the same manner, this gives an actual weight-\(w_F\)
section modulo \(J^{n_0}\). Its restriction is \(b\); Hasse
derivatives of total degree less than \(n_0\) of the difference
vanish on the slice by the product formula. This proves the
lifting requirement at every order for the stated geometric
realization. In particular no effective upper bound on the
degree of the approximating chart numerator is needed in
the pole-length calculation. \label{cor:generic-slice-lifting-2}
\end{proof}


The branch lattices just constructed are the ones used in
\eqref{eq:thick-intersection-lengths}. In particular their
integer orders already include the absolute ramification
in the corresponding completed field.

\paragraph{The branch contribution and the complementary sections.}
This is the local part of the proof of
Proposition~\ref{prop:geometric-run-index}. We first bound
the possible excess order on each branch, and then use the
two spaces of sections that are omitted by the branch
restriction.
\begin{proposition}[The local remainder is independent of the subspace]\label{prop:local-product-colength}
Fix the run, lift and place, and use the short sections,
light columns and auxiliary section of the selected slice.
Use the local lifting condition of
Lemma~\ref{lem:slice-thickening-length}, the depth bound
\eqref{eq:slice-tied-depth}, and the rational containment
\(H_g\subseteq\operatorname{span}(N_1+Z)\) of
Lemma~\ref{hyp:slice-product-lattices}. Define \(r_\lambda\) by
\eqref{eq:geometric-four-terms}. It is independent of the
integer \(h\ge\max(0,T\nu_\lambda)\) and the subspace \(V\)
of the completed cusp space. The resulting colength over the local coefficient DVR satisfies
\[
 \operatorname{length}\!\left(\frac{V\cap M}{V\cap M\cap\Theta}\right)
      \le h\dim V+r_\lambda .
\]
\end{proposition}

\begin{proof}
Let \(h\ge0\) be an integer with \(h\ge T\nu_\lambda\), and
let \(\nu_\zeta\) denote the depth of the light lattice at
the point for \(\zeta\). For every integral weight-\(\kappa\)
section \(F\) the evaluation has nonnegative point valuation.
The branch comparison identifies the shifted order of
\(F/L_0^T\) with \(e_\zeta\) times its point valuation
minus \(e_\zeta T\nu_\zeta\). Its Taylor coefficients on
the intersection are this restriction times the coefficients
of \((1+\sum_l U_{0l}v_l)^{-T}\). The latter are integral.
It follows, also for \(O'\)-linear combinations, that
\(\rho(\pi^h M_{O'})\subseteq A_{\boldsymbol m}\) for
\[
 m_\zeta=\max\bigl(0,\lceil e_\zeta(T\nu_\zeta-h)\rceil\bigr).
\]
On each branch let \(n_\zeta\) be the minimum of the integer
orders of the same scaled columns.

The light pivot in the
normal-chart comparison has that order. The exact order \(n_\zeta=e_\zeta\nu_\zeta\) is an integer,
so the ceiling in the earlier bound makes no change. The
comparison is
\[
 m_\zeta=(Tn_\zeta-e_\zeta h)_+
      =e_\zeta(T\nu_\zeta-h)_+
      \le e_\zeta T(\nu_\zeta-\nu_\lambda)_+ .
\]
The point bound \eqref{eq:slice-tied-depth} and its normal-chart
extension give
\[
 \nu_\zeta\le\nu_\lambda+c_+v_\zeta,\qquad c_+=\max(0,c).
\]
Here \(v_\zeta\) is nonnegative and, before the final lift,
is the value of the selected auxiliary section \(\eta\) at
the same point divided by twice its weight parameter \(i_\eta>0\).
At the final lift \(v_\zeta=0\). Since \(h\ge T\nu_\lambda\),
it follows that \(m_\zeta\le Tc_+e_\zeta v_\zeta\), with no
additive term per branch. Here \(e_\zeta\) is the branch
ramification over the construction field. Equation
\eqref{eq:light-frame-order} has already accounted for the
extension at the algebraic point used to obtain the split chart.
Write

\[
  N_{\eta,\lambda}=\sum_\zeta e_\zeta f_\zeta
\operatorname{val}_{y_\zeta}(\eta)
\]
 before the last lift.

Put \(M'=M_{O'}\), \(H=\operatorname{span}(N_1+Z)\), and
\(U=(\pi^hM'\cap H)+N_1+Z^{\rm sat}\).
Lemma~\ref{lem:slice-thickening-length} identifies the
kernel of \(\rho|_U\). In particular there is an exact sequence
\[
 0\longrightarrow Z^{\rm sat}/Z
  \longrightarrow U/(N_1+Z)
  \longrightarrow \rho U/\rho N_1
  \longrightarrow 0 .
\]
Here \(N_1\cap Z^{\rm sat}=0\), since \(\rho\) is injective
on the span of \(N_1\). The last quotient embeds in
\(A_{\boldsymbol m}/\rho N_1\). Applying the two length
bounds in \eqref{eq:thick-intersection-lengths} and
Lemma~\ref{lem:slice-boundary-saturation} gives
\[
 \operatorname{length}\bigl(U/(N_1+Z)\bigr)
       \le G_\lambda+R I_\lambda
                          +R\sum_\zeta f_\zeta m_\zeta .
\]

The rational containment of \(H_g\) puts the image of
\(\pi^h(M'\cap H_g)\) modulo \(\Theta_{O'}\) in this bound.
The chosen quotient generators \(\theta_i\) and their
\(\pi^h\)-multiples belong to \(\Theta_{O'}\). Thus passing
from \(\pi^h(M'\cap H_g)\) to \(\pi^hM'\) introduces at most
\(L_{Q,\lambda}\) more length. Substituting the branch depth
bound gives
\begin{equation}\label{eq:geometric-four-terms}
 r_\lambda=R\left(I_\lambda+
 {\bf1}_{j<m}\frac{Tc_+}{2i_\eta}N_{\eta,\lambda}\right)
                 +G_\lambda+L_{Q,\lambda},
\end{equation}
where the term involving \(i_\eta\) is absent when \(j=m\).

The extension to \(O'\) preserves lengths of finite lattice
quotients, so \(r_\lambda\) bounds the image of \(\pi^hM\) in
\(M/\Theta\). For \(L=V\cap M\), the image of \(\pi^hL\)
is a submodule of that image. The remaining quotient of
the image of \(L\) is a quotient of \(L/\pi^hL\) and has
length at most \(h\dim V\). This gives the asserted bound
for every \(V,h\) using the same \(r_\lambda\).

\end{proof}

\begin{proposition}[Averaging the four local costs]\label{prop:global-product-colength}
Use the short families at weights \(w_F,w_F+s w_L,\kappa\),
with the branch choices of
Proposition~\ref{prop:local-product-colength}, the ranks of
Proposition~\ref{prop:slice-rank-shares}, the height bounds
\eqref{eq:slice-branch-mean-bounds}, and the determinant bounds
of Propositions~\ref{prop:slice-boundary-determinant-cost}
and~\ref{prop:slice-complement-determinant-cost}. Assume \(T>s\)
and \(0\le Tc_+g_0^*\le1\). If \(R_j>0\), the remainders of
Proposition~\ref{prop:local-product-colength} satisfy
\[
 \frac{\langle r\rangle}{R_j}
 \le H_s+\frac{44s^2}{T}(H_z+P_Q)S_\kappa
\]
Here \(H_s,H_z\) are given in \eqref{eq:slice-height-price-constants}
and \(P_Q\) is the quotient height constant of
Proposition~\ref{prop:slice-complement-determinant-cost}.
At the last lift the auxiliary
and complementary contributions are absent.
\end{proposition}

\begin{proof}
\proofstep{The four terms in the global height bound.}
The global size estimate has a separate role. For a place
quantity \(U_\lambda\) write 
\[
  \langle U\rangle=
\sum_\lambda U_\lambda\log {\rm N}\lambda/[K_0:\Q].
\]
The four contributions enter with weights \(R\), \(RTc_+/(2i_\eta)\),
one, and one in this sum. More explicitly, let \(d_s\) be the
number of short restricted sections, \(r_z\) the number of
boundary columns and \(r_q\) the number of complementary
sections. With 
\[
  S_L=|w_L|,
\]
 \(S_\kappa=|\kappa|\) denoting
weight sizes and \(H(t'_i)\) the mean log supremum over the
embeddings, use the \hyperref[par:shifted-small]{scaled short
families} at the three weights.

Choose the \(t'_i\) from the
first family by independent restriction, choose the boundary
columns among the products of the second family by their generic
rank, and choose independent images of the third family in the
quotient by \(H_g\). These choices precede completion at the
finite places. The new choice keeps all the light columns and
\(\eta\). Thus the depth estimate and the following determinant
estimates apply to the same list, with

\[
  \sum_i H(t'_i)\le d_sG_c|w_F|.
\]
 The discriminant and auxiliary-section contributions to the four-term local index satisfy
\begin{equation}\label{eq:slice-branch-mean-bounds}\begin{aligned}
 2\langle I\rangle &\le d_s\log d_s+2\sum_i H(t'_i)
 +\tfrac83 s(s+4)d_s P_{rw} S_L\\
 &\hspace{5em}+2(s+4)(s+1)d_s(2G_c+6P_H) S_L,\\
 \langle N_\eta\rangle/i_\eta
 &\le g_0^*\big((s+1)d_s(2G_c+6P_H) S_L
              +\tfrac43 s d_s P_{rw} S_L\big)+d_sP_\eta
                      \qquad(j<m).
\end{aligned}\end{equation}
Here \(0\log0=0\). The constant \(G_c\) bounds mean heights per unit weight, \(P_{rw}\) bounds the logarithm of the number of light generators per unit weight, and \(P_H,P_\eta\) bound the projective and auxiliary section heights. The auxiliary weight factor is \(g_0^*\), equal to \(1250r\) at the scale below. For the other two terms the
required height bounds are
\[
 \langle G\rangle\le r_z\big(2G_c+P_{rw}+P_L
                    +\tfrac32(P_d+P_O)\big) S_\kappa,
 \qquad
 \langle L_Q\rangle\le 2r_q P_Q S_\kappa .
\]
Here \(-P_L S_\kappa\) bounds the mean log supremum of a
nonzero integral form from below; \(P_d,P_O,P_Q\) are the
metric, coefficient-field and quotient bounds of the run.
At the last lift the quotient has rank zero.

For the selected
size range \(T>s\), the rank comparisons used in the sum are
\[
 R d_s\le R_j,\qquad
 r_q\le 22s^2 R_j/T,\qquad
 r_z\le r_q+22s^2 R_j/T .
\]
In particular \(r_z\le44s^2 R_j/T\). At the last lift the
coefficient 22 suffices there also. To see the effect on the
first two terms put \(B=2G_c+6P_H\), use
\(H(t'_i)\le G_c|w_F|\), \(\log d_s\le P_d\) for \(d_s>0\),
and \(0\le Tc_+g_0^*\le1\). The last bound is imposed when
choosing the light weight and \(T\); it prevents the auxiliary
divisor from multiplying the slice height by \(T\). Define
\begin{equation}\label{eq:slice-height-price-constants}\begin{aligned}
 H_s={}&\tfrac12 P_d+G_c |w_F|
 +\tfrac43 s(s+4)P_{rw} S_L+(s+4)(s+1)B S_L\\
 &+\tfrac12\big((s+1) B S_L+\tfrac43 s P_{rw} S_L
                           +Tc_+P_\eta\big),\\
 H_z={}&2G_c+P_{rw}+P_L+\tfrac32(P_d+P_O).
\end{aligned}\end{equation}

The first two terms of \eqref{eq:geometric-four-terms} are
bounded after weighting by \(Rd_s H_s\). This uses
\(Tc_+g_0^*\le1\) on the contribution to
\(N_\eta/i_\eta\), while the term involving \(P_\eta\)
keeps its factor \(Tc_+\). The boundary and quotient terms
are bounded by \(r_z H_z S_\kappa\) and
\(2r_q P_Q S_\kappa\), respectively. Thus
\[
 \langle r\rangle
 \le R d_s H_s+(r_z H_z+2r_qP_Q)S_\kappa .
\]
Using \(Rd_s\le R_j\), \(r_z\le44s^2R_j/T\) and
\(2r_q\le44s^2R_j/T\) gives the stated inequality. When
\(R_j=0\), the lattice quotient is zero. At \(j=m\) the
auxiliary term was omitted in the local estimate, and one can
use \(r_q=0\) and \(r_z\le22s^2R_j/T\).
\label{cor:run-index-length-2}

\end{proof}

\paragraph{The same sections over the number field.}
The first term of \eqref{eq:geometric-four-terms} uses a
discriminant at the generic intersection. We identify that
discriminant over the number field, using the same sections
as in the branch computations, before applying the global
height formula.
\hypertarget{source-slice-orders}{}
Put \(F=K(X)\), \(L=\prod_Z L_Z\) and \(\delta=\dim_F L\).

\begin{lemma}[One short basis before localization]
\label{lem:slice-short-basis-before-localization}
With a weight frame \(\tau\) nonzero at every point of the
geometric generic cut,
one can select \(\delta\) short forms from
\(s_F(g_0g_1)^{frw_F.n}\operatorname{Sh}(w_F)\) whose restrictions \(b_i\) give a basis over \(F\),
\[
 b_i=(t'_i/\tau)|_L,\qquad L=\bigoplus_{i=1}^\delta F b_i.
\]

The same list works after extension to every \(K_\lambda(X)\),
with the light columns and \(\eta\) fixed.
\end{lemma}
\begin{proof}
Embed \(K\) in \(\mathbb C\) and specialize \(X\) to algebraically
independent complex numbers over its image. This embeds \(F\) in
\(\mathbb C\). Separability and the chart identify
\(L\otimes_F\mathbb C\) with the product over the \(\delta\) distinct
generic intersection points \(z\). Interpolation of
\hyperref[par:slice-interpolation]{arbitrary values} at weight
\(w_F\), with vanishing low modes, gives the delta functions there.
The bounded rational family spans this section space, so choose
\(\delta\) of its members with \(\det(t'_i(z))_{z,i}\ne0\) in
the chosen nonzero frames. In the \(\tau\)-frame,
\[
 \det(b_i(z))_{z,i}
    =\left(\prod_z\tau(z)^{-1}\right)\det(t'_i(z))_{z,i}\ne0 .
\]
Thus the \(F\)-linear map \(F^\delta\to L\) with columns \(b_i\)
is invertible and remains so over \(K_\lambda(X)\). Clearing
polynomial denominators gives the same assertion for rational
section identities. Only the short restrictions were selected.
\label{cor:slice-orders-1}
\end{proof}

\begin{proposition}[Comparison of the two branch lattices]\label{prop:slice-completed-branch-comparison}
Use the one list of short forms from
Lemma~\ref{lem:slice-short-basis-before-localization}, and suppose
their branch orders are in the pulled-back integral frames.
The integral closures over the two generic coefficient DVRs satisfy
\[
 B'_Z=R_1\otimes_{R_0} B_Z.
\]
 The branches correspond
with the same ramification indices and residue degrees. In
the shifted coordinates of the short forms this gives
the local length inequality and the auxiliary order identity
in \eqref{eq:slice-order-comparison}. The auxiliary identity
is used before the definite last lift.
\end{proposition}

\begin{proof}
\proofstep{Comparing the two integral lattices.}
We compare the two computations of the integral lattice. Put
\[
 R_0=(\mathcal O_{K,\lambda}[X])_{(\pi)},\qquad
 R_1=(\widehat{\mathcal O}_{K,\lambda}[X])_{(\pi)},\qquad
 F_1=\operatorname{Frac}R_1 .
\]
Here \(\pi\) can be taken to generate the maximal ideal of
\(\mathcal O_{K,\lambda}\). Both DVRs have residue field
\(k_\lambda(X)\) and ramification index one over the
coefficient DVR. Let \(B_Z\) be the integral closure of
\(R_0\) in \(L_Z\). It is a finite torsion-free module
over \(R_0\). For example, take an integral field basis.
The trace dual of its span contains every integral element,
by separability and integrality of trace. Each localization
of \(B_Z\) at a maximal ideal is a DVR. Its order \(w\) is
normalized to have image \(\mathbb Z\), with ramification
index \(e_w=w(\pi)\) and residue degree \(f_w\).

The integral closure after base change is exactly

\[
  B'_Z=R_1\otimes_{R_0}B_Z
\]
 in \(F_1\otimes_F L_Z\).
In fact this is finite free over \(R_1\), and its generic
fibre is a product of separable fields. The special fibres identify, because both coefficient rings
have residue \(k_\lambda(X)\):
\[
 (R_1\otimes_{R_0}B_Z)/\pi
       =(R_1/\pi)\otimes_{R_0/\pi}(B_Z/\pi)
       \simeq B_Z/\pi .
\]
This bijects the primes above \(\pi\), even if the generic
fibre splits into several fields. At corresponding
primes localize on both sides; write \(t\) for a uniformizer
of \((B_Z)_w\).

Since \(\pi\) is a unit times \(t^{e_w}\)
there, the quotient of the local ring of \(B'_Z\) by \(t\)
is its residue field. This follows by making that quotient
already in the common localized special fibre. Thus the
maximal ideal is principal. The local ring has dimension
one, since the extension is finite over a DVR and \(\pi\)
is a nonzerodivisor. It is a regular local ring, hence a DVR.
The generic localizations are fields. This proves normality
and the asserted equality with the integral closure, as well
as preservation of both \(e_w\) and \(f_w\).
For corresponding DVRs \(T_w\) and \(T'_{w'}\) on the two models, the branch calculation gives
\[
  e_{w'}=e_w,\qquad f_{w'}=f_w .
\]
 In particular
the branch correspondence counts places even if a field
splits after extending coefficients. \cite{dvr-slice}

\proofstep{Returning to the integer branch order.}
For each branch identified in this way let \(n_w(g)\)
be the integer order of the section \(g\) in an integral
frame of its weight. It is infinite for the zero reading.
To compare with the order on \(L_Z\), restrict the valuation
ring to \(\Omega_Z\). Over this valuation ring the generic
point extends by properness. The induced valuation on
\(\Omega_Z\) is discrete up to a positive integer factor
since \(\pi\) has nonzero order.
Indeed \(v(\Omega_Z^\times)\) is a nonzero subgroup of \(\mathbb Z\) containing \(v(\pi)\). For its positive generator \(b\) this gives
\[
  v(\Omega_Z^\times)=b\mathbb Z,\qquad b\mid v(\pi).
\]

Taking the saturated
flag lattices and a generator of the theta line gives an
integral frame there. Its pullbacks to the two branch
valuations are integral frames. Uniqueness in the valuative
criterion compares the underlying extended points and
the frame changes in the same integral flag have unit
diagonals. Thus for any two nonvanishing sections of the
same weight,
\[
 n_w(g)-n_w(\tau)=w(g/\tau),\qquad
 n_w(g)=e_\zeta\operatorname{val}_{y_\zeta}(g)
\]
at the corresponding branch. This argument fixes the
absolute order of a section as well as the order of a
quotient. It is why the adjusted lattices can be compared
at weights other than the light weight.
\label{cor:slice-orders-2}

\proofstep{Comparing shifted lattices in one trivialization.}
Take \(\tau\) of weight \(w_F\), nonzero on each component,
and choose \(d_Z\in L_Z^\times\) with
\(w(d_Z)=n_w(\tau)\) at all places above \(\lambda\),
by weak approximation. The shifted lattice over \(R_0\)
is \(B(\tau)=\prod_Z d_Z^{-1}B_Z\).
The \(b_i\) belong to it by integrality of the \(t'_i\).
Changing \(\tau\) rescales both the ball and the \(b_i\)
componentwise. Under the coordinate map over \(R_1\)
the local lattice of restrictions with nonnegative point
orders lies in \(R_1\otimes B(\tau)\), and the short
columns map to \(1\otimes b_i\). Hence
\begin{equation}\label{eq:slice-order-comparison}
 I_\lambda\le I^{\mathrm{ff}}_\lambda
 :=\operatorname{length}_{R_0}
       (B(\tau)/\textstyle\sum_i R_0 b_i),\qquad
 N_{\eta,\lambda}=\sum_{Z,w\mid\lambda} f_w n_w(\eta).
\end{equation}
The map on lengths uses flatness and the common residue
field. The second assertion is used only before the final
lift. The integer order already includes the ramification
factor relative to the normalized point valuation.

\end{proof}

\paragraph{A discriminant for the shifted lattice.}

After changing the frame, the discriminant in the index formula is that of the shifted branch lattice. Its change is measured by the determinant of the frame map.

\begin{lemma}[Trace discriminant after the frame shift]\label{lem:shifted-slice-discriminant}
For the branch and shifted lattices of
Proposition~\ref{prop:slice-completed-branch-comparison},
let \(\Delta\) be the discriminant of the trace matrix of the
\(b_i\), and \(t_\lambda\) that of the integral closure in
valuation units. Then the normalized Gauss order satisfies
\[
 v_\lambda(\Delta)+2\sum_{Z,w\mid\lambda}f_w n_w(\tau)
                  =2I_\lambda^{\mathrm{ff}}+t_\lambda .
\]
An empty slice has \(\Delta=1\) and both order contributions
zero.
\end{lemma}

\begin{proof}
Let \(\Delta=\det(\operatorname{Tr}_{L/F}(b_i b_j))\).
The trace pairing of a product of separable fields is
nondegenerate. With an integral basis of \(\prod_Z B_Z\)
let \(t_\lambda\ge0\) be the order of its discriminant.
Write \(C\) for the matrix of \(d b_i\) in an integral
basis of the closure, where \(d=(d_Z)_Z\).
The trace matrix in the new basis is \(C^{\mathsf T}GC\),
with \(G\) the trace matrix of the integral basis. Its
determinant has valuation \(t_\lambda+2I_\lambda^{\mathrm{ff}}\)
by Smith form. The change from \(b_i\) to \(d b_i\)
has determinant \(\mathrm N_{L/F}d\), so the same trace
determinant is \((\mathrm N_{L/F}d)^2\Delta\).

Taking valuations gives
\begin{equation}\label{eq:slice-discriminant-order}
 v_\lambda(\Delta)+2\sum_{Z,w\mid\lambda} f_w n_w(\tau)
     =2I^{\mathrm{ff}}_\lambda+t_\lambda .
\end{equation}
Here the order of the norm is the sum of the orders times
the residue degrees. For an integral \(d\) this follows by
filtering the quotient by \(d\) into the residue fields
of the local DVRs; division gives the formula for
fractional \(d\). The Gauss order of the single rational
function \(\Delta\) thus bounds the first price after the
frame correction. An empty slice has zero lattice, empty
discriminant one and zero contributions in both formulas.
\cite{dvr-slice}\label{cor:slice-orders-3}
\end{proof}

\paragraph{The scalar auxiliary section before the definite lift.}
\label{par:selected-auxiliary-cusp}

The definite last lift is treated separately. Before reaching it one needs an auxiliary scalar section with controlled values for the height calculation.

Let \(\Lambda\) be the rank-\(r\) polarized lattice of the fixed lift. For this construction use the odd integer \(R=5\). Write
\(\Gamma_R=\operatorname{Aut}_{\mathcal O_E}(\Lambda/R^4\Lambda)\).
The module modulo \(R^4\) is free of rank \(r\) over
\(\mathcal O_E/R^4\mathcal O_E\), since \(\Lambda\) is projective of
constant rank and the quotient is a product of Artinian local
rings. The kernel of reduction of the matrix group modulo \(R\)
has order \(R^{6fr^2}\). Its elements are all matrices of the
form \(1+A\) with entries of \(A\) in the nilpotent ideal
\((R)/(R^4)\). The reduction contains \(-1\) of order two.

Consequently 
\[
  |\Gamma_R|=2R^4 i_\eta
\]
 for an integer
\(i_\eta>0\). Also \(|\Gamma_R|\le R^{8fr^2}\) by counting all
matrices. Choose lifts \(x_b\), \(1\le b\le r\), of a basis
modulo \(R\). Index tuples by
\(J=(\Lambda/R\Lambda)^r\). Choose representatives of their
entries and put
\[
             x_{l,b}=x_b+R\widetilde l_b\pmod {R^4\Lambda}
             \qquad(l\in J).
\]
For each \(\gamma\in\Gamma_R\), the elements \(\gamma x_{l,b}\)
span the full module modulo \(R^4\): they generate modulo the
nilpotent ideal \((R)\), then generate by Nakayama's lemma.
There are \(R^{2fr^2}\) tuples.
\label{cor:selected-auxiliary-cusp-1}

\begin{auxlemma}[A nonzero theta value on a torsion translate]
\label{lem:eta-torsion-translate}
Let \(A/\mathbb C\) be a principally polarized abelian variety, let
\(L\) represent its polarization and let \(s\) generate \(H^0(A,L)\).
For every integer \(R\ge2\), the torsion translates have empty common zero locus:
\[
 \forall x\in A(\mathbb C)\quad\exists t\in A[R]\quad s(x+t)\ne0 .
\]
\end{auxlemma}
\begin{proof}
Write \(p=[R]\). Its deck group is \(A[R]\), so the character
decomposition of the finite étale cover and the projection formula give
\[
 H^0(A,p^*L)=
 \bigoplus_{P\in\operatorname{Pic}^0(A)[R]} H^0(A,L\otimes P).
\]
Indeed the character sheaves of \(p_*\mathcal O_A\) are the line
bundles trivialized by \(p\), each with multiplicity one. The
principal polarization identifies the indexing sets by
\[
 A[R]\xrightarrow{\ \sim\ }\operatorname{Pic}^0(A)[R],
 \qquad t\longmapsto T_t^*L\otimes L^{-1}.
\]
Thus each summand is spanned by the pullback of a translate of \(s\).

The line \(p^*L\) is generated by its sections. To see this, the
theorem of the square gives \(p^*L\simeq L^{R^2}\otimes P_0\) for
some \(P_0\in\operatorname{Pic}^0(A)\). At any specified point take
a product of \(R^2\) translates of \(s\), with the sum of the
translation points fixed by \(P_0\). Each individual vanishing
condition cuts out a proper closed subset of that fiber of the
addition map, which is isomorphic to \(A^{R^2-1}\). A product
avoiding them all has the required line bundle and a nonzero value.
Consequently at a lift of \(x\) some character summand is nonzero.
This gives the claimed translate.
\end{proof}

\begin{auxlemma}[Decay along a shortest CM period]
\label{lem:eta-period-decay}
Let \(V/L\) be a complex torus of dimension \(g=fr\) with its
principal Riemann form \(B=\operatorname{Im}H\), Euclidean norm
\(\|v\|^2=H(v,v)\), and polarized \(\mathcal O_E\)-action with
Rosati involution given by conjugation. Let
\(u\in\mathcal O_E\) satisfy \(u+\bar u=1\). Use the theta line
whose semicharacter is \(k\mapsto(-1)^{B(k,uk)}\), and let
\(s\) be a section of Haar \(L^2\)-norm one. Write
\(\lambda=\min_{0\ne k\in L}\|k\|\),
\(\rho=|\operatorname{disc}E|^{1/(2f)}\) and
\(\Omega=2f^{3/2}\rho\). Suppose that
\[
 M\ge\rho^2,\quad \rho\ge16,\quad
 M^{-1/2}\le h\le 1/\Omega,\qquad
 P=\begin{cases}
 2f/(\rho^2\lambda^2),&\lambda<h,\\
 0,&\lambda\ge h.
 \end{cases}
\]
Let \(R\ge2\), and let \(a_1,\ldots,a_r\in V\) satisfy
\(R^2a_b\in L\) with these periods generating \(L/RL\) over
\(\mathcal O_E\). Put
\[
 C_0=r\bigl[3+g\log(2\sqrt{gM})
                  +2f\log(2\sqrt{2f})\bigr].
\]
For the pointwise theta metric,
\begin{equation}\label{eq:eta-period-decay}
 \prod_{b=1}^r\|s(a_b)\|
 \le e^{C_0}\max(1,P)^{rg/2}
                 \exp\!\left(-\frac{\pi P}{2R^4}\right).
\end{equation}
\end{auxlemma}
\begin{proof}
\proofstep{A bound valid on the whole torus.}
Put \(S=\sup_z\|s(z)\|\) and \(b_0=\min(1,\lambda/2)\).
A polydisc of coordinate radius \(b_0/\sqrt g\) in orthonormal
coordinates injects into the torus. Center the Appell--Humbert
frame at the point in question. On this polydisc the squared
metric is \(e^{-\pi\|z\|^2}\) times the squared absolute value.
The principal alternating form makes the period covolume one in
these coordinates, so Euclidean measure on a fundamental domain
is Haar probability measure. Submean on the polydisc gives
\[
 S\le e^{\pi/2}\left(\frac{\sqrt g}{\sqrt\pi\,b_0}\right)^g
 \le e^2 (2\sqrt{gM})^g\max(1,P)^{g/2}.
\]
In fact \(\lambda^{-1}\le\sqrt M\) when \(P=0\), and
\(\lambda^{-1}=\rho\sqrt{P/(2f)}\) otherwise. This already proves
\eqref{eq:eta-period-decay} for \(P=0\).

\proofstep{An isotropic subspace and its frequencies.}
Choose a shortest period \(\xi\) when \(\lambda<h\). There are
integers \(\omega_j\in\mathcal O_E\) spanning \(E\) over
\(\mathbb Q\) with \(\max_\sigma|\sigma\omega_j|\le\Omega\).
Indeed, equip \(E\otimes\mathbb R\) with the squared norm
\(\sum_\sigma|\sigma x|^2\). The dual of \(\mathcal O_E\)
is its codifferent for this metric. Every nonzero member of the
codifferent satisfies
\[
 |N_{E/\mathbb Q}y|\ge |\operatorname{disc}E|^{-1},
 \qquad \sum_\sigma|\sigma y|^2
       \ge 2f |N_{E/\mathbb Q}y|^{1/f}\ge 2f/\rho^2 .
\]
Thus the dual first minimum is at least \(\sqrt{2f}/\rho\).
Euclidean transference gives last minimum of \(\mathcal O_E\)
at most \(2f\) divided by the dual first minimum
\cite{banaszczyk}. This is at most \(\Omega\).

For each \(j\), integrality of the alternating form and the
bound on the operator norm of \(\omega_j\) give
\[
 |B(\xi,\omega_j\xi)|\le\Omega\lambda^2<1,
 \qquad B(\xi,\omega_j\xi)=0.
\]
By linearity and the Rosati identity, \(W=\operatorname{span}_{\mathbb R}
(\mathcal O_E\xi)\) is isotropic. Set \(I=W\cap L\) and
\(n=\dim_{\mathbb R}W\le2f\). The semicharacter is one on \(I\),
since \(uI\subset I\). With \(I^*\) the Euclidean dual in \(W\),
every nonzero \(p\in I^*\) has nonzero integral pairing with some
\(\omega_j\xi\). Hence
\[
 \|p\|\ge(\Omega\lambda)^{-1}\ge1.
\]

\proofstep{The Fourier expansion on the isotropic subspace.}
For \(a\in V\) let \(U(a)\in W\) represent
\(k\mapsto B(a,k)\) by the Euclidean inner product. In the
Appell--Humbert frame centered at \(a\), restrict the section to
\(W+iW\). Multiply its holomorphic representative by
\(\exp(-\pi(z,z)/2)\), where \((\ ,\ )\) is the complex bilinear
extension of the real inner product on \(W\).
The resulting function \(G_a\) satisfies
\[
 |G_a(0)|=\|s(a)\|,\qquad
 |G_a(x+iy)|\le S e^{\pi\|y\|^2}\qquad(x,y\in W).
\]

Its translation multiplier on \(I\) is \(e^{\pm2\pi iB(a,k)}\),
with the sign fixed by the convention for \(H\). For a frequency
\(p\in I^*\pm U(a)\), move the Fourier coefficient integral from
\(W/I\) to \((W-ip)/I\). Opposite faces cancel by the multiplier
law, and the exponential from the Fourier character gives
\[
 |\widehat G_a(p)|\le S e^{\pi\|p\|^2-2\pi\|p\|^2}
                    = S e^{-\pi\|p\|^2}.
\]

These bounds give absolute convergence at the origin. Reversing all
frequencies if necessary, we may therefore sum over \(U(a)+I^*\).

\proofstep{A norm bound for every tuple of frequencies.}
Take arbitrary \(p_b\in U(a_b)+I^*\). Since \(I\) is saturated in
\(L\), every integer functional on \(I\) extends to \(L\). By the
principal polarization \(U(L)=I^*\), so we can write
\[
 R^2p_b=U(k_b),\qquad k_b=R^2a_b+R^2t_b,\quad t_b\in L .
\]
The \(k_b\) still generate \(L/RL\). The integral functional
\(\alpha\mapsto B(k_b,\alpha\xi)\) on \(\mathcal O_E\) has the
form \(\operatorname{Tr}_{E/\mathbb Q}(y_b\alpha)\) for an element
\(y_b\) of the codifferent. At least one \(y_b\) is nonzero.
Otherwise, by the Rosati identity, every element of
\(\sum_b\mathcal O_E k_b+RL=L\) would pair with \(\xi\) in
\(R\mathbb Z\). Perfectness excludes this because a shortest
period is primitive.

The norm of this trace functional for the embedding metric is
\((\sum_\sigma|\sigma y_b|^2)^{1/2}\). Its value at \(\alpha\) is
\((R^2p_b,\alpha\xi)\), and
\(\|\alpha\xi\|\le\|\xi\|(\sum_\sigma|\sigma\alpha|^2)^{1/2}\).
For a nonzero \(y_b\) the preceding codifferent estimate gives
\[
 R^4\|p_b\|^2\lambda^2\ge
 \sum_\sigma|\sigma y_b|^2\ge 2f/\rho^2,
 \qquad \sum_b\|p_b\|^2\ge P/R^4 .
\]

Reserve half of each Gaussian exponent for this last inequality.
The rest is summable uniformly on every translate \(D\) of \(I^*\).
Indeed the balls of radius \(1/2\) about \(p\in D\) have disjoint
interiors. On such a ball \(\|p\|^2\ge\|x\|^2/2-1/4\), whence
\[
 \sum_{p\in D} e^{-\pi\|p\|^2/2}
 \le \frac{e^{\pi/8}}{\operatorname{vol}B_n(1/2)}
                  \int_{\mathbb R^n}e^{-\pi\|x\|^2/4}\,dx
 \le e^{\pi/8}(2\sqrt n)^n .
\]

The last inequality uses the inscribed cube of side \(1/\sqrt n\).
Multiplying the absolutely convergent Fourier bounds for all
\(b\), then using \(n\le2f\) and the bound for \(S\), proves
\eqref{eq:eta-period-decay}.
\end{proof}

\begin{auxlemma}[Parameter bounds for the two theta estimates]
\label{lem:joint-theta-parameters}
Let \(g=fr\), \(D=D_F\) be the differential degree of the light weight,
\(M=g_0g_1\), and \(\rho=\operatorname{rd}(E)\). Suppose
\[
 \begin{gathered}
 f,r,m,D\ge1,\quad g\ge2,\quad m^4\le f,\quad \rho\ge\max(16,f^2),\\
 R_0^2=90(g+D),\qquad M\ge R_0^2\rho^{28}Dr .
 \end{gathered}
\]
With \(d=2fm\), define
\[
 \Omega=2f^{3/2}\rho,\qquad h^2=fR_0^2\rho^{21}/M,\qquad
 B=M/(R_0^2\rho^{28}),\qquad T_0=MB .
\]
Then \(M^{-1/2}\le h\le\Omega^{-1}\), and all six estimates used
in the lower bound for the light determinant hold simultaneously:
\begin{align}
 d^2\Omega h^2&<1, & B&\ge1,\label{eq:joint-theta-smallness}\\
 90m^2d^2\rho&\le M/2, & D&\le M,\label{eq:joint-theta-dimension}\\
 4D(7+\log(MD))+d+1 &\le T_0,\label{eq:joint-theta-log-cost}\\
 g\log2+2d+3+(D+1)(r\log r+3r+3)&\le T_0.\label{eq:joint-theta-rank-cost}
\end{align}
\end{auxlemma}
\begin{proof}
We have \(d\le2\rho\), \(m^2\le\rho\), and
\(f^{3/2}\le f^2\le\rho\). The common level bound gives
\[
 B\ge Dr,\qquad T_0\ge MDr,\qquad R_0^2\ge270.
\]
In particular \(D\le M\) and \(h\sqrt{M}\ge1\). The two
smallness requirements use different powers of \(\Omega\):
\[
 d^2\Omega h^2\le\frac{8R_0^2\rho^{26}}M\le\frac8{\rho^2}<1,
 \qquad (\Omega h)^2\le\frac{4 f^4}{\rho^5Dr}\le1.
\]
Also \(90m^2d^2\rho\le360\rho^4\le M/2\).

For the logarithmic allowance use \(D\le M\) and
\(\log M\le M/32\) for \(M\ge256\), by monotonicity of
\(t/32-\log t\) in that range. Since \(M\ge270\rho^{28}\),
\[
 4D(7+\log(MD))+d+1
 \le D(29+2\rho+M/4)\le MD\le T_0 .
\]
Finally \(r\le g\), \(\log r\le g\), and \(D+1\le2D\) give
\[
 g\log2+2d+3+(D+1)(r\log r+3r+3)
 \le (3g+4\rho+15)Dr\le MDr\le T_0 ,
\]
using \(M\ge90(g+D)\rho^{28}Dr\) in the second inequality.
\end{proof}

\begin{proposition}[The scalar cusp section before the definite lift]
\label{prop:auxiliary-scalar-section}
For each nondefinite lift the scalar construction gives a cusp section
\(\eta\) nonzero on every component with the same finite and
archimedean boundary comparisons used for the light columns. For the
page parameters, one may take the decay constant \(P_2\) and the upper bound \(P_\eta\) for the mean log supremum divided by \(i_\eta\) to satisfy
\[
 \max\{P_2,P_\eta\}\le200000(frL_0)^2,
\]
where \(L_0\) is the run's logarithmic scale.
\end{proposition}

\begin{proof}
\proofstep{The scalar section before descent.}
Over the characteristic-zero torsor of full linear level \(R^4\), denote the
normalized odd-level theta null values by \(A(t)\) for
\(t\in\Lambda/R^4\Lambda\). They are the values without differentiation at the point \(t\) of order dividing \((R^2)^2\)
with the odd-level normalization of order \(R^2\). Use
\begin{equation}\label{eq:selected-eta-product}
 \eta=\prod_{\gamma\in\Gamma_R}
             \left(\sum_{l\in J}m_l
                       \prod_{b=1}^{r}A(\gamma x_{l,b})\right).
\end{equation}
The frame law gives each theta null value central degree one
and zero row degrees. The torsor action permutes the outer
factors, so their product descends to the original coefficient
field and level. Its weight is 
\[
  r|\Gamma_R|u_0=i_\eta g_0^*u_0,
\]
with \(g_0^*=2\cdot5^4r\).
\label{cor:selected-auxiliary-cusp-2}

Choose one frame at the chosen complex embedding on each
component. For a torsion label \(t\), rigidification gives a nonzero frame
at \([R^2]t\), hence
\[
 A(t)\ne0\quad\Longleftrightarrow\quad s([R^2]t)\ne0.
\] At each tested frame and each \(\gamma\),
the shifts in a single column give the following points:
\[
  R^2(R\widetilde l_b)=R^3\widetilde l_b,\qquad R(R^3\widetilde l_b)=0\pmod{R^4}.
\]
As \(\widetilde l_b\) ranges through representatives modulo \(R\),
these last classes cover all of \(A[R]\), including after the
automorphism \(\gamma\). Applying Lemma~\ref{lem:eta-torsion-translate}
at \([R^2]\gamma x_b\) gives a nonzero value for that column. The shifts
of the columns are independent and all occur in \(J\); hence the
linear functional on \((m_l)_l\) in that factor is nonzero.

With \(M_c\) tested frames, the box \(0\le m_l\le M_c|\Gamma_R|\)
contains an integral point on which every such functional is
nonzero. Set \(B=M_c|\Gamma_R|\). A proper hyperplane contains at most \((B+1)^{|J|-1}\) points of the box, so the number left after excluding all \(B\) hyperplanes is at least
\[
  (B+1)^{|J|}-B(B+1)^{|J|-1}=(B+1)^{|J|-1}>0.
\]
 The odd-level test requires
characteristic zero, the full \(R^4\) torsion frame and the
rigidified theta line. All occur in the chosen complex uniformization. For the chosen test point \(x_a\) on component \(a\) the coefficients give
\[
 \prod_{\gamma\in\Gamma_R}
   \left(\sum_l m_l\prod_b A(\gamma x_{l,b})(x_a)\right)\ne0.
\]
Thus \(\eta\) is nonzero on every component. They are used at every finite and infinite
place.
\label{cor:selected-auxiliary-cusp-3}

\proofstep{Finite orders in an integral frame.}
At a finite algebraic test point, pass to a split semiabelian
uniformization in a finite valued extension. Write \(Q\) for
the nonnegative quadratic form of the integral theta expansion in
the degeneration lattice and \(y_t\) for the lattice class of a torsion evaluation at level \(R^4\). In a frame with integral theta
generator, we claim that
\[
       \operatorname{ord} A(t)
       \ \ge\ \frac{1}{R^4}
                   \min_z Q(y_t+R^2z).
\]

To compute the order, let \(m=R^2\), and put
\(B_Q(y,z)=Q(y+z)-Q(y)-Q(z)\). Use the period and character
lattices identified by the principal polarization. In the symmetric
Mumford degeneration data, a period \(y\) contributes
\(B_Q(y,z)\) to the order of character \(z\) and \(Q(y)\) to
the order of its theta multiplier. Indeed the integrally normalized Poincar\'e
biextension has pairing of order \(B_Q(y,z)\). Pullback along
the graph of \(u\) gives for its multiplier at \(y\)
\[
 B_Q(y,uy)=B_Q(y,\bar u y)=\tfrac12B_Q(y,y)=Q(y).
\]

This identifies the quadratic refinement of this particular line.
The coefficient transformation and its compatibility with the
integral metric are given in
\cite[Proposition~3.19 and Corollary~3.21]{frss-nonarch-theta}. Write \(a_z\) for its coefficient of index \(z\).
Translation by the period \(z\) carries the integral abelian-part
section at index zero to that at index \(z\), with multiplier
of order \(Q(z)\). In the integral frame this gives
\[
                 \operatorname{ord}(a_z)\ge Q(z).
\]

Write the original torsion lift as \(t\), with \([m^2]t\) equal
to the period \(y_t\) on the uniformizing group. Character orders
are additive under multiplication. At \([m]t\) the character
\(z\) contributes \(B_Q(y_t,z)/m\). The symmetric torsion frame
contributes \(Q(y_t)/m^2\): its square has valuation
\(B_Q(y_t,y_t)/m^2\) by the symmetric biextension pairing after clearing \(m^2\),
and the rigidified cubical isomorphism \([m]^*L\simeq L^{m^2}\) is integral. Thus
the term of index \(z\) has order bounded below by
\[
 Q(z)+\frac{B_Q(y_t,z)}m+\frac{Q(y_t)}{m^2}
       =\frac{Q(y_t+mz)}{m^2}.
\]
All orders are measured in the uniformizer of the same split chart.
The terms of the convergent theta expansion lie in the closed
valuation ball specified by their common lower bound. Their sum
lies in that ball as well, even when some terms cancel.

For a spanning tuple of torsion labels the classes \(y_t\) span the
degeneration lattice modulo \(R^2\). Indeed the period map from
geometric \(m^2\)-torsion to \(Y/m^2Y\) sends a lift to its
\(m^2\)-multiple modulo periods. It is surjective: each period
has an \(m^2\)-division point in the semiabelian uniformizing
group over the algebraic closure of the fraction field. It
commutes with the \(\mathcal O_E\)-action. The resulting surjections, with \(A\) the degenerating abelian
variety, are \(\mathcal O_E\)-equivariant:
\[
 A[m^2](\overline K)\twoheadrightarrow Y/m^2Y\twoheadrightarrow Y/mY .
\]
Thus their images generate \(Y/mY\). Replacing each image by a minimizing
representative preserves its class in this quotient. Since the lattice is a finite module over
\(\mathcal O_E\), these representatives span after tensoring
with \(E\) by Nakayama at a prime above \(R\).
Write \(v_b\) for these representatives. On balanced integral
symmetric bilinear forms \(C\) use the integral linear functional
\[
 \ell(C)=\sum_b C(uv_b,v_b)
        =\tfrac12\sum_b C(v_b,v_b),\qquad u+\bar u=1 .
\]

The equality follows from symmetry and
\(C(ux,y)=C(x,\bar u y)\). If \(C\ge0\) and \(\ell(C)=0\),
all \(v_b\) lie in the radical, which is \(\mathcal O_E\)-stable.
Their \(E\)-span is the whole rational lattice, so \(C=0\).
Thus \(\ell\) is one of the integral functionals allowed in the
infimum defining \(\Phi\). Evaluating the functional on \(B_Q\) and using the preceding
termwise bound gives
\[
 \Phi\le\ell(B_Q)=\sum_b Q(v_b)
       \le R^4\operatorname{ord}\!\left(\prod_b A(\gamma x_{l,b})\right).
\] The integer coefficients and the
ultrametric inequality give that same lower bound for the
whole factor. Multiplication over \(\Gamma_R\) gives
\(2i_\eta\Phi\). Descent uses the normalized valuation on the
coefficient field: if the extension ramification is \(e\),
both the raw order and the boundary form computed using the
extension's uniformizer acquire a factor \(e\).

If \(V\subset V'\) are the two DVRs, the numerical descent is
\[
  {\rm ord}_{V'}(a)=e\,{\rm ord}_{V}(a),\qquad {\rm ord}_{V}(a)=e^{-1}{\rm ord}_{V'}(a)\qquad(a\in V,\ a\ne0).
\]
In particular \(v_\lambda(\eta)\ge2i_\eta\phi_\lambda\)
in the relative uniformizer units of the finite light
comparison. A zero value has infinite order. The section
is integral at all algebraic tests and therefore belongs
to the integral cusp lattice once cuspidality is established below.
\label{cor:selected-auxiliary-cusp-4}

\proofstep{The complex norm on both halves.}
For a good complex embedding take a unit frame. With \(d\mu\) the Haar probability measure of the torus, the chosen unit frame normalizes the Appell--Humbert section \(\vartheta\) by
\[
  \int\|\vartheta(z)\|^2\,d\mu(z)=1.
\]
The absolute value of the normalized theta null value is the norm
of this section at the torsion point represented by the period
divided by \(R^2\). Multiplying these period representatives
by \(R^2\) produces a spanning family modulo \(R\) for every
tuple in \eqref{eq:selected-eta-product}. In this chart the
polarization has semicharacter \((-1)^{B(k,uk)}\). Indeed the
rigidified theta line is obtained by pulling the Poincar\'e line
back along \(z\mapsto(z,uz)\) using the polarization. Its sign
on a period is the pairing of the two period components modulo two.
The two axis rigidifications determine this sign, as in
Lemma~\ref{lem:placed-row-theta}.

Take \(g=fr\), \(M=g_0g_1\), \(R_0^2=90(g+D_F)\), and
\(\rho=\operatorname{rd}(E)\). The boundary height used on the
selected page is the function \(P\) of Lemma~\ref{lem:eta-period-decay}
at the threshold
\[
 h=\sqrt{f R_0^2\rho^{21}/M}.
\]

The page has \(D_F,m\ge1\), \(fr\ge2\), \(m^4\le f\),
\(\rho\ge\max(16,f^2)\), and
\begin{equation}\label{eq:selected-torus-level}
 M\ge R_0^2\rho^{28} D_F r .
\end{equation}

To verify the level bound on the chosen lift, use the same size
parameter \(\mathcal P=C_2z^{A_2}\) as in the proof of
Lemma~\ref{lem:short-weight-curvature}. The construction takes
\(\mathcal P\ge90(fr+1)\) and \(\mathcal P\ge r\); its prime
inequality is \(\rho^{28}\mathcal P^3 D_F^2\le\min(g_0,g_1)\).
Since \(R_0^2\le\mathcal P D_F\) and \(r\le\mathcal P^2\), we obtain
\[
 R_0^2\rho^{28}D_F r\le\rho^{28}\mathcal P^3D_F^2
                    \le\min(g_0,g_1)\le M .
\]

At the nonterminal signature the smaller block has size
\(m-j\le m\). Lemma~\ref{lem:joint-theta-parameters} now applies
both to the light-family estimate in
Lemma~\ref{lem:scaled-light-metric} and to the decay estimate at
this same threshold. In particular \(M\ge\rho^2\) and
\(M^{-1/2}\le h\le\Omega^{-1}\) as required by
Lemma~\ref{lem:eta-period-decay}.
\label{cor:selected-auxiliary-cusp-5}

Write \(A_0=R^4rg\ge1\). For \(P>1\), concavity of the logarithm
gives \(\log P\le\log A_0+P/A_0\). For \(0\le P\le1\) the same upper
bound holds for \(\log\max(1,P)=0\). Since \(\pi\ge2\),
\eqref{eq:eta-period-decay} yields, for every tuple in
\eqref{eq:selected-eta-product},
\[
 \prod_b|A(\gamma x_{l,b})|
 \le\exp\!\left(C_0+\frac{rg}2\log A_0
                          -\frac{P}{2R^4}\right).
\]
The triangular change to a unit frame preserves this boundary
function, and the flag metric gives exactly the norm of
the value in a unit frame.

The Lie factor is one at a scalar
weight. Using \(|J|\le R^{2fr^2}\) and
\(0\le m_l\le B\le B_1=\max(1,M_c)R^{8fr^2}\) in the triangle
inequality, we get
\begin{align*}
   \log\|\eta\|&\le i_\eta(p_{16}-P),\\
   p_{16}&=2R^4(C_0+\log B_1+2fr^2\log R)+A_0\log A_0.
\end{align*}
The bound applies on each good embedding, and on the others
the norms come from conjugation.

To see cuspidality directly, fix a rational boundary component and
a primitive period in its isotropic lattice. In a Witt basis one
may approach this component with the real split diagonal acting
by \(e^{-t}\) on the isotropic space and by \(e^t\) on its dual.
The unipotent radical fixes that isotropic space pointwise.
For fixed Levi coordinate and uniformly over the compact unipotent
period of integration, the first minimum therefore satisfies
\[
 \lambda_1(t)\le C e^{-t}<h,\qquad
 P(t)\ge\frac{2f}{\rho^2 C^2}e^{2t}\qquad(t\gg1).
\] The change from
the metric frame to the holomorphic Fourier frame has at most
exponential growth in \(t\): the determinant characters of the
scalar weight are fixed powers of the diagonal entries. Hence
\[
 \sup_{\rm period}|\eta|_{\rm Fourier\ frame}
      \le C_1\exp(C_2t-C_3e^{2t})\longrightarrow0
      \qquad(C_3>0).
\]

The coefficient of index zero in the normal Fourier variables is
their period average. It is independent of the imaginary
translation in those variables by holomorphy and periodicity,
so taking this limit makes it zero. This holds at each rational
boundary component. A component with no rational boundary has
no constant-term condition. Together with the finite orders this
puts \(\eta\) in the integral cusp lattice.
\label{cor:selected-auxiliary-cusp-6}

\proofstep{The size bound in the selected parameters.}
For the explicit size, denote the run's logarithmic scale by
\(L_0\) and put \(Z=frL_0\). The page has \(L_0\ge2\),
\(\log r\le L_0\), \(\log f,\log\operatorname{rd}(E)\le 2L_0\),
\(\log M\le70L_0\), \(m\le r\), and \(\log M_c\le41Z^2\).
For \(R=5\), \(r\ge2\), and \(g=fr\), they give
\[
 C_0\le fr^2(8+39L_0),\qquad
 \log A_0=4\log5+\log f+2\log r\le7+4L_0.
\]

Indeed \(\log g\le3L_0\), \(2f\le g\), and
\(\log(2\sqrt{2f})\le2+L_0\). As \(\log5<2\), it follows that
\[
 p_{16}\le1250[41Z^2+fr^2(28+39L_0)]
             +625fr^2(7+4L_0)
        \le 200000 Z^2 .
\]

For the last inequality use \(fr^2L_0^2\le Z^2\) and
\(L_0\ge2\). This supplies both the constant \(P_2\) in the
decay estimate and the constant \(P_\eta\) in the mean
log supremum divided by \(i_\eta\). The same scalar form
supplies the finite order comparison. It is selected before
the small short basis on the incidence field.
\label{cor:selected-auxiliary-cusp-7}
\end{proof}

\paragraph{Integral sections for the height comparison.}
\label{par:selected-height-sections}
\hypertarget{source-selected-height-sections}{}
\begin{established}[Complementary weights and the gaps of the selected rows]
\label{hyp:height-section-weights}\label{cor:selected-height-sections-1}
Fix the coefficient field, the two level primes and either orientation.
Write \(u_0\) for the weight with zero row entries and central
coordinate one in the group of weights. The auxiliary section before
the last lift has weight \(i_\eta g_0^*u_0\). Put
\[
 w_1=2(w_L-u_0),\qquad w_2=(s+4)w_L-w_F,\qquad
 d=s+2,\qquad T'=T+d .
\]
The two central degrees are \(2n_W-2\) and \((s+4)n_W-n_P\).
The run satisfies
\[
 8fr+4=n_P\le n_W,\qquad (s+4)\operatorname{DF}(w_L)+2
       \le\min(g_0,g_1),\qquad s\ge3,\qquad T\ge40d .
\]
Every interior adjacent difference \(a\) in a nonempty block of
\(\kappa\) satisfies \(a+2\ge T'(T'+2r)\). Each terminal entry \(a\)
satisfies \(a\ge2T^2\); its offset is \(d_0=dn_W+2r\) in the
positive block and zero in the negative block, with \(0\le d_0\le T\).
\end{established}
\begin{proof}
The central degrees follow by addition of weights.
Lemma~\ref{hyp:common-slice-scales} gives the degree and order bounds.
Since \(D_L=\operatorname{DF}(w_L)\le 2frb\), \(s,f,r,k\) are
polynomially bounded in \(x\), and the two primes satisfy
\(g_i\ge(kb)^8\), the level inequality follows at all large scales.

We give the row argument, including the margin needed at a block end.
Let \(h=\eta_{\rm size}/m\) be the spacing of the central parameters,
and \(\Delta\) the separation parameter of the outer ones. In the
successive choices of Paragraph~\ref{par:window-successive-bounds}
one can impose
\[
       h\ge4T^2+2t,\qquad
       \Delta\ge(T+D_*+2)(T+D_*+2+2r).
\]
Indeed these lower bounds are fixed polynomials in \(x\) after
Lemma~\ref{lem:rank-window}. In choosing the width and separation
before rounding the nodes, adjoining either polynomial to their lower
bounds costs \(O(\log x)\) on the logarithmic scale; the available
upper bounds there are \(x^{\epsilon/4}\) for \(\log\eta_{\rm size}\)
and \(x^{\epsilon/4}/2\) for \(\log\Delta\). Equivalently the
degree-power lower bounds suffice with exponent three. For
\(N=k+m(t+1)\), Lemma~\ref{hyp:common-slice-scales} gives
\[
    N\ge t\ge T,\qquad
    (T+D_*+2)(T+D_*+2+2r)\le(41T/40)(81T/40)\le N^3 .
\]
These choices give \(G:=T'(T'+2r)\le\min(\Delta,h)\), since
\(d\le T/40\) and \(2r\le T\).

Fix an embedding and number its central nodes \(b_0,\ldots,b_{m-1}\).
Their spacings and the corresponding integral centres \(e_i\) satisfy
\[
     h\le b_{i+1}-b_i\le h+1,\qquad |e_i-b_i-h/2|\le2 .
\]

The two errors come from rounding the nodes and their half-shifts.
The decreasing list for the positive block consists of outer
parameters \(v_0,\ldots,v_{k-2}\) followed by
\(-b_0,\ldots,-b_{j-1}\). Within the outer list consecutive differences
are at least \(\Delta\), by its spacing, and the central differences
are at least \(h\). At the junction \(v_{k-2}\ge0\); the lower endpoint
of the translated windows gives \(e_0\ge t\eta_{\rm size}\). Thus
\[
   v_{k-2}-(-b_0)\ge b_0\ge e_0-\eta_{\rm size}
                   \ge(t-1)\eta_{\rm size}\ge h .
\] In the reversed negative block the consecutive
unshifted differences are again differences of adjacent \(b_i\)'s.
Subtracting the row shift changes an interior difference by one; in
particular \(a+2\ge G\) in either block.

At lift \(1\le j\le m\), substitution of the last row in the two
weight formulas gives respectively
\[
 e_{j-1}-b_{j-1}-(t+1)/2,\qquad
 b_j-e_{j-1}-(t+1)/2 \quad(j<m).
\]
The centre error and the lower bound on node spacing make both at
least \((h-4-t-1)/2\ge2T^2\), since \(t\ge5\).
The positive offset is at most \(T\) by
Lemma~\ref{hyp:common-slice-scales}; the empty last block at \(j=m\)
requires no terminal inequality.
\end{proof}

\begin{proposition}[Bounded complementary sections at the selected weights]\label{prop:height-complementary-sections}
For the run of Lemma~\ref{hyp:height-section-weights}, the weights \(w_1,w_2\) used at the
corresponding lifts are admissible. There are integral
sections \(t_1,t_2\) at these weights, nonzero on every
component, with
\[
  \tfrac12\log\|t_1\|_{\sup,\sigma}\le P_4 sS_L,\qquad
              \log\|t_2\|_{\sup,\sigma}\le P_4 sS_L
\]
with \(P_4=\max\{1,5(1009Z_b+\lambda_*)+162 Z_b^2/s\}\),
where \(Z_b=frL\) is the parameter of Proposition~\ref{prop:slice-evaluation-net} and
\(\lambda_*\) bounds the logarithmic cost of the integralizing scalar
per unit of weight size. The choices use the fixed
coefficient field and light columns. At the definite lift
only \(t_2\) is needed.
\end{proposition}

\begin{proof}
\proofstep{The allowable row-entry shifts.}
Here is the row-entry check. Put \(d=s+2\) and \(T'=T+d\).
At an interior gap in either sign block let \(a\) be the gap of
\(\kappa\), and \(l=\lfloor(a+2)/T'\rfloor\) the corresponding
gap of \(w_L\). The page gives \(a+2\ge T'(T'+2r)\). Write
\(a+2=T'l+\rho\) with \(0\le\rho<T'\). The gap of \(w_2\) is
\[
 (T'+2)l-a=2l-\rho+2>0.
\] At a last entry \(a\ge2T^2\) of
\(\kappa\), put \(d_0=dn_W+2r\) in the positive block and
\(d_0=0\) in the negative block.

Now
\(l=\lfloor(a-d_0)/T'\rfloor\) and \(0\le d_0\le T\).
Because \(T\ge40d\), we have
\[
 T'\left(1+\tfrac32T\right)\le 2T^2-T .
\]
Write \(a-d_0=T'l+\rho\) with \(0\le\rho<T'\). The last entry
of \(w_2\) in this block is positive, since
\[
 l=\left\lfloor\frac{a-d_0}{T'}\right\rfloor\ge\frac{3T}{2},
 \qquad 2l-d_0-\rho>3T-T-T'=T-d>0 .
\]
The rows of \(w_1\) are twice the rows of \(w_L\).
The row and central bounds for \(w_L\), with \(2,s+4\le T\),
give \(w_1,w_2\le\kappa\) componentwise.
For either weight \(u\), the rank and level tests follow from
\[
 \begin{aligned}
 n(u)&\in\{2n_W-2,(s+4)n_W-n_P\}, &n(u)&\ge n_W\ge r,\\
 \operatorname{DF}(u)+2&\le (s+4)\operatorname{DF}(w_L)+2,
 &\operatorname{DF}(u)+2&\le\min(g_0,g_1).
 \end{aligned}
\]
In particular 
\[
  g_0g_1\ge(\operatorname{DF}(u)+2)^2
\ge4\operatorname{DF}(u)+4.
\]
 The size bounds are
\(\lvert w_1\rvert\le2S_L\) and \(\lvert w_2\rvert\le(s+4)S_L\).
\label{cor:selected-height-sections-2}

\proofstep{The two sections from the nonvanishing jet test.}
Fix \(u\in\{w_1,w_2\}\) among the weights needed at this lift.
Work at a good complex embedding. On each component \(a\) choose an
orthonormal frame. The positive last entries and \(n(u)\ge r\)
allow the block determinant harmonic polynomial. Let \(D_u\) denote its nonzero leading differential operator.
At the level just checked, Lemma~\ref{lem:ideal-lattice-jet} gives
\[
 d=\operatorname{ord}(D_u)\le\operatorname{DF}(u),\qquad
        D_u\theta_\nu(0)\ne0\quad\text{for some }\nu .
\]
The corresponding theta evaluation is linear in the polynomial.
The bounded grid of determinant polynomials in the proof of
Proposition~\ref{prop:three-shifted-spanning-sets} spans the same
harmonic space. Fix a translate \(\nu\) on which the jet is nonzero. If every grid
polynomial vanished on this translate, linearity would give zero for the
original polynomial there as well. Choose one grid rule \(F_a\) with
nonzero value at this frame and translate.

Let \(A_u\) be the sum of the flag entries and put \(h_u=frn(u)\).
The grid estimate and the complex theta norm calculation in that
proposition use only \(d\le A_u\), \(r\le n(u)\) and the logarithmic
size bounds at \(u\). These hold by the coordinatewise bounds
\(u\le\kappa\) above. Before applying the integralizing scalar they give
at every embedding \(\sigma\)
\[
 \log\|F_a\|_{\sup,\sigma}
 \le 140rL A_u+10L A_u+200h_uL
 \le 400 Z_b |u| .
\]
Let \(x_a\) be the chosen frame on component \(a\), and put
\(C=\{0,\ldots,M_c\}^{M_c}\).

For fixed \(m_b\) with \(b\ne a\),
the coefficient of \(m_a\) in \(\sum_b m_bF_b(x_a)\) is nonzero.
Consequently
\begin{equation}\label{eq:height-section-grid-avoidance}
 \begin{aligned}
 H_a&=\{m\in C:\textstyle\sum_b m_bF_b(x_a)=0\},&
       |H_a|&\le(M_c+1)^{M_c-1},\\
 \left|C\setminus\bigcup_a H_a\right|
     &\ge (M_c+1)^{M_c-1}.&&
 \end{aligned}
\end{equation}
By \eqref{eq:height-section-grid-avoidance} we can choose one \(m\)
avoiding every zero locus and multiply \(\sum_b m_b F_b\)
by the integralizing scalar \(s(u)\) of the theta module.
This scalar is nonzero and has log
absolute value at most \(\lambda_*|u|\) at every embedding.
The scaled combination lies in the integral cusp lattice at \(u\),
is nonzero on each analytic component and therefore on each
algebraic component. Its log supremum is at most
\[
               (400Z_b+\lambda_*)|u|+2\log M_c .
\]

For this estimate \(M_c\ge1\) and \(\sum_a m_a\le M_c^2\).
The marked-lattice count in Proposition~\ref{prop:slice-evaluation-net}
gives \(\log M_c\le41Z_b^2\), hence also \(\log M_c\le81Z_b^2\).
For \(|w_1|\le2S_L\), \(|w_2|\le(s+4)S_L\) and \(s\ge3\), put
\[
       P_4=\max\{1,5(1009Z_b+\lambda_*)+162 Z_b^2/s\}
\]
Since \(s+4\le5s\) and \(S_L\ge1\), the preceding bounds give
\[
 \tfrac12\log\|t_1\|_{\sup,\sigma}\le P_4 s S_L,\qquad
           \log\|t_2\|_{\sup,\sigma}\le P_4 s S_L .
\]
The two selections take place over the original coefficient field.
Their restrictions to an incidence field are nonzero because
that field contains the corresponding component field. The
selection uses only the section data of this run. In particular
\(t_2\) serves both the thickened weight-\(w_F\) line and the
comparison of its height with the light height.
\label{cor:selected-height-sections-3}
\end{proof}

\paragraph{From the discriminant to a height.}
\hypertarget{source-slice-heights}{}
We give the height calculation before the final lift. Use block
sphere probability measure on the rows of complex coefficients
and the Gauss norm at the finite places. For a polynomial \(q\)
let \(h_0(q)\) denote the sum of its finite log Gauss norms and
the integrals of \(\log|q|\) on the product of spheres, with the
number-field normalization on both sums. On constants the sum
is zero. The height is additive on products, so extends to
rational functions. Let \(N\) be the number of light columns
and \(\ell_N=\log N+1/2\).

\begin{lemma}[The coefficient-sphere height]
\label{lem:coefficient-sphere-height}
If \(q\ne0\) has degrees \(m_i\)
homogeneously in the respective rows, then
\begin{equation}\label{eq:slice-sphere-lower}
 h_0(q)\ge-\ell_N\sum_i m_i .
\end{equation}
\end{lemma}

\begin{proof}
Indeed fix a lexicographic order on its monomials. Condition
at each sphere on the absolute values of its coordinates.
Their phases are independent uniform angles. Successive use
of Jensen's inequality for a univariate polynomial on the
circle bounds the phase integral below by the logarithm of
the absolute value of the lexicographically leading term.
For a uniform complex unit vector the mean log absolute
value of one coordinate is \(-H_{N-1}/2\).

For \(N>1\)
its squared absolute value has density
\((N-1)(1-t)^{N-2}\) on \((0,1)\). Integrating \(\log t\)
gives this formula, including the value zero when \(N=1\).
For \(N>1\) the one-coordinate expectation is
\[
  \mathbb E\log|z_1|=\frac{N-1}{2}\int_0^1(1-t)^{N-2}\log t\,dt=-\frac{H_{N-1}}2 .
\]
The same leading coefficient is available in every Gauss
norm, so the product formula cancels its logarithms.
Finally \(H_{N-1}/2\le\ell_N\). The conditional Jensen
calculation proves integrability because the positive part on
the spheres is bounded and the lower integral is finite.

\end{proof}

\begin{proposition}[The discriminant of the selected slice]
\label{prop:slice-height-discriminant}
For the selected finite reduced slice, the weighted discriminant term satisfies \eqref{eq:slice-height-disc},
\[
 2\langle I\rangle\le\delta\log\delta+2\sum_i H(t'_i)+2s a_F\ell_N+2H_F.
\]
 An integral auxiliary section
on those same branches contributes through its normalized height and
its mean logarithmic supremum, with the same weight line and branch
multiplicities.
\end{proposition}

\begin{proof}
\proofstep{The discriminant in a trivialization of the weight line.}
For a section \(\tau\) nonzero on each component let
\(n_q(\tau)\) be the sum of its integral frame orders over
places of the incidence fields above the divisor \(q=0\)
of \(K[X]\), with residue degrees. A place trivial on the
component field contributes zero in the constant generic
frame. Define \(J_\tau=\prod_q q^{n_q(\tau)}\), and
\(a_\tau=\sum_q n_q(\tau)\deg_1 q\) where row 1 is one
of the \(s\) cut rows. These horizontal divisors have finite support, are independent
of row zero and have equal degree in each cut row for each
contributing irreducible \(q\).

We prove these properties
below by extension of the incidence point to its divisorial
valuation rings and symmetry of the cut. The height of a framed
weight is
\[
 H_\tau=\langle N(\tau)\rangle
       +\operatorname{mean}_{\sigma}\int
            \sum_{P}-\log\|\tau(P)\|_\sigma
       +h_0(J_\tau),\qquad
 N_\lambda(\tau)=\sum_{Z,w\mid\lambda} f_w n_w(\tau).
\]
Here \(P\) runs over complex points of the generic slice
at the coefficient tuple being integrated. Discriminants
and exceptional denominator loci are null sets. The integrability argument below applies to each of the
metric terms on this family.
Changing the frame within a weight adds the orders and
log norms of the field norm of its componentwise ratio.
The product identity for that rational function cancels
the change in the three terms. Likewise heights and
horizontal degrees add upon multiplication of weights
and sections. \label{cor:slice-heights-1}

\proofstep{Horizontal orders and the incidence divisor.}
Form 
\[
  U=\Delta J_\tau^2.
\]
 At a horizontal prime apply
the trace calculation of \eqref{eq:slice-discriminant-order}
in the corresponding DVR. At the contributing places the
short sections have nonnegative orders in the integral
frame. At the other places a nonzero constant component
reading has order zero. Hence the local order of \(U\)
is nonnegative at every irreducible of \(K[X]\). Unique
factorization makes \(U\) a polynomial.
In reduced form \(U=A/B\) any irreducible \(q\mid B\) would have
\[
  \operatorname{ord}_q U=\operatorname{ord}_q A-\operatorname{ord}_q B=-\operatorname{ord}_q B<0 .
\]
 Its row degrees
are \((0,2a_\tau,\ldots,2a_\tau)\) because the equal-weight
coordinates of \(b_i\) depend only on the cut points, so \(\Delta\) is invariant under
independent scalar changes of the rows. For example the
corresponding substitutions on the incidence fields fix the
\(b_i\) and intertwine the traces. All the degree therefore
comes from the horizontal factor.

\proofstep{The trace determinant in the archimedean average.}
For almost every tuple the trace discriminant is the square
of the determinant of the matrix \(b_i(P)\). Multiply each
row by the reading of \(\tau(P)\) in a unit norm frame.
Hadamard on columns bounds the squared absolute determinant
of this new matrix by \(\delta^\delta
\prod_i\|t'_i\|_{\sup,\sigma}^2\).
In the finite height of \(\Delta\) the log Gauss norm has the
opposite sign to \(v_\lambda(\Delta)\log\mathrm N\lambda\).
Additivity of \(h_0\) gives the exact identity
\[
 \langle v(\Delta)+2N(\tau)\rangle
 = 2H_\tau-h_0(U)
  +\operatorname{mean}_{\sigma}\int
       \left(\log|\Delta|+2\sum_P\log\|\tau(P)\|_\sigma\right).
\]
Combining the preceding inequalities yields
\begin{equation}\label{eq:slice-height-disc}
 2\langle I\rangle\le \delta\log\delta+2\sum_i H(t'_i)
                     +2s a_F\ell_N+2H_F .
\end{equation}
Here the subscript \(F\) means the weight \(w_F\).
For an integral section \(g\) with nonnegative horizontal
orders and bounded norm the same height formula gives
\[
 0\le\langle N(g)\rangle\le H_g+
    \delta\operatorname{mean}_\sigma\log\|g\|_{\sup,\sigma}
                       +s a_g\ell_N .
\]
Apply this to \(\eta\) and divide by \(i_\eta\).
Thus both vertical prices in \eqref{eq:geometric-four-terms}
are bounded by heights of the same global family.

\end{proof}

\paragraph{The degree and the auxiliary height.}

The preceding height contains contributions of different degrees. Divide by the incidence degree before estimating the auxiliary term.

For the following comparisons set \(n_1=i_\eta g_0^*\).
Write \(H_e=H_\eta/n_1\), \(a_e=a_\eta/n_1\).
Use the integral sections \(t_1,t_2\) of weights
\(2(w_L-w_\eta/n_1)\) and \((s+4)w_L-w_F\).
\(H_1,a_1\) denote the corresponding quantities for
\(t_1\) divided by two. Choose a constant light combination
nonzero on every component. Such combinations form the
complement of finitely many proper subspaces over an infinite
field. Let \(a_L,H_L\) be its degree and height.
By the weight identities,
\[
 a_L=a_1+a_e,\qquad (s+4)a_L=a_F+a_2 .
\]
All degrees on the right are nonnegative by the integral
section comparison. For \(a_F\) one can take a short
section nonzero on each particular component and use frame
independence.

The normalized moving part of the light
system has degree \(\delta\) because adding a general zeroth
hyperplane to the \(s\) cutting hyperplanes yields the
incidence form of the light image, with its generic field
multiplicity on each component. Equivalently the norm of
the ratio of two such hyperplanes is the ratio of their
products over the \(\delta\) slice points. After removing
the minimum order of the light columns at each place, its
specialization gives the divisor of that moving system.
The actual degree is 
\[
  a_L=\delta+\sum_q E_q\deg_1 q,
\]
where \(E_q\) is the sum of the removed minima with residue
degrees. On the horizontal depth estimate
\(0\le E_q\le c\,n_q(\eta)/(2i_\eta)\) we get
\[
 a_L\le\delta+\tfrac12 c g_0^* a_e
      \le\delta+\tfrac14 a_L .
\]
In particular \(a_L\le4\delta/3\) and
\(a_F\le4(s+4)\delta/3\).

\begin{proposition}[Heights of the chosen incidence sections]
\label{prop:slice-normalized-heights}
Under the degree, metric and Chow bounds used below, for \(s\ge3\) and
\(D_0=\delta S_L\) the pre-final heights satisfy
\[
 H_e\le(s+1)D_0(2G_c+6P_H),\qquad
 H_F\le(s+4)(s+1)D_0(2G_c+6P_H).
\]
The definite lift satisfies the same upper bound for \(H_F\), using
the baseline finite depths alone.
\end{proposition}

\begin{proof}
\proofstep{The normalized height of the moving incidence divisor.}
For precision, let \(H_{\rm mov}\) be the height of the
moving light system, computed at finite places by subtracting the minimum of the
light orders, and at a complex point divide light norms by
\((\sum_l\|w_l\|^2)^{1/2}\). Its incidence height expression
is
\[
 H_{\rm mov}=\sum_Z h_0(F_Z)+\delta H_{N-1}/2 ,
\]
where \(F_Z\) is the Chow form of the light image raised to
the generic covering degree. In fact the zeroth row in
\(F_Z\), with the other rows generic, factors up to scalar
as the product of the light linear forms at the slice
points.

At a Gauss place each factor has the minimum
column order. The mean log absolute value of each factor
on the complex unit sphere is the log norm of its vector
minus \(H_{N-1}/2\), by unitary invariance.
For a nonzero coefficient vector \(v\), the sphere average is therefore
\[
  \int_{\|z\|=1}\log|\langle z,v\rangle|\,d\sigma(z)=\log\|v\|-\tfrac12 H_{N-1}.
\]
 Comparing
with a fixed zeroth row at which the product is nonzero
gives the displayed formula by additivity of the sphere
and Gauss heights. This calculation uses the specialization
of the horizontal divisor at that row.

\proofstep{Using the same integral sections in the height comparison.}
To get the height bounds, form \(H_{\rm mov}\) using exactly the
light-weight columns retained in the new choice of sections. Write \(L=S_L\) and
\(\beta=cg_0^*\le1/2\).
Let \(P_0,P_N,P_5,P_2,P_4\ge1\) bound respectively
the integral-section price in the Chow estimate, the
finite depths, the two terms comparing the archimedean light
metrics, and the norms of the two complementary sections.
Put 
\[
  p=P_0+P_N+P_5+P_2+P_4+P_{rw},
\]
and take \(P_H\ge 2p+G_c\). Use the following inputs in the normalized height calculation.\begin{align*}
 \sum_Z h_0(F_Z)&\le (s+2)\delta
              (P_0L+G_cL+\log N+1),\\
 H_L&\le H_{\rm mov}
    +\delta(P_N+P_5+P_2)L
    +s\beta a_e\ell_N+\beta H_e,\\
 H_e&\le H_L+\delta P_4s L+s a_1\ell_N,\\
 H_F&\le(s+4)H_L+\delta P_4s L+s a_2\ell_N .
\end{align*}

\proofstep{The finite contribution of the auxiliary section.}
For the second inequality, let \(B_\lambda\) be the sum of the
minimum light orders with residue degrees at a finite place
\(\lambda\). At each branch its baseline is
\(e_w\nu_\lambda\). The decomposition identity
\(\sum_{Z,w\mid\lambda}e_w f_w=\delta\) gives
\[
 B_\lambda\le\delta\nu_\lambda+
                   \frac{c}{2i_\eta}N_\lambda(\eta).
\]
Use the weighted baseline bound
\(\langle\nu\rangle\le P_N L\) on the run. Since the
auxiliary orders are nonnegative, one can use the larger
coefficient \(c/i_\eta\) for this finite contribution.
At a complex slice point the change from the moving metric
contributes \(-\log(\sum_l\|w_l\|^2)^{1/2}\). The paired light
and auxiliary metric bound pays \(P_5D_F+cP_2\) and the same
coefficient of \(-\log\|\eta\|\).

Horizontally, let
\(J_{\min}=\prod_q q^{E_q}\) in the frame of the light system.
The effective difference with coefficients
\((c/i_\eta)n_q(\eta)-E_q\) has sphere height at least
\(-s(c/i_\eta)a_\eta\ell_N\), by
\eqref{eq:slice-sphere-lower}. This also holds for rational
or real exponents by summing the individual bounds with
nonnegative multipliers. Therefore
\[
 h_0(J_{\min})\le (c/i_\eta)h_0(J_\eta)+s\beta a_e\ell_N .
\]
In the three terms of \(H_L-H_{\rm mov}\) the resulting
auxiliary coefficient is exactly \(c/i_\eta=\beta/n_1\).
The last two height inequalities follow by weight additivity
and the lower bound for the height of an integral section.
Use the bound \(P_4sL\) for the two log suprema in
Paragraph~\ref{par:selected-height-sections}, with the factor
one half for \(t_1\).

\proofstep{Bounding the coefficient-sphere integrals.}
For example put \(D_0=\delta L\) in this calculation.
Use \(\log N+1,\ell_N,H_{N-1}/2\le P_{rw}L\).
The first two inequalities imply
\[
 H_L\le D_0\{(s+2)(p+G_c)+p
                   +\tfrac23sP_{rw}\}+\beta H_e .
\]
Indeed the harmonic and depth terms not in the first
parenthesis are paid by \(p\). The light-count term in
the first parenthesis leaves further slack. Since
\(H_e\le H_L+D_0 s(P_4+\tfrac43P_{rw})\),
substitution and division by \(1-\beta\) give
\[
 H_L\le D_0\{2(s+2)(p+G_c)+2p+
                            s(P_4+\tfrac83P_{rw})\}.
\]
If \(H_L\) is negative the same upper bound follows
directly. Adding \(D_0s(P_4+\tfrac43P_{rw})\)
costs at most \(D_0s(2P_4+4P_{rw})\) in total
in the last \(s\)-term. Hence, for \(s\ge3\),
\[
 H_e\le(s+1)D_0(2G_c+6P_H),\qquad
 H_F\le(s+4)(s+1)D_0(2G_c+6P_H).
\]
The second inequality uses \(a_2\le (s+4)4\delta/3\).

At the definite final lift use the baseline depth bound alone.
The light minimum has zero horizontal contribution, so
\(a_L=\delta\) and \(a_F\le(s+4)\delta\). The first-minimum
function is zero and the lower bound for the light vector is
uniform over the complex frames. In the height comparison this
gives
\[
 H_L\le H_{\rm mov}+\delta(P_N+P_5)S_L .
\]
The complementary section \(t_2\) gives
\(H_F\le(s+4)H_L+\delta P_4sS_L+s a_2\ell_N\). With \(a_2\le(s+4)\delta\), the Chow bound and
\(\ell_N\le P_{rw}S_L\) give
\[
 H_F\le(s+4)(s+1)\delta S_L(2G_c+6P_H).
\]
Indeed \(H_{\rm mov}\) contributes at most
\(\delta S_L((s+2)(P_0+G_c+P_{rw})+P_{rw})\).

The other terms after dividing by \((s+4)\delta S_L\)
are at most \(P_N+P_5+P_4s/(s+4)+sP_{rw}\).
Their sum is within the stated bound by \(P_H\ge2p+G_c\).
The light and complementary section norms at this lift are
integrable using the same rational ratios and the bounded light
vector above and below. Thus the discriminant argument uses
the same upper bound at \(j=m\) with no auxiliary order term.
\label{cor:slice-heights-2}

Together with \eqref{eq:slice-height-disc} and

\[
  \operatorname{mean}\log\|\eta\|_{\sup}/i_\eta\le P_\eta
\]
these are the two slice-height estimates displayed above.
We next derive the Chow bound and transfer the metric estimates
to the places used here. We then justify horizontal specialization and
integrability on this incidence correspondence.
\label{cor:slice-heights-3}
\end{proof}

\paragraph{Bounding the incidence form by an integral section.}
We detail the Chow estimate used above. The integral lower bound
needed is at one specified multiple of the light weight. Let
\(M=(s+2)N\), \(t=1+\lceil P_O\rceil\), and suppose \(\delta>0\).
Put
\[
 h=8M^2(s+1)\delta t,\qquad h'=8hM(s+1)\delta .
\]
For every component \(Z\) the input is
\begin{equation}\label{eq:chow-lower-section}
 \operatorname{mean}_{\sigma}\log\|g|_Z\|_{\sup,\sigma}
                    \ge-h P_0 S_L
\end{equation}
whenever \(g\) is an integral linear combination of products
of \(h\) light columns and \(g|_Z\ne0\). A lower bound on the
supremum over the union of all components alone would not
give this inequality. The integer \(h\) is fixed before applying it.

\begin{proposition}[The component Chow form]
\label{prop:integral-chow-bound}
With the componentwise input \eqref{eq:chow-lower-section} at the
chosen \(h\), the incidence forms of the selected family satisfy
\[
 \sum_Z h_0(F_Z)
 \le(s+2)\delta(P_0 S_L+G_c S_L+\log N+1).
\]
The sum is over the same components as the slice degree.
\end{proposition}

\begin{proof}
\proofstep{The incidence polynomial under a linear substitution.}
Write \(Y_i=\sum_l U_{il}w_l\), \(0\le i\le s+1\), with
all \(M\) entries of \(U\) independent. Let \(q(U,Y)\) be
the primitive irreducible equation of the projected image of
the given light component in \(\mathbb P^{s+1}_{K(U)}\).
Set \(d_0=\deg_Y q\), \(u=\deg_U q\), and let \(m_Z\)
be the covering multiplicity used in \(F_Z\). We use the
incidence construction with
\[
 1\le m_Zd_0\le\delta_Z,\qquad u\le(s+1)\delta_Z,\qquad
                    h_0(F_Z)\le m_Z h_{\mathbb T}(q).
\]
Here \(h_{\mathbb T}\) uses the same Gauss norms and the
mean log absolute value on unit \emph{circles} in all
variables. To see the incidence comparison, let \(F_0\) be
the Chow polynomial of the light image, so \(F_Z\) is its
\(m_Z\)-th power up to scalar.

On \(Y_0\ne0\) substitute in the \(s+1\) rows of the Chow polynomial the incidence coordinates
\[
  U'_{i l}=Y_0U_{i+1,l}-Y_{i+1}U_{0l}.
\]
After inverting \(Y_0\) this is a change of independent row
variables followed by multiplication by units. The resulting
equation is irreducible up to powers of \(Y_0\); removing
those powers gives \(q\). Geometrically it tests equality
of a projected image point and \([Y_0:\cdots:Y_{s+1}]\).
For a general line in the target its inverse image is cut by
\(s\) light hyperplanes. Counting those points bounds
\(m_Zd_0\) by \(\delta_Z\); the row substitution gives the
stated bound for \(u\). The homogeneous equations vanishing
after 
\[
  Y_i=\sum U_{il} w_l
\]
 form the principal prime \((q)\):
over \(K(U)\) the image has dimension \(s\) (a general linear
projection preserves dimension), and primitivity descends
divisibility to \(K[U,Y]\). Equivalently the differential
argument for the generic rows proves this dimension statement.

In the substitution into \(F_Z\) the coefficient of the
highest power of \(Y_0\) is the original polynomial in the
other rows of \(U\). Its log Gauss norm and torus mean do not
exceed those of the substituted polynomial, by extraction
of a leading coefficient and Jensen on the \(Y_0\) circle.
In particular for \(P(Y_0)=\sum_{j=0}^d a_jY_0^j\) the local coefficient estimate is
\[
  \|a_d\|_{\rm Gauss}\le\max_{0\le j\le d}\|a_j\|_{\rm Gauss}=\|P\|_{\rm Gauss}.
\]
The mean on spheres of \(F_Z\) is at most its torus mean:
condition on the radii and apply monotonicity of circle
means of a subharmonic function, one variable at a time.
Gauss norms and torus means are multiplicative, in logarithms
additive; the scalar and the power of \(Y_0\) have global
height zero.
For a product of nonzero polynomials the logarithmic torus mean satisfies
\[
  \int_{\mathbb T^n}\log|FG|=\int_{\mathbb T^n}\log|F|+\int_{\mathbb T^n}\log|G| .
\]
 This proves the comparison.

\proofstep{The coefficient estimate when the height is positive.}
There is nothing to prove when \(h_0(F_Z)\le0\). Otherwise
\(H=h_{\mathbb T}(q)>0\). Consider coefficient vectors of
polynomials homogeneous of degree \(h\) in \(Y\) and degree
at most \(h'\) in \(U\). Their dimension and the dimension of
the multiples of \(q\) are, respectively,
\[
 A=\binom{h+s+1}{s+1}\binom{h'+M}{M},\qquad
 B=\binom{h-d_0+s+1}{s+1}\binom{h'-u+M}{M}.
\]
Indeed both degrees add exactly in a product in a polynomial
domain, and the quotient of a homogeneous polynomial by a
homogeneous factor is homogeneous. Thus the shifted monomials
times \(q\) give a basis for the second space. We need
\[
 \frac{B}{A-B}\ge \frac h{(s+2)d_0},\qquad
                    \log A+P_O\le h .
\]

For the first assertion use the product expressions for \(B/A\).
The sum of the deficits of their factors is at most
\[
 (s+1)d_0/h+Mu/h'\le(s+1+1/8)d_0/h.
\]
Since \(d_0/h\le 1/(8M^2(s+1))\), division by one minus
this last deficit gives 
\[
  A/B-1\le(s+2)d_0/h.
\]
For the second assertion bound \(A\le(h+1)^{s+2}(h'+1)^M\)
by counting exponents separately. Now \(h'\le h^2\),
\(h\ge32M^2\) for \(s\ge3\), and \(\log(h+1)\le\sqrt h\).
Hence \(\log A\le3M\sqrt h\le3h/\sqrt{32}\), leaving
more than \(P_O\) by the definition of \(h\).

\proofstep{The metric on the quotient of the coefficient lattice.}
Give the coefficient spaces the Euclidean coefficient norms.
The normalized Arakelov degree of the saturated sublattice of
multiples in \(\mathcal O_K^A\) is at most \(-BH\).
Here is a direct determinant calculation. At a finite place, multiplicativity of the Gauss norm identifies the cofactor lattice by
\[
  \|qg\|_{\rm Gauss}\le1\quad\Longleftrightarrow\quad \|g\|_{\rm Gauss}\le\|q\|_{\rm Gauss}^{-1}.
\]
At an embedding, order the cofactor monomials lexicographically.
Gram--Schmidt on their multiples by \(q\) gives at each step
the norm of \(q\) times a polynomial with leading coefficient
one, measured in the coefficient \(L^2\) norm of the product.
For the cofactor multiple \(H\), Parseval followed by Jensen gives
\[
  \Big(\sum_\alpha |[X^\alpha]H|^2\Big)^{1/2}=\|H\|_{L^2(\mathbb T^n)}\ge\exp\!\left(\int_{\mathbb T^n}\log|H|\right).
\]
 The mean
logarithm of the monic factor is nonnegative by iterated
Jensen. Thus each orthogonal length is at least the
exponential of the mean log absolute value of \(q\). Multiplying
these bounds and inserting the finite Gauss balls gives
the claimed degree with the indicated sign.

The orthogonal quotient lattice therefore has normalized
degree at least \(BH\) in rank \(A-B\). Minkowski's convex
body theorem applied to the lattice of real dimension
\([K:\mathbb Q](A-B)\) supplies an integral coefficient
vector \(Q\) with nonzero image such that its quotient norms
satisfy
\[
 \operatorname{mean}_\sigma\log\|[Q]\|_\sigma
    \le-\frac{BH}{A-B}+\tfrac12\log A+P_O .
\]
In detail, the real covolume in the conjugate-compatible
sum metric is \(|\operatorname{disc} K|^{(A-B)/2}\)
times the exponential of minus the unnormalized degree.
The Euclidean unit ball in dimension \(a\) contains the
cube of side \(2/\sqrt a\). Applying Minkowski and then
arithmetic--geometric mean to the embedding norms cancels
\(\sqrt{[K:\mathbb Q]}\), leaving \(\sqrt{A-B}\)
and half the log root discriminant. The bound above enlarges
both allowances. The quotient is the image of the integral
coefficient lattice, so its nonzero vector lifts to such a
\(Q\).

\proofstep{Applying the integral section bound at the prescribed weight.}
After substituting 
\[
  Y_i=\sum_l U_{il}w_l,
\]
 some coefficient
\(g\) of the resulting polynomial in \(U\) is nonzero on
\(Z\) by the kernel assertion. This is a single integral
combination of \(h\)-fold products. To bound it at an embedding
one may subtract the orthogonal projection of \(Q\) onto the
space of multiples of \(q\) over \(\mathbb C\), since the
substitution of the difference has the same value on \(Z\).
Its coefficient one-norm is at most \(\sqrt A\|[Q]\|_\sigma\).

If \((c_i)_{1\le i\le A}\) are the coefficients of the orthogonal representative, Cauchy--Schwarz gives
\[
  \sum_{i=1}^{A}|c_i|\le\sqrt A\left(\sum_{i=1}^{A}|c_i|^2\right)^{1/2}=\sqrt A\,\|[Q]\|_\sigma .
\]
Each substituted \(Y\)-monomial contributes at most
\(N^h e^{h m_\sigma}\) in the pointwise metric, by the
triangle inequality for the expansion. The same upper bound
holds for every extracted coefficient times the one-norm.
Comparing with \eqref{eq:chow-lower-section} gives
\[
 \frac{hH}{(s+2)d_0}\le
 h(P_0 S_L+G_c S_L+\log N)+\log A+P_O .
\]
The two monomial counts and \(m_Zd_0\le\delta_Z\) prove the
Chow bound for \(F_Z\) and summation proves the required
bound for the family. Although \(h,h'\) grow with the number
of rows and columns, both cancel in this estimate. The
remaining arithmetic input in this step is the componentwise
lower bound at precisely \(h w_L\).
\label{cor:slice-chow-1}
\end{proof}

\paragraph{The lattice estimate for the CM approximations.}

The lower bound for the Chow estimate uses an integral section evaluated near a CM frame. Hermitian transference controls how closely the required directions can be approximated.

\begin{lemma}[Hermitian lattice transference]
\label{lem:cm-lattice-transference}
Let \(E\) be the CM field, \(r_0\ge1\) an integer, and \(L\subset E^{r_0}\) a full \(\mathcal O_E\)-lattice
self-dual for a Hermitian form \(H\). Take one embedding from
each conjugate pair and write \(f=[E:\mathbb Q]/2\).
At these embeddings let \(Q_\tau\) take \(H_\tau\) to its
diagonal signature matrix \(J_\tau\). Write
\(q_\tau(z)=\|Q_\tau\tau(z)\|^2\) for the corresponding
positive definite majorants. Suppose

\[
  \max_\tau q_\tau(z)\ge\lambda^2>0
\]
 for every nonzero
\(z\in L\). Put \(d_E=|\operatorname{disc}E|\),
\(t_*=fr_0\), and use the real Euclidean space
\(\prod_\tau\mathbb C^{r_0}\) of dimension \(2t_*\).
Then the image lattice under \(z\mapsto(Q_\tau\tau(z))_\tau\)
has covering radius at most \(4t_*^3d_E/\lambda\). There is
an \(E\)-basis of \(E^{r_0}\) in \(L\) with all Euclidean
image norms at most the same bound.

\end{lemma}

\begin{proof}
\proofstep{Comparing the rational and Hermitian duals.}
Indeed let \(L^\dagger\) be the rational dual of \(L\) for
the pairing \(\operatorname{Tr}_{E/\mathbb Q}H\). For
\(z\in L^\dagger\), pairing with all \(\mathcal O_E\)-multiples
of a vector of \(L\) puts each \(H(z,y)\) in the inverse different.
Thus \(d_E z\in L\) by Hermitian self-duality. Combining a CM type
with its conjugate shows that the image of \(L^\dagger\) under
\(z\mapsto(2J_\tau Q_\tau\tau(z))_\tau\) is the Euclidean dual.

In fact the Euclidean pairing of these two images is
\(2\sum_\tau\operatorname{Re}\tau H(z,y)=\operatorname{Tr}H(z,y)\).
The factor two identifies a real dual basis with the image of the
rational trace dual basis. For \(0\ne z\in L^\dagger\), the element
\(d_Ez\) is a nonzero vector of \(L\), whence
\[
 \Big(\sum_\tau\|2J_\tau Q_\tau\tau(z)\|^2\Big)^{1/2}
 \ \ge\ \frac{2}{d_E}\max_\tau\sqrt{q_\tau(d_Ez)}
 \ \ge\ \frac{2\lambda}{d_E}.
\] Euclidean transference gives \(2t_*\)
independent primal vectors of norm at most
\((2t_*)^2/(2\lambda/d_E)=2t_*^2 d_E/\lambda\).
Their integer span is a sublattice.

For an arbitrary point write \(x=\sum_{j=1}^{2t_*}a_j v_j\) in
this real basis, and choose \(n_j\in\mathbb Z\) with
\(|a_j-n_j|\le1/2\). Then
\[
 \Big\|x-\sum_j n_jv_j\Big\|
 \le\tfrac12\sum_j\|v_j\|
 \le \tfrac12(2t_*)\,2t_*^2d_E/\lambda=2t_*^3 d_E/\lambda .
\]
The rounded point belongs to the original lattice. The same family spans \(E^{r_0}\)
over \(\mathbb Q\) after pulling back, so an \(E\)-basis can be
selected from it. Both conclusions fit the stated bound.
This argument is used in the full dimension and on the
complementary lattice when it has the same self-duality and
majorant hypotheses. The quantitative choice of the frame
uses the resulting covering radius at that frame's cusp times.
\cite{banaszczyk}\label{cor:slice-cm-transference-1}
\end{proof}

\paragraph{Degrees of the CM lines.}

The approximation determines the CM lines. Their contribution to the section height involves both the finite ideal and the norms of its chosen frame at infinity.

For a metrized line \(L\) over a number field let
\(I(L)\) be its fractional ideal and \(q_\sigma(L)\) its
metric on the chosen coordinate at \(\sigma\). We use the
normalized degree
\[
 \operatorname{ndeg}L
   =-\frac{\sum_w\operatorname{ord}_w I(L)\log\mathrm Nw+
                   \sum_\sigma\log q_\sigma(L)}
              {[K':\mathbb Q]} .
\]
The product formula makes this independent of the coordinate.
For the CM object of dimension \(g\), use its Riemann
metric on the Lie lattice \(\mathrm{Lie}\) and the integral
theta line \(\mathcal T\) with the normalized \(L^2\) metric.
\begin{lemma}[The principal cubical theta degree]\label{lem:slice-principal-theta-degree}
For this CM Lie lattice and the normalized principal theta line,
the degree comparison is
\begin{equation}\label{eq:slice-cm-theta-degree}
 \operatorname{ndeg}(\mathcal T^\vee)
      =-\tfrac12\operatorname{ndeg}(\det\mathrm{Lie})
                         +\tfrac{g}{4}\log(2\pi).
\end{equation}
\end{lemma}

\begin{proof}
The theta generator comes from the rigidification of
the principal symmetric cubical line. The analytic length
is the mean of its squared Gaussian norm on a period
parallelepiped, with its square root taken, and the Lie
determinant length is the square root of the Riemann
determinant in the same frame. To see the comparison, write
\(\mathcal D=\det\mathrm{Lie}\) and use the powered cubical invariant
of the rigidified symmetric principal line. For the chosen theta
and Lie coordinates it is a scalar \(k\in (K')^\times\) for an
integer power parameter \(a>0\) depending on the dimension.
Multiplying the theta generator by \(u\) and changing the Lie
frame by \(M\) multiplies \(k\) by \(u^{-8a}\det(M)^{4a}\).
The invariant in integral theta and Lie generators is a unit, by
the integral cubical identity.

Comparing those generators with
the chosen coordinates gives
\[
 (k)=I(\mathcal T)^{8a}I(\mathcal D)^{-4a}.
\]
This is an equality of fractional ideals. One may compare over a
valued extension, since both sides acquire the same factor of the
ramification index. At a complex embedding \(\sigma\), the same
invariant in a theta uniformization satisfies
\[
 |\sigma(k)|\,q_\sigma(\mathcal T)^{8a}
       q_\sigma(\mathcal D)^{-4a}=(2\pi)^{2ga}.
\]
Here the complex cubical formula uses the mean square of the theta
generator, and the period quotient has volume
\(\pi^g |\det W|^2/q_\sigma(\mathcal D)^2\) in the corresponding
complex uniformization frame \(W\). Thus the volume factor uses
the Lie determinant of the same frame. Let \(d=[K':\mathbb Q]\). Summing the complex identities and
using the product formula for \(k\) replaces
\(\sum_\sigma\log|\sigma(k)|\) by the finite log norm of the
displayed ideal. Thus
\[
      -8ad\,\operatorname{ndeg}(\mathcal T)
             +4ad\,\operatorname{ndeg}(\mathcal D)
             =2gad\log(2\pi).
\]
Divide by \(8ad\) and use duality for the degree of a
metrized line. \label{cor:slice-integral-norm-1}
\end{proof}

For an abelian CM differential line \(\ell_k\), put \(\Omega=2\operatorname{ndeg}\ell_k\). The CM differential degree formula gives the same real number for every abelian line in both constructions. The toric differential lines on the arc have degree zero.
\label{cor:slice-integral-norm-2}

\begin{lemma}[Degree of a CM coefficient line]\label{lem:cm-coefficient-line-degree}
With the theta degree identity \eqref{eq:slice-cm-theta-degree}
and the abelian letter metrics above, the value line of an arc
coefficient satisfies \eqref{eq:slice-cm-line-degree}. For the coefficient line \(\mathcal L_\mu\), the bound is
\[
 \operatorname{ndeg}\mathcal L_\mu\le P_\Lambda(\mathcal A+frn),\qquad P_\Lambda=\tfrac12(1+|\Omega|+\log\Lambda).
\]

\end{lemma}

\begin{proof}
\proofstep{Compensating the norms of the CM letters.}
The abelian letters \(u_k\) form a basis over \(K'\). Each
scaled line \(I(\ell_k)u_k\) is inside the completed Lie
lattice. Thus their exterior product is inside its determinant
lattice. Over each complex embedding their lengths are at
most \(\Lambda q_\sigma(\ell_k)\).
To combine the two estimates, identify the exterior line of
the \(u_k\) with \(K'\) via that exterior product. The
determinant ideal contains \(\prod_k I(\ell_k)\), and
Hadamard bounds the determinant metric by
\(\Lambda^g\prod_k q_\sigma(\ell_k)\).

Taking degrees with the indicated sign gives
\[
 \operatorname{ndeg}(\det\mathrm{Lie})
             \ge g\,\Omega/2-g\log\Lambda .
\]
By \eqref{eq:slice-cm-theta-degree},
\(\operatorname{ndeg}(\mathcal T^\vee)\le g P_\Lambda\),
where
\[
             P_\Lambda=\frac{1+|\Omega|}{2}
                                  +\tfrac12\log\Lambda .
\]
Here \(\log(2\pi)<2\). Every dual letter has degree
at most \(P_\Lambda\) since \(\Lambda\ge1\).
The arc has abelian dimension \(f(r-2)\le fr\).
The line \(\mathcal L_\mu\) of the arc coefficient therefore
has
\begin{equation}\label{eq:slice-cm-line-degree}
    \operatorname{ndeg}\mathcal L_\mu
           \le P_\Lambda(\mathcal A+frn).
\end{equation}
This is exactly the degree of the line in the finite and
complex comparisons. All three bounds for the arc
coefficient now have their dependence on \(\Lambda\), \(n'\),
\(\Omega\), and the actual cutoff specified.
\label{cor:slice-integral-norm-3}

\end{proof}

\paragraph{Uniformity of the radius and torsion order.}

The detecting CM point may depend on the form. The approximation radius and the allowed torsion order must nevertheless be bounded before that point is chosen.

Put \(w_0=S_\kappa+r\), \(N_A=d_E=|\operatorname{disc}E|\),
\(N_C=\mathrm N\mathfrak c_K\), and \(M_0=g_0g_1\).
The split approximation compares the two cusp times using
a positive time lower bound \(e_{\rm sp}\), and uses
\(\rho_{\rm sp}=1/(12\pi(c_1+w_0))\) for the approximation.
With
\begin{align*}
 N_0&=M_0d_E^2,\qquad N_2=4f,\\
 H_0&=N_C^2(2+4f^4\operatorname{rd}(E)^2 N_A^2),\\
 K_0'&=1+N_0H_0(2c_1+1),\\
 R_0&=w_0+N_2\log(2K_0' N_2)+N_2+2 ,
\end{align*}
the time bound is
\[
 e_{\rm sp}=\min\{1,f^2/(N_0^2N_C^2c_1R_0)\},
 \qquad e_{\rm sp}^{-1}\le\max\{1,N_0^2N_C^2c_1R_0\}.
\] The coefficients in the bounds
for the split complementary frame are
\begin{align*}
 B_{\rm ind}&=f(2f(r-2))^3 d_E N_A
                     \sqrt{N_C^2\operatorname{rd}(E)e_{\rm sp}^{-1}},\\
 B_{\rm rad}&=2\cdot4\sqrt{r-2}\cdot
         4(f(r-2))^3d_E\sqrt{e_{\rm sp}^{-1}}\cdot40\rho_{\rm sp}^{-1}.
\end{align*}

The choice is \(n'=1+\lceil1000B_{\rm ind}\rho_{\rm sp}^{-1}\rceil\)
and \(\Lambda=\max(1,B_{\rm rad}^{2q_j-1})\).
Fix this order before choosing the detecting frame.
The split approximation theorem then chooses a frame
where multiplication by this order makes the abelian
projection of the cusp lattice integral. The constants
are the same for all choices of arc on a lift.

\begin{lemma}[Uniform order and radius for the detecting arcs]\label{lem:uniform-cm-arc-size}
Use the split approximation with the selected weight and
cutoff \(c_1\), and the fixed parameters above. For \(z=frL\),
the order \(n'\) and radius bound \(\Lambda\) satisfy
\[
 \max(\log n',\log\Lambda)\le200(K_*z)^2,\qquad K_*=20+8|C_A|+|C_K|.
\]
 These choices precede selection
of a detecting arc. The three coefficient allowances in
the proof each have the bound
\((17+2K_*+200K_*^4)z^4S_\kappa\).
\end{lemma}

\begin{proof}
\proofstep{Uniform bounds for the conductor and radius.}
Here is the quantitative consequence used in the division
budget. Use \(z=frL\) as in the run and
\(K_*=20+|C_A|+7|C_A|+|C_K|\).
The size inequalities put \(f,r,d_E,M_0,N_C,c_1,w_0\)
below 
\[
  U=e^{K_*z}.
\]
 For \(c_1\) this uses
\(Tc\le1\) and \(q_j\le r\). For \(w_0\) use
\(S_\kappa\le C_Kkfb\), \(k\le r\), \(\log b=L\).
The size comparison also gives
\(|\Omega|\le |C_A|L\le K_*z\).
The bound for \(\log d_E\) is \(|C_A|z\), and the bound
for \(\log M_0\) is \(18L+2\).

Let \(U\ge2\), count exponents of \(U\) by addition
under multiplication, and charge a sum of two terms by
one plus the larger exponent. Since \(\operatorname{rd}(E)\le d_E\)
and \(\log v\le v\) for \(v>0\), this gives
\(N_0\le U^3\), 
\[
  H_0\le U^{13},
\]
 \(K_0'\le U^{20}\),
\(R_0\le U^{30}\), and \(e_{\rm sp}^{-1}\le U^{50}\).
Also \(\rho_{\rm sp}^{-1}\le U^8\).
Substitution gives \(n',B_{\rm rad}\le U^{100}\), including
the rounding in \(n'\). Thus \(\log n'\le100K_*z\),
while
\[
 \log\Lambda\le(2q_j-1)\max(0,\log B_{\rm rad})
                         \le 200rK_*z .
\]
Because \(K_*z\ge r\) and \(K_*z\ge1\), both are at most
\(200(K_*z)^2\).

Combining them with the coefficient and degree bounds
shows that \(C_{\rm arch}\), \(B_{\rm deg}\), and the
normalized mean of the finite baseline each cost at most
\((17+2K_*+200K_*^4)z^4 S_\kappa\).
It is the upper bound on \(\log\max(1,c_1)\), together
with the fixed frequency restriction, that is used here.
\label{cor:slice-integral-norm-4}

\end{proof}

\paragraph{The lower norm of an integral form of arbitrary weight.}
\label{subsec:deep-cm-interior}
In this paragraph the ambient abelian dimension is \(fr\).
Keep the coefficient-field and level data of the run fixed,
and write \(Z=fr+L+\log d_E+|\Omega|\) in the normalization
of the preceding CM estimates.

\begin{proposition}[Lower norm at the actual weight]\label{prop:slice-lower-integral-norm}
Use the coefficient field and the two-prime level \(g_0g_1\)
of an orientation in Proposition~\ref{prop:linked-comparison-rows}.
Set \(Z=fr+\log b+\log|\operatorname{disc} E|+|\Omega|\).
For a nonzero integral cusp form \(F\) on this run of dominant
weight \(v_0\) with \(|v_0|\ge1\), the normalized supremum norms satisfy
\[
 \operatorname{mean}_\sigma\log\|F\|_{\sup,\sigma}
 \ge -3000Z^2(1+\log|v_0|)|v_0| .
\]
The coefficient \(3000\) is independent of the run, the
weight and the form.
\end{proposition}

\begin{proof}
For the lower norm one uses an interior CM frame of the full
abelian dimension \(g=fr\). Let \(F\) be nonzero in the integral
cusp lattice at a dominant weight \(v_0\), and put \(S=|v_0|\ge1\).
The normalized supremum at the distinguished good embedding
is finite and positive. Take a frame where the reading exceeds
half the supremum. The quantitative CM frame selection starts
with this frame. Write \(p_0\le r\) for the positive rank of
its Hermitian space in the signed frame convention and put
\[
 d_E=|\operatorname{disc}E|,\qquad
 M=g_0g_1\,\mathrm N\mathfrak c_K\,d_{\rm lev},\qquad
 e_0=\min\left\{1,\frac{\pi}{M(fS+2)k_0\log 2}\right\}.
\]
Here \(d_{\rm lev}\) is the positive level discriminant in the
transvection parameter and \(k_0>0\) is the polarization scale of
the component lattice. The principal level gives the integral
real transvections with period \(M\).

At a half-maximal frame,
the nonnegative Fourier modes along these transvections
force at least one majorant \(e_0\) at every nonzero lattice
vector in the reference lattice. Indeed, if all the majorants
\(H_\tau\) are less than \(e_0\le1\), the integral Hermitian
norm and the bound of each of its conjugates by \(H_\tau\)
force that norm to be zero. Along the isotropic flow the
domain condition is \(1+\operatorname{Im}\zeta\,k_0 H_\tau>0\).
Set 
\[
  t_0=1/(2k_0e_0).
\]
 The flow section is then holomorphic
on a strip containing \(\operatorname{Im}\zeta=-t_0\),
is \(M\)-periodic, and tends to zero at the cusp.

On
\(\operatorname{Im}\zeta=-t_0\) each of the factors in the
majorant change lies between \(1/2\) and \(1\). The
weight-spread sum is at most \(S\), so the raw section
has norm at most \(2\cdot 2^{S/2}\) times its modulus at
the starting frame. This uses the half-supremum bound
for the normalized section. Under \(u=e^{2\pi i\zeta/M}\)
the function extends holomorphically across \(u=0\)
with value zero. Apply the maximum principle to its
quotient by \(u\) in the disc of radius \(e^{2\pi t_0/M}\).
At \(\zeta=0\) it yields

\[
  1\le 2^{1+S/2}\exp(-2\pi t_0/M).
\]
The definition of \(e_0\) gives
\[
 2\pi t_0/M\ge(fS+2)\log 2>(1+S/2)\log 2,
\]
a contradiction.

\proofstep{Moving to an interior frame.}
Orthonormalize the starting flag with its order and signs
unchanged. The normalized reading and the positive majorants
are preserved. Put
\[
 \rho_0=(10\sqrt{rS})^{-1},\qquad \rho=3\rho_0/5,\qquad
 \mu=4(fr)^3d_E\sqrt{f/e_0},\qquad A=4\sqrt r\,\mu/\rho.
\]
The majorant lower bound implies the hypothesis of the previous
lattice estimate with \(\lambda=\sqrt{e_0/f}\). Thus \(\mu\)
bounds both the covering radius and the lengths of an independent
spanning family. Approximate the first \(p_0\) axes, scaled by \(A\),
simultaneously at all embeddings by lattice vectors. A matrix with
these vectors as columns has top block \(A I+E_+\) and bottom
block \(E_-\), with \(\|E_\pm\|\le\sqrt r\,\mu\le A\rho/4\).

The columns of this matrix span the positive graph of 
\[
  Z=E_-(A I+E_+)^{-1},
\]
 and
\(\|Z\|\le\rho/3=\rho_0/5\) at every embedding. The absence
of a positive block gives the zero graph. Taking the graph as
the positive part and its Hermitian orthogonal as the negative
part defines the common CM structure on these lattice subspaces.

For the complementary letters let \(G\) be the Gram matrix
over \(E\) of the positive columns. It is positive definite
at every inducing embedding. The orthogonal projection onto
the positive space is given by \(G^{-1}\) in the coordinates
of those columns. Multiplying the complementary projection
by \(\det G\) therefore carries the short lattice vectors into
the original lattice by the adjugate formula and integral
Hermitian pairings.
The matrix identity clearing the projection is
\[
  (\det G)G^{-1}=\operatorname{adj}(G).
\]
 The resulting vectors span the orthogonal
complement. Select an independent family there.

At each
embedding the graph has norm less than \(1/3\), and the
complementary projection in the original Euclidean norm
has norm at most \((1+\|Z\|^2)/(1-\|Z\|^2)\le2\).
This follows by writing the negative graph as \((Z^*y,y)\)
and inverting its Gram matrix \(-(I-ZZ^*)\) for the projection.
Hadamard on the positive Gram matrix gives
\(|\tau(\det G)|\le(2A)^{2p_0}\). Hence the positive
and negative letters all have absolute Hermitian lengths
at most \((2A)^{2p_0+1}\). In the absence of one block
the bound follows directly for the other block by the
same approximation or short basis. Substitution of the approximation radius gives
\[
 \Lambda=X^{2p_0+1},\qquad
 X=\frac{1600}{3}\,r\sqrt f\sqrt S\,(fr)^3 d_E\,e_0^{-1/2}.
\]
The upper estimate on \(\Lambda\) uses any larger bound
for the chosen letters as well.

\proofstep{Retaining a nonzero value through the move.}
To see that the new frame still detects \(F\), use the
holomorphic chart with positive graph \(tZ\) and unchanged
negative coordinate axes. For each signed leading Gram determinant \(\Delta_i\) in this frame, the chart gives
\[
  (1-\|tZ\|^2)^r\le|\Delta_i|\le1.
\]
For leading blocks confined to the positive axes this is the
Gram matrix \(I-(tZ)^*(tZ)\) restricted to those axes.

When negative axes are included, the Schur complement removes
their projection from \(tZ\), and gives the same bound.
The full determinant has absolute value one. Thus for
\(|t|<5\) the change in the log norm factor, with weight
spread at most \(S\), has absolute value at most
\(4Sr\|tZ\|^2\le1/25\), by summing the log minor bounds.
The half-maximal hypothesis bounds the raw reading on
that disc by \(2e^{1/25}<3\) times its modulus at zero.
If the raw reading \(F_0(t)\) vanished at \(t=1\), apply
Schwarz's lemma to

\[
  H(u)=F_0(1+4u)/(3|F_0(0)|)
\]
 on the unit disc.
We have \(H(0)=0\), so
\[
      \tfrac13=|H(-1/4)|\le\tfrac14 .
\]
This contradiction proves that the moved frame still
detects the section.
The graph's majorant is within a factor
\((1+\|Z\|)/(1-\|Z\|)\le2\) of the original. Indeed the
matrix with the positive and negative graph columns is
\(I+J_Z\), where \(J_Z\) is self-adjoint with off-diagonal
blocks \(Z,Z^*\). Consequently the new majorant, represented by \(H_{\rm new}\) in the old frame, satisfies
\[
  H_{\rm new}=(I-J_Z)(I+J_Z)^{-1}.
\]
 This proves the
majorant comparison without changing the polarization.

\proofstep{Polynomial evaluation on the algebraic CM model.}
The arithmetic CM realization uses the lattice of this
basis, its two CM types and the original marking. The
tangent coordinates and the normalized theta generator
are algebraic over a finite extension \(K'\); all finite
comparisons use the corresponding integral models.
\label{cor:slice-integral-norm-5}

The section at this CM point gives a homogeneous polynomial
\(G(X')\) of frame degree \(\mathcal A\), in the entries of the
block matrices \(X'\).
With \(X'\) the frame variables, homogeneity means
\[
  G(cX')=c^{\mathcal A}G(X') .
\]
 Since the section is nonzero in this frame, \(G\ne0\). Choose a nonzero coefficient \(a_\mu\).
At every finite place normalize the theta line and use an
integral eigenbasis of the Lie lattice. The defining
integrality test on \(F\) bounds the normalized polynomial
by one on integral changes of that basis with unit block
determinants. Reduction over an extension with infinite
residue field, or over arbitrarily large finite residue
extensions, gives integral coefficients by the same
determinant test as for the arc. The \(I(\ell_k)\)-scaled
CM letters are integral in this basis. Substituting their
integral coordinate matrices preserves this integrality.
Undoing the scaling gives
\(a_\mu\in I(\mathcal L_\mu)\) at every place, with

\[
  \mathcal L_\mu=(\mathcal T^\vee)^n
\otimes\bigotimes_{\tau,i,k}(\ell_{\tau k}^\vee)^{\mu_{\tau i k}}.
\]
The finite order defect of the toroidal arc is absent in
this interior calculation.
\label{cor:slice-integral-norm-6}

\proofstep{Comparing the archimedean and integral coefficient bounds.}
At \(\sigma\) take the polydisc of radii
\(q_\sigma(\ell_{\tau k})^{-1}\) in the CM matrices.
Each resulting row has length at most \(r\Lambda\) in the
Riemann metric, by the triangle inequality. For invertible
blocks, orthogonalize the rows in the ordering of the flag.
The triangular factor acts by the weight character.
Each diagonal length of that factor is at most the
length of the original row, and all its flag exponents
are nonnegative.

In the orthonormal frame the supremum
bound applies with the fixed theta metric. Consequently
\[
 |G(X')|q_\sigma(\mathcal T)^{-n}
    \le \|F\|_{\sup,\sigma|K}(r\Lambda)^{\mathcal A}.
\]
The bound extends to singular blocks by continuity.
Cauchy's formula in the entries of \(X'\) then bounds
\(|\sigma a_\mu|q_\sigma(\mathcal L_\mu)\) by the same
right side. The coefficient is integral in its line, so
\[
 -\operatorname{ndeg}\mathcal L_\mu
       \le \operatorname{mean}_\sigma
                      \log(|\sigma a_\mu|q_\sigma(\mathcal L_\mu)).
\]
This is the product formula for \(a_\mu\), using
\(\operatorname{ord}_w a_\mu\ge\operatorname{ord}_w I(\mathcal L_\mu)\).
Apply \eqref{eq:slice-cm-line-degree} with \(g=fr\) and
cancel the multiplicity of extensions of each embedding of
\(K\). It follows that
\begin{equation}\label{eq:slice-cm-point-norm}
  \operatorname{mean}_\sigma\log\|F\|_{\sup,\sigma}
       \ge -P_\Lambda(\mathcal A+frn)
                   -\mathcal A\log(r\Lambda).
\end{equation}
This bound is independent of the degree of the CM field over which the chosen data are defined.

\proofstep{Summing the dependence on the weight.}
For the norm budget set \(Z=fr+L+\log d_E+|\Omega|\).
The chosen principal level gives
\(\log(Mk_0)\le 19L+3+\log d_E\), using the two level
prime bounds and the upper bound for \(k_0\).
Since \((fS+2)\log2\le 3fS\) and \(\pi>3\),
\[
 -\log e_0\le 19L+6+\log d_E+f+\log S .
\]

Taking the logarithm of \(X\) gives
\(0\le\log X\le600 Z+\log S\), using
\(\log(1600/3)<7\), \(Z\ge fr\ge1\) and
\(\log fr\le fr\). Thus
\(\log\Lambda\le1800 Z^2(1+\log S)\).
The weight convention gives
\(\mathcal A+frn\le S\) and \(\mathcal A\le S\).
Inserting these bounds into \eqref{eq:slice-cm-point-norm}
yields
\begin{equation}\label{eq:slice-integral-norm}
        \operatorname{mean}_\sigma\log\|F\|_{\sup,\sigma}
                    \ge -3000 Z^2(1+\log |v_0|)|v_0|.
\end{equation}
The absolute coefficient \(3000\) works on both runs,
also when \(v_0\) is the large multiple of the light weight
used in the Chow estimate. The estimates for \(X\) and the
letter metric are the necessary dependence on the weight
of that multiple. They permit the choice of the nonzero
coefficient to vary with \(F\).
\label{cor:slice-integral-norm-7}
\end{proof}

\paragraph{The integral light lattice in the marked chart.}
\label{par:slice-light-integral}
Fix a lift and either of the two runs, with construction field
\(K_0\), half-level ideal \(\mathfrak m_b\), and full level
\(M=g_0g_1\). Write \(a_{\tau,i}\) for the entries of a weight
\(w\) on the \(\tau\)-eigenspace and \(d_\tau=\sum_i a_{\tau,i}\).
Use \(\mathcal T_w\) for the theta rules with these entries and
the prescribed number of copies. A rule uses an integral harmonic
polynomial \(P\), of the specified flag weight, and marked level
points on the copies.

For each marked point take an \(M\)-division
point killed by \(\mathfrak m_b^2\). The normalized multiplication
isomorphism and the symmetric theta lift give a translated section
\(H_b\) of the \(M^2\)-th power of the rigidified theta line on
copy \(b\). The rule reads
\begin{equation}\label{eq:light-rule-value}
 M^{-\deg P}\sum_{\boldsymbol\mu} P_{\boldsymbol\mu}
        \prod_b \nabla^{\boldsymbol\mu_b} H_b(0).
\end{equation}
Here \(\boldsymbol\mu_b\) uses the weight-labelled tangent
directions on copy \(b\), and \(\nabla\) is the covariant
derivative in the given frame. The translation and trivialization
laws give this expression after a faithfully flat cover carrying
the marked sections. The harmonic and theta transformation laws
identify it with the extension of the original rule value for
any admissible choice of the \(H_b\). The power of \(M\) in
\eqref{eq:light-rule-value} is part of the rule throughout the
comparison. \label{cor:slice-light-integral-1}

\begin{proposition}[Integral evaluation of the scaled lattice of theta rules]
\label{prop:scaled-integral-light-reading}
After scalar multiplication by \(a_W\) in \eqref{eq:light-global-scalar},
the chosen light family generates the image of the integral span
of \(\mathcal T_{w_L}\). Let \(w_l\) be the scaled light generators and \(b_T\) the scaled value of a rule \(T\). At every allowed integral framed point \(y\), in the fixed orientation and coefficient embedding,
\[
 \min_l\operatorname{ord}w_l(y)=\inf_{T\in\mathcal T_{w_L}}\operatorname{ord}b_T(y).
\]

\end{proposition}

\begin{proof}
\proofstep{Integral frames and division by the level scalar.}
Take an algebraic integral framed point \(y\) over a complete
DVR \(V\) above \(\lambda\), with relative ramification \(e\).
Its residue field is finite. Choose one of its integral models,
with polarized semiabelian group \(G\). The split base-change
theorem for the semiabelian model over this complete DVR gives
a finite local extension \(V'/V\) with finite residue field and a
split uniformization over \(V'\). The integral-model comparison
transports that uniformization and its torsion identification to
every integral model of the extended point.

Choose the
integral formal-group chart at the pulled-back frame. Its
polarization generators are compatible under jet truncation
over \(V'\).
For example, for jet orders \(a\le b\), the compatible reductions of the chosen formal parameters give
\[
  V'[\![X]\!]/(X)^{b+1}\longrightarrow V'[\![X]\!]/(X)^{a+1}.
\]
 The framed base-change map pulls back the
underlying point, all the labelled tangent vectors and the
theta-line generator. The incidence map respects this frame
extension, so the new point lies over the same point in the
boundary comparison. \label{cor:slice-light-integral-2}

Let \(\iota:\operatorname{Frac}V\hookrightarrow\operatorname{Frac}V'\)
be the extension. The normalized reading commutes with \(\iota\).
Indeed, for a rule at the new frame choose a field over
\(\operatorname{Frac}V'\) carrying the translated sections. Such
a field is obtained by taking a residue field at a maximal ideal of
the faithfully flat coefficient algebra and pulling back its
sections. Over this field, the composite frame map from the old
point and the frame map from the new point give the same
expression \eqref{eq:light-rule-value} using the same sections.
Injectivity of the field map proves

\[
  f(y')=\iota(f(y)).
\]
 This works for each rule and each weight
on the same \(V'\). The auxiliary field used to test the equality
can vary with the rule. If \(r\) is the ramification of
\(V'/V\), then
\begin{equation}\label{eq:light-frame-order}
 \frac{\operatorname{ord}_{V'} f(y')}{re}
        =\frac{\operatorname{ord}_{V} f(y)}{e}.
\end{equation}
Both sides are infinite when the reading is zero. All the
integral germs and the local interpolation estimates below may
thus be taken on the split chart, with its own ramification.
\label{cor:slice-light-integral-3}

\proofstep{Using the half-ideal of the selected orientation.}
For the actual orientation choose the elements \(\pi_i\) of
the fixed imaginary quadratic subfield used by the run, with
\(\pi_i\bar\pi_i=g_i\), and put \(\alpha=\pi_0\pi_1\).
The oriented markings give
\(\mathfrak m_b=(\bar\alpha)\) and 
\[
  \alpha\bar\alpha=M.
\]
In fact each \(\bar\pi_i\) divides a generator of the half
level. Their ideals are coprime because \(g_0,g_1\) are distinct,
so their product divides that generator. The relative norm of the
quotient is one, since the generator and \(\bar\alpha\) have the
same relative norm \(M\). Thus the quotient is a unit. Away
from the level one can equivalently use any generator with this
relative norm.

At the level use \(\alpha\) with the indicated
orientation. Let \(\iota_0\) denote the coefficient embedding of
the fixed imaginary quadratic field defining the plus type.
Each plus embedding agrees with \(\iota_0\) there and each
minus embedding is its conjugate. Put
\(D^\pm(w)=\sum_{\tau\text{ of type }\pm}d_\tau\). The weight
determinant of the action of \(\alpha\) on the labelled frame is
\begin{align}
 \mathcal D_w(\alpha)
     &=\prod_\tau\tau_K(\alpha)^{d_\tau} \notag\\
     &=\prod_{i=0}^1
         \iota_0(\pi_i)^{D^+(w)}\,
         \iota_0(\bar\pi_i)^{D^-(w)},
   \label{eq:light-oriented-determinant}
\end{align}
where \(\tau_K\) are the embeddings used in the coefficient
field. The first formula is the action on each row degree of
\(P\). The second follows by summing those degrees in the two
signature blocks. In particular \(D_F=D^+(w_L)+D^-(w_L)\).
\label{cor:slice-light-integral-4}

\proofstep{The minimum valuation at the smaller weight.}
At the smaller weight, the scalar multiplying the lattice is
\begin{equation}\label{eq:light-global-scalar}
 a_W=\mathcal D_{w_L}(\alpha)\,
       Q_S^{D_F}q^{\lceil D_F/2\rceil}d^{\lceil D_F/2\rceil}.
\end{equation}
The local value of \(a_W\) is obtained from this coefficient
scalar by the same embedding as the rule. The light lattice is
\(a_W\operatorname{span}_{\mathcal O_{K_0}}\mathcal T_{w_L}\),
using the cusp inclusion at this weight. This equality follows
by taking the integral span of the rules and applying the scalar
map to the span. The chosen family \((w_l)\) generates exactly
this image, and the new choice of short sections retains it.

In the fixed integral frame, an integral combination has
valuation at least the least valuation of its terms. Hence the
least valuation of the scaled evaluations is attained on a generator \(w_l\). Equations \eqref{eq:light-frame-order}
and \eqref{eq:light-oriented-determinant} apply at every allowed
weight, including the two weights in the full and truncated
evaluation comparisons. The choice of chart depends on the
point and marking. It introduces no weight-dependent scaling
into these evaluations. \label{cor:slice-light-integral-5}
\end{proof}

\paragraph{Values of the light and auxiliary sections at finite places.}
\label{par:slice-finite-depth}

A bound for the light lattice over the coefficient field gives no local order by itself. The light and auxiliary sections must be compared in the integral frames at each finite place.

Fix either run and an allowed lift \(j\). At a coefficient place
\(\lambda\mid p\) put \(e_2=\operatorname{ord}_{\lambda}(p)\).
For an evaluation in a complete discrete valuation ring above
\(\lambda\) of relative ramification \(e\), use

\[
  v_\lambda=\operatorname{ord}/e=e_2v_p,
\]
 where \(v_p(p)=1\).
The chart depth is the infimum of \(v_\lambda\) of the values of the covariant differential operators on integral formal germs and integral harmonic
differential polynomials of the light weight. One takes the
infimum over charts at the point, the supremum over algebraic
integral framed points, and then the maximum with zero. Write
\(e_2\nu_W(\lambda)\) for this last quantity. Chart independence
identifies the infimum over charts at the extended point with
that on the split chart. The extended point is among the
algebraic points in the defining supremum. Its relative
ramification is the denominator used in the depth.
For ramification index \(e\), the order to use after base change is normalized by
\[
  \frac{\operatorname{ord}_{V'}(a)}{e}=\operatorname{ord}_{V}(a)\qquad(a\in V\setminus\{0\}).
\]
The \hyperref[par:unscaled-light-depth]{unscaled estimates}
use this normalization. \label{cor:slice-finite-depth-1}

Let \(D_L=D_F\) be the differential degree of \(w_L\) and
\(M=g_0g_1\) the product of the two distinct level primes.
The valuation of the integer correction in
\eqref{eq:light-global-scalar} is \(e_2 b_p\). The level
primes, the two primes \(q,d\), and the primes dividing
\(Q_S\) are disjoint. Thus
\[
 b_p=
 \begin{cases}
 D_L\operatorname{ord}_p(Q_S),&p\text{ one of the designated small primes},\\
 \lceil D_L/2\rceil,&p=q,d,\\
 0,&\text{otherwise},
 \end{cases}
\]
The baseline uses
\[
 \nu_*=
 \begin{cases}
 D_L/(p-1),&p=g_0,g_1,\\
 \nu_W(\lambda)+b_p,&\text{otherwise}.
 \end{cases}
\]
The normalization tie is
\begin{equation}\label{eq:slice-finite-baseline}
                     \nu_\lambda=e_2\nu_* .
\end{equation}
The exponent used for the index is the rounding of
\(T(\nu_\lambda+e_2{\bf1}_{\rm bad})\) in the construction
field. An extension to a place \(v\) of relative ramification
\(e_1\) multiplies this exponent by \(e_1\). Thus there is
only one rounding before the place comparison.
\label{cor:slice-finite-depth-2}

\begin{proposition}[Finite depths on the selected light lattice]
\label{prop:slice-finite-readings}
For the baseline in \eqref{eq:slice-finite-baseline}, every allowed
integral framed point before the definite lift at which the value of \(\eta\) has finite valuation satisfies
\[
 \min_l v_\lambda(w_l(y))\le
 \nu_\lambda+\frac{c}{2i_\eta}v_\lambda(\eta(y)).
\]
The auxiliary valuation is nonnegative, allowing \(+\infty\) if that
section vanishes. At the definite lift the upper bound is
\(\nu_\lambda\). These comparisons use \(v_\lambda=\operatorname{ord}/e\)
in every allowed coefficient extension.
\end{proposition}

\begin{proof}
\proofstep{The finite boundary term at a framed point.}
Let \(\phi_\lambda(y)\ge0\) be the finite boundary function of
an algebraic integral framed point \(y\). The local
uniformization of \(y\) has lattice rank \(r_0\le N_0=2fm\).
Its short basis for the integral quadratic form of the
uniformization supplies an upper bound \(\Phi\) for the basis
values and a cross-term bound \(J_0\) in the translates.
The short-basis comparison gives
\[
 0\le\Phi\le eR_0N_0^3\phi_\lambda(y),\qquad
 \frac{R_0(N_0^4 n_W+2N_0^3)D_L^2}{g_i}\le c\quad (i=0,1).
\]
Here \(R_0=(\operatorname{rd}E)^2/(4f)\)
is the discriminant factor of the chart.

The bounds also
give \(D_L J_0<g_i\), \(4D_L+4\le g_i\), and
\(n_W\ge2(p_j+q_j)\) with \(p_j,q_j\) the signature ranks.
At the definite lift \(\Phi,J_0,\phi_\lambda\) are zero.
These hypotheses use the actual polarized chart, the level
marking with conjugate coprime half-level ideals, and the
theta coefficient expansion for its integral germs.
\label{cor:slice-finite-depth-3}

In the following comparisons \(e\) is the ramification of
the split chart over the construction field. If \(p\nmid M\),
the half-level elements are units, since each is integral and
its product with its conjugate is a unit. The local interpolation
on the uniformized copies uses this unit condition, the short
basis above, 
\[
  4D_L+4\le M,
\]
 and \(D_LJ_0<M\). On the finite
list of monomials of degree at most \(D_L\), harmonicity, row
degrees and unipotent invariance impose coefficient equations.
After clearing their denominators, let \(A\) be the matrix
of these homogeneous equations on the finite coefficient
module \(P\) over \(\mathcal O_{K_0}\).

If \(O'\) is the local
coefficient ring, torsion-freeness over the Dedekind domain
gives the exact kernel identity
\[
 \ker(A\otimes O')=(\ker A)\otimes_{\mathcal O_{K_0}}O' .
\]
Thus every local jet polynomial is an integral linear
combination of the original harmonic coefficient polynomials.
For any integral
germ whose jet has finite valuation, the translated section estimate
gives a term with order at most the order of the jet plus
\(D_L^2\Phi/M\). The term uses one of the original integral
polynomials. After translation the value has the form
\(u\mathcal D_{w_L}(\alpha)f(y')\) for \(f\in\mathcal T_{w_L}\)
and \(\operatorname{ord}(u)=0\), by the normalized translation
identity. This identity keeps the marking and the factor
\(M^{-\deg P}\) of \eqref{eq:light-rule-value}.

At the level
prime the same identity will include a nonunit determinant.
Here the determinant is a unit. Denote the least unscaled valuation at the point by \(b_{\rm loc}\), using the local order divided by \(e\). Taking the infimum over the integral germs and polynomials gives
\[
  b_{\rm loc}\le e_2\nu_W(\lambda)+\frac{D_L^2\Phi}{eM}.
\]
Equations \eqref{eq:light-frame-order} and
\eqref{eq:light-global-scalar} give the bound
\(\nu_\lambda+c\phi_\lambda(y)\) on the original light columns,
using \(M\ge g_i\) and the short-basis estimate.
\label{cor:slice-finite-depth-4}

\proofstep{The calculation at a level prime.}
If \(p=g_i\), set \(g'=g_{1-i}\).
The half-level generator for \(p\) and its conjugate have
product \(p\) and are coprime. Here the marked factor is
\(\bar\pi_i\), and in the translated evaluation the scalar acts by \(\pi_i\pi_{1-i}=\alpha\). At each conjugate pair
of embeddings exactly one reads \(\pi_i\) as a nonunit.
Their product is a nonunit and their B\'ezout identity
forbids two nonunits. Allocate disjoint copy blocks to the two resulting types.
For a label \(\tau\), write \(\Delta_{\tau,s}\) for the leading
minor of size \(s\) in its block. Set the next exponent after its
last row to zero and take the polynomial
\[
 H(X)=\prod_{\tau,s}\Delta_{\tau,s}(X)^{a_{\tau,s}-a_{\tau,s+1}},
 \qquad
 \deg_{\tau,i}H=\sum_{s\ge i}(a_{\tau,s}-a_{\tau,s+1})
                    =a_{\tau,i}.
\]
The total degree is \(D_L\), and the upper flag operations fix
each \(\Delta_{\tau,s}\). Every conjugate contraction pairs
derivatives in a single copy.

In each such copy at least one
of the two conjugate types is absent, so every mixed partial
in that contraction kills \(H\). This proves harmonicity over
the original coefficient ring.
The level Hodge--Tate jet comparison in this chart
applies when \(D_L<p-1\), \(4D_L+4\le g'\),
\(D_LJ_0<g'\), \(n_W\ge2(p_j+q_j)\), and the
level marking is split. In the chart's valuation scale the Hodge--Tate
inclusion for the polarized split semiabelian model gives integral
multiplicative characters whose differentials span a module
containing \(c_0\) times the full tangent dual for

\[
  \operatorname{ord}(c_0)=\operatorname{ord}(p)/(p-1).
\]
Use a \((p-1)\)-st root of \(p\) in the valued coefficient
extension for \(c_0\). The split eigenspace projectors modulo
\(p\) and the integral dual frame give the same inclusion
after projecting the differentials to each label. Let \(P_s\) be projection onto the first \(s\) rows and
let \(J_s\) be the ideal of the \(s\)-minors of the projected
character matrix. It is the determinant ideal of the image
lattice in those rows. The inclusion of \(c_0 O'^s\) gives
\[
                 c_0^s O'\subset J_s,\qquad
                 \operatorname{ord}(J_s)\le s\operatorname{ord}(c_0).
\]
One can realize the minimum with nested column pivots.

In
the first row choose a coefficient of minimum valuation and
subtract integral multiples of that column from the others.
With first pivot \(m\) these operations put the top \(s\) rows
in block form \(\begin{pmatrix}m&0\\ *&B_s\end{pmatrix}\).
An \(s\)-minor omitting the first column vanishes, and one
containing it is \(m\) times an \((s-1)\)-minor of \(B_s\).
Thus
\[
        J_s=m\,I_{s-1}(B_s)
\]
where \(I_{s-1}\) denotes the maximal-minor ideal. Repeat on the
remaining rows, using integral column operations that fix the pivots
already chosen. This
gives the stated leading-minor bounds in one nested choice.
Different labels of one type may assign pivots to the same copy.

\proofstep{Evaluation on the chosen finite eigenspaces.}
For each label pull its chosen character back by a CM element
acting modulo \(p\) as the projector onto that label. On a
given copy take the product of these pullbacks over all its
assigned labels. For the assigned characters \(\chi_\tau\) on one copy, their product has logarithmic differential
\[
  d\log\!\left(\prod_\tau\chi_\tau\right)=\sum_\tau d\log\chi_\tau .
\]
On projection to any one label the other summands vanish
modulo \(p\), and the chosen pivot survives. Indeed the copy
indices used within that label are distinct. Thus a single
character on each copy realizes all the prescribed minor
matrices modulo \(p\).

In the block where \(\pi_i\) is a nonunit the powered minors
therefore cost at most its degree times
\(\operatorname{ord}(p)/(p-1)\).
For powered minors of degrees \(d_b\), summing their individual costs gives
\[
  \sum_b d_b\,\frac{\operatorname{ord}(p)}{p-1}\le D_L\,\frac{\operatorname{ord}(p)}{p-1}.
\]
 The comparison with the
covariant eigencharacter jets is modulo \(p\). The bound
\(D_L<p-1\) puts this determinant product strictly below
that order, so the error cannot change its valuation.
\label{cor:slice-finite-depth-5}

The local eigencharacter section on each of the \(p_j+q_j\)
nonunit-block copies costs at most \(r_0\Phi/(4g')\) in addition.
On the unit block the jet comparison produces a unit value before multiplication by the section factors. If the required jet degree on a copy
is \(d_b\), its section factor costs at most \(d_b^2\Phi/g'\).
Since the degrees \(d_b\) are nonnegative and sum to at most \(D_L\),
the total section cost on these copies satisfies
\[
 \sum_b d_b^2\,\Phi/g'\le\left(\sum_b d_b\right)^2\Phi/g'
                         \le D_L^2\Phi/g'.
\] The covariant evaluations of
the two blocks multiply, since the differential polynomial on
each block uses only that block's variables. The expansion into
marked translates at the companion level costs at most another
\(D_L^2\Phi/g'\) for the selected term. The term is chosen by
the ultrametric inequality after expanding the covariant differential expression with integral coefficients.

The selected term reads \(u\mathcal D_{w_L}(\alpha)f(y')\) with unit-valued \(u\) and
\(f\in\mathcal T_{w_L}\). The determinant equals the scalar's
determinant by \eqref{eq:light-oriented-determinant}, and
the integer correction in \eqref{eq:light-global-scalar}
is a unit here. Thus the term has the normalized valuation
of \(a_W f\) at \(y\), by \eqref{eq:light-frame-order}.
After division by \(e\) the first cost becomes
\(e_2D_L/(p-1)\). The remaining cost is at most
\[
       \frac{((p_j+q_j)r_0/4+2D_L^2)\Phi}{e g'}
       \le \frac{R_0(N_0^4n_W+2N_0^3)D_L^2}{g'}
                         \phi_\lambda(y).
\]
Here \(p_j+q_j\le n_W\), \(r_0\le N_0\), and \(D_L\ge1\).
The least valuation of a scaled section is again at most
\(\nu_\lambda+c\phi_\lambda(y)\).
\label{cor:slice-finite-depth-6}

\proofstep{Adding the auxiliary cusp section before the final lift.}
Before the last lift the auxiliary cusp comparison gives
\(2i_\eta\phi_\lambda(y)\le v_\lambda(\eta(y))\)
with \(i_\eta>0\) for the same frame, point and boundary
function. We therefore have, at points where the auxiliary section has finite valuation,
\begin{equation}\label{eq:slice-tied-depth}
 \min_l v_\lambda(w_l(y))\le
 \nu_\lambda+\frac{c}{2i_\eta}v_\lambda(\eta(y)),
 \qquad v_\lambda(\eta(y))\ge0 .
\end{equation}
Nonnegativity also holds when the auxiliary section vanishes and its valuation is infinite. At the definite lift the local
bound is simply \(\nu_\lambda\). Rechoosing the short and
complementary bases keeps the light columns and the
auxiliary section fixed. Thus the inequalities hold on
the actual two runs. They are independent of the
framed point and the degree and ramification of its
coefficient extension. The \hyperref[par:slice-valued-orders]{normal-chart passage} uses these same
sections even when the branch field is transcendental.
\label{cor:slice-finite-depth-7}
\end{proof}

\paragraph{Bounding the unscaled depth by covariant differential operators.}
\label{par:unscaled-light-depth}
All orders in this paragraph are in units of \(v_p(p)=1\).
They give bounds for \(\nu_W(\lambda)\) by the
\hyperref[par:slice-finite-depth]{definition of chart depth}.
A covariant differential operator on an integral germ may be evaluated over a complete algebraically closed valued extension of the chart field. Its value depends on a finite jet of the germ. Over the extension,
these jets are integral linear combinations of coefficient
monomials. The integral harmonic coefficient lattice
extends as well. Indeed, in the fixed multidegree the
harmonicity equations are linear over the original
valuation ring. Their kernel extends under flat scalar
extension, since the new valuation ring is torsion-free
over the original DVR.

If \(H\) is the matrix of the harmonicity equations and \(V\subset V'\) the extension of valuation rings, this says
\[
  (\ker H)\otimes_V V'=\ker(H\otimes_V V').
\]
 The covariance operators commute
with coefficient extension. By bilinearity and the
ultrametric inequality, each lower bound on the original valuations also holds after scalar extension, in the same units. Thus any extended test giving a finite valuation bounds the original infimum from above.
\label{cor:slice-unscaled-depth-1}

\begin{proposition}[Upper bounds for the unscaled chart depth]
\label{prop:primitive-light-covariants}
At every allowed finite place, primitive columns in the local
eigenspaces bound the unscaled depth \(\nu_W(\lambda)\) used in
Proposition~\ref{prop:slice-finite-readings}. Use the
\hyperref[par:slice-finite-depth]{weighted chart depth} with \(v_p(p)=1\).
At the small primes the eigenflag and colength estimate is
\eqref{eq:light-small-prime-depth}; at \(q,d\) the integral
Hodge--Tate estimate is \(\nu_W\le D_L/(p-1)\). On the good charts
the bounds use monomials or alternating pairs of Frobenius orbits. In the balanced level-one conjugating case write \(C_{\rm lead}\) for the leading flow cost. The degree contributions \(D_a\) enter as
\[
 C_{\rm lead}=\frac1{p^2-1}\sum_a D_a.
\]

\end{proposition}

\begin{proof}
\proofstep{Choice of primitive columns in each eigenspace.}
For each label \(\tau\) of rank \(r_\tau>0\), let
\(a_{\tau,1}\ge\cdots\ge a_{\tau,r_\tau}\ge0\) be the
light row weights, and put \(a_{\tau,r_\tau+1}=0\).
Choose \(r_\tau\) columns of integral characters on that
frame, with disjoint column sets for distinct labels.
The chosen number of rows permits this. If their derivative
matrix has leading minors \(\Delta_{\tau,i}\), use the
product of the corresponding determinant polynomials to
the powers \(a_{\tau,i}-a_{\tau,i+1}\).
Row \(j\) occurs precisely in the minors of sizes \(i\ge j\).
Its exponent is therefore
\[
 \sum_{i=j}^{r_\tau}(a_{\tau,i}-a_{\tau,i+1})=a_{\tau,j}.
\] A harmonic
contraction pairs conjugate types in the same column.
Each column belongs to a single type, so every contraction
kills the polynomial. Its coefficients are integral.

In formal logarithm coordinates the character derivatives
evaluate the determinant product with no factorial loss.
If \(v_p(\Delta_{\tau,i})\le h_\tau\) for each \(i\), this
product has valuation at most 
\[
  \sum_\tau a_{\tau,1}h_\tau.
\]
The integral covariant character-germ comparison replaces the
ordinary reading by this product plus an error of strictly
higher valuation. It uses the integral pairing and chooses
auxiliary eigenvalues congruent to one to accuracy greater
than the order of the nonzero product. Thus the covariant
reading has the same valuation.
\label{cor:slice-unscaled-depth-2}

For the bounds on the minors, start with a finite collection
of differential columns attaining the chart colength.
In the fixed first flag row choose an entry of least valuation
and use its column to clear that row in every other column
by integral multiples.
For a pivot \(a\ne0\) of least valuation and another entry \(b\) on its row, the clearing step uses
\[
  v_p(b/a)\ge0,\qquad b-(b/a)a=0 .
\]
 At each step take the minimum in the
next fixed flag row of the residual columns. Each residual pivot comes from an original
column after earlier column subtractions.

For the nested
sequence the leading determinants are unchanged by these
subtractions and have valuations equal to the partial sums
of pivot valuations. Since each pivot is integral, the
partial sum is at most the full sum, which is the valuation
of the determinant ideal.
If \(p_1,\ldots,p_n\) are the successive nonzero integral pivots, the cost of any initial segment is bounded by
\[
  \sum_{j=1}^{i}v_p(p_j)\le\sum_{j=1}^{n}v_p(p_j)\qquad(i\le n).
\]
 This choice of integral characters over the original ring gives the required leading minors.
\label{cor:slice-unscaled-depth-3}

\proofstep{The Hodge and colength inputs at the small primes.}
At the designated small primes use the split auxiliary
quadratic field. The selected CM half is partitioned by the
host places of the split blocks. If a block has \(n_w\)
embeddings and discriminant coefficient \(d_w\), the CM
character comparison on the integral chart bounds the
product of its absolute determinant lower bounds below by
\(|p|^{n_wq_j(1/(p-1)+d_w/2)}\).
Taking logarithms and dividing by \(\log|p|<0\), then using
\(\sum_w n_w=f\) and \(\sum_w n_w d_w=fv_d(p)\), gives
\[
        \sum_{\tau\text{ in the selected half}} h_\tau
            \le fq_j(1+v_d(p)/2).
\]
On the complementary half the integral Hodge--Tate
character inclusion contains \(p\) times the tangent
dual, since \(v_p(p)=1\ge1/(p-1)\).

Its restriction projects onto \(p\) times the dual of the
integral rank-\(q_j\) eigendirection frame, with determinant
cost at most \(q_j\) for each label.
The small-prime page bounds the first light row of the
selected half by \((2b+S_c)/T\), and the sum of the first
rows of the complementary half by \(f\eta_0/T\).
With \(S_e=S_c+\eta_0\), \(q_j\le m\) and
\(v_d(p)\le v_d\), the determinant product gives
\begin{equation}\label{eq:light-small-prime-depth}
              T\nu_W\le m(2+v_d)(2fb+fS_e).
\end{equation}
This uses the summed colength once on the selected half.
The global comparison may use \(2fb\) as a majorant for
the weight sum of any lift. At \(q,d\) the crude integral
Hodge--Tate bound, obtained by using the degree of each
label in the determinant reading, is

\[
  \nu_W\le D_L/(p-1).
\]
The CM colength and Hodge--Tate inclusions are local inputs
on the given uniformization and eigendirection frames.
\label{cor:slice-unscaled-depth-4}

\paragraph{Depth bounds at good primes and odd shifts.}
On the unramified chart, order the embedding labels by their
residue Frobenius. Let \(C_\tau\) be the character colength at label \(\tau\). If \(\tau_h\) denotes its \(h\)-th successor, the unramified estimate uses the ranks of the conjugate eigenspaces
\[
  C_\tau\le\sum_{h\ge1}p^{-h}r_{\overline{\tau_h}}.
\]
 The row systems use this same
labelled orbit. The characteristic-zero frame and the
integral chart give the character and monomial blocks with
the prescribed types.
\label{cor:slice-unscaled-depth-5}

\proofstep{The diagonal monomial on a split block.}
For a monomial block use formal curve coordinates with
identity linear part on the rows of the indicated half.
The product of the leading minor polynomials has
diagonal monomial with coefficient one. For each minor,
the identity permutation uniquely maximizes the sum of
the products of the row and column indices, by the
rearrangement inequality.
For increasing row and column sequences \((r_i),(c_i)\), strictness for a nonidentity permutation reads
\[
  \sum_i r_ic_{\sigma(i)}<\sum_i r_ic_i\qquad(\sigma\ne1).
\]
 This identifies its diagonal
term uniquely even inside the product. If \(b_i\) are the diagonal monomial exponents, the
prescribed germ derivative reads \(\prod_i b_i!\) because
the coefficient of this monomial is one and the higher
total-degree terms give zero. The single-half covariant
factor is one. Legendre's formula gives the loss
\[
         \sum_i v_p(b_i!)=\sum_i\sum_{a\ge1}\lfloor b_i/p^a\rfloor
                 \le \frac{\sum_i b_i}{p-1},
\]
with zero loss when all \(b_i<p\).

If a first-row
exponent is at most \(s_0(p-1)\), its integral weight gaps
can be distributed among \(s_0\) row slots of capacity
\(p-1\), by filling the slots successively and splitting
gaps at the boundaries. Later rows use only a subset of
the gaps. Different labels of one type reuse the slots
using their labelled coordinates. When 
\[
  4rf s_0\le n_W,
\]
there are enough rows of each type to obtain a monomial of valuation zero on the whole chart. In particular, \(p>D_L\) permits \(s_0=1\).
In the split case the character block of exponent \(u\)
costs \(u m/(p-1)\) on the distinguished part and the
complementary monomials use the remaining degree.
\label{cor:slice-unscaled-depth-6}

For an alternating pair of orbits the selected inert-prime
hypothesis gives \(p\ge5\). Put \(n=k-1\ge1\) and let
\(d_0\) be its orbit length. An odd shift \(1\le i<d_0\) uses block
exponents \(p^i y,y\).
On that pair the total degree consumed is
\[
  p^iy+y=(p^i+1)y .
\]
 The adapted chain begins with at least
\(n-((i-1)/2+2)(m-j)\) independent residue directions,
loses at most \(m-j\) at each of \((i-1)/2\) successive
semilinear steps and loses at most \(j\) in the final
pairing. Hence at least \(n-(i+1)m\) rows remain.

An integral completion of the opposite-side vectors puts
their retained pairing in negative identity form. The
integral gauge estimates and the odd-level jet theorem on
this chart then give admissible \(y\) in every window of
the specified length, as well as \(y=0\). For positive
\(y\), write \(\gamma=p/(p^2-1)\). The block cost is at most
\[
  n\gamma(p^i+1)y-ny+n(1+\log_p n)
                 +\frac{(i+1)m}{p-1}(p^i+1)y .
\]
The term with \(m\) bounds the factorials of the lost
rows. At \(y=0\) the cost is zero. This uses the odd-level
jet estimate at the stated row ranks, normal form of its
mixed gauge coefficients and integral curve inputs.
\label{cor:slice-unscaled-depth-7}

\proofstep{The cost of the alternating orbit.}
To see how the blocks give the orbit price, write \(W_a\)
for the distinguished exponents at the light weight and
\(w_a>0\) for the target exponents. The degree comparison
gives \((1-\theta)w_a\le T W_a\le w_a+C_1S_e\) with
\(0\le\theta\le1\) and \(n\theta\le5m\).
Set \(V_a=(1-\theta)w_a/T\).
The directed odd shifts allocate nonnegative variables
\(x_{a,i}\) with constraints
\[
         \sum_i p^i x_{a,i}+\sum_i x_{a-i,i}\le nV_a .
\]

On demands \(D_a=nV_a\), define the cost
\(F_p(D)=\gamma\sum_a D_a-\max\sum_{a,i}x_{a,i}\).
Summing the constraints proves that the maximum is at
most \(\sum_a D_a/(p+1)\). The cost is nonnegative,
homogeneous and invariant under cyclic shifts.
If levels \(i\ge I=i_0(k,p)\) are discarded from a
maximizer, the lost flow is at most \(p^{-I}\sum_aD_a\).
Round each remaining \(x_{a,i}/n\) down to an allowed
integer \(y_{a,i}\), with loss at most the window size
\(S\) in this quotient. Then
\[
 u_a=W_a-\sum_i p^i y_{a,i}-\sum_i y_{a-i,i}\ge0
\]
is an integer character-block exponent.

In the total cost
the leading character term and leading odd terms sum to
\(n\gamma\sum_a W_a\), by the displayed equality and
cyclic reindexing. The saving from the rounded flow loses
at most \(d_0 I n S\). The positive odd factorial costs
are at most \(2I m\sum_a V_a/(p-1)\) since
\(p^i+1\le2p^i\) and \(i+1\le I\) on the retained levels.
The character cost beyond its leading term is at most
\(m\sum_aW_a/(p-1)\) and the additive odd cost is at most
\(d_0 I n(1+\log_p n)\).
The chart range with these odd blocks has
\[
 I\le5\log k,\qquad S+1+\log_p n\le707 n^2\log k.
\]
Multiplying the rounding and additive window costs by \(T\) gives the combined upper bound
\[
     3535 d_0 Tn^3(\log k)^2 .
\]

The discarded levels cost at most
\(p^{-I}k\sum_a w_a\).
Replacing \(V_a\) by \(W_a\) and paying the remaining
monomial degrees adds at most a fixed multiple of
\((m(\sum_a w_a)\log k+d_0 kS_e)/p\) on the light-weight
and row-budget comparisons. For the sum over orbits, the total
length is at most \(f\). Use the uniform weight bound
\(\widehat W=2fb\). Below the threshold
\(p_0=mb/(Tf n^3\log b)\),
\[
 Tf n^3(\log k)^2
 < \frac{mb(\log k)^2}{p\log b}
 \le \frac{m\widehat W\log k}{p}.
\]
Thus the absolute window loss fits the weight allowance with
its fixed numerical multiplier. At \(p\ge p_0\) the typed row
supply on the selected chart gives the zero-depth monomial
allocation, including equality at the threshold.
\label{cor:slice-unscaled-depth-8}

\proofstep{Vanishing of the depth at large nonsplit primes.}
At large nonsplit primes the leading cost has a direct
bound. If \(d_0\ge2\) and all demand ratios are at most \(p\), the system
\(p x_a+x_{a-1}=D_a\) has the cyclic solution
\[
 x_a=\frac{\sum_{i=0}^{d_0-1}(-p^{-1})^iD_{a-i}}
                {p(1-(-p^{-1})^{d_0})}.
\]
Pair successive terms in the numerator. Each pair is
nonnegative and in odd length the leftover term is
positive for positive demands. This realizes the upper
bound on the total flow using only level one. Thus
\[
 F_p(D)=\sum_a D_a/(p^2-1).
\]
The selected target ratios give the required balance for
\(p>X_2\), so the leading cost is at most
\(2k\sum_a w_a/p^2\).

At smaller primes the forced content is the orbit average
of the orbit programmes. Cyclic invariance gives their sum
over representatives in the place formula. When a
Frobenius power is conjugation, choose the first
admissible odd level on the same typed chart and the
programme comparison for that orbit. For Frobenius equal to
conjugation choose level one directly and write

\[
  W=(p+1)y+u
\]
 with \(y\) in the jet window and
\(0\le u\le(p+1)S\). This gives the explicit
\(1/(p^2-1)\) leading term in that case.
The cost comparison for the first conjugating odd level
and the block supplies in the respective charts are used
with these hypotheses.
\label{cor:slice-unscaled-depth-9}
\end{proof}

\paragraph{The integral lower bound on a single component.}
\hypertarget{source-slice-metrics}{}
Keep the integer \(h\) of \eqref{eq:chow-lower-section}. A product
of \(h\) integral light forms, and any integral sum of those
products, satisfies the integral framed-point test at weight
\(h w_L\). This follows by multiplying their values in each framed
valuation ring. Positive weight and cuspidality persist under this
multiplication. Suppose the resulting form \(g\) restricts
nontrivially to \(Z\). On the generic fibre multiply it by the
idempotent of \(Z\). The resulting section \(g_Z\) is algebraic
and cuspidal. In every framed DVR test its generic point belongs
to exactly one component. The value tested there equals the
original value or zero. Thus \(g_Z\) belongs to the same global
integral lattice and is nonzero. This proof also covers components
whose closures meet in a special fibre.

\begin{proposition}[The lower norm bound on a single component]\label{prop:integral-chow-component-height}
For the fixed exponent \(h\) of \eqref{eq:chow-lower-section}
let \(g\) be an integral sum of products of \(h\) integral
light sections with \(g|_Z\ne0\).
Use the constant \(P_0\) of \eqref{eq:component-chow-lower-constant},
chosen from the common column-size and weight data before \(g\) and \(Z\).
On every such component,
\[
 \operatorname{mean}_\sigma\log\|g|_Z\|_{\sup,\sigma}\ge-hP_0S_L.
\]

The integral lifting uses the generic component idempotent
and the original cusp lattice at weight \(h w_L\).
\end{proposition}

\begin{proof}
\proofstep{Bounding the supremum by algebraic evaluations.}
At an embedding \(\sigma\), the supremum norm of \(g_Z\) is bounded
by the supremum of the norms of \(g\), with the same metric, under embeddings \(\Omega_Z\to\mathbb C\) extending
\(\sigma\). Indeed, work on an open affine flag chart over \(K\)
on which the frame has been trivialized. Its ring is countable.
For a nonzero element of the ring the set of complex points
where it vanishes has empty interior in each analytic component.
Each geometric component dominates the integral chart over \(K\),
and on the smooth complex chart the assertion follows from
the identity theorem. Baire's theorem therefore gives a dense
set of points where the ring embeds in \(\mathbb C\); the map
at such a point extends to its fraction field.

Equivalently the evaluation map on the chart ring \(A\) is injective at each point of the dense set
\[
  \bigcap_{a\in A\setminus\{0\}}\bigl\{x:a(x)\ne0\bigr\}.
\]
Let \(F\) be the coefficient of the section in an adapted frame and
let \(J>0\) be the norm of that frame. The section has norm \(|F|J\).
Under a change of frame with weight character \(\chi\), the coefficient
and frame norm transform inversely, so
\[
  |\chi|\,|F|\cdot |\chi|^{-1}J=|F|J.
\]
Consequently the section norms on the dense set agree
with the generic norm. The metric is continuous across the
chart and the corresponding flag cells are dense in the component.
At conjugate embeddings one uses the conjugate uniformization
with its prescribed metric. This proves the supremum comparison.

Apply \eqref{eq:slice-integral-norm} for nonzero integral
\(F\) of the required weight \(v\). In the notation of that
estimate, put \(P_*=3000Z^2\) with \(Z=fr+L+\log d_E+|\Omega|\).
Then
\[
  \operatorname{mean}_\sigma\log\|F\|_{\sup,\sigma}
        \ge-P_* (1+\log|v|)|v| .
\]
Here \(|v|\) is the weight size. Apply the bound to \(g_Z\)
with \(v=h w_L\) and divide by \(hS_L\). Use \(1\le |h w_L|\le hS_L\)
and
\[
 \log h\le \log\!\big(8(s+2)^3(2+P_O)\big)
            +2\log N+\log\delta .
\]
The family bounds are \(\log N\le n_B\) and
\(\log\delta\le1041z^3\), where \(n_B\) bounds the logarithm
of the column size. One may take
\begin{equation}\label{eq:component-chow-lower-constant}
 P_0=\max\{1,P_*[1+\log S_L+\log(8(s+2)^3(2+P_O))
                                         +2n_B+1041z^3]\}.
\end{equation}
Thus the lower bound of Proposition~\ref{prop:slice-lower-integral-norm}
at the single weight \(h w_L\) gives \eqref{eq:chow-lower-section}.
The integer \(h\) and the bound for \(P_0\) are common to all components. \label{cor:slice-metrics-1}

\end{proof}

\paragraph{From algebraic integral points to the orders of the slice.}
\label{par:slice-valued-orders}
The depth test available on algebraic points is a bound on the minimum valuation of the light sections in an adapted integral frame. We explain how
it gives orders at the nonalgebraic places used by the height.
The following extension lemma is the input to the branch calculation.

\begin{lemma}[Integral tests at the closed fibre]\label{lem:slice-valuative-extension}
Let \(O\) be a complete
DVR of a number field completion, with uniformizer \(\pi\).
Let \(A\) be a flat normal integral algebra of finite type over
\(O\), and let \(g\in\operatorname{Frac}A\). Suppose that in
every map \(A\to V\) to a finite extension DVR of \(O\) whose
generic point belongs to a fixed dense open where \(g\) is
defined, the value of \(g\) lies in \(V\). Then the simultaneous local conclusion is
\[
 g\in\bigcap_{\substack{q\in\operatorname{Spec}A\\ q+(\pi)\ne A}}A_q.
\]


\end{lemma}

\begin{proof}
A pole of \(g\) at a height-one prime \(P\subset q\)
gives a point \((P,\infty)\) on the closure of its graph
in \(\operatorname{Spec} A\times\mathbb P^1\), by evaluating
on the DVR at \(P\). Taking the closure gives a closed point
\((y,\infty)\) on the special fibre, since 
\[
  (q,\pi)\ne A.
\]
An integral curve through that point with generic point in
the prescribed open can be chosen as follows.

In the local
domain of the graph at this point, of dimension \(d\), choose
a nonzero \(b\) divisible by \(\pi\) such that \(D(b)\) is in
the original graph over that open and \(1/g\) is nonzero
there. The case \(g=0\) is immediate. Complete \(b\) to a
system of parameters \(b,a_2,\ldots,a_d\). Write \(B\) for this local domain and \(I=(a_2,\ldots,a_d)\).
Killing one element of the maximal ideal lowers the dimension
by at most one \cite{slice-dimension-drop}. Hence
\[
 \dim B/I\ge d-(d-1)=1,\qquad \sqrt{I+(b)}=\mathfrak m_B .
\]
If every minimal prime over \(I\) contained \(b\), these two facts would
give \(\dim B/I=0\). Choose such a prime \(\mathfrak p\) avoiding \(b\).
In the local domain \(B/\mathfrak p\) the nonzero image of \(b\) generates a
maximal-primary ideal. The same dimension inequality now gives
\(\dim B/\mathfrak p=1\).

Let \(C\) be the closure of this branch. It dominates \(\operatorname{Spec} O\)
since \(\pi\mid b\). At the chosen closed point its local dimension is one
and its residue field is finite over the residue field of \(O\).
The dimension formula over the universally catenary DVR \(O\) gives
\[
 1=1+\operatorname{trdeg}_{\operatorname{Frac}O}K(C)
       -\operatorname{trdeg}_{O/(\pi)}\kappa(y).
\]
Thus the finitely generated field \(K(C)\) is finite over
\(\operatorname{Frac} O\) \cite{dim-slice}. A local ring of the
normalization above the closed point is a DVR dominating \(O\).
The extension of the valuation of a complete DVR to this finite
extension is unique and its valuation ring is finite over \(O\)
\cite{dvr-slice}. On the resulting test the generic point lies in
\(D(b)\), whereas \(1/g\) has strictly positive valuation.
This contradicts the assumption on \(g\).
Consequently \(g\) has no pole at any height-one prime of \(A_q\).
Normality now gives
\[
 g\in\bigcap_{\operatorname{ht}_{A_q}\mathfrak p=1}
               (A_q)_{\mathfrak p}=A_q
\]
by the intersection theorem for Noetherian normal domains
\cite{slice-normal-intersection}.\label{cor:slice-metrics-2}

\end{proof}

\begin{proposition}[Transfer of the framed depth and metric tests]\label{prop:slice-chart-metric-readings}
Use the selected light columns and, before the final lift,
the auxiliary \(\eta\) with its stated pointwise depth
inequality. On the chosen normalized charts the inequality
holds at the generic slice branches, using the coefficient
baseline on vertical branches and without that term on
horizontal branches. At a horizontal branch \(q\) its incidence exponent obeys
\[
 0\le E_q\le\frac{c\,n_q(\eta)}{2i_\eta}.
\]

Evaluation in the adapted complex frames gives the paired light and
auxiliary metric bound in the proof. At the definite lift
only the finite baseline is used.
\end{proposition}

\begin{proof}
\proofstep{Transferring the valuation test to the flag chart.}
Apply this assertion on integral framed charts which record a
smallest light column \(w_b\). Such a chart adjoins the regular
ratios \(w_l/w_b\) and uses the normalization of the corresponding
blowup chart. For a valued point with valuation \(v\), choose a column with least
finite order. Then
\[
 v(w_l/w_b)=v(w_l)-\min_j v(w_j)\ge0
\]
for every \(l\), so all the adjoined ratios are in this valuation ring. The required cover uses the
integral model, with frames chosen over its affine opens. To see that these charts suffice for the centre criterion, first
take a vertical DVR. Complete it. Its map from the coefficient
DVR extends to the completion, and its order on the original
field stays unchanged. Use the proper model at this prime
and its flag bundle as in the construction of \(y_\zeta\).

At each elimination step in the flag matrix, choose a nonzero eligible pivot \(a\) of minimum valuation among the eligible entries \(a_j\). Comparability of principal ideals in the valuation ring \(V\) gives
\[
        v(a_j)\ge v(a)\quad\Longrightarrow\quad a_j/a\in V .
\]
The elimination coordinates therefore give an integral point on this flag chart.
The theta generator and these flags may be chosen on finitely
many trivializing charts. After adjoining the ratios for a
smallest light column, the generic fibre over the original
nonboundary locus is still covered: the columns there
generate the unit ideal in a trivialization.

Inverting the
pivot therefore gives that chart on the generic fibre.
After passing to the completed coefficient ring take the closures
of the generic components and normalize. These normalizations
are finite, since the base is excellent, and the valuative
criterion for the resulting finite maps lifts the point with
its generic lift. On the open generic flag charts the
normalization agrees with the smooth model. The resulting
affines are normal and flat over the DVR by torsion-freeness.
In particular if \(A\) is such an affine algebra over the coefficient DVR, its flatness follows from
\[
  0\longrightarrow A\xrightarrow{\ \pi\ }A\longrightarrow A/\pi A\longrightarrow0.
\]

\proofstep{Horizontal orders and the line-frame correction.}
For a horizontal valuation there are two orders to keep track
of. Let \(T\) be its DVR with a copy of \(K\) on which its
valuation is trivial. One can extend it, with ramification
index one, to a DVR \(V\) containing \(K_\lambda\) as a
field of units together with zero. To see this, complete \(T\)
and choose a coefficient field containing \(K\). Explicitly,
lift a transcendence basis of the residue field \(k/K\).
The embedding of the rational function field so obtained
extends across the separable algebraic residue extension by
Hensel lifting, uniquely with the chosen residues. Thus

\[
  \widehat T=k[[u]]
\]
 for a uniformizer \(u\).
Embed \(k\) and \(K_\lambda\) over \(K\) in a common field \(k'\), using the fraction field of a prime quotient of \(k\otimes_K K_\lambda\). The resulting inclusions are
\[
       T\longrightarrow\widehat T=k[[u]]
          \longrightarrow V=k'[[u]],\qquad K_\lambda\subset k'.
\]
The first nonzero coefficient of a series stays nonzero after extending \(k\) to \(k'\), so \(u\)-adic orders are unchanged. Every nonzero element of \(K_\lambda\) is a constant unit of \(V\).

Let \(R'\subset k'\) be a valuation ring extending the ring of \(K_\lambda\).
The composite valuation has ring and maximal ideal
\[
 W=\{a\in V:\bar a\in R'\},\qquad
 \mathfrak m_W=\{a\in V:\bar a\in\mathfrak m_{R'}\}.
\]
Thus \(\mathfrak m_V\subset\mathfrak m_W\). The
coefficient uniformizer \(\pi\) is a unit of \(V\), but its
residue in \(k'\) lies in \(\mathfrak m_{R'}\), so it belongs
to \(\mathfrak m_W\).
In particular the coefficient uniformizer distinguishes the two valuation rings by
\[
  \pi\in V^\times,\qquad \pi\in\mathfrak m_W .
\]
The extension of valuations used here follows by domination of
a local ring by a valuation ring, or equivalently the extension
theorem for Krull valuations.

Apply proper extension, flag pivots, light pivots and the normal
charts first over \(W\). The morphism to the fine valuation ring
evaluates the chart in \(W\subset V\). Pullback to \(V\) gives
the original horizontal point after the unramified extension.
The integral frames differ there by units on the diagonals:
the two underlying flag points agree by separatedness, and both
frames are bases over \(V\). In this chart the inverse image of
\(\mathfrak m_W\) contains both the inverse image of
\(\mathfrak m_V\) and \(\pi\).
Writing \(A\) for the chart ring, the two inverse images thus satisfy
\[
  (\mathfrak m_V\cap A)+\pi A\subseteq\mathfrak m_W\cap A .
\]
 This proves the specialization condition for the chosen model. An arbitrary neighborhood of a horizontal centre need not satisfy it.
The good-level compactification used here is the same at each
run and at all its places. \label{cor:slice-metrics-3}

\proofstep{The corresponding least orders of the selected sections.}
The algebraic test \eqref{eq:slice-tied-depth}
in the light and auxiliary frames gives, with
\(v(\pi_\lambda)=1\) for a uniformizer of \(K_\lambda\),
\[
 0\le v(\eta),\qquad
 v(w_b)\le\nu_\lambda+c'v(\eta),\qquad c'=c/(2i_\eta).
\]
Take \(v(\eta)\) finite by excluding its zero divisor in the
open used to test regularity. It is nonzero on each component.
For each integer \(l\ge1\) apply the assertion to
\[
 g_l=\pi^{\lceil e_O l\nu_\lambda\rceil}
           \eta^{\lceil l c'\rceil}w_b^{-l},
 \qquad e_O=\operatorname{ord}_O(\pi_\lambda).
\]
Use the readings in the chosen trivializations. The algebraic
test and \(v(\pi)=1/e_O\) make every such value integral.
Apply the same criterion to \(\eta\). Evaluate the resulting
local-ring expressions at the centre of a slice valuation,
where the generic readings are defined. After evaluation at the branch, the inequalities are
\[
 l\,v(w_b)\le
   \lceil e_Ol\nu_\lambda\rceil\,v(\pi)
      +\lceil lc'\rceil\,v(\eta).
\]
Divide by \(l\), keeping the branch fixed. The rounding
errors tend to zero. On a vertical branch
\(e_Ov(\pi)=1\) in the \(v(\pi_\lambda)=1\) normalization,
which retains the coefficient baseline.

On the horizontal
branch the coefficient valuation is trivial and
\(v(\pi)=0\), so only \(c'v(\eta)\) remains. Multiplying
by the ramification factor on a vertical branch converts to
its integer order. Summing with residue degrees gives exactly
the auxiliary coefficient in the finite and horizontal
comparisons for \(H_L\). Indeed 
\[
  E_q\le c'n_q(\eta)
\]
and each \(n_q(\eta)\ge0\). For a horizontal polynomial
factor \(q\) its sphere height may be negative. Applying
\eqref{eq:slice-sphere-lower} to the nonnegative difference
\((c/i_\eta)n_q(\eta)-E_q\) incurs altogether at most
\(s(c/i_\eta)a_\eta\ell_N\). This is the term
\(s\beta a_e\ell_N\) above. At the last lift the
test has only its baseline and the power of \(\eta\) is omitted.

\proofstep{The complex metric in the adapted frame.}
For the complex metric, choose the adapted frame attached to a
generic embedding of \(\Omega_Z\). The invariant norms of the
light forms then agree with the chart norms. For the auxiliary
section the generic norm is at most the chart norm. The chart
estimates with first-minimum function \(P\) are
\[
 -\log\big(\sum_l\|w_l\|^2\big)^{1/2}\le P_5 D_F+cP,
 \qquad \log\|\eta\|\le-i_\eta(P-P_2).
\]
Since \(c\ge0\) and \(i_\eta>0\), substitution gives at each
generic reading
\[
 -\log\big(\sum_l\|w_l\|^2\big)^{1/2}
  \le P_5D_F+ (c/i_\eta)(-\log\|\eta\|)+cP_2 .
\]
The auxiliary reading is nonzero under a field embedding. These
arguments reduce the metric comparisons to the algebraic-point
bound and the two analytic first-minimum
bounds on the selected forms. \label{cor:slice-metrics-4}
\end{proof}

\paragraph{Specialization of the incidence form.}
\hypertarget{source-slice-integrals}{}
The finite orders describe the vertical contribution. For the horizontal term in the moving height we must identify the divisor obtained by specializing the incidence form. Suppress row zero, put \(A=K[X_{ij}:1\le i\le s]\),
\(B=\operatorname{Frac} A\), and let \(L/B\) be the incidence
field of one component with that row suppressed. Its degree
is \(\delta_Z\). Adjoining the independent zeroth row to
both fields gives the incidence extension used above. In a
generic light frame choose \(w_b\ne0\), and put \(r_j=w_j/w_b\).
The polynomial
\[
       U(Y)=\operatorname{Norm}_{L/B}\big(\sum_j Y_jr_j\big)
             \quad\in B[Y]
\]
has degree \(\delta_Z\).

Multiplication by one scalar of \(B^*\)
makes it primitive over the factorial ring \(A\), unique up
to \(K^*\). This primitive polynomial is \(F_Z(Y,X)\).
Indeed over the geometric generic cut it is the product of
the linear forms of its points, so its zero set for the last
hyperplane is exactly the locus cut out by the Chow polynomial.
The multiplicity at a general point of the light image is
the degree of the component field over that image. Both
polynomials have that multiplicity: on the norm side the
points over the image appear as the factors of the separable
norm. The irreducible Chow polynomial depends on row zero
and is primitive in that row over \(A\). Gauss's lemma
identifies its indicated power with the primitive polynomial.

\begin{proposition}[Specialization and integrability on the incidence slice]
\label{prop:integrable-incidence-heights}
On the selected horizontal cut for a component \(Z\), write \(L/B\) for its incidence extension of degree \(\delta_Z\), and \(r_j\) for the light-coordinate ratios in a nonzero light frame. The moving incidence divisor is computed by specializing the primitive representative of
\[
 \operatorname{Norm}_{L/B}\!\left(\sum_j Y_jr_j\right),\qquad [L:B]=\delta_Z.
\]
 The weight metrics
used to compute \(H_L\) and \(H_F\) are integrable on those charts.
At the definite lift the same discriminant comparison applies with
the baseline finite-depth term and the full cusp space.
\end{proposition}

\begin{proof}
\proofstep{The contribution over a horizontal prime.}
For an irreducible \(q\) of \(A\), write \(w\mid q\) for the
places of \(L\) above its order. The order of the content of
\(U\) at \(q\) is
\[
           \sum_{w\mid q} f_w\min_j w(r_j).
\]
Indeed extend the base valuation to the independent \(Y\)'s
by the Gauss rule. Its extensions from each \(w\) still have
the same ramification index and residue degree. This follows
by tensoring the local normalization with the Gauss DVR and
localizing: the old residue fields acquire independent
transcendentals \(\bar Y\), so remain fields and the old
uniformizers still generate the local maximal ideals.

In these extensions the order of the sum is the minimum.
The norm valuation formula follows by taking the determinant
of multiplication on the normalization, or by counting the
successive residue field quotients of each local principal
ideal. Replacing \(Y\) by a constant tuple \(x\) for which
\(w_x=\sum_jx_j w_j\ne0\) on each component gives
\[
 v_q(F_Z(x,X))=
   \sum_{w\mid q} f_w\big(w(w_x/w_b)-\min_j w(r_j)\big).
\]
Fix a \(B\)-basis of the finite algebra. If \(M_l\)
is multiplication by the \(l\)-th column ratio, the norm in
the moving row is
\[
             \det\!\left(\sum_l X_{0l}M_l\right).
\]
Its specialization is the determinant for multiplication
by the specialized row, since determinant is a polynomial
in the matrix entries over \(B\). Hence no denominator
depending on the specialized row has been introduced.

At a place trivial on the component field both terms vanish.
At any other contributing place \(\zeta\), compare in an
integral frame. If \(a_\zeta=\min_l v_\zeta(w_l)\), then
\[
       v_\zeta(w_x)-v_\zeta(S_0)=v_\zeta(w_x)-a_\zeta .
\]
Here the generic zeroth row realizes the minimum by the
Gauss rule. For each horizontal prime \(q\), the pushforward coefficient is
\[
 \sum_{\zeta\mid q} f_{\zeta}
            \bigl(v_{\zeta}(w_x)-a_{\zeta}\bigr).
\]
This is precisely the coefficient of the moving divisor at \(q\).
The specialized polynomial is nonzero. Holding the other cut rows
generic, its factors as a polynomial in the remaining row are
readings at the \(\delta_Z\) points of the complementary cut,
counted with degree. Each such factor is linear in that row. The
degree is therefore \(\delta_Z\) per row, and summing over components
gives \(\delta\).
Adding the row degrees of the component contributions gives
\[
  \sum_Z\delta_Z=\delta .
\]

\proofstep{Line weights and changes of cut coordinates.}
For the horizontal orders of any rational section of a weight
line there are only finitely many contributing \(q\). Here
is a geometric way of seeing this on the present proper
flag model. Normalize the affine space of cut rows in \(L\).
Write \(\pi:Y\to\operatorname{Spec}K[X]\) for this finite
normalization and \(\mathcal F\) for the proper flag compactification.
At a codimension-one point \(\zeta\) of \(Y\), properness extends the
generic flag reading to \(\operatorname{Spec}\mathcal O_{Y,\zeta}\).
The finitely many coordinates in an affine chart about its closed
image extend on an open neighbourhood of \(\zeta\). Thus the rational
flag reading is a morphism \(\theta:Y^\circ\to\mathcal F\) away from
a subset of codimension at least two.

For a rational weight section
\(u\), the coefficient at \(q\) is
\[
 \bigl[\pi_*\operatorname{div}_{Y^\circ}(\theta^*u):q\bigr]
         =\sum_{\zeta\mid q} f_\zeta\,v_\zeta(\theta^*u).
\]
The Weil divisor on the left has finite support because \(\theta^*u\)
is a rational section of an invertible sheaf on the noetherian open
\(Y^\circ\). Each order on the right is exactly the flag-and-theta
frame order at \(\zeta\).
The same reasoning applies to the minimum of finitely many
column orders. Adjoining the independent row zero pulls
back these divisors, without acquiring any further order.

A change of basis of the \(s\) cut rows leaves the generic
point on the flag variety fixed under the induced isomorphism
of incidence fields. Thus its pushed-down divisors are invariant
under \(\mathrm{GL}_s\).
For a pushed-down horizontal divisor \(D\) this invariance is
\[
  g^*D=D\qquad(g\in{\rm GL}_s).
\]
 Each of their finitely many irreducible
components is invariant: the image of the product of the
connected algebraic group and a geometric irreducible component
is irreducible, contains that component, and lies in the
finite union. Descent gives the statement over \(K\).

If \(F_D\) generates the corresponding prime in \(K[X]\),
each \(F_D(gX)\) is another generator.
Since the units of \(K[X]\) are constants, the two prime generators satisfy
\[
  F_D(gX)=c(g)F_D(X),\qquad c(g)\in K^\times\quad(g\in{\rm GL}_s(K)).
\]
 Its ratio to \(F_D(X)\)
is a scalar, and the scalars form an algebraic character of
\(\mathrm{GL}_s\). Write it as \(\det^m\). Restriction to the
diagonal torus gives
\[
 F_D\bigl(\operatorname{diag}(t_1,\dots,t_s)X\bigr)
          =(t_1\cdots t_s)^m F_D(X).
\]
Thus every nonzero term has degree \(m\) in each cut row. This justifies
using one cut-row degree in \eqref{eq:slice-height-disc}.
\label{cor:slice-integrals-1}

\paragraph{Integrability of the weight metrics.}
The logarithm of the absolute value of a nonzero polynomial
in the rows is integrable on the product of spheres. The
phase argument for \eqref{eq:slice-sphere-lower} applies
also to a polynomial with terms of unequal degrees. Its
positive part is bounded and its leading term gives a finite
lower integral. In particular a nonzero polynomial has zero
set of sphere measure zero. If \(r\in L^*\), consider the
characteristic polynomial of multiplication by \(r\) and the
one for \(r^{-1}\). Their coefficients are rational
functions of the cut rows.

For every generic complex
point of this finite cover the two Cauchy root bounds give
\[
          |\log |r(P)||\le A_r(X)
\]
where \(A_r\ge0\) is integrable on the product of spheres:
one can take the sum of two copies of \(\log 2\) and the
positive logarithms of the absolute values of all nonzero
coefficients. These logarithms are dominated by sums of the
absolute logarithms of polynomial numerators and denominators.
The bound holds simultaneously at all sheets over a tuple.

\proofstep{Integrability of the logarithms of norms on the flag chart.}
Use a nonzero component light column \(w_b\) and set
\[
 x=-\log\|w_b(P)\|,\qquad
 y=-\log\|\eta(P)\|/n_1,\qquad z_1=-\log\|t_1(P)\| .
\]
The bounded suprema give constant lower bounds for these
three quantities. Write \(B_L=(\sum_j\|w_j(P)\|^2)^{1/2}\).
The rational ratios \(r_j\), with \(r_b=1\), show that
\(0\le\log(B_L/\|w_b(P)\|)\) is bounded by an integrable
function, by the preceding root bounds applied to the
nonzero ratios. The archimedean inequality just proved
therefore implies \(x\le A(X)+\beta y\) with an integrable
majorant \(A\) common to the sheets of this component.

The first weight identity makes
\[
 h=\frac{t_1^{n_1}\eta^2}{w_b^{2n_1}}
\]
a nonzero rational function on this component. Away from its zeros and
poles, \(\log|h|=2n_1(x-y-z_1/2)\). The preceding bounds for
\(\log^+|h|\) and \(\log^+|h^{-1}|\) give an integrable majorant
for \(\bigl|\log|h|\bigr|\). Hence

\[
  y=x-z_1/2+O(A'(X)),
\]
 where \(A'\) is integrable.
In particular \(y\le x+A''(X)\) since \(z_1\) is bounded
below. As \(\beta\le1/2\), substitution bounds \(x\)
above by an integrable function. This bounds \(y\) above,
then \(z_1\) as well. For any nonzero short section of
weight \(w_F\) on this component use its product with \(t_2\)
divided by \(w_b^{s+4}\). The logarithm of this rational
function and the bound for \(x\) control the sum of the
two negative log norms; both summands are bounded below.

Finally a change to any other rational frame within a
weight adds the logarithm of a nonzero rational function.
The additional term is integrable by the two log bounds for a rational function \(h\ne0\) and its inverse, since
\[
  \big|\log |h|\big|=\log^+|h|+\log^+|h^{-1}|\qquad(0<|h|<\infty).
\]
This proves integrability for all the section and moving
metrics used in the height argument. Measurability of the
sheet sums follows on the finite \'etale open by locally
labeling its sheets and using the continuous weight metrics;
the excluded algebraic set is null. Here the integrable
majorants may depend on the lift. Their sizes never enter
the final inequality, which uses the uniform numerical
prices of the height calculation. \label{cor:slice-integrals-2}

\paragraph{At the definite last lift.}
For the terminal, negative-definite lift the first-minimum term in the analytic height vanishes. The resulting integral can therefore be estimated directly. On the negative-definite space the invariant majorant
is the definite Hermitian form with its sign changed, independent
of the negative subspace; the run's first-minimum threshold is
below its lower lattice bound. At the last index the displayed function therefore vanishes. We have

\[
  -\log B_L\le P_5D_F
\]
 and the local test is simply
\(v(w_b)\le\nu_\lambda\) relative to the uniformizer of
\(K_\lambda\). Its horizontal version gives minimum order
at most zero. Integrality gives the opposite inequality.
Thus \(E_q=0\), \(a_L=\delta\), and
\[
 H_L\le H_{\rm mov}+\delta(P_N+P_5)S_L .
\]
For integrability, \(-\log\|w_b(P)\|\) is bounded below
and the upper bound follows directly from \(-\log B_L\)
and the integrable ratio bound. The calculation with
\(t_2\) then gives integrability at \(w_F\) and at the
complementary weight. The Chow comparison gives \(a_2,a_F\le(s+4)\delta\), and the discriminant contribution gives
\[
 H_F\le(s+4)H_L+\delta P_4 s S_L+s a_2\ell_N .
\]

For example the Chow bound gives
\(H_L/(\delta S_L)\le
(s+2)(P_0+G_c+P_{rw})+P_{rw}+P_N+P_5\).
With \(P_H\) as above, \(s\ge3\), the resulting upper bound
on \(H_F\) fits \((s+4)(s+1)\delta S_L(2G_c+6P_H)\).
Only the discriminant term is used at this lift. These
simplifications give the definite height inequality without
a bound for the auxiliary order. \label{cor:slice-integrals-3}
\end{proof}

\paragraph{Counting the short and boundary sections.}
The orders on the slice control its finite contribution. Its archimedean rank estimate now requires a count of the sections that remain short and of those lost to the boundary. On a connected unfolded complex component put
\(S_i=L_i\) for the \(s\) cutting sections and write
\(w_0=L_0\) for the nonvanishing reference at their
common zeros. The flat-section hypotheses are that these
zeros consist, modulo the uniformizing group, of
\(\delta_k\) distinct framed points, that
\(u_i=S_i/w_0\) give analytic local coordinates there, and
that the chosen short weight provides interpolation of the
values at those points. The interpolation with the required
vanishing of low Fourier coefficients will be used also for
the lower bound on the selected subspace. Both interpolation conditions use the same point count \(\delta=\sum_k\delta_k\) as the heights.

\begin{proposition}[Ranks of the full and short evaluation maps on the slice]
\label{prop:slice-rank-shares}
Assume the simultaneous cut consists on component \(k\) of
\(\delta_k\) distinct framed points at which \(w_0\) is nonzero
and the \(S_i/w_0\) are local coordinates. Assume interpolation
of values at all these points in the short weight with the
prescribed low modes zero. Put \(\delta=\sum_k\delta_k\). For the full rank \(R_j\), the rank deficit \(r_q\) of the rational good space, and the rank \(r_z\) of evaluation on the boundary,
\[
 R_j\le\delta\binom{T+2s+6}{s},\qquad \max(r_q,r_z)\le\frac{22s^2}{T}R_j\quad(s\ge2,\ T\ge11s^2).
\]
 The definite final lift
uses the full space as its good space.
\end{proposition}

\begin{proof}
\proofstep{An upper bound for the entire rank.}
To bound the full rank from above put \(e=T+s+6\).
The jets of order \(\le e\) at the \(\delta_k\) points have
dimension \(\delta_k\binom{e+s}{s}\). A nonzero section
\(v\) in the kernel vanishes to order \(e+1\) there.
We give the maximum argument which excludes it. In the
indefinite chart, use \(b_2,b_1,\xi\) of respective weights
\((s+4)w_L-w_F,\ 2w_L-w_d,\ n w_d\), where
\(n=i g_0^*/2>0\) in this analytic normalization.
All three are nonzero on the component. Put \(S=\sum_i
\|S_i\|^2\), the sum over the \(s\) cutting sections,
and consider
\[
 \Phi=\frac{\|v\|^2\|b_2\|^2\|b_1\|^2
                     \|\xi\|^{2/n}}{S^{e+1/2}} .
\]

The bounded section and the auxiliary factors satisfy the
chart estimates for the first minimum. Near a cut point \(S\gg|u|^2\) and \(v=O(|u|^{e+1})\);
the other numerator factors are bounded there. Thus
\[
 \Phi\ll \frac{|u|^{2(e+1)}}{|u|^{2(e+1/2)}}=|u|.
\] Away from neighborhoods of the cut
orbits the chart estimate gives \(S\gg e^{-2cP}\), and the
numerator is \(O(e^{-4P/g_0^*})\). The hypotheses \(s+6+1/2<T\) and \(Tcg_0^*\le1\) give the strictly positive net decay exponent
\[
  \frac{4}{g_0^*}-2c(e+1/2)>0.
\]
 As the first minimum
goes to zero, \(P\) tends to infinity. Thus every positive
superlevel has compact representatives, away from the cut
points.

On a definite component compactness itself supplies
this step with the bounded auxiliary factors in its chart.
A common nonvanishing locus for the numerator and at
least one cutting section exists by the identity theorem
on the connected domain. In the case \(\delta_k=0\) the
denominator has no common zero and the same argument applies.

A positive maximum contradicts the Levi form. In a local
holomorphic frame the weight of the numerator terms totals
\(e w_L\), leaving a positive half of the strictly positive
light-weight potential in \(\log\Phi\). The log absolute
values of its individual holomorphic factors are
pluriharmonic at a positive maximum.
For a holomorphic factor \(h\) nonzero at the maximum this vanishing reads
\[
  {\cal L}(\log|h|)=0\qquad\text{locally on }\{h\ne0\}.
\]
 The remaining log
sum of squares contributes the pullback of a
Fubini--Study form by the map \([S_1:\cdots:S_s]\).
This maps an \(s\)-dimensional tangent space to one of
dimension \(s-1\).

Choose a nonzero vector in the kernel
of its differential. The Levi form of the log sum is zero
on this vector and that of the remaining potential is
strictly positive. At a maximum their sum must be
nonpositive. This proves injectivity of the jet readings.
Restriction of the global space to all components is
injective, hence summing gives
\[
                      R_j\le\delta\binom{T+2s+6}{s}.
\]

\proofstep{The rank of the complementary sections.}
For the complementary rank take forms \(v_j\) of weight
\(w_F\) interpolating the point delta functions on one
component, with their required low Fourier coefficients
zero. For \(|\alpha|\le T\) the products
\(v_j w_0^{T-|\alpha|}\prod_i S_i^{\alpha_i}\)
belong to the prescribed analytic subspace. This uses
stability of the low-mode vanishing condition under
multiplication by the bounded light forms, whose Fourier
indices are nonnegative. Independence follows by first
taking the least total degree appearing in a relation,
then taking that homogeneous Taylor coefficient at each
point after dividing by \(w_0^T\). Its monomials are
distinct and the value matrix of the \(v_j\) is invertible.
One can extend the families separately by zero to the
other components. The \hyperref[par:shifted-good]{descent of the low-mode space} gives

\[
  \dim H_g\ge\delta\binom{T+s}{s}.
\]
 This identifies
\(r_q=R_j-\dim H_g\) after rationally spanning the cusp
lattice. At the definite final lift the comparison uses
all sections.

\proofstep{Short restrictions and the boundary rank.}
Likewise the \(|\alpha|<T-s\) products from the short
restrictions are independent modulo the span of all
boundary products. Apply the least-degree argument to
a relation of degree less than \(T-s\) against sections
vanishing to that order. There are \(R\delta\) of these
independent classes. The boundary products can be tested
by their rank on the generic coefficient functions: after
extension of scalars their minors vanish precisely when
the corresponding generic complex minors vanish. Taking
a basis out of their short generating family gives

\[
  R\delta+r_z\le R_j.
\]
 This calculation also deals with a
zero-dimensional vector space of readings.

For \(s\ge2\) and \(T\ge11s^2\), put \(B=\binom{T+2s+6}{s}\).
The fraction \(R/B\) is the product of the \(s\) factors
\(1-(2s+7)/(T+2s+6-i)\) for \(0\le i<s\).
A product of \(1-x_i\), \(0\le x_i\le1\), is at least
\(1-\sum x_i\). Consequently
\[
 0\le 1-R/B\le s(2s+7)/T\le22s^2/T .
\]
Since \(R_j/B\le\delta\), both deficient ranks are at most
\(22s^2R_j/T\). This implies the weaker pre-final bound
\(r_z\le r_q+22s^2R_j/T\) used in the slope sum. It uses
no assumption that different primes admit independent
choices of the short family. \label{cor:slice-shares-1}
\end{proof}

\paragraph{Choosing the cut on the unfolded domain.}
\label{par:slice-interpolation}
\hypertarget{source-slice-interpolation}{}
We derive the flat-section facts used in the rank estimates.
We use the algebraic realization of the listed forms in the
complex uniformization, including their common weight lines.
\begin{lemma}[One transverse cut on the unfolded domain]
\label{lem:transverse-selected-cut}
A common choice of rows gives the reduced finite cut on all the
components. Writing \(w_0\) for row zero and \(S_1,\ldots,S_s\) for the cut rows, at every point \(y\) of the cut one has
\[
 w_0(y)\ne0,\qquad d(S_1/w_0)\wedge\cdots\wedge d(S_s/w_0)|_y\ne0,
\]
 and away from the selected
zero neighborhoods their squared norms control the squared norm of
the light vector as used below.
\end{lemma}

\begin{proof}
\proofstep{The generic incidence map.}
Take all the cut-row coefficients algebraically independent
over the coefficient field. Their exceptional sets below
are null for Lebesgue measure. The Zariski closure of the
projective light image of one component has dimension \(s\).
Thus \(s+1\) generic linear forms have no common zero on
that closure. On its compact complex projective points the
sum of their squared absolute values in a unit representative
is bounded below. With the notation \(S_i\) of the rank
calculation this says
\[
       \sum_l\|w_l\|^2\le C\sum_{i=0}^s\|S_i\|^2 .
\]
Let \(X\) be the smooth open flag chart, \(A\) the affine
space of cut coefficients and \(L\) the light line. The incidence
space is the zero scheme of \((S_1,\ldots,S_s)\) on \(X\times A\).
At a common zero choose a light column \(w_b\ne0\).
Differentiation in the \(s\) variables \(X_{ib}\) gives
\[
 (\dot X_{1b},\ldots,\dot X_{sb})
               \longmapsto(\dot X_{1b}w_b,\ldots,\dot X_{sb}w_b)
                        \in L_x^{\oplus s}.
\]
This is onto, so the incidence scheme is smooth there.

The generic slice computation gives a generically finite
map to \(A\), so generic smoothness gives reduced
zero-dimensional fibres. The complement in the
incidence variety of the dense open on which its map is
finite \'etale has proper image after restricting the base
to a suitable dense open. Indeed its dimension is smaller
than the dimension of the coefficient space. Equivalently
one may first compactify the graph and remove its boundary
and the branch divisor. Thus the complex points at a generic
cut are the \(\delta_Z\) points of the incidence field and
give the required isolated orbits under uniformization.
Over a complex point outside the branch and boundary loci, a finite separable component algebra of degree \(\delta_Z\) has fibre algebra
\[
  \underbrace{\mathbb C\times\cdots\times\mathbb C}_{\delta_Z\ \mathrm{factors}} .
\]
With \(w_0=S_0\), their transverse coordinates are
\(S_i/w_0\).


\proofstep{The lower bound off the zero neighborhoods.}
There is also an estimate off neighborhoods of these
orbits. Here are the details before the last lift, where
the domain need not be compact. In the analytic weight
conventions let 
\[
  n=i_\eta g_0^*/2
\]
 and let \(b_1\)
be the bounded nonzero complementary section of weight
\(2w_L-w_d\). The section \(a=\eta b_1^n\) has weight
\(2n w_L\). On each component it is nonzero, and the
first-minimum bounds give, for \(W=\sum_l\|w_l\|^2\),
\[
                  \|a\|^2\le C' W^{2n+1}.
\]
In fact \(W\gg e^{-2cP}\), \(\|\eta\|^2\ll e^{-2i_\eta P}\),
and \((2n+1)c\le i_\eta\) by the chosen margin
\(Tcg_0^*\le1\), \(T\ge2\), \(P\ge0\).

Consider the projective map with coordinates \(a\) and
all the degree-\(2n\) monomials in the \(w_l\).
The closure of its image has dimension at most \(s\)
and no component lies entirely in \(a=0\). If \(Y\) denotes this image closure, the hyperplane therefore cuts its dimension by at least one:
\[
  \dim\big(Y\cap\{a=0\}\big)\le\dim Y-1\le s-1.
\]
In that hyperplane the remaining coordinates lie in the
Veronese image of light directions. Pulling back by the
Veronese isomorphism onto its image gives a closed set
of light directions of dimension at most \(s-1\). It
misses the generic linear space cut by rows \(1,\ldots,s\).
If along a sequence \(W\to0\) and

\[
  \sum_{i=1}^s |S_i|^2/|S_0|^2\to0,
\]
 pass to a limiting
light direction. The displayed estimate shows that the
normalized coordinate \(|a|/(\sum_l |w_l|^2)^n\) tends
to zero. In this expression all factors from the common
weight metric cancel. The limiting direction lies both
in the forbidden set and in the generic linear space.
This contradiction bounds \(W\) away from zero when the
ratio is sufficiently small.

Each light column is cuspidal, with normalized norm tending
to zero when the first minimum tends to zero.
For the finite light family the individual decay has the simultaneous consequence
\[
  \lambda_1(y_n)\longrightarrow0\quad\Longrightarrow\quad\max_l\|w_l(y_n)\|\longrightarrow0 .
\]
 This uses the
cuspidal decay estimate for these bounded sections. Hence
in the small-ratio region there are compact representatives
by reduction theory. Away from prescribed invariant open
neighborhoods of the cut points, the continuous ratio
\(W/\sum_{i=1}^s\|S_i\|^2\) is bounded on these representatives.
Where the ratio to \(|S_0|^2\) is large, the first displayed
inequality already bounds it. At the definite lift take
compact representatives from the start. This proves the
off-neighborhood bound used in the maximum argument.
Only existence of these chart-dependent compactness constants
enters the rank, allowing one choice of rows to work on the
finitely many components. \label{cor:slice-interpolation-1}

\end{proof}

\paragraph{Curvature at the selected short weight.}
\hypertarget{source-slice-short-analytic}{}

Control on the chosen cut must extend to a neighborhood. The curvature of the short weight determines the local norm comparison used in this extension.

Write \(q\) for the length of the first signature block,
\(p=r-q\), and \(s=fr(r-1)/2\). At an orthonormal flag,
use a lower unipotent coordinate matrix \(I+Z\). The
Gram matrix in each embedding is \((I+Z)^*J(I+Z)\), where
\(J\) is diagonal with \(q\) signs of one kind followed by
\(p\) of the other. The potential of a weight \(v\) is
\[
 \varphi_v=\sum_\tau\sum_{i=0}^{r-1}
               v_{\tau,i}\log\left|
                 \frac{\Delta_{\tau,i+1}}{\Delta_{\tau,i}}\right|,
\]
where \(\Delta_{\tau,k}\) is the leading minor of
order \(k\) and \(\Delta_{\tau,0}=1\).

In the quadratic expansion of this minor's logarithm,
an entry \(Z_{ji}\), \(i<j\), contributes only when
\(i<k\le j\). The first \(k\) rows give a unit lower
triangular factor whose determinant is constant. The
remaining rows contribute the quadratic trace with sign
\(J_j/J_i\). For a tangent connecting positions \(i<j\), the coefficient of the Levi form after summing the weight differences is
\[
  \begin{cases}v_i-v_j&i,j\text{ in the same block},\\v_j-v_i&i,j\text{ in opposite blocks}.\end{cases}
\]
 At another
point take its orthonormal frame by triangular
orthogonalization. The induced change on tangents is
invertible and the logarithm of the change in the weight
character is pluriharmonic. This transports the diagonal
calculation to that point.

In this chart the invariant density is
\(D=\prod_{\tau,k<r}|\Delta_{\tau,k}|^{-2}\).
The full Gram determinant is constant on the unipotent
cell, so

\[
  \mathcal L(\log D)=\mathcal L(\varphi_{k_{\rm vol}})
\]
with \((k_{\rm vol})_{\tau,i}=2i-(r-1)\). Write \(d\)
for the sign weight that is \(-1\) on the first block
and \(+1\) on the second. The radius potential \(R\)
is \(\varphi_d/2\) plus a pluriharmonic function, since
the only nonzero interior weight difference of \(d\)
occurs at \(q\). In particular \(\mathcal L(R)\) has unit
coefficients in the cross-block directions and zero
coefficients within blocks in these coordinates.
\label{cor:slice-short-analytic-1}

\begin{lemma}[Curvature for the short-weight interpolation]
\label{lem:short-weight-curvature}
In the flat cut chart with the given signatures, the short weight
has the positive Levi directions used in the local interpolation
estimate. For the selected cutoff on each fixed zero neighborhood, put \(G=w_F-k_{\rm vol}-s w_L\). With \(i,j\) in opposite blocks in the indicated order, the signature margin is
\[
 G_j-G_i\ge2((w_L)_j-(w_L)_i)\ge19T/5.
\]

\end{lemma}

\begin{proof}
\proofstep{Positivity at the chosen short weight.}
Here is the weight comparison used on both runs.
Put \(G=w_F-k_{\rm vol}-s w_L\). Within either block the
selected weight inequalities give
\(G_i-G_{i+1}\ge2((w_L)_i-(w_L)_{i+1})>0\).
Across the blocks they give
\((w_L)_j-(w_L)_i\ge19T/10\) and
\(G_{r-1}-G_0\ge2((w_L)_{r-1}-(w_L)_0)\).

Summing the adjacent inequalities from the outer
indices to \(i,j\) gives
\[
 G_j-G_i\ge2((w_L)_j-(w_L)_i)\ge19T/5 .
\]
These are the bounds for the actual shifted short weight.
In taking consecutive differences of \(w_F-s w_L\),
the volume weight contributes
\((k_{\rm vol})_{i+1}-(k_{\rm vol})_i=2\). The light weight and
\(2w_L-d\) are strictly positive in the Levi formula.
For \(0\le u\le T\), the weight \(G-u d-w_L/2\)
has nonnegative Levi form. Its cross-block gap is at
least \(3(19T/10)/2-2T>0\), and the within-block
difference of \(d\) is zero.

\proofstep{The weighted curvature comparison.}
The invariant potential \(\psi_0\) is bounded above and
has \(\mathcal L(\psi_0)\ge-(r/2)\mathcal L(\varphi_d)\).
Its cusp coefficient and the selected cutoff are
\[
 a_1=\frac{1}{2\cdot5^4|\operatorname{disc}E|^{1/f}},
 \qquad
 c_1=2(\operatorname{rd}E)^7 T c\,(fq)^2 q .
\]
Here \(q\ge1\) on the indefinite lift and \(c>0\)
is the light first-minimum coefficient. To check the allowance, put \(d_E^{\rm root}=\operatorname{rd}E\), \(\mathcal P=C_2z^{A_2}\), and write \(Q\) for the denominator of the light coefficient.
The numerical inequalities on this run are
\[
 c=(d_E^{\rm root})^7\mathcal P D_F^2/Q,\qquad
 (d_E^{\rm root})^{28}\mathcal P^3D_F^2\le Q,\qquad
 1250\,rfm^3\le\mathcal P^2,\qquad
 2fm\le d_E^{\rm root}.
\]
Thus \(c\mathcal P^2(d_E^{\rm root})^{21}\le1\).

Since \(q\le m\) and \(d_E^{\rm root}\ge1\),
\[
 \frac{r c_1}{a_1}
    =Tc\, (1250\,rfq^3)\,(2f(d_E^{\rm root})^9)
    \le Tc\mathcal P^2(d_E^{\rm root})^{21}\le T .
\]
With \(r\le T\) this gives \(r(1+c_1/a_1)\le2T\).
Put \(\psi=(1+c_1/a_1)\psi_0\). The preceding
comparison with \(u=r(1+c_1/a_1)/2\) gives
\[
 \mathcal L(\varphi_F-\log D-s\varphi_L+\psi)
                    \ge \tfrac12\mathcal L(\varphi_L).
\]
The cross-block estimate also gives the independent
lower bound \(m_{\rm int}\mathcal L(R)\) on this expression,
with \(m_{\rm int}=19T/5-2T\).

For division the potential
omits the \(\psi\) term and the subtraction of
\(s\varphi_L\). The margin there is
\(m_{\rm div}=19T/5+s(19T/10)\).
Set \(\beta=5fr^2\) and \(u_0=f\min(p,q)\). Then
\[
 4fr^2<\beta<m_{\rm int},
 \qquad \beta^2 u_0+\beta < m_{\rm div}.
\]
Indeed \(f\ge1\), \(r\ge8\), \(2u_0\le fr\), and
\(T\ge11s^2\). These give
\(\beta\le20s^2/49\) and
\(\beta^2 u_0\le20s^3\) by \(5r^2\le(r-1)^3\).
The additional \(19T/5\) and \(19sT/10\) cover the
two terms in the division margin. For a definite
signature \(R=0\) and \(\psi_0=0\); the within-block
calculation supplies the nonnegative Levi form and the
same inequalities with arbitrary margins in the zero
directions of \(\mathcal L(R)\).
\label{cor:slice-short-analytic-2}

\end{proof}

\paragraph{The Poincar\'e estimate on coordinate balls.}

The curvature estimate is local. To obtain the norm bound after Poincar\'e summation one must compare the coordinate balls about the translated frames.

Let \(x_*\) be the orthonormal frame of the seed cell.
At an orthonormal frame \(x\) write the relative
isometry \(x_*^{-1}x\) at each embedding in signature
order \(p,q\), with bottom-right block \(D_0\).
Then \(R(x)=2\sum_\tau\log |\det D_{0,\tau}|\).
For a general admissible frame, subtract the logarithm
of its Gram minor of order \(q\) in each embedding.
This formula is invariant under the upper triangular
change of frame.

\begin{lemma}[Local integrals and the Poincaré estimate]
\label{lem:short-poincare-summation}
For the coordinate balls and weights below, the estimate for the
Poincaré sum follows from the local integral of the lifted section
and the bound on the overlap of the translated balls. In particular, for a seed of finite weighted integral \(I=\int|P|^2W_u e^{\beta R}D\), its integral \(I_\gamma\) on a translated ball, before the factor \(e^{\beta R}D\), satisfies
\[
 I_\gamma\le C I e^{-\beta R(\gamma x)}.
\]
The constants may depend on the fixed cut chart.
\end{lemma}

\begin{proof}
\proofstep{Bounding the overlap of translated balls.}
Use discreteness of the group in the product of the
archimedean matrix groups. Two distinct relative
isometries \(g,h\) from \(\Gamma x\) have an entry of
\(hg^{-1}-I\) bounded below in absolute value by a
positive constant depending on the group and \(x_*\).
This is independent of \(x\) because the factor \(x\)
cancels in the quotient. If all entries of both
isometries have absolute value at most \(A\ge1\),
the entries of \(g^{-1}=Jg^*J\) are at most \(A\).
Since 
\[
  hg^{-1}-I=(h-g)g^{-1},
\]
 at least one entry of
\(h-g\) is separated by a constant times \(A^{-1}\).
Partitioning the box of real and imaginary matrix
entries into cubes of that sidelength bounds its
number of isometries by \(C A^{4fr^2}\).

For a single signature matrix \(g=\left(\begin{smallmatrix}
A_0&B_0\\ C_0&D_0\end{smallmatrix}\right)\) write \(s_i\) for the singular values of \(B_0\). The identity
\(D_0^*D_0=I+B_0^*B_0\) gives
\[
 |\det D_0|^2=\prod_i(1+s_i^2)\ge1+\sum_i s_i^2
 =1+\|B_0\|_{\mathrm{HS}}^2.
\]
The other block identities bound the square of
each entry of \(g\) by \((r+1)|\det D_0|^2\).
Thus a dyadic shell of maximal entry \(A\) contributes
\(O(A^{4fr^2-2v})\) to the sum of \(e^{-vR}\).
Consequently
\[
          \sup_{x\ {\rm orthonormal}}\
               \sum_{\gamma\in\Gamma}e^{-vR(\gamma x)}<\infty
                 \quad(v>2fr^2).
\]

The same bound controls local integrals. On a small
fixed unipotent coordinate ball about an orthonormal
centre \(y\), the Gram minors are close to their identity
values, uniformly in \(y\). In signature order, right
multiplication of the relative isometry by the nearby
coordinate frame \(N\) changes its bottom-right
determinant by the factor
\(\det(N_{22}+D_0^{-1}C_0N_{12})\).
The relation \(D_0D_0^*=I+C_0C_0^*\) implies that \(D_0\) is
invertible and gives the operator inequality
\[
 (D_0^{-1}C_0)(D_0^{-1}C_0)^*
 =I-D_0^{-1}(D_0^{-1})^*\le I.
\]
Hence \(\|D_0^{-1}C_0\|_{\rm op}\le1\). On a ball fixed in
these coordinates the determinant of the factor is at
least \(1/2\), independently of the relative isometry.
The Gram-minor formula gives \(R(yN)\ge R(y)-3f\log2\).
The density \(D\) is bounded below there.

\proofstep{Weights in the local integral estimate.}
Let \(W_u\) be a squared-reading weight transforming
inversely to the squared weight character and invariant
under \(\Gamma\). Suppose the seed has finite integral
\(I=\int |P|^2 W_u e^{\beta R}D\) on the entire flag
cell. Write \(I_\gamma\) for the integral without
\(e^{\beta R}D\) in the coordinates of the translated
ball. The invariant measure and the uniform bounds
give 
\[
  I_\gamma\le C I e^{-\beta R(\gamma x)}.
\]
Here the change to the seed cell preserves the
integral even when the small ball crosses its chart
divisor, which has measure zero. Summing square roots
is allowed since \(\beta/2>2fr^2\).
The holomorphic submean inequality gives the bound on
the normalized Poincar\'e sum by a constant times
\[
       \sup_{z\ \text{in the ball}}
              \left(e^{-\varphi_v(z)}/W_u(z)\right)^{1/2}.
\]

On a smaller neighborhood of any centre the same
submean inequality supplies a common summable
majorant for all evaluation points, provided
\(W_u\ge c_*e^{-\varphi_v}\) for some \(c_*>0\).
Triangular orthogonalization at that centre changes
this majorant by its nonzero weight character. This
proves locally normal convergence on the frame domain.
All constants in the ball bound can depend on the
lift, weight, cut and seed. They are uniform in the
orthonormal centre and its orbit, which is the
uniformity needed for the Fourier test.
\label{cor:slice-short-analytic-3}

\end{proof}

\paragraph{Extracting Fourier coefficients along the cusp flow.}
On an indefinite component the potential is the
logarithm of the sum of normalized squared theta
values over ordered \(r\)-tuples of points of order
dividing 25 whose multiples by 25 generate the
period lattice modulo 5 over \(\mathcal O_E\).
Each factor is divided by the theta squared
\(L^2\) norm. The theta estimate used here assumes an
integral CM element \(\omega\) with
\(\omega+\bar\omega=1\), the even trace comparison for
multiplication by \(\omega-\bar\omega\), an isotropic
lattice vector, \(\operatorname{rd}E\ge\max(f^2,16)\),
and a rank parameter \(h\) with \(h^4\le f\) and
\(\min(p,q)\le h\). In the trace comparison the difference
of the two trace pairings on any two integral period
vectors is in \(2\mathbb Z\). The chosen host supplies
\(\omega\) and this comparison.

On lift \(j\) the
signature has 
\[
  q=m-j>0,
\]
 so one can use \(h=m\).
The invariant smooth potential satisfies the lower Levi
bound used above and, for \(0<\varepsilon<1/2\),
\[
 \psi_0(x)\le C_\varepsilon
    +A_\varepsilon\log(1+\lambda_1(x)^{-1})
    -8\pi(1-2\varepsilon)a_1\min_\tau t_\tau(x,\xi)
\]
whenever all \(t_\tau(x,\xi)\ge1\) for an isotropic
lattice vector \(\xi\). The bound with just the constant
and logarithm also holds on the entire domain.
Here \(\lambda_1^2\) is the minimum of the sum of the
positive majorants on nonzero lattice vectors, and
\(t_\tau(x,\xi)\) is the reciprocal of the corresponding
individual weighted majorant. In the thin region a
vector with sum of the majorants at most
\(2\lambda_1^2\) is isotropic by integrality and all its
times are at least \((2\lambda_1^2)^{-1}\). Using it in
the displayed inequality bounds \(\psi_0\) above even
as \(\lambda_1\) tends to zero.
\label{cor:slice-short-analytic-4}

\begin{proposition}[Fourier vanishing and the cusp decay estimate]
\label{prop:slice-cusp-gap-decay}
At the prescribed short weight the vanishing test on each complete rank-one period forces its low Fourier--Jacobi coefficients \(a_\nu\) to satisfy
\[
 a_\nu=0\qquad\left(\sum_\tau\nu_\tau\le c_1\right).
\]
 Here \(\tau\) ranges over the chosen CM half, so
\(\sum_\tau\nu_\tau=\tfrac12\operatorname{Tr}_{E/\mathbb Q}(\nu)\).
The resulting bound on a section with those terms zero
is the uniform first-minimum decay used in the division norm.
The cross-block gap gives the stated limiting constant-coefficient
test also for \(w_F+k w_L\) with \(k\ge0\).
\end{proposition}

\begin{proof}
\proofstep{The local estimate for coefficient tests.}
For the coefficient test, a small coordinate ball
about any orthonormal flag changes each positive
majorant by factors in \([1-\delta,1+\delta]\), for
any fixed \(\delta>0\) after decreasing the radius.
Indeed the majorant matrices in that ball are given by
orthogonalizing the Gram matrix of the unipotent
coordinate frame. In orthonormal coordinates these
matrices vary continuously from the fixed positive
form at the origin. The comparison holds simultaneously
for all vectors, so it gives the same comparison for
the squared minima and a reciprocal comparison for
each time. Put \(a=(1+c_1/a_1)a_1=a_1+c_1\).

For every \(a'<a\) with \(a'>0\), choose the radius
and \(\varepsilon\) so that
\((1-2\varepsilon)a/(1+\delta)>a'\).
The normalized ball bound for the Poincar\'e sum then
gives
\[
 \|P_{\rm aut}(x)\|_v
     \le C(1+\lambda_1(x)^{-1})^A
               e^{-4\pi a'\min_\tau t_\tau(x,\xi)}
\]
when all times at the centre are at least 2. Triangular
normalization carries it to general frames with the
same majorants. The exponent \(4\pi\) comes from the
square root of the metric bound.

\proofstep{Decay along the integral cusp flow.}
On the holomorphic flow of \(\xi\), translate each
complex parameter to \(u_\tau+iY\), with the real
parameters \(u_\tau\) in a fixed period parallelepiped.
Each time becomes \(t_\tau(x,\xi)+Y\). This follows
from the rank-one isotropic shear. The Gram matrix at
time \(Y\) is a rank-one Hermitian update of the
original Gram matrix, since the quadratic term of the
shear is a multiple of the norm of \(\xi\) and vanishes.
Each leading minor divided by its value at zero is
\(1+Y b_{\tau,k}\) with \(b_{\tau,k}\ge0\).
Writing \(\Delta_{\tau,k}(Y)\) for the moving leading minor, the rank-one calculation gives
\[
  \frac{\Delta_{\tau,k}(Y)}{\Delta_{\tau,k}(0)}=1+Yb_{\tau,k}\ge1\qquad(Y\ge0).
\]

The signs of the minors persist for all \(Y\ge0\).
This gives polynomial upper and inverse bounds for
the diagonal factors of the weight metric. The entries
of the frame and of its inverse grow at most polynomially
in \(Y\) for bounded \(u\), since the shear is linear
in the flow parameter and its rank-one term squares
to zero. Orthogonalization using these minors shows
that the inverse of the positive majorant matrix grows
at most polynomially as well. Hence
\(1+\lambda_1^{-1}\) and the change from the normalized
norm to the holomorphic reading both have polynomial
bounds in \(1+Y\) on this slice.
In particular along this slice there are constants \(C,N\), depending on the initial data, for which
\[
  1+\lambda_1(Y)^{-1}\le C(1+Y)^N .
\]
 These bounds may depend
on the initial flag and on \(\xi\).

The holomorphic reading has character frequency
\(2\nu\) on the real period lattice. By definition the
coefficient is zero when \(2\nu\) is outside the dual
period lattice. For an actual character, shifting the torus
integral to height \(Y\), by holomorphy and periodicity,
gives
\[
 |F_{\xi,\nu}(x)|e^{-4\pi Y\sum_\tau\nu_\tau}
     \le \sup_u |P_{\rm aut}(\operatorname{flow}_\xi(u+iY)x)|.
\]
If \(\sum_\tau\nu_\tau<a\), take \(a'\) strictly between
\(\max(0,\sum_\tau\nu_\tau)\) and \(a\). The right hand
side is bounded by a polynomial times
\(e^{-4\pi a'Y}\). Letting \(Y\) grow gives the required
vanishing. The strict inequality at \(a=a_1+c_1\) kills
the mode of total frequency exactly \(c_1\).
\label{cor:slice-short-analytic-5}

\paragraph{Uniform decay of a section with zero low modes.}
\label{par:slice-cusp-decay}

The extracted low Fourier coefficients describe the part that must be removed. Once those coefficients vanish, the cusp flow bounds the remainder uniformly.

Let \(P\) be a bounded normalized section whose coefficients
of total frequency at most \(c_1\) vanish on the rank-one
periods. The moderate-growth coefficient estimate, using
one quarter of its exponential rate as slack, gives
\[
 \|P_{\xi,\nu}(x)\|_\kappa
 \le C \exp\!\left(-3\pi\sum_\tau\nu_\tau
                                      t_\tau(x,\xi)\right)
\]
when the majorants of the integral isotropic vector \(\xi\)
are at most one. The constant is independent of the vector,
the frequency and the flag \(x\) on the specified quotient.
The support has \(\nu_\tau\ge0\), and the nonzero terms are
admissible frequencies. The frequency-sum bound on the
period lattices gives a uniform finite majorant for the sum
of \(e^{-(3\pi/2)\sum_\tau\nu_\tau}\) over their nonnegative
frequencies. These coefficient and summability bounds apply
to the bounded weight-\(\kappa\) sections here.

\proofstep{Using the first minimum to choose the flow.}
For \(\lambda_1(x)^2\) sufficiently small, choose an integral
vector with sum of its weighted majorants

\[
  S<3\lambda_1(x)^2/2
\]
 and \(S<1\).
At this threshold the vector is isotropic, by integrality
of its Hermitian norm and the bounds under all embeddings.
Its times are at least \(1/S\). Each surviving \(\nu\) has
\(\sum_\tau\nu_\tau>c_1\). Half of the exponential rate
therefore bounds the term by \(C e^{-\pi c_1/\lambda_1^2}\),
and the other half leaves the summable frequency majorant
above. The period expansion and the triangle inequality
give \[
 \|P(x)\|_\kappa\ll e^{-\pi c_1/\lambda_1(x)^2}.
\]
On the rest of the quotient the positive lower threshold
for \(\lambda_1(x)\) and boundedness give the same estimate
with a larger constant. The squared norm bound is uniform
in \(x\) at the given lift. Its constant can depend on the
lift and section since it is used for integrability.
\label{cor:slice-short-analytic-6}

\proofstep{The norm needed by division.}
For the division coefficients only the zero mode must
vanish. A bounded holomorphic section of a weight \(v\)
with nonincreasing weights inside each block and

\[
  v_{\tau,r-1}-v_{\tau,0}\ge1
\]
 at one embedding \(\tau\)
is already cuspidal. Shift the imaginary flow at this
embedding alone. Let \(\rho_k\) be the ratio of the
leading Gram minor at the new time to its original
value. The isotropic rank-one update has
\(\rho_0=\rho_r=1\) and \(1\le \rho_k\le\rho_q\):
the signed squared lengths in the isotropic vector's
orthogonal-frame coordinates accumulate to an extreme
at the cut between signs. Also 
\[
  \rho_q=1+Y/t_\tau.
\]
Summation by parts yields
\[
 \sum_{i=0}^{r-1} v_{\tau,i}
       (\log\rho_{i+1}-\log\rho_i)
     \le (v_{\tau,0}-v_{\tau,r-1})\log\rho_q .
\]
The coefficients of \(\log\rho_k\) for \(k\ne q\) are
nonnegative and the term at \(q\) keeps its exact value.

The metric conversion now bounds the holomorphic
reading on the translated period torus by a constant
times \((1+Y/t_\tau)^{-1/2}\).
The other embeddings retain their metric factors.
The torus integral for the zero coefficient tends
to zero. These weight conditions hold for \(w_F+k w_L\)
with \(k\ge0\). In a block the integral light gap is
at least one, so the two light gaps in \(G\) absorb
the successive change of two in \(k_{\rm vol}\). The
outer cross-block gap is at least one by the bounds
on \(G\), \(w_L\) and the increasing \(k_{\rm vol}\).
\label{cor:slice-short-analytic-7}
\end{proof}


\paragraph{Interpolation with a Fourier cutoff.}
The short weight \(w_F\) interpolates arbitrary values at
the cut with all modes through the selected cutoff zero. Write \(\varphi_F,\varphi_L\) for the weight
potentials, so the squared norm is the squared modulus of the holomorphic coefficient times \(e^{-\varphi}\). Let \(D\) be the invariant measure density in the open flag cell and \(R\) its radius potential.

Take the invariant
\(\psi=(1+c_1/a_1)\psi_0\), where \(a_1>0\) is the decay
coefficient of \(\psi_0\) on the cusp flows, \(c_1\ge0\)
the required Fourier cutoff. Use \(m=m_{\rm int}\), \(\beta=5fr^2\) and the two
curvature bounds just proved.
\[
 \mathcal L(\varphi_F-\log D-s\varphi_L+\psi)
        \ \ge\ m\mathcal L(R),\qquad
 \mathcal L(\varphi_F-\log D-s\varphi_L+\psi)
        \ \ge\ 2\epsilon_0\mathcal L(\varphi_L)
\]
for some \(\epsilon_0>0\), on the centred cells.
The light form has positive curvature, \(R\) is plurisubharmonic,
and \(\psi\) is smooth and bounded above. Subtract \(\beta R\)
from the expression in parentheses to obtain a potential
\(V\). Its Levi form is at least
\((1-\beta/m)2\epsilon_0\mathcal L(\varphi_L)\). In particular
it is strictly positive on a small closed coordinate ball
at the chosen cut point.

\begin{proposition}[Prescribing cut values with the low modes zero]
\label{prop:cut-interpolation-with-cusp-gap}
At the specified short weight on a nondefinite lift, let \(S_0\) be the space of cusp forms whose Fourier--Jacobi modes through the selected cutoff vanish. If \(Z\) is the common cut and \(\mathcal L_y\) is its value line at \(y\), then evaluation is surjective
\[
 \operatorname{ev}_Z:S_0\twoheadrightarrow\bigoplus_{y\in Z}\mathcal L_y,\qquad f\longmapsto(f(y))_{y\in Z}.
\]
 Interpolation
is at this particular cut and weight. The definite case uses the
full cusp space.
\end{proposition}

\begin{proof}
\proofstep{A local interpolating section with the frequency cutoff.}
Take a smooth cutoff \(\chi\) equal to one near that point
and zero outside the ball, with its derivative supported
in an annulus free of other zeros. In the matrix coordinates
of the big cell, the function \(R(z)+\log(1+\|z\|^2)\)
is a plurisubharmonic exhaustion. In the matrix coordinates of the CM factors the exhaustion is
\[
  R=-\sum_\tau\log\det(I-Z_\tau^*Z_\tau).
\]
The summands are plurisubharmonic, by the Levi formula for
these determinants and for the log of a squared norm of
holomorphic coordinates. Each determinant lies in \((0,1]\).
A sublevel therefore bounds the coordinates and keeps the
determinants away from zero. Its closure stays in the cell
because the limiting semidefinite matrices are then positive
definite.

Removing the zero hypersurface of one nonzero
cutting section preserves the exhaustion after adding the
squared absolute value of its holomorphic reciprocal.
\label{cor:slice-interpolation-2}
There apply the
weighted \(L^2\) estimate for \(\bar\partial\) with weight
\(V+s\log\sum_i|S_i|^2\) and right hand side
\(\bar\partial\chi\).
On the complement of the removed divisor the correction satisfies
\[
  \bar\partial u=\bar\partial\chi,\qquad \bar\partial(\chi-u)=0 .
\]
 The log sum of squares has nonnegative
Levi form. On the supporting annulus the inverse Levi cost
is bounded, since \(V\) is strictly positive in curvature
there and the sum of squares is bounded away from zero.
Equivalently, pair the datum with a test form, use the
weighted basic estimate for its Levi norm and then apply
Hilbert space duality to the adjoint equation. One obtains
a solution \(u\) with
\[
           \int |u|^2 e^{-V}
                 \big(\sum_i |S_i|^2\big)^{-s}<\infty .
\]

By exhaustion this proof also applies to a weight whose
Levi form has zero directions off the support. The
holomorphic \(F=\chi-u\) extends across the removed
hypersurface by local square integrability. The singular
factor still controls \(F-\chi\) almost everywhere after
extension. At a cut point the simple coordinates make
this factor comparable with \(|z|^{-2s}\), which has
divergent integral in dimension \(s\). Continuity forces
\(F\) to equal the cutoff there. Thus it reads one at the
chosen point and zero at other zeros in the cell.
\cite{hormander-slice}\label{cor:slice-interpolation-3}

\proofstep{The estimate needed for Poincaré summation.}
For passage to the automorphic section one uses the integral
on the whole cell, including the factor \(e^{\beta R}\) in
\(e^{-V}\). In an adapted frame the invariant density for
that integral is
\[
 D e^{-(\varphi_F+\psi)} e^{\beta R}
          \big(\sum_i\|S_i\|^2\big)^{-s}.
\]
All \(S_i\) are bounded in the normalized light norm. Boundedness of the \(S_i\) gives a positive lower bound for the inverse power of their squared norm sum. The weighted estimate for \(u\), compact support of \(\chi\) and \(F=\chi-u\) then give the bound needed for summation
\[
  \int_{\rm cell}|F|^2D e^{-(\varphi_F+\psi)}e^{\beta R}<\infty.
\]
 The frame extension uses LU decomposition
on the big cell and the weight character on its triangular
factor. Local square integrability in the invariant metric
on compact flag neighborhoods gives removable singularities
for this extension across the missing chart divisor.

Near a
zero outside the original chart the same change of coordinates
applies to the singular integral, with the cutoff extended
by zero there: its compact support in the original chart
is disjoint from a sufficiently small such neighborhood.
The smooth local densities and frame factors are bounded
above and below on the smaller neighborhoods. A nonzero limiting value would give, up to a positive constant, the lower bound
\[
  \int_0^\delta r^{2s-1}r^{-2s}\,dr=\int_0^\delta\frac{dr}{r}=+\infty.
\]
Local integrability therefore forces the value to vanish. This gives a Borel-equivariant
holomorphic seed, with the specified values on the translates
of all cut orbits. The good level acts without point stabilizers
on the flags in this comparison.

\proofstep{Returning the sum to the arithmetic cusp space.}
The Poincar\'e sum with inverse character factors makes this
seed automorphic. The estimate on coordinate balls gives locally normal
convergence and bounds the sum's normalized norm by a
constant times the supremum of \(e^{\psi/2}\) on balls of
any sufficiently small fixed radius. The cusp-flow test
then kills every frequency below \(a_1+c_1\), including
all those at most \(c_1\).
Thus a frequency of total size at most \(c_1\) satisfies
\[
  \sum_\tau\nu_\tau\le c_1<a_1+c_1.
\]
At the chosen representative just one summand is nonzero,
with its nonzero weight character. At all other representatives
the sum vanishes. Rescaling and adding these sections now
interpolates arbitrary values with the stated Fourier
vanishing. The interpolation coefficients and convergence constants
are used here only for existence and rank.

The subsequent \hyperref[par:shifted-good]{trace descent} applies
this vanishing to the coefficient tests at \(\kappa\). It uses the
\hyperref[par:shifted-theta]{full cusp readings} at the selected
level and every extension of the coefficient embedding. At the
definite end the test space consists of all sections.
\label{cor:slice-interpolation-4}
\end{proof}

\paragraph{The filtered product decomposition.}

Interpolation gives values on the cut. To estimate the cokernel of multiplication we need to express a section through products of short and light sections in successive degrees.

Write \(a=T-s\), and let \(t'_i\) be the chosen short
rational forms of weight \(w_F\) whose restrictions
are a basis on the cut. For a section \(f\) of weight
\(\kappa=w_F+Tw_L\) with the low modes through \(c_1\)
zero, match its jets of order less than \(a\) using
the products
\[
            t'_i L_0^{T-|\alpha|}\prod_j S_j^{\alpha_j},
           \qquad |\alpha|<a .
\]
Divide by \(L_0^T\) locally at every cut point and use
the coordinates \(S_i/L_0\). Let \(\mathcal E\) be the evaluation matrix of the \(t'_i\) on
the entire cut.

If \(v_\alpha\) is the residual Taylor coefficient
on the cut for \(|\alpha|=j<a\), solve
\[
       \mathcal E\,b_\alpha=v_\alpha,\qquad
       g_{j+1}=g_j-\sum_{|\alpha|=j}\sum_i(b_\alpha)_i\,
             t'_i L_0^{T-j}\prod_m S_m^{\alpha_m},
       \qquad g_0=f .
\]
The new terms lie in the \(j\)-th power of the cut ideal, so
\(\mathcal E\) inverts the degree-\(j\) coefficient without
changing any coefficient of smaller degree. This leaves a section \(g\) vanishing
to order at least \(a\) along the whole cut. The goal
is to divide \(g\) by the degree-\(a\) monomials
with coefficients of weight \(w_F+s w_L\), bounded
and cuspidal on the quotient.

Put \(W_L=\sum_l\|w_l\|^2\) and
\(S=\sum_i\|S_i\|^2\) in the normalized light metric.
Take the quotient integral of \(|g|^2\) with the
weight
\[
     W_{\rm div}=e^{-\varphi_\kappa}
                           W_L^{-1/2}S^{-(T-1/2)} .
\]
Near a cut point use the complex dimension \(s\) when integrating
the bound from the order-\(a\) vanishing:
\[
   \int_0^\epsilon r^{\,2a-2(T-1/2)}r^{\,2s-1}\,dr=\epsilon .
\]
On the complement \(W_L\ll S\) bounds the singular factor
by a constant times \(W_L^{-T}\).
The \hyperref[par:slice-cusp-decay]{cuspidal decay estimate}
on complete rank-one periods at total cutoff \(c_1\) is
\(\|f(x)\|_{\kappa}^2\le C e^{-2\pi c_1/\lambda_1(x)^2}\).
The selected cutoff has \(TcP\le c_1/\lambda_1^2\),
where \(P\) is the first-minimum penalty of the light
family.

In detail \(P\le2 f/\lambda_1^2\) and
\(c_1\ge 2fTc\). The lower bound \(W_L\gg e^{-2cP}\) gives uniform control
on the complement of the cut in the cuspidal region:
\[
   \|f\|_\kappa^2 W_L^{-T}
    \ll \exp(2TcP-2\pi c_1/\lambda_1^2)
    \le \exp\bigl(-2(\pi-1)c_1/\lambda_1^2\bigr).
\]
Each correcting product has squared norm \(\ll W_L^T\):
the \(t'_i\) are bounded, and the squared norm of each light
factor is bounded by a constant times \(W_L\). The quotient has finite invariant
volume. For the definite lift compactness and
nonvanishing of the light family give this integrability
with the full cusp space.

\paragraph{Lifting before dividing.}
Let \(k=e^{-\beta R}\), using the exponent \(\beta\) of Lemma~\ref{lem:short-weight-curvature}.
Poincar\'e summation maps the holomorphic functions
of finite norm
\(\int |F|^2 W_{\rm div}D/k\) on the open flag cell to
a subspace of the quotient weighted Hilbert space.
After automorphy factors are absorbed into the norms, the
pointwise Cauchy--Schwarz inequality is
\[
 \Bigl|\sum_\gamma F(\gamma x)\Bigr|^2
 \le \Bigl(\sum_\gamma k(\gamma x)\Bigr)
       \sum_\gamma |F(\gamma x)|^2/k(\gamma x).
\]
Integration on the quotient unfolds the second sum to the
whole-cell integral, and the first sum is uniformly bounded.
The ball bound gives holomorphy and boundedness because
\(W_{\rm div}\gg e^{-\varphi_\kappa}\). The weight-gap
cuspidality criterion at \(\kappa\) places the image in
the finite-dimensional cusp space, so the image is closed
in the quotient weighted norm. We use the same criterion
for bounded holomorphic sections at \(w_F+s w_L\).
The within-block and extreme cross-block comparisons
in the zero-mode calculation apply at both weights.

\begin{proposition}[Analytic lifting and division on the common cut]
\label{prop:lift-and-divide-at-slice}
Let \(f\) be a section of weight \(w_F+Tw_L\) whose low modes
through \(c_1\) vanish, or a section in the full cusp space on the
compact definite lift. Assume the \(t'_i\) of weight \(w_F\) give
a basis on the whole cut, that the \(S_j/L_0\) are coordinates
there, and use the division margin
\(\beta^2u_0+\beta<m_{\rm div}\) of
Lemma~\ref{lem:short-weight-curvature} and the cell estimate of
Lemma~\ref{lem:short-poincare-summation}. Put \(a=T-s\).
Let \(J\) denote the section matching the jets through degree \(a-1\). The difference lies in the \(a\)-th power of the cut ideal and has an expansion
\[
 f-J=\sum_{|\alpha|=a}S^\alpha h_\alpha,\qquad \operatorname{wt}(h_\alpha)=w_F+s w_L,
\]
with every \(h_\alpha\) a bounded cusp form.
\end{proposition}

\begin{proof}
\proofstep{Surjectivity of the automorphic Poincaré sum.}
The next calculation proves surjectivity onto cusp
forms with finite quotient norm for this singular
weight. Suppose \(g_0\) in that space is orthogonal
to the image. Unfolding with the summable kernel gives
\[
          I=\int k|g_0|^2 W_{\rm div}D<\infty
\]
on the cell. In these coordinates write
\(\Psi=-\log W_{\rm div}-\log D\). The formula for
\(W_{\rm div}\), the nonnegative Levi forms of the
two logs of sums of squares, and the division
margin give 
\[
  \mathcal L(\Psi)\ge m_{\rm div}\mathcal L(R)
\]
where the cut functions do not vanish simultaneously.
Work on the pseudoconvex cell after removing the
zero hypersurface of one nonzero cut function.

\proofstep{Transferring the radius to the unfolded cell.}
The radius is pulled back from the bounded matrix
domain with potential \(-\log\det(I-Z^*Z)\) at each
embedding. Its gradient squared in its Levi form is
\(\operatorname{tr}(Z^*Z)\): at diagonal singular
values \(z_i\) the gradient entries are
\(z_i/(1-z_i^2)\) and the corresponding Levi
coefficients are \((1-z_i^2)^{-2}\).
This is at most \(\min(p,q)\) since \(\|Z\|_{\rm op}<1\).
Cauchy--Schwarz across the embeddings and pullback
to the flag cell give

\[
  |\partial R(\xi)|^2\le u_0\mathcal L(R)(\xi)
\]
for every tangent vector \(\xi\).
Apply the weighted \(\bar\partial\) estimate with
potential \(\Psi-\beta R\) to the derivative of
\(kg_0\). Its datum is \(-\beta k g_0\bar\partial R\),
and the inverse Levi cost is bounded using
\((m_{\rm div}-\beta)\mathcal L(R)\).
This gives a holomorphic \(F\) such that
\[
 \int |kg_0-F|^2\frac{W_{\rm div}}{k}D
       \le \frac{\beta^2u_0}{m_{\rm div}-\beta}I .
\]
The Levi form can be semidefinite in the compact flag
directions because the datum vanishes in those directions.
Use exhaustion and the directional estimate to apply the
weighted Hilbert space argument. Since the scalar on the
right is strictly less than one, \(F\) has finite seed
norm. It extends across the removed hypersurface by local
square integrability, using the local positive lower
bound for the metric. It extends to equivariant readings
on the other frame charts in the same way as the
interpolating seed.

Orthogonality to the sum of this \(F\) implies

\[
  \int F\overline{g_0}W_{\rm div}D=0
\]
 by unfolding.
The integrals are absolutely convergent by the seed norm
and \(I<\infty\). Thus the left hand side in the displayed
inequality equals \(I+\int |F|^2 W_{\rm div}D/k\).
It follows that \(I=0\), so \(g_0=0\) by the positive
weight almost everywhere and holomorphy. This shows that
\(g\) has a seed in the specified whole-cell norm.
The stronger inequality 
\[
  \beta^2u_0+\beta<m_{\rm div}
\]
enters here.
\cite{hormander-slice}\label{cor:slice-short-division-1}

\paragraph{Iterating division and returning to the quotient.}

After lifting once, division leaves a section of lower degree. The norm change at each repetition determines the cost on returning to the quotient.

In a local holomorphic frame write
\(S_0'=\sum_i|S_i|^2\), \(W_0'=\sum_l|w_l|^2\).
The seed norm for \(g\) is
\[
  \int |F|^2 (S_0')^{-(T-1/2)} e^{-\Phi},\qquad
     \Phi=\varphi_F-\log D+\tfrac12\log W_0'-\beta R .
\]
The light family has no common zero, and the division
margin makes \(\Phi\) plurisubharmonic. Use Skoda's
\(L^2\) division theorem with \(s\) holomorphic
generators in dimension \(s\). For an exponent
\(H>s\), finite norm with weight
\((S_0')^{-H}e^{-\Phi}\) gives division by the \(S_i\)
with coefficients of finite norm with exponent \(H-1\).
In the theorem's notation \(q=s-1\) and

\[
  H=\alpha q+1
\]
 with \(\alpha>1\).
Here the last input exponent is \(s+1/2\), so all
\(a=T-s\) divisions are allowed. Sum the finitely
many ordered products with the same exponent tuple.
Each resulting seed has the weight with exponent
\(s-1/2\). It extends across the removed zero
hypersurface by local square integrability.

\proofstep{Convergence and descent after division.}
The \(S_i\) and \(w_l\) are bounded in their normalized metric. Consequently each resulting seed \(u\) has finite norm for the regular metric, with the indicated density,
\[
  \int|u|^2 e^{-\varphi_{w_F+s w_L}}\frac{D}{k}<\infty.
\]
Extend each seed across the frame chart divisor
and Poincar\'e sum it in weight \(w_F+s w_L\).
The resulting forms are bounded by the ball estimate
and are cuspidal by the shifted-weight criterion.
Multiplication by each automorphic monomial
\(\prod_i S_i^{\alpha_i}\) commutes with summation
in the corresponding weight. Absolute convergence
therefore yields the division equation for \(g\)
on the whole quotient. The construction also works
in the definite case: \(R=0\), the gradient cost
is zero, and the compact quotient has no cusps.
\cite{skoda-slice}\label{cor:slice-short-division-2}

\end{proof}

\paragraph{Analytic realization of theta sections at the shifted weights.}\label{par:shifted-theta}
Fix one of the two runs and its embedding \(\sigma_0:K\to\C\).
Write \(C(u)\) for the space of rational cusp forms at weight \(u\).
We need three consequences of the arithmetic realization. The analytic realizations of \(C(w_F+k w_L)\) span the full space of complex cusp families for
every \(k\ge0\). At \(\kappa=w_F+T w_L\) the subspace with the
specified low modes zero is spanned by the analytic images of the rational good space. Finally at \(w_F,w_F+s w_L,\kappa\) there are spanning
sets in the prescribed scaled theta lattices with a bound on the mean
log supremum of each member. The last assertion is a size estimate.

\begin{proposition}[Theta spanning on marked components]\label{prop:shifted-component-span}
Use the marked level and two runs of
Proposition~\ref{prop:linked-comparison-rows}. On one run let
\(u=w_F+k w_L\) for an integer \(k\ge0\). Assume \(u\) has
even central degree at least \(4[E:\mathbb Q]r+4\) and
positive last weights in the respective blocks, as furnished
by the chosen \(w_F,w_L\). Let \(S_{\rm an}(u)\) be the space of analytic cusp families on the marked components, and \(\Theta_\alpha(u)\) the restrictions of the classical theta series on the common lattice family, with the prescribed scalar on each component. Then
\[
 S_{\rm an}(u)=\operatorname{span}_{\mathbb C}\{\Theta_\alpha(u)\}_\alpha.
\]

\end{proposition}

\begin{proof}
\proofstep{Apply the theta spanning theorem at the shifted weight.}
Write \(n=u.n\) for the central degree and
\(r=p_j+q_j\). The central degree of \(w_F\) is even and at least

\[
  8fr+4=4[E:\Q]r+4.
\]
 The last positive row weights of \(w_F\)
are positive. Adding \(k w_L\) preserves both conditions.
Each analytic component has the prescribed CM type, Hermitian form,
signature and polarization scalar. Use one component as reference.
Move the integral period lattices with their markings into the
self-dual lattice family of the reference space. Write \(g_\gamma\) for
the move from component \(\gamma\), \(L_\gamma\) for its moved lattice,
and \(\Gamma_D\) for the level group. The level data include the
principal polarization, the split odd level, the marked character
tuple, and the auxiliary type prime distinct from three. There are
two properties of this marked cover that matter for spanning.
Writing \([L]_{\Gamma_D}\) for the orbit of a lattice, the marking distinguishes the components by
\[
  [L_\gamma]_{\Gamma_D}=[L_{\gamma'}]_{\Gamma_D}\quad\Longrightarrow\quad \gamma=\gamma'.
\]

On the
stabilizer of \(L_\gamma\) the residual vacuum multiplier is one, and the
group transported back to the component \(\gamma\) acts trivially on
its section. Here the multiplier is what remains after the finite
character of the reduction is taken out of the Weil action.
We explain the use of these properties before giving the comparison
of the markings.

\proofstep{Comparison of theta functions in a common frame.}
Let \(P_\gamma\) denote the function representing the analytic cusp family in the chosen frame on component \(\gamma\). Move it
to the reference frame by \(g_\gamma\), dividing by the nonzero component
normalization \(C_\gamma\) in the theta function. Extend this function to
the entire orbit of \(L_\gamma\) by the inverse change of frame and the
vacuum multiplier in the automorphy law. Indeed two moves \(g_1,g_2\) leading to the same lattice differ by a stabilizer element, since
\[
  g_1L_\gamma=g_2L_\gamma\quad\Longrightarrow\quad g_2^{-1}g_1\in\operatorname{Stab}(L_\gamma).
\]
 The stabilizer properties and the cocycle law give
the same extension. Set the function to zero on every remaining
orbit.

The resulting function is bounded in the invariant norm since the multipliers
are unitary and there are finitely many chosen components.
Section~\ref{subsec:theta-row-rank} now applies to this one family.
Its inverse overlap operator and reproducing kernel give a finite
sum of theta series on the lattice family with coefficients in the
harmonic highest-weight polynomial spaces. Restricting the sum to
the \(L_\gamma\) gives all the component readings with one list of
lattices and polynomials.
\label{cor:slice-shifted-theta-1}

\end{proof}

\input{slice-marked-theta-cover}

\begin{proposition}[Rational sections at the shifted weights]\label{prop:shifted-rational-cusp-span}
Let \(\mathcal S(u)\) be the rational short sections and \(K\) the coefficient field of \(C(u)\). For each \(k\ge0\),
\[
 u=w_F+k w_L,\qquad \operatorname{span}_K\mathcal S(u)=C(u).
\]
Their analytic realizations span all analytic cusp families on the run.
\end{proposition}

\begin{proof}
In the resulting classical series write \(\Lambda\) for a source
lattice and \(N_\Lambda\) for an integer exhibiting its admissible
level. The integral harmonic polynomials at this level span its
complex highest-weight space. A concrete spanning set is constructed
below. The classical evaluation of a short rule attached to an
integral polynomial \(P\), on component \(\gamma\), is the corresponding
classical theta reading times \(C_\gamma u_{\Lambda,N_\Lambda}\), with
\(u_{\Lambda,N_\Lambda}\ne0\) independent of the component \(\gamma\). Write \(\Theta_{\Lambda,\gamma}(P)\) for the classical theta
reading on component \(\gamma\), linear in the harmonic polynomial.
The common classical expansion has coefficients \(P_\Lambda\) with

\[
  C_\gamma^{-1}P_\gamma=\sum_\Lambda\Theta_{\Lambda,\gamma}(P_\Lambda).
\]
Replace each polynomial by \(P_\Lambda/u_{\Lambda,N_\Lambda}\)
and express the result in the spanning integral polynomials.
The corresponding short rules then have total reading \(P_\gamma\)
on every component, using the same coefficients on them all. Both sides can be
extended by zero off the frame domains. For an empty component set
the assertion holds with an empty sum. The transformation law
gives rational sections and the last-positive-weight condition
gives cuspidality at every good embedding. Faithfulness after
extension of \(K\) follows from the independent evaluation rows of
Section~\ref{subsec:selected-cusp-reading}. In particular the rational
span of the short rules is \(C(u)\), since a proper rational
subspace stays proper after extending scalars.
\label{cor:slice-shifted-theta-2}

\end{proof}

\paragraph{Descent of the low-mode space.}\label{par:shifted-good}
Suppose \(j<m\). For each admissible CM arc over \(K'/K\), use the
coefficient functional \(\phi:C(\kappa)\to K'\) of the specified
frame monomial and order. The arc calculation at the half-trace
cutoff \(c_1\) identifies the intersection of their kernels over
\(K\) with \(H_g\). For a complex linear combination with all low
modes zero, the coefficient equation holds at \emph{each} extension
\(\iota:K'\to\C\) of \(\sigma_0\). Indeed at that embedding the
conjugate arc realizes the same Fourier--Jacobi test on the unfolded
components, with the same cutoff. Extracting the coefficient from
the sum commutes with its finite complex linear combination.

\begin{proposition}[The rational good space and its complex extension]
\label{prop:shifted-low-mode-descent}
At \(\kappa\) the complex families with the prescribed low modes
zero are exactly the scalar extension of the rational kernel \(H_g\)
of the coefficient tests. Thus if \(S_{\rm low=0}\) denotes the complex subspace defined by those coefficient tests,
\[
 S_{\rm low=0}=H_g\otimes\mathbb C,\qquad H_g=C(\kappa)\quad\text{at the definite lift}.
\]

The rational kernel spans the same good families used in the
slice-rank and division arguments.
\end{proposition}

\begin{proof}
\proofstep{Descent by annihilation of the low modes.}
To see why these equations suffice, take a full-span expansion
of the analytic family in rational forms and choose a
\(\sigma_0(K)\)-basis \(b_1,\ldots,b_h\) of the span of its complex
coefficients. Regroup the rational forms into \(F_1,\ldots,F_h\)
so the expansion has coefficients \(b_i\). For each arc,
\(\sum_i b_i\iota(\phi(F_i))=0\) for every \(\iota\) above
\(\sigma_0\).

For any \(a\in K'\), multiply by \(\iota(a)\) and
sum over these embeddings. The result is
\[
 \sum_i b_i\,\sigma_0\!\left(
       \operatorname{Tr}_{K'/K}(a\phi(F_i))\right)=0 .
\]
Independence and the nondegenerate trace pairing give
\(\phi(F_i)=0\) for every \(i\). The same \(F_i\) work for every
arc, so all belong to \(H_g\). Conversely the readings of \(H_g\)
have zero low modes by the detector. At \(j=m\) the Hermitian
form is definite, there is no isotropic line, and \(H_g=C(\kappa)\).
Thus in both cases the interpolation, rank and division arguments
use precisely the scalar extension of the rational good space.
\label{cor:slice-shifted-theta-3}

\end{proof}

\input{slice-short-heights}


\paragraph{Descent of the containment.}

The three spans have now been defined over the coefficient field. The containment needed in the index estimate concerns their saturated intersections with the integral row.

\begin{proposition}[The saturated containment on the selected row]
\label{prop:descent-of-slice-division}
With the rational good space of Proposition~\ref{prop:shifted-low-mode-descent},
the containment provided by Proposition~\ref{prop:lift-and-divide-at-slice}
holds over the generic coefficient field \(K(X)\). Writing \(M_g\) for the integral good submodule, and taking saturation inside the completed cusp lattice, one has
\[
 \operatorname{Sat}(M_g)\subseteq\operatorname{Sat}(N_1+Z).
\]

Finite indices inside these saturations require separate bounds.
\end{proposition}

\begin{proof}
A rational form of the selected good space has the required
low modes zero at \(\kappa\) by the maximal cusp comparison.
The preceding division writes its reading in the complex
span of the leading products and the boundary products.
Work first over \(K(X)\), the coefficient field with the
cutting coefficients indeterminate. The products are
polynomials in \(X\) in the rational cusp space.

At the
chosen algebraically independent complex tuple, the
full reading on that space is faithful after scalar
extension. Thus the ranks of the columns with and
without the rational form are exactly the ranks of the
original matrices over \(K(X)\), as polynomial minors
vanish at this tuple precisely when they are zero
polynomials.
Writing \(A\) for the column matrix and \(v\) for the rational form to be tested, membership is the minor test
\[
  v\in\operatorname{im}A\quad\Longleftrightarrow\quad \operatorname{rank}[A\mid v]=\operatorname{rank}A .
\]
 This gives the containment over \(K(X)\).
Extend scalars to the fraction field of the Gauss DVR
at each finite place. Intersecting with the completed cusp
lattice then gives containment of the saturated good
lattice in the saturation of \(N_1+Z\). The indices inside
those saturations still require the finite-place comparisons.
\label{cor:slice-short-division-3}
\end{proof}

\paragraph{The coefficient tests at the prescribed cutoff.}
We keep the two level primes, the CM field \(E\) of degree \(2f\),
and the lift \(j<m\). Write \(\mathcal A\) for the sum of the
nonnegative flag exponents of the section being tested and \(n\)
for its central weight. For weight \(\kappa\) these numbers satisfy
\(\mathcal A+frn\le S_\kappa\). Let \(q_j\) be the negative signature rank. Define the
Fourier--Jacobi cutoff at this lift by
\[
 c_1=\tfrac12\operatorname{rd}(E)^{C_R}(2fq_j)^2q_j\,Tc,
 \qquad C_R=7 .
\]

The weight and cutoff fix the torsion order \(n'\) and the letter
radius bound \(\Lambda\) before the tests are chosen.
The split approximation uses weight allowance \(S_\kappa+r\),
level \(g_0g_1\), and the two integer norms
\(|\operatorname{disc}E|\) and \(\mathrm N\mathfrak c_K\).
The radius is enlarged to be at least one. The split construction
is used only at \(j<m\), where \(q_j\ge1\) and \(c_1>0\).
Thus the two signature cases entering the split construction are separated by
\[
  q_j=m-j\ge1\quad(j<m),\qquad q_m=0 .
\]
At \(j=m\) there is no complementary quotient.
\label{cor:slice-cm-coefficients-1}

An arc has a character lattice of rank one over \(\mathcal O_E\), a semiabelian model,
a normalized theta generator, and tangent letters in each
embedding block. Let \(t\) be its parameter and let \(X=(X_\tau)\)
be independent square frame matrices. Replace each letter in
row \(i\) by the combination with coefficients \((X_\tau)_{ik}\)
of the original letters.
Writing \(\ell_{\tau k}\) for the original letters, the variable frame has rows
\[
  \ell_{\tau i}(X)=\sum_k(X_\tau)_{ik}\,\ell_{\tau k}.
\]
 The section reading is a polynomial
\(P_F(X,t)\) in the entries of \(X\), with coefficients in
\(K'((t))\), over the number field of the arc. Homogeneity follows
by scaling all the tangent letters and applying the flag character.
In particular
\[
 P_F(X,t)\text{ has total frame degree }\mathcal A,
 \qquad \phi_{\mu,J}(F)=[X^\mu t^J]P_F(X,t).
\]
It is linear in \(F\). The same definition applies to a section
with complex scalar coefficients, by extending each coefficient
reading linearly along the chosen embedding of \(K'\).

\paragraph{Keeping the selected mode under a change of frame.}

To construct the CM arc, follow one nonzero Fourier coefficient through the change of flag.

Fix the component and a nonzero low coefficient \(F_\nu\) on a
rational isotropic line \(E\xi\). Let \(a_\tau(x)>0\) be the
cusp time at a flag frame \(x\), normalized as the reciprocal of
\(c_\tau\operatorname{maj}_{x,\tau}(\xi)\), with \(c_\tau\) the
positive flow constant of the component. The complex transvection
with parameter \(s=(s_\tau)_\tau\) sends these times to
\(a_\tau(x)+\operatorname{Im}s_\tau\). Its domain is the product
of half-planes given by positivity of the translated times.
Real transvections act on the Fourier coefficient by

\[
  \exp(4\pi i\sum_\tau\nu_\tau s_\tau).
\]
The product of the flowed coefficient and the reciprocal
exponential is holomorphic on these half-planes, by integrating
the section over its compact period torus. It is constant under
real translations. Its complex derivatives in those same
directions vanish, so the product is constant throughout the
domain. In particular transporting a nonzero reading to all
times equal to one keeps it nonzero.

For a flag write \(\widehat N(x)\) for the supremum of the
normalized point norms of \(F_\nu\) over its frames. This is
finite by the Fourier coefficient bound on the compact real
period torus, since those real transvections preserve the point
metric. It is positive at some flag of time one. Put

\[
  S=\sup_{a(x)=1}\widehat N(x)>0.
\]
 Take a flag of time one
whose supremum is greater than \(S/2\), then a frame on it
whose point norm is still greater than \(S/2\).
Signed Gram--Schmidt is a right upper triangular change of frame.
The weight character of the section and the norm factor cancel
under this change, and the majorant depends only on the signed
flag. We may therefore take an orthonormal frame \(x_0\) of the
same flag, with point norm between \(S/2\) and \(S\).

\begin{proposition}[A split frame detecting the selected mode]
\label{prop:split-frame-low-mode}
Fix the torsion order before the definite lift. If \(a_\nu\) is a nonzero selected low coefficient, it is detected by a bounded split approximation,
\[
 a_\nu\ne0\quad\Longrightarrow\quad \exists\,x_{\rm split}\quad a_\nu(x_{\rm split})\ne0.
\]

The construction keeps the component and isotropic line fixed.
It chooses the quotient complement and the frame after the coefficient
and order are fixed, using the metric and radius bounds below.
\end{proposition}

\begin{proof}
\proofstep{A motion that preserves a large coefficient.}
Put \(\rho=(12\pi(c_1+w_0))^{-1}\).
For a frame with times in
\([(1-\rho)/(1+\rho),(1+\rho)/(1-\rho)]\), move its times to one
as above. A bound for the ratio of the norm factors follows
from the leading Gram minors under that imaginary transvection.
The change in the Gram matrix under a complex transvection
with imaginary time \(s_\tau\) is a Hermitian rank-one
update. Its coefficient has the sign of \(s_\tau\).

For the \(i\)-th leading block let \(G_i(0)\) and \(G_i(s_\tau)\) be its Gram matrices before and after the update. The determinant lemma gives
\[
  \frac{\det G_i(s_\tau)}{\det G_i(0)}=1+s_\tau b_i.
\]
 Here \(b_i\) is the signed squared projection
of the line onto that block in the flow normalization.
Since the line is isotropic, the leading coefficients start
and end at zero, increase on the positive block and decrease
on the negative block, with maximum \(1/a_\tau\).
Thus the ratios lie between \(1\) and \(1+s_\tau/a_\tau\)
in either order.
More explicitly, with \(a_\tau=a_\tau(x)\), the norm factor at
\(x\) divided by the norm factor after transport is bounded by

\[
  \prod_\tau\max(a_\tau,a_\tau^{-1})^{s_\tau/2},
\]
 where
\(s_\tau\) is the spread of the weight at \(\tau\). Both time
ratios are at most \(1+3\rho\). This contributes at most
\(\exp(3\rho w_0/2)\le e\).
The quotient of the raw coefficient norms is
\(\exp(4\pi\sum_\tau \nu_\tau(1-a_\tau))\).
Since \(\nu_\tau\ge0\), \(\sum_\tau\nu_\tau\le c_1\) and
\(1-a_\tau\le2\rho\), the quotient is at most

\[
  \exp(8\pi\rho c_1)\le e.
\]
 Thus
\begin{equation}\label{eq:slice-split-near-sup}
  \widehat N(x)\le e^2 S .
\end{equation}
The movement and the comparison apply to every frame of the flag,
which explains the supremum on the left.

\paragraph{The first minimum on the isotropic quotient.}
Use \(L'\) for the Hermitian self-dual lattice of the component,
and \(\mathcal L\) for its period lattice. The fixed factors satisfy
\(N_A\mathcal L\subset L'\) and \(N_C L'\subset\mathcal L\). The
majorant-orthogonal complement of the hyperbolic plane at \(x_0\)
represents \(\xi^\perp/E\xi\). Write \(v_\tau\) for the image there
of a putative short nonzero class represented by \(v\in L'\cap\xi^\perp\).
Let \(N_0,H_0,K_0',R_0\) be the split parameters defined below, and
\(e_{\rm sp}=\min(1,f^2/(N_0^2 N_C^2 c_1 R_0))\).

We claim that
\[
 \max_\tau\operatorname{maj}_{x_0,\tau}(v_\tau)\ge e_{\rm sp}.
\]
If the maximum is smaller, choose \(0<\epsilon_0<e_{\rm sp}\)
bounding it. The Hermitian self-pairing of \(v\) is integral
and is the pairing of \(v_\tau\) at every embedding. It has
absolute value less than one everywhere, so it is zero.
Shift \(v\) by \(N_A b\xi\), \(b\in\mathcal O_E\), to make its
majorant projection coefficient \(b_\tau'\) onto \(\xi\) satisfy

\[
  |b_\tau'|\le 2f^2\operatorname{rd}(E)N_A
\]
 at all places.
Indeed the Euclidean dual of the embedded \(\mathcal O_E\)
is the conjugate codifferent scaled by two, by the trace formula.

For each nonzero element of the codifferent the field norm has
absolute value at least \(1/d_E\), by the inclusion of its
principal ideal in the codifferent. The arithmetic--geometric
mean inequality gives Euclidean dual minimum at least
\(2\sqrt f/\operatorname{rd}(E)\).
Let \(\rho_2\) be the Euclidean covering radius and \(\lambda_1^*\) the dual minimum. Banaszczyk's bound in dimension \(2f\) gives
\[
  \rho_2\le\frac{2f}{2\lambda_1^*}\le\frac{\sqrt f\,\operatorname{rd}(E)}2.
\]
The maximum norm radius is therefore at most
\(2f^2\operatorname{rd}(E)\). \cite{banaszczyk}
The shift leaves \(v_\tau\) and the assumption of shortness
unchanged. The two vectors \(\xi,N_Cv\) are now in \(\mathcal L\)
and span a totally isotropic plane.

\proofstep{The average along the transvection torus.}
Average along the real transvection torus for this plane. Its
characters are Hermitian two-by-two frequency matrices
\(\chi_\tau=\tau(\beta)\). The principal level torus contains the
translations obtained from integral Hermitian matrices after
multiplying by the level. Trace duality, including the off-diagonal
trace pairing, gives 
\[
  N_0\beta_{ij}\in\mathcal O_E,
\]
 with
\(N_0=(g_0g_1)d_E^2\). Taking the \(\xi\)-coefficient first retains
just \((\chi_\tau)_{00}=\nu_\tau\). The surviving frequencies are
positive definite at every embedding.

In fact translating in a
rank-one positive direction with negative frequency would give
exponential growth against the polynomially varying norm factor.
This gives positive semidefiniteness. If one determinant vanishes,
the algebraic determinant vanishes. A rational vector in the
kernel gives a rational isotropic line on which the frequency
is zero. The corresponding constant-term projection is zero by
cuspidality. Translation averages commute on this abelian torus,
so such a frequency has zero coefficient.

Normalize the direction \(v_\tau\) to have cusp time one, writing
it as \(a_{\tau,0}\xi+a_{\tau,1}N_Cv\) after this rescaling.
The normalization gives
\[
 |a_{\tau,1}|^2
   = (N_C^2 c_\tau\operatorname{maj}_{x_0,\tau}(v_\tau))^{-1},
 \qquad 0<c_\tau\le2 .
\]
Write \(d_\tau=\det\chi_\tau>0\), let \(R_\tau\) be the value of
this frequency on the normalized direction \(a_\tau\), and put
\(T=2\sum_\tau R_\tau\). Completing the square in the first
coordinate gives \(\nu_\tau R_\tau\ge d_\tau|a_{\tau,1}|^2\).

The algebraic integer \(\det(N_0\beta)\) is nonzero, so the
product of \(N_0^2d_\tau\) over the complex half is at least one.
Since \(\sum_\tau\nu_\tau\le c_1\), two applications of the
arithmetic--geometric mean inequality give
\[
 \sum_\tau d_\tau/\nu_\tau\ge \frac{f^2}{N_0^2c_1},
 \qquad T\ge\frac{f^2}{N_0^2N_C^2c_1\epsilon_0}> R_0 .
\]
The first inequality uses
\(\prod_\tau(f\nu_\tau/c_1)\le1\). In particular all the
\(\nu_\tau\) under consideration here are positive.

\proofstep{The first minimum at the same frequency.}
The height count is on the same frequencies. In each embedding
write the second standard vector as \(q a_\tau+p(1,0)\).
Then \[
 |q|^2=N_C^2 c_\tau\operatorname{maj}(v_\tau)\le2N_C^2,\qquad |p|^2=N_C^2|b_\tau'|^2.
\]
Positivity and Cauchy--Schwarz bound the second diagonal and
the off-diagonal frequency entries by \(H_0(c_1+T/2)\).
For the off-diagonal use

\[
  |(\chi_\tau)_{01}|^2\le\nu_\tau(\chi_\tau)_{11}.
\]
The first diagonal is prescribed, so the other two entries
determine the frequency. In the maximum of the complex norms
the integer lattice has minimum at least one, by the field norm.
Disjoint translates of the open ball of radius \(1/2\) give at
most \((1+2B)^{2f}\) lattice points in a ball of radius \(B\).
For \(t\ge1\) use \(B=N_0 H_0(c_1+t/2)\) for each of the two
entries. Counting both entries therefore bounds the number of Hermitian frequency matrices \(\beta\) by
\[
  \#\{\beta:T(\beta)\le t\}\le(K'_0t)^{4f}\qquad(t\ge1).
\]

The normalized \(v_\tau\)-transvection commutes with the
\(\xi\)-flow. The time-one flag, and every frame in it, can be
moved back by \(i/2\) in the normalized \(v_\tau\) direction
at each place. Its time becomes \(1/2\), and the \(\xi\)-time
stays one, since the two directions are majorant-orthogonal.
The real period averages of the rank-two group at that moved
point therefore bound the normalized coefficient by \(S\).
Transport back multiplies the raw coefficient by \(e^{-\pi T}\)
and changes the norm factor by at most \(2^{w_0/2}\).
Group the resulting bounds by \(m=\lceil T\rceil\). For
\(m>R_0=w_0+4f\log(8fK_0')+4f+2\), the inequality
\(\log(m/(8f))\le m/(8f)-1\) gives

\[
  (K_0'm)^{4f}\le e^{2m-w_0-2}.
\]
Thus the sum of the coefficient bounds, divided by \(S\), is
at most
\[
 2^{w_0/2} e^{\pi-w_0-2}
 \frac{e^{4(2-\pi)}}{1-e^{2-\pi}}<\tfrac12 .
\]
The absolute bound justifies recovering the section from its
Fourier series on the compact torus. Taking the supremum over
the original flag contradicts the choice above \(S/2\). This
proves the quotient minimum with the cutoff fixed at \(c_1\).
\label{cor:slice-split-frame-1}

\paragraph{A signed rational complement.}

The quotient on the isotropic line controls a first minimum. To turn that estimate into a CM frame we need a rational complement with the prescribed signs.

The quotient of \(L'\cap\xi^\perp\) by \(L'\cap E\xi\) is
a self-dual Hermitian lattice on \(\xi^\perp/E\xi\).
One can check this locally. A generator of the saturated line
is primitive, and self-duality supplies a partner pairing to one.
If these vectors are \(\xi_0,\eta_0\), the Gram
matrix on their plane is
\[
 \begin{pmatrix}0&1\\1&H(\eta_0,\eta_0)\end{pmatrix}
\]
and its determinant is \(-1\). Its inverse is integral,
so subtracting the plane projection preserves the lattice.
Inside this plane the perpendicular to \(\xi_0\) is
the line it spans. Thus the quotient lattice is the
integral orthogonal complement. Its signature in the selected convention
is \((q_j-1,r-q_j-1)\).

\proofstep{Choosing the rational complement by transference.}
Use the transference bound for this lattice with lower majorant
\(e_{\rm sp}\). In real dimension \(2f(r-2)\) it gives covering
radius and an independent short family at

\[
  \mu=4(f(r-2))^3d_E/\sqrt{e_{\rm sp}}.
\]
Use the same lattice construction as for the full interior frame
with approximation allowance \(\rho/40\). Namely, approximate
the positive axes scaled by \(4\sqrt{r-2}\mu/(\rho/40)\),
and multiply the orthogonal projections of the short spanning
family to the negative space by the positive Gram determinant.
The positive block errors have operator norm at most
\(\sqrt{r-2}\mu\), so its span is a graph of norm at most
\(\rho/120\). The negative projection is given by the Gram
adjugate and is integral.

The leading positive Gram determinant
and the Gram-adjugate formula give the length bound
\(B_{\rm rad}^{2q_j-1}\) on the resulting letters, with
\(B_{\rm rad}=2\cdot4\sqrt{r-2}\mu/(\rho/40)\).
Indeed, for a vector \((u,v)\) the negative projected norm
squared is
\[
 \|(1-AA^*)^{-1/2}(v-Au)\|^2\le (1+\|A\|^2)\|(u,v)\|^2/(1-\|A\|^2)\le2\|(u,v)\|^2.
\]
Raising the common bound
for the positive lengths to twice the number of positive
axes pays the determinant. If \(q_j=1\) there is no positive
axis in the quotient, and the short negative family needs
no projection. Lift the letters to \(L'\cap\xi^\perp\).
Their pairings and their quotient basis are independent
of the choices of lifts.

In the signed orthonormal basis put
\(C_+=(1-A^*A)^{-1/2}\) and \(C_-=(1-AA^*)^{-1/2}\).
The block matrix with the required positive graph is
\[
                 U_A=\begin{pmatrix}C_+&A^*C_-\\
                                        AC_+&C_-\end{pmatrix}.
\]
For the form with matrix \(\operatorname{diag}(1,-1)\),
the two column blocks are orthogonal: their mixed matrix
is \(C_+(A^*-A^*)C_-=0\). Their squared Gram blocks are
\[
 C_+(1-A^*A)C_+=1,\qquad C_-(AA^*-1)C_-=-1.
\]
Thus \(U_A\) is unitary. Extend it by the identity on the
hyperbolic plane.

The
resulting frame \(x_1\) has the selected letters in the two
eigenspaces on the complement, and its \(\xi\)-time is still
one. The positive plane in the chart of \(x_0\) has graph
\(Z_1\) of norm at most \(\rho/100\). With
\(t=(1-\rho/100)/(1+\rho/100)\), the majorant at \(x_1\) is
between \(t\) and \(1/t\) times that at \(x_0\). The quotient
lower bound for \(x_1\) is therefore \(t e_{\rm sp}\).

\paragraph{Clearing the projection at the prescribed order.}
Let \(W_\tau\) denote the complex complement at \(x_1\),
and write \(P_\tau(\mu)\) for the coefficient of \(\xi\)
in the projection of \(\mu\in\mathcal L\cap\xi^\perp\).
The coefficient projection comes from the exact sequence
\[
 0\longrightarrow \mathcal L\cap E\xi
 \longrightarrow \mathcal L\cap \xi^\perp
 \longrightarrow
   (\mathcal L\cap\xi^\perp)/(\mathcal L\cap E\xi)
 \longrightarrow 0.
\]
The quotient is torsion-free because the submodule is the
intersection with an \(E\)-line. It is finitely generated over
the Dedekind ring \(\mathcal O_E\), hence projective, so the
sequence splits over \(\mathcal O_E\). Write \(l(\mu)\xi\) for an integral
projection from such a splitting. Its coefficient extends
by scalar extension to an \(E\)-linear form with 
\[
  l(\xi)=1.
\]
Hence the difference \(P_\tau-l_\tau\) is a linear form on
the quotient.

\proofstep{Simultaneous clearing at the prescribed torsion order.}
We approximate \(n'(P_\tau-l_\tau)\) simultaneously by the
linear forms of one lattice element. In the real tangent
space \(\prod_\tau W_\tau\) with the Euclidean majorant metric,
the image of \(\mathcal L\cap\xi^\perp\) is a lattice of
dimension \(d=2f(r-2)\). The quotient minimum is at least
\(\sqrt{t e_{\rm sp}}/N_A\): multiplication by \(N_A\) sends
the projected lattice to the one used for the first-minimum bound.
The marked Riemann pairing is
\(\operatorname{Tr}_{E/\mathbb Q}(-\xi_0 H(-,-))\) on period
vectors, where 
\[
  2|\tau(\xi_0)|=k_0.
\]
This pairing is perfect on the projected lattice, since it is
perfect on the period lattice and the isotropic sublattice is
saturated.

To see this, extend a basis of the primitive
isotropic sublattice by symplectic partners, obtained by lifting
their dual functionals under the perfect pairing. The resulting
even-rank block is unimodular, and the perpendicular quotient
is its integral orthogonal complement.
Relative to the majorant, the real operator associated to the
pairing is \(k_0\) times an orthogonal map.
Writing \(T=k_0 U\) with \(U\) orthogonal for the majorant gives
\[
  \|Tx\|=k_0\|x\|,\qquad T^*T=k_0^2I .
\]
 Indeed it is the
product of multiplication by a purely imaginary scalar,
the signature involution, and the factor two from the CM trace.

The similarity \(T\) takes the projected lattice to its Euclidean dual.
Thus the dual minimum is at least \(k_0\sqrt{t e_{\rm sp}}/N_A\).
Transference and rounding in \(d\) independent vectors give
covering radius at most
\(d^3N_A/(2k_0\sqrt{t e_{\rm sp}})\) for the projected lattice.
Via Riesz representation, the linear forms
\[
 u\longmapsto d_E(-\tau(\xi_0)) H_\tau(u,\mu_{0,\tau}),
 \qquad \mu_0\in\mathcal L\cap\xi^\perp ,
\]
come from the projected lattice by an operator of norm
\(d_Ek_0/2\) times an isometry. The covering estimate
therefore approximates the tuple \(n'(P_\tau-l_\tau)\)
with error in its representing vector bounded by
\(d_Ed^3N_A/(4\sqrt{t e_{\rm sp}})\).
The majorant of \(\xi\) is \(1/k_0\).

Multiplying the error
norm by \(1/\sqrt{k_0}\) for its value on the line is
covered by \(B_{\rm ind}(t e_{\rm sp})\), the index allowance
below evaluated with lower bound \(t e_{\rm sp}\).
For the discriminant factor, the pairings \(H(L',L')\)
generate the unit ideal by Hermitian self-duality.
Since 
\[
  N_C L'\subset\mathcal L,
\]
 integrality of the trace
pairing puts \(N_C^2(-\xi_0)\) in the codifferent. All its
embedding norms equal \(N_C^2 k_0/2\). The ideal norm lower
bound \(1/d_E\) therefore gives
\(N_C^2 k_0\operatorname{rd}(E)\ge2\), as used in the index
allowance.
All these comparisons concern the projected lattice of the
full period module, independent of the particular spanning
letters just chosen.

\proofstep{Returning to the Hermitian frame.}
Dividing by \(n'\) gives a correction \(\delta_\tau\) on \(W_\tau\)
of operator norm, as a map into the line with its majorant metric,
at most \(\kappa=B_{\rm ind}(t e_{\rm sp})/n'\).
It has the required integral effect. For \(\mu\in\mathcal L\cap\xi^\perp\)
the corrected projection multiplied by \(n'\) has coefficient
\[
 n'l(\mu)+d_E(-\xi_0)H(\mu,\mu_0).
\]
The second summand is integral since its value before multiplication
by \(d_E\), paired against each \(a\in\mathcal O_E\) by the trace,
is \(\operatorname{Tr}_{E/\mathbb Q}(-\xi_0 H(a\mu,\mu_0))\in\mathbb Z\).
Thus it is in the trace dual. The first summand belongs to the
line coefficient module by the integral splitting. The vector
with the displayed coefficient is in \(\mathcal L\cap E\xi\).

At each embedding choose \(\xi'\) in the hyperbolic plane with
\(H_\tau(\xi,\xi')=1\), of norm squared the reciprocal of
\(\operatorname{maj}_{x_1,\tau}(\xi)\), and isotropic.
Take \(w_\tau\in W_\tau\) with \(H_\tau(u,w_\tau)=\delta_\tau(u)\).
The unitary fixing \(\xi\) acts by
\[
 u\longmapsto u+\delta_\tau(u)\xi,\qquad
 \xi'\longmapsto
 \xi'-w_\tau-\tfrac12 H_\tau(w_\tau,w_\tau)\xi .
\]
Preservation of all pairings follows by expanding them on the
three summands. In this moved frame the line projection is
\(P_\tau-\delta_\tau\) and the eigenspaces on the quotient are
unchanged. The two linear off-diagonal terms have operator
norm at most \(\kappa\) each and the quadratic term at most
\(\kappa^2/2\), so the move differs from the identity by at
most \(3\kappa\) in the \(x_1\) majorant when \(\kappa\le1\).

In the \(x_0\) majorant the bound is \(\eta=3\kappa/t\).
With \(z_1=\rho/100\), the transformed graph has norm at most
\[
 \frac{z_1+\eta(1+z_1)}{1-\eta(1+z_1)}<\rho/21 .
\]
Indeed the new upper block on \(\binom{1}{Z_1}\) differs from
the identity by at most \(\eta(1+z_1)\), and the lower block
changes by at most the same amount. The inequality on \(n'\)
gives 
\[
  \kappa\le\rho/(1000t),
\]
 and \(\rho\le1/10\) gives the
strict numerical bound. This construction uses the \emph{chosen}
positive integer \(n'\). Clearing denominators first and then
taking an arbitrary multiple would give a different quantifier.
\label{cor:slice-split-frame-2}

\paragraph{The small-chart nonvanishing test.}
Let \(Z\) be the lower-block chart parameter of a final moved
flag in the orthonormal frame \(x_0\). Use the operator norm on
each block, and the maximum over the complex places. For
\(\|W\|\le\rho\) the points \(x(W)\) of this chart are in the
domain. For each singular value \(s\) of \(W\), the majorant in the
associated signed two-plane has eigenvalues \((1+s)/(1-s)\)
and its reciprocal relative to \(x_0\).
The two relative eigenvalues are reciprocal, as can be read from
\[
  \frac{(1+s)^2}{1-s^2}=\frac{1+s}{1-s},\qquad \frac{(1-s)^2}{1-s^2}=\frac{1-s}{1+s}.
\]
 This follows from the
unitary graph matrix above, whose two-by-two block is
\((1-s^2)^{-1/2}\left(\begin{smallmatrix}1&s\\s&1\end{smallmatrix}\right)\).
The remaining eigenvalues are one. Thus the majorants change by
factors between \((1-\rho)/(1+\rho)\) and its reciprocal.

For a dominant weight
of spread at most \(w_0\), the norm factor \(J(x)\) satisfies
\[
  |\log J(x(W))-\log J(x_0)|\le 4w_0 r\|W\|^2\le \tfrac14 .
\]
For leading minors ending in the positive block one takes
determinants of \(1-W_1^*W_1\), with \(W_1\) the relevant columns
of \(W\). For larger minors, eliminate the retained negative
identity block. The remaining factor is \(1-W_2^*W_2\), with
\(W_2\) the rows of \(W\) after deleting that block. The full
determinant has absolute value one.

For each of the other positive absolute determinant ratios \(q_i\), the block bound gives
\[
  |\log q_i|\le-r\log(1-\|W\|^2).
\]
Summation by parts in the consecutive minor ratios charges
only the weight spread, giving the stated quadratic estimate.
This and the majorant comparison put the chart in the range of
\eqref{eq:slice-split-near-sup}. If \(\|Z\|<\rho/21\), put
\(R=\rho/\max(\|Z\|,\rho/22)>21\).
The function \(g(t)=J(x_0)F_\nu(x(tZ))\) on \(|t|\le R\)
has \(|g(t)|\le e^{9/4}S\), and \(|g(0)|>S/2\).
If \(g(1)=0\), precompose with the automorphism of this disc
that exchanges 0 and 1. Schwarz's lemma gives

\[
  |g(0)|\le e^{9/4}S/R<S/2,
\]
 a contradiction.
This is the form of the nonvanishing argument used in the split
detector. It allows the frame move to depend on the coefficient
and the prescribed torsion order.
\label{cor:slice-split-frame-3}
\end{proof}


\paragraph{Isolation of the frequency on one arc.}

Nonvanishing on a small chart still permits several Fourier frequencies. An arc on which the chosen frequency can be separated gives a single coefficient to estimate.

Let \(E^+\) be the fixed field of complex conjugation and
write \(\operatorname{Tr}=\operatorname{Tr}_{E/\mathbb Q}\).
Let \(N\) be the transvection level, \(Y\) the rank-one
character lattice and \(\xi:Y\to\mathbb Z\) its marked
nonzero functional. Choose \(D>0\) clearing the trace dual
of \(\mathcal O_E\).

\begin{lemma}[The frequency lattice at a split cusp]\label{lem:slice-cm-frequency-lattice}
Every occurring frequency \(\nu\) at the selected split cusp
satisfies the archimedean positivity and finite trace conditions
\[
 \sigma(\nu)\ge0\quad(\sigma:E^+\hookrightarrow\mathbb R),\qquad \operatorname{Tr}_{E^+/\mathbb Q}(ND\nu\,\mathcal O_{E^+})\subseteq\mathbb Z.
\]

\end{lemma}

\begin{proof}
The principal-level transvections for the selected isotropic
line have trivial weight character. The full transvection
period lattice is rational in \(E^+\). Thus Fourier characters
are trace pairings with elements of \(E^+\). For a nonzero
coefficient, sending one cusp time to infinity with the other
times fixed multiplies its raw term by the corresponding Fourier
exponential. For a nonzero coefficient \(a_\nu\), with the
other times fixed, the modulus of that Fourier reading is a
constant multiple of
\(|a_\nu|\exp(-2\pi t\,\tau(\nu))\).
It is at most the supremum on the real period torus.
The norm change on this time family grows at most polynomially
in \(t\) in both directions. If \(\tau(\nu)<0\) the displayed
exponential violates that bound as \(t\to\infty\).
This proves total nonnegativity.

Trace duality gives
\(p:Y\to E\) with 
\[
  \operatorname{Tr}(a p(y))=\xi(ay)
\]
 for
\(a\in\mathcal O_E\). The map is \(\mathcal O_E\)-linear and
injective since it is nonzero on the torsion-free rank-one
lattice. For \(t_0\in\mathcal O_{E^+}\) the
bilinear form
\[
  (y,y')\longmapsto
  \operatorname{Tr}\bigl(2NDt_0 p(y)\overline{p(y')}\bigr)
\]
is divisible by \(N\), since \(D\overline{p(y')}\) is integral.
Under the unimodular alternating period pairing this divisibility
says that the transvection with parameter \(NDt_0\) moves each
lattice vector by an element of \(N\) times the lattice. One sees
this by pairing its displacement against every lattice vector.
The transvection is therefore a period with trivial weight
character. Its Fourier character gives
\(\exp(2\pi i\operatorname{Tr}(ND\nu t_0))=1\) for every \(t_0\)
in the integer ring of \(E^+\), which is the asserted dual
lattice condition.
\end{proof}

\begin{proposition}[Isolating a low mode by arc coefficients]
\label{prop:cm-low-mode-detector}
For \(j<m\), let \(a_\nu\) denote extraction of the Fourier--Jacobi coefficient of frequency \(\nu\). On the selected source space, the allowed arc coefficient tests detect precisely the modes of half-trace at most \(c_1\),
\[
 \bigcap_{\mu,J,\ {\rm allowed\ arcs}}\ker\phi_{\mu,J}=\bigcap_{\nu:\,\sum_\tau\nu_\tau\le c_1}\ker a_\nu.
\]

Here \(\tau\) ranges over the chosen CM half, so
\(\sum_\tau\nu_\tau=\tfrac12\operatorname{Tr}_{E/\mathbb Q}(\nu)\).
The statement holds after extension along the chosen complex
embedding because each nonzero low coefficient has a detecting arc
and frame. The detecting field may depend on that coefficient.
\end{proposition}

\begin{proof}
\proofstep{Choosing the frequency on the arc.}
Put \(j_0=\operatorname{Tr}(ND\nu_0)\), a nonnegative integer.
The totally nonnegative elements of this trace dual of trace
\(j_0\) are finite. Their real conjugates lie in a compact
simplex and the dual embeds as a lattice. Choose
\(b\in\mathcal O_{E^+}\) such that
\(\operatorname{Tr}((\lambda-ND\nu_0)b)\ne0\) for all other
elements of the slice. The nondegenerate trace pairing makes
each excluded set a proper rational hyperplane, and an
integral lattice cannot be covered by finitely many such
hyperplanes. Set \(M=\max_\tau |\tau(b)|\) and take an
integer \(R>M(2j_0+1)\). The element \(u=R+b\) is totally
positive.

If 
\[
  \operatorname{Tr}(\lambda u)=
\operatorname{Tr}(ND\nu_0 u)
\]
 and \(i=\operatorname{Tr}\lambda\),
then \(R|i-j_0|\le M(i+j_0)\). If \(i\ne j_0\) then \(|i-j_0|\ge1\) and
\(i+j_0\le |i-j_0|+2j_0\). Combining the two inequalities gives
\[
 R|i-j_0|\le M(i+j_0)\le M(2j_0+1)|i-j_0|,
\]
contrary to the choice of \(R\).
Equality of the perturbed traces in the slice now forces
\(\lambda=ND\nu_0\).

\proofstep{The period and its Fourier variable.}
Use \(s_E=NDu\) in the period form above. This form is
positive definite and level-divisible. Conversely it determines
\(s_E\) uniquely by trace duality against the multiples of
one nonzero \(p(y)\overline{p(y)}\). The rank-one toroidal
uniformization therefore gives an arc with the times
\(y\tau(s_E)\) at height \(y\). Its exponent for \(\nu_0\) is
\[
 J=2\sum_{\tau\text{ in the CM half}}
              \tau(\nu_0)\tau(s_E).
\]
All occurring frequencies have integral exponents by the level
condition, and \(\nu_0\) is the unique one with exponent \(J\).
Along the deep part of the flow the Fourier expansion converges
normally. Integrating the circle of \(t=e^{2\pi i w}\) against
\(t^{-J}\) therefore selects exactly the \(\nu_0\)-term.

For invertible \(X\), the resulting identity has the form
\begin{equation}\label{eq:slice-arc-coeff}
     [t^J]P_F(X,t)\,t^J
          = \kappa_X\,F_{\nu_0}(x_X),
          \qquad \kappa_X\ne0 .
\end{equation}
Here \(F_{\nu_0}\) is measured at the corresponding complex
frame and height, and \(\kappa_X\) includes the change of
theta section and tangent frame. The normalized base change
identifies the arc theta generator with the complex generator
up to a nonzero scalar. Formula \eqref{eq:slice-arc-coeff}
holds on the complex chart and hence computes its Laurent
coefficient, even if there is a pole at the boundary. Rescaling
\(t\) by a nonzero constant changes both sides of the coefficient
comparison by the matching power.
\label{cor:slice-cm-coefficients-2}

\proofstep{The kernel of the coefficient tests.}
These identities give both directions of the kernel assertion.
Every allowed arc has
\(\tfrac12\operatorname{Tr}_{E/\mathbb Q}\nu_0\le c_1\).
If every mode of half-trace at most \(c_1\) vanishes, write \(Q_J(X)=[t^J]P_F(X,t)\). On the punctured arc, \eqref{eq:slice-arc-coeff} gives
\[
 Q_J(X)=t^{-J}\kappa_X F_{\nu_0}(x_X)=0
       \qquad (X\ \text{invertible}).
\]
The product of the general linear groups is Zariski dense in the product of the matrix spaces. Hence the polynomial \(Q_J\) vanishes identically, which kills every coefficient test on this arc. Conversely take a nonzero
coefficient on one component with

\[
  \sum_\tau\nu_\tau=\tfrac12\operatorname{Tr}\nu\le c_1.
\]
The split construction just proved keeps this coefficient nonzero
and gives the bounded CM basis on the complement of the hyperbolic
plane. Use the resulting frame and the above positive form
to construct the arc. The abelian period lattice is obtained
from the orthogonal to the isotropic line by quotienting by
the line. Its complex structure is diagonal on the split CM
letters. The principal Riemann form on the original lattice
identifies the rank-one quotient with the dual of the isotropic
lattice.
Writing \(I\) for the isotropic lattice and \(L\) for the ambient period lattice, this is the perfect duality induced by the principal alternating form:
\[
  L/I^\perp\simeq I^\vee .
\]
 These are the two toric parts of the period lattice.
The complementary Riemann form is principal on the abelian
part. The bounded split approximation also gives the prescribed
\(n'\) which clears the projection to the isotropic line.

Write \(V_{\rm ab}/\Lambda_{\rm ab}\) for the abelian torus in the period decomposition and \(\widetilde y_{\rm ab}\) for the abelian component of a lift of \(y\in Y\). Clearing the projection gives
\[
       n'\widetilde y_{\rm ab}\in\Lambda_{\rm ab}
       \quad\Longrightarrow\quad
       n'[\widetilde y_{\rm ab}]=0\quad\text{in }V_{\rm ab}/\Lambda_{\rm ab}.
\]
Thus the gluing points have the prescribed torsion bound. The rank-one period
comparison keeps the complex flow on this component. The
canonical trivialization and algebraic descent are described
below. The identification of tangent spaces respects the embedding blocks,
so a block change \(X\) of these letters reads the prescribed
split frame.

Changes along the cusp flow multiply that mode by
a nonzero exponential. Thus \eqref{eq:slice-arc-coeff} is
nonzero at one such \(X\), and some monomial of the polynomial
is detected. The field of definition of the split CM data may depend on the coefficient being tested. Taking all such tests gives the rational kernel. After scalar extension, the tests annihilate every complex family whose corresponding low modes vanish.
\label{cor:slice-cm-coefficients-3}

\end{proof}

\paragraph{The normalized degeneration and its fields of definition.}
On the complementary period lattice let \(H_B\) be the Hermitian
form, normalized so that its imaginary part is the integral Riemann
form. Let \(C_y\) be the abelian component of the period belonging
to \(y\in Y\). The toric periods are affine in the arc coordinate
\(w\), with linear term given by the positive form of the preceding
construction. The constant term in their pairing can be chosen to be
\begin{equation}\label{eq:slice-cm-normal-pairing}
 T(y,y')=-\frac{H_B(C_y,C_{y'})}{2i}
                 +\operatorname{Im}H_B(\alpha C_y,C_{y'}),
 \qquad \alpha+\bar\alpha=1,
\end{equation}
where \(\alpha\) is the marked integral element of the CM host.
The Rosati relation and the displayed equation for \(\alpha\)
make \(T\) symmetric. They also give

\[
  T(y,\alpha y')+T(y',\alpha y)=T(y,y').
\]

Consequently \(\exp(2\pi i T(y,\alpha y))\) is the even
trivialization with polarization given by \(T\).
Replacing the old symmetric constant pairing by \(T\) is a
translation along the complex cusp flow. Indeed the difference
is conjugate balanced over \(\mathcal O_E\) on \(Y\), so over
the trace pairing it is a complex linear combination of real
parameters of that flow. The slope and abelian projection are
preserved. At large imaginary \(w\) the positive slope absorbs
the fixed translation in each embedding. A Fourier--Jacobi
coefficient under that translation receives its trace character,
which is never zero.

In the normalized model \(n'C_y\) belongs
to the abelian period lattice. Thus
\begin{equation}\label{eq:slice-cm-tau-normal}
  \exp\bigl(2\pi i T(n'y,n'y')\bigr)
       \exp\bigl(\pi H_B(n'C_y,n'C_{y'})\bigr)=1.
\end{equation}
Indeed its remaining exponent divided by \(2\pi i\) is the
integral Riemann pairing of \(\alpha n'C_y\) and \(n'C_{y'}\).
On this sublattice the two gluing points are the identity of the
abelian variety. The exponential product is the comparison of
the polarization of the trivialization with the rigidification
there. This is the canonical torsion normalization.
\label{cor:slice-cm-arc-model-1}

\begin{proposition}[Comparison of the two evaluations on the normalized CM arc]
\label{prop:marked-cm-arc-base-change}
On the marked CM degeneration with pairing
\eqref{eq:slice-cm-normal-pairing}, let \(s,a_\nu\) and \(s',a'_\nu\) denote the positive slope and the Fourier--Jacobi coefficient before and after the change of constant period. Then
\[
 s'=s>0,\qquad a'_\nu=\gamma_\nu a_\nu,\quad\gamma_\nu\ne0.
\]
 Over the fields with the prescribed marking, the
complex and finite-place evaluations use the same toric parameter and
the same normalized theta and frame lines, up to the indicated constant
frame changes.
\end{proposition}

\begin{proof}
\proofstep{Realizing and descending the marked period.}
The marked CM period realization applies to this complementary
lattice with its principal alternating form, split CM type and
bounded basis. It gives a polarized abelian scheme \(B\) after a
finite number field extension and an identification of its complex
periods and CM action with the above ones. In this construction one
keeps the isometry of marked period lattices, so the same
identification carries the Riemann matrix and the Weil pairings on
the gluing points.

Toroidal uniformization for the isotropic
quotient uses \(B\), the character lattice \(Y\), the gluing
homomorphism \(c:Y\to B[n']\), the positive period form and a
trivialization in the pulled-back cubical line. The norm and the
trivialization are compared using the rigidification at zero and
the symmetric primitive theta characteristic. Writing \(E_B\) for the integral period pairing, the symmetric characteristic evaluates on a period \(a\) as
\[
  (-1)^{E_B(\alpha a,a)}.
\]
 The equality of
the period pairings gives equality of these signs, which is needed
for the two Appell--Humbert readings to agree.
\cite{faltings-chai}\label{cor:slice-cm-arc-model-2}

\proofstep{Preserving the toric parameter in descent.}
To descend the normalized model, first adjoin the \(n'\)-torsion
of \(B\). The characteristic-zero torsion scheme is finite, so
all points of \(c\) are now defined over a number field \(k\).
Choose algebraic generators for the one-dimensional cubical
fibres at \(c(y)\). Compare them with the normalized analytic
generators by nonzero scalars \(q(y)\). In the quotient \(\mathbb C^\times/k^\times\), divide the
analytic and algebraic symmetry and cube identities. They
give an even function \(\bar q\) with bimultiplicative
polarization
\[
 b(y,y')=\frac{\bar q(y+y')}{\bar q(y)\bar q(y')}.
\]
The normalized product in \eqref{eq:slice-cm-tau-normal} is
the rigidification over \(k\) on \(n'\)-multiples, so
\[
 b(y,y')^{(n')^2}=b(n'y,n'y')=1.
\]
Bimultiplicativity also gives \(\bar q(0)=1\).
Since \(b(y,-y)=\bar q(y)^{-2}\), one has

\[
  \bar q(y)^{2(n')^2}=1.
\]
 Thus \(q(y)\) satisfies a
power equation over \(k\) and is algebraic.
For a finite integral generating family of \(Y\), it suffices to
adjoin the scalars on the generators and on their pairwise sums.
In the quotient by the enlarged field all polarizations vanish
by bimultiplicativity, and then all the scalars vanish in the
quotient group by the character law. This is one finite extension
for the entire trivialization. The remaining level data are points
in finite torsion schemes of the bounded level and can be adjoined
in a further finite extension.
If \(k_i/k\) are the finitely many fields of definition of the residual level records, their compositum is finite since
\[
  [k_1\cdots k_t:k]\le\prod_{i=1}^t[k_i:k].
\]
 The rigidified line, the marked
action and the normalized toric parameter are transported together.
\label{cor:slice-cm-arc-model-3}

\paragraph{Finite-place realization of the same arc.}

The archimedean construction has now specified the arc and its field. Its evaluation must also extend to the completed model at every finite place needed for the height.

Use a finite place \(w\) of the field \(K'\) just constructed,
write \(V\) for its complete DVR and \(F_V\) for its fraction
field. Base change the abelian part \(B\) to \(V\). By properness
every gluing point in \(B(F_V)\) extends uniquely to \(B(V)\).
The extended points form a homomorphism \(c_V:Y\to B(V)\):
both sides of 
\[
  c_V(y+y')=c_V(y)+c_V(y')
\]
 extend the same
generic point, as do the sides of the CM equivariance identity.
For the rigidified cubical line \(\mathcal T\) on \(B/V\),
the datum retains each generic generator in
\[
             (c_V(y)^*\mathcal T)\otimes_V F_V .
\]

In particular one retains its valuation relative to
\(c_V(y)^*\mathcal T\), even if the generator is not a basis
over \(V\). This is
the datum used in the integral toroidal model. In particular it
has the same generic fibre as the original selected arc. The
level points in the degeneration datum may be kept as generic
level records. Extension of a generic level point as a point of
an abelian torsion group alone would not specify its toric record.

\proofstep{Integral frames and tangent coordinates on the chart.}
An integral framed chart has constant tangent letters in each
CM embedding block and a generator of the theta line on the
integral abelian part. The comparison of two such descriptions
is made over \(F_V\) by base change of the same degeneration
datum.
Thus two ordered tangent bases \(e,e'\) on a common CM block over the generic fibre differ by
\[
  e'=eA,\qquad A\in{\rm GL}_n(F_V),
\]
 Each block of algebraic tangent letters is related to the
integral block by a constant matrix. Constancy follows by taking
invariant tangent vectors of the semiabelian extension and
identifying their pullbacks via the common generic group. The
theta generators there span the same one-dimensional space.

Let
\(\gamma\) be the scalar expressing the algebraic section in the
integral generator.
Thus if \(\vartheta_{\rm alg}\) and \(\vartheta_{\rm int}\) are the indicated generators, their comparison and its valuation are
\[
  \vartheta_{\rm alg}=\gamma\,\vartheta_{\rm int},\qquad v_w(\vartheta_{\rm alg}/\vartheta_{\rm int})=v_w(\gamma).
\]
 These identifications determine the frame
comparison over the entire arc. A change of the torsion
normal form multiplies the reading by a root of unity, as the
torsion sections differ by their level character. This leaves the
finite order of \(\gamma\) unchanged.
\cite{faltings-chai}\label{cor:slice-cm-arc-model-4}

The generator of the global theta ideal at \(w\) is a local
coordinate for the theta line of this same polarized model.
Carry the generator and a Lie frame through the polarized
pullback. The integral theta generator differs from the former
by a unit in \(V\), since both generate the same rank-one module
and the isomorphism is rigidified. If \(t_w\) generates the theta ideal in this trivialization,
then 
\[
  \gamma^{-1}=u t_w
\]
 with \(u\in V^\times\).

Passing to the dual
ideal changes the sign:
\[
 v_w(\gamma)=-v_w(t_w)=v_w I(\mathcal T^\vee).
\] The determinant
of the Lie comparison is that of the same model. This identifies
the central factor of the coefficient line used at infinity as
well as at finite places. In the finite comparison below, the
integral toroidal evaluation is made on smaller punctured discs
over valued extensions. The input is the framed chart extending
the normalized degeneration, with the valuation-compatible
extension of the coefficient field. Base change makes the two
evaluations agree after the constant frame change. The bound on
the integral reading uses the integral lattice of sections on
that chart.
\label{cor:slice-cm-arc-model-5}
\end{proof}


\paragraph{Finite order of an arc coefficient.}
At a place \(w\) of \(K'\) over \(\lambda\) of \(K\), take
the completed integral semiabelian model. Let
\(\mathcal T\) be the theta line with its normalized generator
and \(\ell_{\tau k}\) the line of the letter in column \(k\)
of the \(\tau\)-block. A toric letter has the unit line.
Write \(\gamma\) for an inverse local generator of the ideal of \(\mathcal T\)
and \(d_{\tau k}\) for the generators of the ideals of the
letter lines. Use \(\operatorname{ord}_w\) with integer orders.
The integral comparison supplies a matrix \(B_\tau\) from
an integral frame to the coordinates normalized by these
generators. In its inverse \(A_\tau\), every entry has order
at least \(-b_w^0\), where
\[
 b_w^0=e_{w/\lambda}
             \max\{\operatorname{ord}_\lambda|\operatorname{disc}E|,
                         \operatorname{ord}_\lambda n'\}.
\]

For a toric letter the infinitesimal character on \(y\in Y\)
is the labelled conjugate of \(p(y)\). A torus tangent vector
over the dual numbers is determined by its evaluations
\(1+\varepsilon h(y)\) on the characters, with \(h\) additive.
The trace-dual condition gives \(d_Ep(y)\in\mathcal O_E\),
so the multiplied tangent vector is integral by this
character description on the same constant torus. For an
abelian letter the local generator \(d_{\tau k}\) gives an
integral tangent vector on \(B\). The classifying extension
has \(n'\)-torsion gluing homomorphism and the chosen
trivialization on \(n'Y\). The pullback under \([n']\) is split,
with these trivializations specifying a splitting compatible
with the generic fibre. Applying its tangent map to the
integral vector on \(B\) extends precisely the generic
semiabelian letter multiplied by \(n'\).

On invariant tangent vectors the abelian multiplication map acts by
\[
  d[n']_0(v)=n'v .
\]
 The local unit
needed to compare a differential model with this frame
is invertible over \(V\). These are separate columns, and
the same bound on the entries uses the larger of the two
orders. The theta generator is the integral rigidified
generator up to a local unit after rescaling by its line.
These comparisons are made on the integral chart of the
toroidal model. An integral section is bounded by one when
evaluated there in any integral frame with unit determinant.
\label{cor:slice-cm-coefficients-4}

\begin{lemma}[Finite and complex bounds for one coefficient]
\label{lem:cm-coefficient-metrics}
Use the theta and tangent lines of the selected arc. For an integral
source form the coefficient bounds in the following calculation hold
on the ideal defined by those lines for the same \(\phi_{\mu,J}\)
at the finite and the complex places. For the finite reading on an integral source form \(F\), the coefficient ideal of the value line gives
\[
 \operatorname{ord}_w\phi_{\mu,J}(F)-\operatorname{ord}_w I(\mathcal L_\mu)\ge-\mathcal A b_w^0.
\]
The complex bound uses \(\max(1,c_1)\) and the frame radius bound \(\Lambda\), independently of the length of the isolating direction.
\end{lemma}

\begin{proof}
\proofstep{The nonarchimedean coefficient estimate.}
We give the coefficient argument because it also fixes the finite
offset in the quotient. In the variable
\(Y=(Y_\tau)\) for the integral frames put
\[
 R(Y,t)=\gamma^{-n}P_F((Y_\tau B_\tau D_\tau^{-1})_\tau,t),
 \qquad D_\tau=\operatorname{diag}(d_{\tau k}^{-1}),
 \qquad A_\tau=B_\tau^{-1}.
\]
The finite chart comparison gives convergence and
\(|R(Y,t)|_w\le1\) for the integral \(Y\) with unit block
determinants and the allowed \(t\) with \(0<|t|_w<1\).
One may evaluate over the completed algebraic closure with
infinite residue field. The allowed radii are dense in
\((0,1)\) and the readings of \(t\) are stable under multiplication
by units with arbitrary nonzero residues in an infinite
subfield.
In particular unit variation on any allowed circle preserves its radius:
\[
  |t_0u|_w=|t_0|_w\qquad(u\in V^\times).
\]
 Such extensions and parameters in the integral arc
comparison leave the value line on \(K'\) fixed.

For a fixed \(Y\) and allowed radius \(|t_0|_w\), the maximum
of \(\big|[t^i]R(Y,t)t_0^i\big|_w\) is attained if it is nonzero.
Convergence makes only finitely many of the terms attain it.
Divide by one such term and reduce the sum at \(t=t_0u\).
If \(a_i=[t^i]R(Y,t)\), the normalization is by one term attaining
\[
  |a_jt_0^j|_w=\max_i |a_it_0^i|_w>0 .
\]
The remaining terms of maximal norm give a nonzero Laurent
polynomial in the residue of \(u\). A unit with a residue
outside its finitely many roots gives a value of precisely
the maximal norm. Hence
\[
 \big|[t^J]R(Y,t)\big|_w\,|t_0|_w^J\le1.
\]
Let the radii tend to one from below. This works for every
integer \(J\).

\proofstep{Clearing denominators of the chosen coefficient at finite places.}
For the fixed polynomial \([t^J]R(Y,t)\) divide again by a
coefficient of maximal absolute value. Its reduction times
the product of the block determinants is nonzero. A tuple
in the infinite residue subfield detects this polynomial,
so lifts of the tuple give unit block determinants and
detect the original maximal coefficient norm. It follows
that every coefficient of \([t^J]R(Y,t)\) is integral.
The reverse substitution 
\[
  Y_\tau=X_\tau D_\tau A_\tau
\]
uses \(\mathcal A\) factors from the matrices \(A_\tau\)
in every term. Their contribution to a coefficient has
order at least \(-\mathcal A b_w^0\), with no additional
loss from the integer multinomial coefficients or the sums.
Evaluation after the remaining diagonal rescaling multiplies
the \(\mu\)-coefficient by exactly
\(\prod_{\tau,i,k}d_{\tau k}^{-\mu_{\tau i k}}\).

Consequently the finite estimate is
\[
 \operatorname{ord}_w\phi_{\mu,J}(F)-
       \operatorname{ord}_w I(\mathcal L_\mu)
       \ge-\mathcal A b_w^0,
 \qquad
 \mathcal L_\mu=(\mathcal T^\vee)^n
                  \otimes\bigotimes_{\tau,i,k}
                         (\ell_{\tau k}^\vee)^{\mu_{\tau i k}}.
\]
This gives \(s_{iw}=-\operatorname{ord}_w I(\mathcal L_\mu)\)
in the quotient calculation below. A test of the wrong
total frame degree vanishes by homogeneity. For the other
tests the common baseline has the factor \(\mathcal A\).
\label{cor:slice-cm-coefficients-5}

\paragraph{Archimedean bound for the same coefficient.}

The local order accounts for the finite part of the height. The archimedean contribution must be evaluated on the same coefficient and the same arcs.

Use a good embedding \(\sigma\) of \(K'\) and denote the
positive metric factors of \(\mathcal T,\ell_{\tau k}\) there
by \(q_\theta,q_{\tau k}\). Put \(v=\max(1,c_1)\) and
\(\epsilon_*=v^{-1}\). Move the frame in
\eqref{eq:slice-arc-coeff} along its cusp flow to all CM
times equal to \(\epsilon_*\). The shift adds
\(\epsilon_*-y\tau(s_E)\) to each old time. The Fourier
factor at time \(y\tau(s_E)\) is recovered by multiplying by
\[
 \exp\left(4\pi\sum_{\tau}\tau(\nu_0)
                         (\epsilon_*-y\tau(s_E))\right)
 = |t|^J\exp\left(4\pi\epsilon_*\sum_\tau\tau(\nu_0)\right).
\]
Here \(y=\operatorname{Im}w\) and the sum is over the CM half.

The total frequency bound and the common shifted time give
\[
  \epsilon_*\sum_\tau\tau(\nu_0)\le\frac{c_1}{\max(1,c_1)}\le1,\qquad \exp\!\left(4\pi\epsilon_*\sum_\tau\tau(\nu_0)\right)\le e^{4\pi}.
\]
 The \(|t|^J\) cancels the factor
from taking the modulus of \eqref{eq:slice-arc-coeff}.
Thus the initial height and the size of the isolating direction
\(s_E\) add no cost to this estimate. In particular the proof of
isolation needs no uniform bound on \(R\). The transvection torus integral for the coefficient
at the new time bounds it by the supremum of the section
divided by its pointwise normalizing factors.

\proofstep{Archimedean normalization on the same arc.}
For clarity, the remaining normalizing factors in this
integral can be computed before coefficient extraction in
\(X\). Let \(H_X\) be the positive Riemann form in the
shifted frame and \(h_{\tau i}=\pi (H_X)_{\tau i,\tau i}\).
The complex comparison of the theta and frame
normalizations gives the factor
\[
 \left(C v^f\right)^{n/2}\prod_{\tau,i}h_{\tau i}^{a_{\tau i}/2},
 \qquad C\le q_\theta^2\operatorname{rd}(E)^f .
\]
Indeed the central normalization is initially
\((C\prod_\tau \mathrm{time}_\tau^{-1})^{n/2}\)
times the central-weight norm in the old frame.

The central-weight norm times this product is invariant
under the move. Dividing by the norm of the total weight
at the new frame leaves the diagonal factors displayed
above. In the normalized theta trivialization, comparison
with the period volume gives \(C/q_\theta^2\le
\sqrt{|\operatorname{disc}E|}\).
These comparisons use the Riemann form on the same
semiabelian degeneration as the trace formula and the
abelian CM metric on its central fibre.
\label{cor:slice-cm-coefficients-6}

\proofstep{Bounding the coefficient on the polydisc.}
Give \(X_{\tau i k}\) radius \(q_{\tau k}^{-1}\) in a closed
complex polydisc. At new time \(\epsilon_*\), the individual
CM letter satisfies

\[
  \sqrt{\pi H_X(e_{\tau k},e_{\tau k})}\le
\Lambda\sqrt v\,q_{\tau k},
\]
 where the pairing on the left
is the pullback of the shifted Riemann form to the letter.
The same inequality also holds for the toric letter with
its unit line; its squared length is bounded by the
reciprocal time. The row is a sum of at most \(r\) such
letters with the indicated coefficients. The triangle
inequality in that Hermitian norm gives
\(\sqrt{h_{\tau i}}\le r\Lambda\sqrt v\).
Put
\[
 p_{\rm arch}=\log(r\Lambda)
       +\tfrac12\log v+\log\operatorname{rd}(E).
\]
The product above is at most
\(q_\theta^n\exp(p_{\rm arch}(\mathcal A+fn))\).
These bounds first apply where the block matrices are
invertible.

Since the reading is a polynomial in \(X\),
continuity extends the bound to the full polydisc.
The multivariate Cauchy coefficient formula gives
\[
 |\sigma\phi_{\mu,J}(F)|\,\|\mathcal L_\mu\|_\sigma
 \le \|F\|_{\sup,\sigma|K}\,
       \exp\big(4\pi+p_{\rm arch}(\mathcal A+fn)\big).
\]
In fact the product of the Cauchy radii is exactly the
differential metric factor of the dual value line. The
division by \(q_\theta^n\) gives its remaining factor.
At the other embeddings use the conjugate good
uniformization, with the same metric by complex
conjugation. This gives the archimedean offset for precisely
the ideal in the finite calculation.
\label{cor:slice-cm-coefficients-7}
\end{proof}

\paragraph{The height of the boundary columns.}
The arithmetic proof of Proposition~\ref{prop:geometric-run-index}
must use one bound for each of the two determinant costs
in \eqref{eq:geometric-four-terms}. The boundary determinant
below is evaluated in the trace dual. The complementary
determinant is evaluated by CM coefficient functionals.
\hypertarget{source-slice-budget}{}
\begin{lemma}[Short trace-dual tests for boundary columns]
Let \(r_z>0\), \(d=R_j\) and \(n=[K:\mathbb Q]\). For every
\(B>d\sqrt n\exp(P_LS_\kappa+P_d/2)\) the \(K\)-dual of the
column space has a spanning family of functionals \(\phi\) satisfying the two bounds
\[
 \phi(M)\subseteq\mathfrak D_K^{-1},\qquad |\sigma\phi(v)|\le B\|v\|_{2,\sigma}\quad(\text{every }\sigma).
\]

\end{lemma}
\begin{proof}
Equip \(M\) with the Euclidean norm from the sum of the
conjugate-compatible squared \(L^2\) norms. Its real dimension
is \(nd\). The height bound \(H(v)\ge-P_LS_\kappa\) for integral
\(v\ne0\), the inequalities
\(\|v\|_{2,\sigma}\le\|v\|_{\sup,\sigma}
 \le e^{P_d/2}\|v\|_{2,\sigma}\), and arithmetic--geometric mean give
\[
 \lambda_1(M)\ge\sqrt n\,e^{-P_LS_\kappa-P_d/2},\qquad
 \lambda_{nd}(M^*)\le \frac{nd}{\lambda_1(M)}<B .
\]
The second inequality is Banaszczyk's transference theorem.
By nondegeneracy of trace the rational functional attached to
an integral dual vector is \(v\mapsto{\rm Tr}_{K/\mathbb Q}\phi(v)\).
Testing on \(av\) for every \(a\in\mathcal O_K\) gives the
codifferent condition. Each Hermitian component of its Riesz
vector has norm at most its full Euclidean norm. Applying these
facts to the spanning short dual vectors gives both assertions.
\cite{banaszczyk}\label{cor:slice-budget-1}
\end{proof}

\begin{proposition}[The boundary determinant cost]\label{prop:slice-boundary-determinant-cost}
For the selected \(r_z\) boundary columns and the norm bounds
above, the finite costs \(G_\lambda\) have finite support and satisfy
\[
 \langle G\rangle\le r_z H_z S_\kappa,\qquad H_z=2G_c+P_{rw}+P_L+\tfrac32(P_d+P_O).
\]

If there are no columns, the cost is zero.
\end{proposition}

\begin{proof}
\proofstep{The determinant of the boundary columns.}
The boundary columns are independent over \(F\), so some
minor of their values on these functionals is a nonzero
polynomial. Choose one nonzero coefficient \(A\in K\).
At each prime, multiplying all its rows by an element of
order \(v_\lambda(\mathfrak D_K)\) makes the functionals
integral there. The definition of \(G_\lambda\) gives
\[
 0\le G_\lambda\le v_\lambda(A)+r_z v_\lambda(\mathfrak D_K).
\]
This proves finite support and supplies one coefficient for
all primes, though the minimizing coefficient may vary.

For each embedding, expand that coefficient multilinearly
in the coefficient vectors \(z_{i,\gamma}\) of the columns.
Every determinant is a pairing of exterior products, bounded
by \(B^{r_z}\prod_i\|z_{i,\gamma_i}\|_{2,\sigma}\)
by Cauchy--Schwarz and Hadamard. Summing bounds over
\emph{all} tuples of indices therefore gives
\[
 \log |\sigma A|\le r_z\log B+
       \sum_i\log\left(\sum_\gamma\|z_{i,\gamma}\|_{2,\sigma}\right).
\]
For column \(i\) write \(a=T-s\). Its coefficient sum is
at most \((N e^{m_\sigma})^a\|t_i^0\|_{\sup,\sigma}\),
by expanding the light products, where \(m_\sigma\) bounds
all the light log suprema. Its mean is bounded in logarithm
by \(a(\log N+G_c S_L)+G_c S_0\). The weight size convention
gives \[
 \begin{gathered}
 aS_L+S_0=S_\kappa+a, \\ a\le S_\kappa, \\ a\log N\le P_{rw}a(S_L-1)\le P_{rw}S_\kappa.
\end{gathered}
\]
Use the product formula for \(A\) and
\(\mathrm N\mathfrak D_K=|\operatorname{disc}K|\).

Let \(B\) decrease to the transference bound. Since
\(\log d\le P_d\), \(\log n\le P_O\) and
\(\log|\operatorname{disc}K|/n\le P_O\), this proves
the asserted \(\langle G\rangle\le r_z H_z S_\kappa\).
For no boundary columns the determinant is one and the price
is zero.

\end{proof}

\paragraph{The complementary quotient and its CM determinant.}

The boundary-column estimate leaves the determinant of a complementary quotient. Its torsion-free lattice and CM evaluation give the second determinant contribution.

Write \(Q=M/(M\cap H_g)\) and \(\rho=r_q\). It is torsion-free
over \(\mathcal O_K\). Choose \(\theta_1,\ldots,\theta_\rho\)
out of the short spanning sections of weight \(\kappa\) to be
a rational basis modulo \(H_g\). Here the short-section price
for each mean log norm is at most \(H_\theta\). We explain
the coindex of their images. At the definite end \(\rho=0\)
and the price is zero.

\begin{proposition}[The complementary determinant cost]\label{prop:slice-complement-determinant-cost}
For the chosen \(r_q\) basis sections of the complementary
quotient, assume the CM coefficient and line bounds below.
For the stated \(P_Q\), the weighted quotient length satisfies
\[
 \langle L_Q\rangle\le2r_qP_Q S_\kappa.
\]
 The coefficient argument uses the same value-line
coordinates at finite and archimedean places. At the definite
lift the quotient and its cost are zero.
\end{proposition}

\begin{proof}
\proofstep{Coefficient functionals on the complementary quotient.}
Before the last lift, the functionals used in the coefficient
test extract a monomial in formal frame letters and a Laurent
coefficient on a CM arc. Evaluation is linear in the section.
The frame polynomial is homogeneous of degree \(A\), the sum of the nonnegative flag exponents, so a nonzero test has that total monomial degree. Its value is
measured in a one-dimensional space with the dual theta
generator to power \(n_\kappa\) (the central weight) and the dual differential
generators to the powers in this monomial. Thus an arc
functional \(\phi_i\) has offsets \(s_{iw}\) at finite places
and logarithmic norm offsets \(a_{i\sigma}\) at infinity.
The finite offsets have finite support. Over its coefficient
field the three CM estimates used are
\begin{align*}
 w(\phi_i(v))+s_{iw}&\ge-e_{w/\lambda} b_\lambda,\\
 \log|\sigma\phi_i(v)|+a_{i\sigma}
    &\le \log\|v\|_{\sup,\sigma|K}+C_{\rm arch},\\
 \bigg(\sum_w s_{iw}\log {\rm N}w
       -\sum_{\sigma} a_{i\sigma}\bigg)/[K':\mathbb Q]
    &\le B_{\rm deg}.
\end{align*}

In this formula \(v\in M\) has nonzero reading, and \(\sigma\)
runs over embeddings. The last inequality concerns the
degree of the value line. In particular its two terms must
use the same differential and theta coordinates. These
bounds apply to any conjugate of the test and persist
after a number-field extension: finite orders and their offsets
multiply by relative ramification, and the archimedean
data pull back along embeddings. The common baseline for
the allowed torsion order \(n'\) is
\[
 b_\lambda=A\max\{v_\lambda(|\operatorname{disc}E|),
                               v_\lambda(n')\}.
\]
It is nonnegative and its mean is at most
\(A(\log|\operatorname{disc}E|+\log n')\) by norming the
two principal ideals. A change of generator of the value
line changes its two offsets by a principal divisor, so
the degree inequality is unaffected.

The CM annihilation and detection of low coefficients give
common rational kernel \(H_g\) in the rational cusp space. To select independent rows on its quotient, include
all coefficient-field conjugates of the tests. Over \(K\)
the coordinate readings of a test in any basis of its finite
value field are linear combinations over a normal closure
of the conjugate readings, by the invertible embedding
matrix.
For a value field of degree \(n\), a basis \(b_1,\ldots,b_n\) and its distinct embeddings \(\sigma_i\), the inversion uses
\[
  \det\bigl(\sigma_i(b_j)\bigr)_{i,j=1}^n\ne0 .
\]
 The coordinate readings have trivial common kernel
on the rational quotient. Select a finite spanning subfamily
of its dual and put the finitely many fields involved inside
one normal extension \(K'/K\). There are then \(\rho\)
independent conjugate readings \(\phi_i\) with the same
bounds as above. This selection requires no bound for
\([K':K]\).

\proofstep{The determinant of the chosen CM evaluations.}
Set 
\[
  d_\theta=\det(\phi_i(\theta_j))\ne0.
\]
For a local basis of \(Q_\lambda\) lift its elements to
\(M_\lambda\). Clearing denominators away from \(\lambda\)
preserves the finite inequality. The determinant on this
basis has order at least \(-\sum_i s_{iw}-\rho e_{w/\lambda}
b_\lambda\). Changing to the \(\theta_j\) adds precisely
\(e_{w/\lambda} L_{Q,\lambda}\) by the elementary divisor
theorem. Averaging over the places above \(\lambda\) gives
\[
 L_{Q,\lambda}\le n_\lambda+\rho b_\lambda,\qquad
 n_\lambda=
 \frac{\sum_{w\mid\lambda} f_{w/\lambda}
                    (w(d_\theta)+\sum_i s_{iw})}{[K':K]} .
\]
This signed quantity has finite support. Bounding only its
positive part would spoil the ensuing product formula.

At \(\sigma\) factor \(e^{-a_{i\sigma}+C_{\rm arch}}\) from
row \(i\) in the determinant estimate, then use the column
norm bound \(\sqrt\rho\|\theta_j\|_{\sup,\sigma|K}\).
Hadamard followed by the product formula for \(d_\theta\)
therefore gives
\[
 \langle n\rangle\le
 \rho\big(H_\theta+C_{\rm arch}+B_{\rm deg}
                              +\tfrac12\log\rho\big).
\]
When \(\rho=0\) use the empty determinant. This gives the
length on completing \(Q_\lambda\), since flatness preserves
the kernel and the finite-length quotient. In particular
both the baseline and the degree price are taken at this
lift. There is no induction on earlier local coindices.

\proofstep{The final cost in the selected rank.}
For the run under consideration \(A\le S_\kappa\), and
the selected radius and torsion estimates bound
\(C_{\rm arch}\), \(B_{\rm deg}\) and \(\langle b\rangle\)
each by \(C_0 z^4 S_\kappa\). The bounds use \(r\ge3\),
the good level, \(\log {\rm N}\mathfrak c_K\le L\) in
the size data, the positive central and last weights,
\(S_\kappa\le C_K kfb\), \(T,c>0\) and \(Tc\le1\).
For the cutoff parameter fixed at seven one may use
\(K_*=20+|C_A|+7|C_A|+|C_K|\) and

\[
  C_0=17+2K_*+200K_*^4.
\]
 The selected dimension bound
gives \(\log\rho\le1041z^3\) when \(\rho\ge1\).
With \(H_\theta\le z^4 S_\kappa\) the two estimates above
give \[
 \langle n\rangle\le\rho P_Q S_\kappa,\qquad \langle b\rangle\le P_Q S_\kappa,
\] with
\(P_Q=(1+3C_0+1041)z^4\). Their sum gives the required
\(2r_q P_Q S_\kappa\) for the fourth price. The degree uses the same line as the two coefficient bounds.
The numerical comparison for the radius and torsion order
is at the chosen weight and cutoff on this lift. \label{cor:slice-budget-2}

\end{proof}

\paragraph{Absorbing all four terms at the selected scale.}

The four bounds can now be compared at the scale fixed for both orientations. The numerical inequalities below show how their sum fits the allowed remainder.

Put \(P=C_2z^{A_2}\) with \(z=fr\log b\), and put \(P_D=H_z,\ P'_Q=2P_Q\), or increase these to the nonnegative majorants for the comparison. The numerical conditions used here are
\begingroup\small
\begin{equation}\label{eq:slice-remainder-scales}\begin{aligned}
 T S_L&\le2 S_\kappa,&
 |w_F|&\le(s+2)S_L+5fr^2T,&Tc_+g_0^*&\le1,\\
 60G_c+180P_H+44P_D+22P'_Q&\le P^2s,&
 5G_cfr^2+P_d&\le P^3,& P_\eta&\le P,\\
 S_\kappa&\le C_K kfb,&4C_K&\le P^2 f,&
             P^3 T^2&\le b,\\
 P^4s^3kf&\le Tm\delta_2,&
 1\le m&\le k,&f,T,g_0^*&\ge1 .
\end{aligned}\end{equation}
\endgroup
Here \(0<\delta_2\le1\) and \(P\ge87\). In the pre-final
case \(c_+=c\); omit the auxiliary term in the final case.
The coefficient bounds are polynomial bounds in the same
\(z\). For example the cutoff used for the field, integral
section and quotient constants includes \(z\ge
C_L+|C_K|+|C_q(C_A,C_K)|\) and \(z\ge10^9\).
It controls \(C_q z^4\) by \(z^5\). Once these fixed
constants are chosen the cutoff is independent of the run
place, component, lift, short sections, \(V\) and \(h\).
\label{cor:slice-budget-3}

\begin{established}[One order and precision for the two orientations]
\label{hyp:common-slice-scales}\label{par:slice-scales}
\hypertarget{source-slice-scales}{}
Take the selected rank of Lemma~\ref{lem:rank-window}. Write
\(L=x^\epsilon,\ b=e^L,\ z_0=frL,\ D_*=f\binom r2\).
With the exponents \(A_2=12,\ C_p=1\) of
Lemma~\ref{lem:window-parameter-choice}, put \(P=C_2z_0^{A_2}\)
for the fixed precision constant \(C_2\). On both runs, at every lift
and all sufficiently large selected scales, the order \(T=\lceil
x^{377}\rceil\) satisfies
\[
 \frac{P^4D_*^3kf\,x^{C_p}}{m\delta_2}\le T,\qquad
 P^3T^2\le\sqrt b,\qquad
 (D_*+2)n_W+2r\le T,\qquad 40(D_*+2),\,11D_*^2\le T .
\]
The cut has dimension parameter \(s=D_*\), and
\(n_P=8fr+4\le n_W\).
\label{cor:slice-scales-1}
\end{established}
\begin{proof}
For \(p+q=r\) the unitary base has dimension \(fpq\); the flags in
the two summands contribute \(f\binom p2+f\binom q2\).
Their sum is \(D_*\), including at \(q=0\).
Lemma~\ref{lem:rank-window} chooses
\[
 n_W=8\left\lceil\frac{10r}8
         \left(f+\frac{f^2(k-1)^3L^2}{m}\right)\right\rceil,
 \qquad r+t=n_P+Tn_W .
\]
In particular \(n_W\ge10fr\ge n_P\). No order is chosen anew after
the rank \(t\) has been fixed.

For an explicit comparison put \(y=x^{1/100}\). The bounds in that
lemma give \(f\ll y^{216}\), \(k,r\ll y^{470}\), \(m\ge y^{49}\),
\(L\le y^2\), and \(n_W\ll y^{2267}\). The nonempty set of larger
indices has density \(\delta_2\ge 1/n\), where \(2n=f\). Thus
\(z_0\ll y^{688}\), \(D_*\ll y^{1156}\), and
\[
 \frac{P^4D_*^3kf\,x}{m\delta_2}
 \ll_{C_2} y^{48\cdot688+3\cdot1156+470+216+100-49+216}
 =y^{37445}\ll x^{377}.
\]
The remaining lower bounds for \(T\) have still smaller exponents.
Finally \(\log(P^3T^2)=O_{C_2}(\log x)\), whereas
\(\log\sqrt b=x^\epsilon/2\). All quantities in this comparison are
fixed before the orientation and lift are chosen.
\end{proof}

\begin{lemma}[Size bounds in both orientations]
\label{lem:common-slice-scales}
For the data of Lemma~\ref{hyp:common-slice-scales}, both orientations satisfy the finite boundary estimate, the allowance conditions in \eqref{eq:slice-remainder-scales}, and the inclusion at every weight coordinate
\[
 T(w_L)_{\tau,i}\le\kappa_{\tau,i}\qquad\text{for all }\tau,i.
\]

For fixed choices of the constants in the polynomial \(P\), the
threshold on \(x\) is common to all lifts in the two runs.
\end{lemma}

\begin{proof}
\proofstep{Compatibility of the weight inclusions.}
The weight inclusion is coordinatewise:
\(T\) times the light rows is bounded by the selected
\(\kappa\) rows, so \(TD_L\le D_\kappa\).
The weight size has the differential degree, a central term
\(fr\mu_j n\), and a constant one, with \(1\le\mu_j\le r\).
The three central parameters at \(w_L,w_F,\kappa\) are
\(n_W,n_P,Tn_W+n_P\). The constant term in the weight size gives \(TS_L\le S_\kappa+T\).
Using \(T\le b\) and the head-weight bounds \(b\le S_\kappa\le7kfb\),
\[
 TS_L\le S_\kappa+T\le S_\kappa+b\le2S_\kappa .
\]

For the short weight, the row-gap
remainder after subtracting \(T w_L\) is at most
\((s+2)\) times the light gap plus \(3T\). For the
terminal gap of the large block the adjustment
\((s+2)n_W+2r\) is at most \(T\), giving that same remainder
bound. Summing the gap comparisons from the last row in
each block bounds the differential degree by
\((s+2)D_L+3fr^2T\). The short central term is at most
\(fr^2 n_P\le fr^2T\). Including the constant term gives
\(|w_F|\le(s+2)S_L+5fr^2T\).
\label{cor:slice-scales-2}

\proofstep{Room for the finite boundary term.}
For the finite boundary cost write \(Q=\min(g_0,g_1)\),
\(R_7=(\operatorname{rd}E)^7\). The record uses
\[
             c=R_7P D_L^2/Q,
             \qquad g_0^*=1250r .
\]
The level-size input is \(Q\ge(kb)^8\).
To see the large margin in this comparison, let
\(B=\sqrt b\). The rank and weight data give
\(r\le2k^2\), \(f,m\le k\) and \(D_L\le2frb\le4k^2B^2\),
where \(f\le k\) and \(r\le2k\) give the second bound.
Also \(R_7\le B\), \(P^3T^2\le B\), \(TP\le B\)
and \(B\ge10^{100}\). Thus
\[
\begin{split}
 R_7^4P^3D_L^2 &\le16 k^4 B^9,\\
 Tg_0^*R_7PD_L^2 &\le40000 k^6 B^6,\\
 2fm^3 R_7^3 P D_L^2 &\le32 k^8 B^8 .
\end{split}
\]
Each is at most \(k^8 B^{16}\le Q\).

The middle bound
gives \(Tc g_0^*\le1\), and the last pays the auxiliary
cusp allowance at the same order. The short-basis
comparison for the \hyperref[par:slice-finite-depth]{finite readings} uses
\(R_0\le R_7\) and \(N_0^4n_W+2N_0^3\le P\).
These give its two pointwise inequalities for \(c\).
Also \[
 J_0\le2\sqrt{N_0}\le3k,\qquad D_LJ_0\le12k^3 b<Q.
\] The degree estimate
\(2D_\kappa+2<g_i\), \(TD_L\le D_\kappa\) and
\(T\ge40(D_*+2)\) give the level conditions for the
Hodge--Tate comparison. The inequality \(4D_L+4\le g_i\)
also follows directly from \(D_L\le4k^2b\) and
\(g_i\ge(kb)^8\).
\label{cor:slice-scales-3}

\proofstep{Checking the four coefficient allowances.}
The coefficient bounds for the four prices use the
polynomial domination at this same \(z_0\), including
the auxiliary, height and quotient constants at their
fixed exponents. Its threshold is chosen once after
those constants have been chosen. The record supplies
\(z_0\ge x\) from \(x\le fr\) and \(L\ge1\), so a
single lower bound on \(x\) ensures the required
polynomial domination for every place and for all
lifts on either run.
In particular the single threshold \(x_*\) chosen for the comparison has the implication
\[
  x\ge x_*\quad\Longrightarrow\quad z_0\ge x_*.
\]
 We use the construction in the specified data on the common family, including the
level interval and positive density of the indices with larger weights, as a
size input here. This argument derives the two-run
inequalities used by the index on that record.
\label{cor:slice-scales-4}
\end{proof}


\begin{proposition}[The remainder at the selected scale]\label{prop:scaled-geometric-remainder}
On either selected run with \(0<\delta_2\le1\), \(P\ge87\),
the bounds \eqref{eq:slice-remainder-scales}, and the common
choice in Lemma~\ref{lem:common-slice-scales}, the weights
and branch lattices in Proposition~\ref{prop:local-product-colength}
satisfy
\[
 \langle r\rangle\le R_j fmb\,\delta_2.
\]

The choice of scale precedes the place and the tested subspace.
\end{proposition}

\begin{proof}
To check the consumption of these inequalities, use
\(R\delta\le d\) and the rank bounds in the four prices.
The light part of \(H_s\) above, after separating
\(5G_c fr^2T\), is at most
\((s+4)(s+1)(3G_c+9P_H) S_L\).
For instance the coefficient of \(G_c\) before this
simplification is
\((s+2)+2(s+4)(s+1)+(s+1)\).

For the \(P_H\) terms use \(P_{rw}\le P_H\) and
\(\tfrac43s(s+4)+\tfrac23s+3(s+1)
 \le3(s+4)(s+1)\). With \(s\ge2\),
\((s+4)(s+1)\le10s^2\). Division by \(d>0\) now gives
\begin{align*}
 \langle r\rangle/d\le{}&
 (60G_c+180P_H+44P_D+22P'_Q)s^2 S_\kappa/T\\
 &+(5G_cfr^2+P_d)T+\tfrac12 P_\eta T c_+ .
\end{align*}
For the first summand, combine the allowance
\(P^2s\) with \(S_\kappa\le C_K kfb\) and obtain
\[
 \frac{P^2s^3S_\kappa}{T}
 \le\frac{C_K}{P^2}\,mb\delta_2
 \le\tfrac14 fmb\delta_2 .
\]
The first inequality uses \(P^4s^3kf\le Tm\delta_2\);
the second uses \(4C_K\le P^2f\).
For the second summand use its coefficient bound \(P^3\)
and \(P^3T^2\le b\).

Since \(Tf m\delta_2\ge4\), the
resulting allowance \(P^3T\le b/T\) is at most
\(\tfrac14 fmb\delta_2\). Finally \(Tc_+g_0^*\le1\)
and \(P_\eta\le P\) give
\(\tfrac12P_\eta T c_+\le P^3T\).
Adding these three estimates proves the asserted remainder bound.
Rank zero has zero lengths. This yields the weighted
remainder in \eqref{eq:geometric-run-index} for either
selected run with its own degree normalization. In
particular the paired trace calculation can use its
\(12fmb\) allowance without a factor for the number of
branches, primes or choices of a subspace.
\end{proof}



\begin{lemma}[Extension of the run field and restriction to a subspace]
At \(v\mid\lambda\) let \(e_1=e(v/\lambda)\),
\(e_2=\operatorname{ord}_\lambda(p)\), and put
\[
 h=\max\{0,\lceil e_2T(\nu_*+\mathbf1_{\rm bad})\rceil\},\qquad E=e_1h.
\]
For every subspace \(W\) of the row over the completed coefficient
field, write \(M_W=M\cap W,\ B_W=B_v\cap W\). Then
\begin{equation}\label{eq:run-cut}
 \operatorname{length}_v(M_W/B_W)\le E\dim W+r_v,\qquad r_v=e_1r_\lambda .
\end{equation}
Writing \(\omega_v=\log{\rm N}v/[K:\Q]\), the selected remainders
satisfy \(\sum_{v\in S}\omega_v r_v\le2R_jfmb\delta_2\)
for every finite \(S\), including its conjugate.
\end{lemma}

\begin{proof}
Over the completed ring at \(\lambda\), put \(N_0=M_0\cap\Theta\).
It has finite colength by Proposition~\ref{prop:geometric-run-index}.
Take a Smith basis with exponents \(d_i\). The span of the basis
vectors for which \(d_i>h\) is an admissible test space:
\(\nu_\lambda=e_2\nu_*\) by \eqref{eq:slice-finite-baseline}.
Subtract \(h\) times its dimension in \eqref{eq:geometric-run-index}:
\[
 \sum_i(d_i-h)_+\le r_\lambda .
\]
Flat extension preserves \(N_0\), computed as the kernel of the map
\(M_0\to (M_0+\Theta)/\Theta\). Its extension lies in \(B_v\) and
has exponents \(e_1d_i\) in \(M\). For a uniformizer \(\pi_v\),
\[
 \operatorname{length}_v((\pi_v^EM+B_v)/B_v)
 \le\sum_i(e_1d_i-E)_+\le e_1r_\lambda .
\]
In \(M_W/B_W\) the image of \(\pi_v^EM_W\) embeds in this quotient.
The quotient by that image is a quotient of \(M_W/\pi_v^EM_W\),
of length \(E\dim W\). This proves \eqref{eq:run-cut} for \(W\)
defined over the larger field as well.

Finally \(\sum_{v\mid\lambda}e_1\omega_v=\log{\rm N}\lambda/[K_0:\Q]\).
Each cover therefore contributes at most \(R_jfmb\delta_2\) over
all places. Restrict each of its nonnegative sums to the places
assigned to it and add. The cut uses the same core charts and bad
primes on both covers. \label{cor:paired-row-indices-3}
\end{proof}
\begin{lemma}[Length on a mirabolic character space]
Let \(P\) be the finite mirabolic group, \(U\) its upper unitriangular
subgroup, \(\psi\) its nondegenerate character and \(I=[P:U]\).
Suppose \(M,B_v\) are \(P\)-stable. For the good subspace \(W\)
in the \((U,\psi)\)-space of dimension \(\dim W=D_0\),
\[
 I\,\operatorname{length}_v(M_W/B_W)
     \le E I D_0+r_v+R_j v(|U|).
\]
\end{lemma}
\begin{proof}
Over one residue field use the mirabolic fixing the last row
and the additive character of the sum of the first superdiagonal.
Index the rows starting at zero. If
\(g\) is outside the upper unitriangular group, choose the last
row \(s\) where its diagonal or lower entries differ from those of
the identity, and choose \(l\le s\) with \(g_{sl}\ne\delta_{sl}\).
The fixed last row gives \(s+1<r\). With column vectors denoted
by \(e_i\), put
\[
 u=1+a(ge_l)e_{s+1}^{\mathsf T},\qquad
 u'=1+ae_l(e_{s+1}^{\mathsf T}g).
\]
Both are upper unitriangular by the choice of \(s\), and
\(ug=gu'\). Their superdiagonal sums differ by
\(a(g_{sl}-g_{s+1,l+1})=a(g_{sl}-{\bf1}_{s=l})\).
This ranges over the field, so a choice of \(a\) distinguishes
the characters. Rescaling nonzero simple-root coefficients
amounts to diagonal conjugation, and inversion preserves the test.

For \(g_i\in P/U\) (one representative per coset) put
\[
 e_i=g_i\left(|U|^{-1}\sum_{u\in U}\psi(u)^{-1}u\right)g_i^{-1}.
\]
On \(g_iW\) this is the identity. On \(g_kW\), \(i\ne k\), an
element of the intersection of the two conjugate subgroups acts by
different characters on \(g_kW\) and on the image of \(e_i\).
It commutes with \(e_i\), so \(e_i(g_kW)=0\). For a product
group, take this element in a component in which the cosets differ.
It follows that \(W_*=\bigoplus_i g_i W\) has dimension \(I D_0\).
For \(X=M,B_v\) put \(\Lambda_X=\bigoplus_i(X\cap g_iW)\).
Since \(|U|e_i\) preserves \(B_v\),
\[
 |U|(B_v\cap W_*)\subseteq\Lambda_{B_v}\subseteq B_v\cap W_* .
\]
The quotient on the right has length at most \(R_jv(|U|)\).
The inclusion \(\Lambda_M\subseteq M\cap W_*\) and
\eqref{eq:run-cut} now give
\[
 \begin{aligned}
 I\,\operatorname{length}_v(M_W/B_W)
 &=\operatorname{length}_v(\Lambda_M/\Lambda_{B_v})\\
 &\le \operatorname{length}_v(M_{W_*}/B_{W_*})+R_jv(|U|)\\
 &\le E I D_0+r_v+R_jv(|U|).
 \end{aligned}
\]
\label{cor:paired-row-indices-4}
\end{proof}

To pass from the unipotent character to the chosen twist we must
also account for the scalar \(\sigma_v\) in the theta lattice.

\begin{auxlemma}[The unnormalized character trace]\label{lem:unnormalized-row-character}
Let \(U\) be a finite group acting by \(\rho\) on a module \(V\) over
a commutative ring \(R\), and let \(\psi:U\to R^\times\) be a character.
Put \(V_\psi=\{y:\rho(a)y=\psi(a)y\ \text{for all }a\in U\}\).
The integral operator
\[
 T_\psi=\sum_{u\in U}\psi(u)\rho(u)^{-1}
\]
maps \(V\) into \(V_\psi\) and restricts to multiplication by \(|U|\) there.
\end{auxlemma}

\begin{proof}
For \(a\in U\) substitute \(v=ua^{-1}\). Then
\[
 \rho(a)T_\psi=\sum_{u\in U}\psi(u)\rho(ua^{-1})^{-1}
              =\sum_{v\in U}\psi(va)\rho(v)^{-1}
              =\psi(a)T_\psi.
\]
If \(y\in V_\psi\), each term \(\psi(u)\rho(u)^{-1}y\) is \(y\).
\end{proof}

\begin{lemma}[The integral unipotent trace on a scaled row]
\label{lem:scaled-row-unipotent-trace}
Fix the common scaled theta lattice \(B\) and its scalar
\(\sigma_v\) in the level comparison of
Lemma~\ref{lem:selected-form-level-tie}. For the finite unipotent
group \(U\) and the assigned character \(\psi\) at that tuple,
let \(W\) be the good subspace of the \(\psi\)-space. Let
\(J_{\mathrm{tr},v}\) be the local lattice generated by traces of
unscaled short theta orbit families. The comparison in the integral coordinates is
\[
 \frac{|U|}{\sigma_v}(B\cap W)\subseteq J_{\mathrm{tr},v}.
\]

\end{lemma}

\begin{proof}
Apply Lemma~\ref{lem:unnormalized-row-character} to the level action.
Write \(L_{\mathrm{short}}\) for the span of the unscaled short theta generators.
We claim \(T_\psi L_{\mathrm{short}}\subseteq J_{\mathrm{tr},v}\). Indeed for each \(u\) right-translate its
level coordinate and take the series of that translated datum.
The level action carries it back to the original series.
These series are cusp forms by positivity of the last
weight. The translation identity has trivial rational twist
on this run, and its determinant factor is one because the
unipotent matrices have all diagonal entries one at the
half-level places. Thus their character-weighted sum is
precisely the trace of the generator in cusp coordinates.
Linearity proves the claim. For \(y\in B\cap W\), the scaling
puts \(y/\sigma_v\) in \(L_{\mathrm{short}}\), and the character law gives
\[
 (|U|/\sigma_v)y=T_\psi(y/\sigma_v)\in J_{\mathrm{tr},v}.
\] These steps use the weight and level
readings of the chosen run.\label{cor:paired-twist-index-1}

\end{proof}

\paragraph{The integral comparison away from the distinguished prime.}
\label{par:off-prime-row-trace}
\hypertarget{source-paired-off-prime-trace}{}
Fix the run used at \(v\) and keep its base weight, base
character and half level. Denote its distinguished rational
prime by \(p_\circ\). In the argument let \(\lambda\) be any
finite place of \(K\) of residue characteristic different
from \(p_\circ\). The character of the finite trace is the
one on the arithmetic row, namely \(\psi_0^{-1}\) in the
base tuple convention. A generator of the source trace
lattice is
\[
 t=\sum_{u\in U}\psi(u)\operatorname{cusp}(f_u).
\]
Here \((f_u)_{u\in U}\) is one orbit family of short theta
rules.

The orbit family has a common unimodular source lattice, an
admissible odd level denominator, a common harmonic
polynomial and a common residue datum. The labels of the
rules are obtained from that datum by the prescribed level action. The target module \(J\) is the integral span, in
these cusp coordinates, of the base kernel traces. An
admissible kernel trace uses a finite split-supported
source translate and a coefficient in its integral
harmonic lattice. Its form coordinate is specified by
evaluation at the lowest target types. We prove that
\(t\in J_\lambda\).
\label{cor:paired-off-prime-trace-1}

\begin{proposition}[Integral kernel trace away from the distinguished prime]
\label{prop:off-prime-trace-inclusion}
At each place of residue characteristic different from
\(p_\circ\), the orbit trace lattice maps into the kernel theta
lattice in the common form coordinates. For the fixed coefficient character on all finite vectors, \eqref{eq:off-prime-trace-identity} reads
\[
 \mathcal A(t')(\mathbf w)=\sigma(c)\,T_\psi^{\rm ad}S_{\Phi_{\rm fr}}(h,y;\mathbf w)\quad(\text{all }\mathbf w),\qquad v(c)=0.
\]
 The inclusion also holds after the good
projector, on the same local row.
\end{proposition}

\begin{proof}
\proofstep{Integral trivialization on the adapted orbit.}
First adapt the lattice of the orbit to \(\lambda\). The
congruence adjustment of a self-dual lattice makes it
standard locally over the residue characteristic of
\(\lambda\). It preserves the standard lattice at the
old bad and level primes and is the identity modulo the
level there. In the second adjustment, write its finite
adelic lattice translate as \(\gamma h k\). Here \(\gamma\)
is rational unitary, \(h\) has split support away from
the forbidden primes and is identity at infinity, and
\(k\) preserves the standard lattice and its congruence
data. For this factorization use rational unitary
approximation with a common denominator supported
only at split primes outside the prescribed set.
Write \(g\) for the adjusted finite matrix after removal
of \(\gamma\).

Let \(e\) be the finite-adelic idempotent of
its allowed split support, including both conjugate halves.
Set
\[
 h=eg+(1-e)1,\qquad k=e1+(1-e)g.
\]
Since \(e(1-e)=0\) and \(\bar e=e\), the products split into
their two idempotent components. Writing \(*\) for the unitary involution,
\[
 hk=eg+(1-e)g=g,\qquad
 h^*h=e\,g^*g+1-e=1,\qquad k^*k=e+(1-e)g^*g=1.
\] Off the allowed support \(g\) lies in the integral
congruence stabilizer, so \(k\) preserves its lattice and
level marking.
This factorization is used for the conjugated
lattice because the coefficient convention contracts
through conjugate source coordinates. Thus the adjusted
polynomial lattice is the conjugate of the lattice read
at \(h\). The representative has a common denominator
prime to the old bad and level primes and to the residue
characteristic of \(\lambda\). The two rational unitary
changes push forward the polynomial and keep the residue
label at the level.

The invariance of the harmonic
equations and the short theta transformation with even
source rank show that all the short rules change by the
same root of unity \(\zeta\in K_0\).
The common scalar leaves the coefficient valuation unchanged, since
\[
  \zeta^n=1\ \text{for some }n\ge1\quad\Longrightarrow\quad v_\lambda(\zeta)=0.
\]
 The common factor
comes from the normalized theta multiplier, which is
independent of the translated residue label. The multiplier \(\zeta\) is in the base field \(K_0\), so it factors out of the trace on the whole orbit. For the adapted orbit and its trace one obtains
\[
        f'_u=\zeta f_u,\qquad t'=\zeta t,\qquad
        \mathcal O_{K,\lambda}\,t'=\mathcal O_{K,\lambda}\,t .
\]
The last equality uses that both \(\zeta\) and its inverse are integral. The adjustment can
increase the level denominator away from the indicated
primes, where it has no effect on the local inclusion.
\label{cor:paired-off-prime-trace-2}

For the adapted polynomial \(P\), choose a rational integer
\(b\) prime to the residue characteristic of \(\lambda\)
that carries each of the adapted and standard lattices
into the other. If \(D_P\) bounds its total degree, then
\(b^{D_P}P\) has integral coefficients in the standard
letters. In degree \(a\) the lattice inclusions in the two directions give, on the tensors in the source,
\[
 b^a L_{\rm ad}^{\otimes a}\subseteq L_{\rm std}^{\otimes a},
 \qquad b^a L_{\rm std}^{\otimes a}\subseteq L_{\rm ad}^{\otimes a}.
\]
Dualizing reverses the inclusions and again costs at most \(b^a\) on coefficients. The same coefficient substitution in the
polynomial realizes these inclusions after imposing the
harmonic equations. For homogeneous degree \(a\le D_P\), a denominator cleared by \(b^a\) is also cleared by the same power for the whole polynomial, since
\[
  b^{D_P}=b^{D_P-a}b^a,\qquad D_P-a\ge0.
\]
 Homogeneity and the harmonic
equations are preserved. The prescribed highest-weight coefficient map gives a lower Fock vector for this
polynomial in the standard integral model.

The highest-weight coefficient map is defined over the common row field, identified with the
coefficient field of the lower vector through its chosen
embedding. Transport into the adapted coefficient lattice
can cost another power of \(b\). To see the support of
this denominator, at every place over the residue
characteristic the two source lattices coincide, so the
integral symmetric tensors and their differential
conditions coincide there as well. Outside those places
only finitely many primes require clearing. Descent by
the coefficient isomorphism and the integral highest-weight
map then supply a lower vector \(y\) for \(P\) and an
integer \(s_1\notin\lambda\) with
\[
                 s_1 y\in\mathcal C(h).
\]
Here \(\mathcal C(h)\) is the prescribed integral coefficient lattice, written in the lower Fock coordinates. The
product of the two denominator clears is allowed in
\(s_1\). This step uses the harmonic lattice on the
adapted source, and the coefficient lattice obtained from
the chosen highest-weight model.
\label{cor:paired-off-prime-trace-3}

\proofstep{The adelic theta evaluation.}
We describe the adelic evaluation used for this generator.
Let \(\chi_{\rm tuple}\) be the run's tuple character and
\(\chi_{\rm aux}\) its finite-order auxiliary character.
For the base kernel the splitting character is
\[
                    \chi_{\rm split}=\chi_{\rm tuple}\chi_{\rm aux}^{-1}.
\]
It is trivial on real-field ideles and has the prescribed values of the tame determinant character over the selected half of \(p_\circ\).
The type residues are the prescribed full-rank ones. The half
places are disjoint from their conjugates and from the
type places. The self-dual frames give the base kernel
as a product of a vacuum and the indicated residue
functions. Write \(\Phi_{\rm fr}\) for the finite vector
with free lower block at the twist places and the
specified full-rank residue matrices at the type places.
These full-rank residues are determined by the adapted
orbit's label. At the twist places the upper \(r\)-by-\(r\)
block is invertible and receives the inverse determinant
character.
\label{cor:paired-off-prime-trace-4}

For a finite vector \(\Phi\), let
\(S_\Phi(h,y;\mathbf w)\) be its single adelic theta series
with coefficient 
\[
  y_{\C}=c'_0(y)
\]
 and the tuple \(\mathbf w\) of
lowest target vectors. Here \(c'_0\) is the coefficient realization for this choice of levels. It sends the algebraic harmonic vector
to the complex coefficient representation using the fixed embedding.
Thus \(S_\Phi(h,y;\mathbf w)\) reads \(y_{\C}\) in the series,
and clearing \(s_1 y\) at \(h\) gives the coefficient
\(s_1 y_{\C}\) in the realized lattice. Denote the adelic realization of the coefficient row by \(\mathcal A(-)(\mathbf w)\). Let \(f\) be the single adapted short rule. With the
fixed choice of \(\Phi_{\rm fr}\) the extremal comparison
is the identity of adelic forms
\[
          \mathcal A(f)(\mathbf w)=c\,S_{\Phi_{\rm fr}}(h,y;\mathbf w),
          \qquad c\in K,\quad s_2c\in O_{K,\lambda},\quad s_2\notin\lambda .
\]
The coefficient and target tuple in this identity are the
ones selected above.
In this comparison one unfolds the geometric class sum
using the fixed adelic lift.

At the extremal tuple the
coefficient map differentiates exactly the letters of
the short polynomial, with the polarization sign and the
common frame and cover factors. The base-tuple class sums
span the cusp space, so the frame comparison applies to the sections defined by short theta rules on the whole row. To obtain that
span, choose a residue point using the surjective row
pairing at each split half-level place and use the
Chinese remainder theorem.
For the disjoint prime-power half-level conditions, choosing the local residue point amounts to using the surjection
\[
  {\cal O}_E\longrightarrow\prod_{\mathfrak p\ {\rm used}}{\cal O}_E/\mathfrak p^{n_{\mathfrak p}} .
\]
 The stable short span can
then be read at that point. Each raw term is moved to a
split-supported self-dual class by the short theta
transformation. The one-class kernel formula there gives
the geometric class sum times the frame factors. These
factors are nonzero, and linear independence of the cusp forms
descends the resulting span to the construction field.
\label{cor:paired-off-prime-trace-5}

\proofstep{Correct the level character in the adelic evaluation.}
The central level translation absorbs the difference
between the two frame normalizations. The base phase is
one, and the quotient of the two cover values and the
remaining central normalization are roots of unity in the field of values.
Thus if \(\xi_1,\xi_2\) denote the indicated normalization factors, their contribution at the chosen place is
\[
  v_\lambda(\xi_1\xi_2)=v_\lambda(\xi_1)+v_\lambda(\xi_2)=0 .
\]
 Their multiplication preserves the
local integrality test on \(c\). When the class
scalar is nonzero its expression, up to the preceding
units, is
\[
    B(\Lambda)^{-1}\Psi(\mathfrak a)
       [\Lambda:\Lambda\cap\Lambda_{\rm st}]^{r/2}.
\]
Here \(B(\Lambda)\) is the base genus scalar,
\(\Lambda_{\rm st}\) is the standard lattice, and
\(\Psi(\mathfrak a)\) is the unramified ideal value
of the combined tame character at the moved class.
Choose the determinant representative with
\(\mathrm N\mathfrak a\) prime to the residue
characteristic, using the local equality of the
adapted and standard lattices there.

In fact the determinant
lattice is the quotient of \(\mathfrak a\) by its conjugate,
and the two integral ideals are coprime. Order zero in
the determinant lattice at all places over the residue
characteristic forces order zero in both ideals there.
The local genus comparison gives \(B(\Lambda)\) order zero
at the same places, by comparing the theta normalization
in the common integral frame.
Consequently both directions of the genus-scalar comparison at such a place have order zero:
\[
  v_w\big(B(\Lambda)^{-1}\big)=-v_w\big(B(\Lambda)\big)=0 .
\]
 Its inverse becomes integral
after multiplication by an integer prime to that
characteristic. For the tame ideal value, take a power
of \(\mathfrak a\) that is generated by an element
congruent to one modulo the character conductor, using
the finite ray class group. Its character value is then
a product of integral powers of the embeddings of that
generator.

The generator divides a power of
\(\mathrm N\mathfrak a\) as an algebraic integer, so a
power of the norm clears the negative exponents.
Integrality descends from a positive power by
transitivity of integrality.
In particular, if \(a^n\) is integral, then \(a\) satisfies the monic equation
\[
  X^n-a^n=0.
\]
 The square of the positive
index factor is a rational integer. Their product is
therefore integral after multiplication by one integer
\(s_2\notin\lambda\). The computation takes place after extension to the chosen complex embedding, and integrality descends to \(K\) by
the injective field embedding. If the coefficient
is zero use \(s_2=1\).
The determinant-one condition on \(U\) makes the
remaining finite weight character trivial on the level
translates. Thus the same \(c\) applies to each term
\(f'_u\), with the unitary level action by \(u^{-1}\) on
the target series.
\label{cor:paired-off-prime-trace-6}

\proofstep{Extending the identity to all finite vectors.}
Both sides of this identity are multilinear in the
lowest target vectors and equivariant separately at
each archimedean place. The compact-group translates of
the nonzero extremal vector span the lowest type at
that place, by irreducibility of the lowest type.
Translate and extend linearly in each variable. Taking
the \(\psi\)-weighted sum gives
\begin{equation}\label{eq:off-prime-trace-identity}
 \mathcal A(t')(\mathbf w)
    =\sigma(c)\, T_{\psi}^{\mathrm{ad}}
                   S_{\Phi_{\rm fr}}(h,y;\mathbf w)
 \quad\text{for all }\mathbf w,
\end{equation}
where the adelic trace uses the same action and summation
as the row trace. The embedding \(\sigma\) here is the
one fixed for the adelic evaluation. The form coordinates
are determined by these tests. Indeed the smooth
archimedean module is generated by its lowest types, and
their translates on the target recover the full
equivariant form. The coordinate map sends its basis to linearly independent forms, hence is injective. In particular the vanishing test for the form coordinates reads
\[
  c=0\quad\Longrightarrow\quad t'=0.
\]
\label{cor:paired-off-prime-trace-7}

Finally turn the free-block vector into a translate
of the base kernel \(\Phi_{\rm b}\). At a chosen twist
place scale the lower \(n_V-r\) rows in the split model
by an element of valuation one, with the inverse adjoint
on the conjugate half. The upper \(r\) rows are fixed.
The original kernel imposes zero lower residues. After
the substitution this is automatic on integral entries.
Conversely, if the scaled matrix is integral with the
required zero lower residues, every entry before scaling
is integral.
For a lower-block entry \(a\) and local uniformizer \(\varpi\), the support equivalence is
\[
  \varpi a\in\varpi{\cal O}_v\quad\Longleftrightarrow\quad a\in{\cal O}_v .
\]
 The upper block determinant and its
character are preserved. This proves the equality of
the two support tests including the integrality test.

At a type prime, complete the prescribed \(r\) columns to an
invertible \(n_V\)-by-\(n_V\) matrix \(B\) over the residue field.
Since \(n_V>r\), column \(r+1\) is among those added in the completion.
Multiplying that column by \((\det B)^{-1}\) gives determinant one and
preserves all \(r\) prescribed columns. Lift the change to the split unitary
stabilizer by using the inverse adjoint at the conjugate
place. It fixes the vacuum and the coefficient lattice
there. The type and twist operations have disjoint
supports. They give the identity of finite vectors
\[
      \Phi_{\rm fr}=\sigma(c')\,\omega(k,1)\Phi_{\rm b}
\]
with a finite source translate \(k\) and nonzero \(c'\in
K\). In the split Weil formula, the determinant factors
of the twist operation use the local values of
\(\mu_0\chi\eta^{-1}\) at the chosen uniformizers.

\proofstep{The local determinant valuation.}
For each unit of determinant valuation in the forward
scaling, the norm factor is the inverse \(r\)th power
of the square root of the residue norm. The normalized
base character reading supplies the other inverse
square root. For residue cardinality \(Q\), the two scalar factors therefore combine as
\[
  Q^{-r/2}Q^{-1/2}=Q^{-(r+1)/2},\qquad (r+1)/2\in\mathbb Z.
\]
The translate uses the inverse scalar. The tame
character values and their inverses are integral
after multiplication by a power of \(p_\circ\), by
the same unramified ideal argument for the base
characters and finite order for the auxiliary
character.

At \(w\nmid p_\circ\) the two clearing conditions on a nonzero tame value \(\alpha\) imply
\[
  v_w(\alpha)\ge0,\qquad v_w(\alpha^{-1})\ge0\quad\Longrightarrow\quad v_w(\alpha)=0 .
\]
 The supported residue norms are powers
of \(p_\circ\). The type operation has determinant one.
The coefficient lattice also changes only at
\(p_\circ\), where a large enough power carries the
required coefficient lattices into one another. Hence
there is an \(m\ge0\) such that \(p_\circ^m c'\) is
integral and \(p_\circ^m s_1 y_{\C}\) is admissible at
\(hk\). Increasing \(m\) preserves both properties.
\label{cor:paired-off-prime-trace-8}

By the Weil representation law, the series with
\(\omega(k,1)\Phi_{\rm b}\) at \(h\) is the series of
\(\Phi_{\rm b}\) at \(hk\). For \(c\ne0\) let \(e\) be the
form coordinate of the latter kernel trace with
coefficient \(p_\circ^m s_1y_{\C}\). Equation
\eqref{eq:off-prime-trace-identity} and the determining
property of the tests identify it with
\[
               e=\frac{p_\circ^m s_1}{cc'}\,t'.
\]
It belongs to \(J\) by the admissibility just proved.

In the original orbit coordinates the identity reads
\[
       \bigl(p_\circ^{2m}s_1s_2\bigr)t
        =\bigl(\zeta^{-1}(s_2c)(p_\circ^m c')\bigr)e .
\]
The coefficient on the right is integral, and the
factor on the left is a unit at \(\lambda\). This proves
\(t\in J_\lambda\), also in the zero case. The argument
is uniform in the sense needed for the index comparison:
it incurs zero local order at every \(\lambda\) of
residue characteristic different from \(p_\circ\).
The clearing powers may depend on the generator and
the place. There is no such charge at the coefficient
place \(v\) because the run pin chooses \(p_\circ\) from
the other rational prime. Pull back along the common
form-coordinate isomorphism to get the inclusion on
the preprojected row. Applying the good projector
gives the inclusion on the good row.
\label{cor:paired-off-prime-trace-9}
\end{proof}


\begin{lemma}[The traced image after good projection]
\label{lem:projected-orbit-trace}
At a place assigned off its own prime, the projected kernel
\(eJ_v\) is a full lattice in the good space \(W\).
Writing \(T_v\) for the orbit trace lattice, the equality before taking signed indices is
\[
 e(J_v+T_v)=eJ_v.
\]
Their signed indices therefore agree.
\end{lemma}

\begin{proof}
Let \(e\) be the good projector. It fixes \(W\).
The localized inclusion also supplies the lower rank bound
for \(eJ\). The
unscaled short lattice is full in the ambient cusp space.
For each \(y\) in the rational character space, fullness supplies a
nonzero integral scalar \(a\) with \(ay\) in the short lattice. By
Lemma~\ref{lem:unnormalized-row-character},
\[
 y=(a|U|)^{-1}T_\psi(ay)
 \quad\text{in the coefficient fraction field}.
\] Therefore the trace lattice
spans the character space, as does \(J\) by the local
inclusion and the character condition on its generators.
The good projector is an idempotent commuting with the
level action. Its image on the character space is exactly
the intersection of that space with its good range.

Consequently \(eJ\) spans \(W\) and is a lattice there as
soon as \(J\) is finitely generated.
If \(V\) denotes this character space, the generic-span test is
\[
  K(eJ)=e(KJ)=eV=W .
\]
 At the selected level
the kernel trace is determined by the finite list of source
classes, and the span of its admissible coefficient vectors
on each class is finitely generated. Linearity gives finite
generation of the kernel traces and of their projected images.
All these modules use the coordinates of
Paragraph~\ref{par:selected-form-coordinates}.
The target projector is the pullback of the retained
finite-group idempotent by the form map. The source
coordinates at \(v\) are those at the assigned tuple \(\mu(v)\),
where the form tie reads the source projector and the class-sum
map with phase one.

Write \(\operatorname{Tr}_{v}\) for the
orbit trace lattice and \(J_v\) for the kernel theta lattice
on the untwisted target. The local unit gives
\((\operatorname{Tr}_{v})_v\subseteq (J_v)_v\), so
\[
       \big(e(J_v+\operatorname{Tr}_{v})\big)_v=(eJ_v)_v .
\]
Localization commutes with sums and with the image of a linear
map. Thus the global lattice sum used to make the containment
integral has the original signed index where it is read.
The conjugate place uses the same tuple by \(\mu(\bar v)=\mu(v)\)
and the same run by equality of residue characteristics.
\label{cor:paired-twist-index-2}
\end{proof}

\begin{proposition}[Twist comparison at the assigned tuple]
\label{prop:paired-twisted-row-index}
For the common good space on the assigned tuple at \(v\), the
projected kernel lattice \(eJ\) and the projected twisted lattice
\(\mathcal X_3\) satisfy \eqref{eq:paired-twisted-row-length}. The excess is bounded over
the lifts with the same ramification normalization; for each conjugate pair over \(p\),
\[
 \sum_j(\epsilon_{wj}+\epsilon_{\bar wj})\le 2/p^2.
\]

The local integral image comparison passes to the good projector
at that same tuple.
\end{proposition}

\begin{proof}
\proofstep{Estimating the excess growth under twisting.}
We now compare \(eJ\) with the projected twisted lattice
\(\mathcal X_3\). For full lattices \(A,B\) in one space over
a DVR define the signed length
\[
 [A:B]=\operatorname{length}(A/(A\cap B))-
        \operatorname{length}(B/(A\cap B)).
\]
For three full lattices choose \(N\subseteq A\cap B\cap C\) of finite
colength and write \(\ell\) for length. The common quotient cancels:
\[
 [A:B]+[B:C]
 =\bigl(\ell(A/N)-\ell(B/N)\bigr)+\bigl(\ell(B/N)-\ell(C/N)\bigr)
 =[A:C].
\] The row is mapped semilinearly into its realization over the fraction field of a theta DVR
\(O'\) over the twisted coefficient ring. The images of a
coefficient-field basis are independent there.

This map intertwines \(e\) with the good projector in the theta
space, and the \(O'\)-spans of the images of the untwisted and twisted generators are the spans of the corresponding
theta series of the same source classes. Hence taking images
under the projector gives the extended lattices of projected
theta vectors in both cases. This assertion also holds after
pulling back by the row-coordinate isomorphism, by surjectivity
and compatibility with the projector. These are exactly the
generator and linear-independence hypotheses for the row realization.\label{cor:paired-twist-index-3}

Let \(e'\) be the order of a uniformizer of the coefficient
DVR at \(v\) in \(O'\). Since the embedding induces \(v\),
\(e'>0\) and 
\[
  \operatorname{ord}_{O'}(p)=e'e_v.
\]
 Signed
lengths on the two extended spans also multiply by \(e'\).
For a pair of full lattices take a common sublattice and
diagonalize each inclusion into the two lattices. The read
basis remains independent and a diagonal factor of order \(a\) has order \(e'a\) upstairs.
Writing \(A',B',N'\) for the extended lattices gives
\[
 [A':B']=\ell(A'/N')-\ell(B'/N')
        =e'\bigl(\ell(A/N)-\ell(B/N)\bigr)=e'[A:B].
\]
Equivalently one can first establish commensurability upstairs.
Each relative quotient commutes with extension to \(O'\)
by flatness of a torsion-free module over a DVR.
Each finite-length quotient downstairs contributes \(e'\)
times its length, and arbitrarily long chains downstairs
remain arbitrarily long upstairs by faithful flatness.
Thus finiteness of both lengths descends as well.

\proofstep{Summing the determinant comparison over the lifts.}
Apply the row determinant estimate of
the \hyperref[subsec:retained-twist]{simultaneous twist construction} at the chosen power of the
ray character. The twisted projected span contains the
specified twisted generators and spans the prescribed good
space. The untwisted index supplies its baseline. Their
difference upstairs divided by \(e'e_v\) is therefore at
most the row excess \(\epsilon_{vj}\). By the preceding
calculation and additivity, for the reference lattice \(M_g\)
on the good space we get
\begin{equation}\label{eq:paired-twisted-row-length}
 [M_g:\mathcal X_3]_v\ \le\ [M_g:eJ]_v+e_v\epsilon_{vj}.
\end{equation}
For this application the initial theta row has full good
rank, and its twisted integral realization and the projected generators use the same character as the determinant test.
The character construction bounds a nonnegative family of
excesses for all \(m\) rows and both sides by \(1/p^2\)
in total. The selected rows map injectively into
\(\{1,\ldots,m\}\), and taking their prescribed side
costs at most that total.

At the conjugate place the
required signed character is among the conjugates of the same specialization. Its excesses have the same
bound. Consequently
\[
             \sum_j(\epsilon_{wj}+\epsilon_{\bar wj})\le 2/p^2 .
\]
We give the integral direction of the image comparison.
Write \(M_t\) for the source level lattice at the chosen tuple,
\(\theta\) for the rational class-sum map and \(e_X,e\) for
the two good projectors. The source coordinates identify the
rational level module with the character space of class forms.
Indeed the level idempotent has this range. Integral
coordinates are tested by membership of the class values in
their prescribed coefficient lattices. The finite class set
and injectivity of evaluation give a full lattice on the
character space.

\proofstep{The integral map at each actual tuple.}
For one tuple generator take its translated finite representative
and its integral source polynomial there. In the class-sum
formula use the inverse representative to place the polynomial
on its class, then trace over its stabilizer to make the value
covariant.
For a stabilizer \(G_c\) the resulting class value is invariant by reindexing the finite sum
\[
  h\sum_{\gamma\in G_c}\gamma P=\sum_{\gamma\in G_c}h\gamma P=\sum_{\gamma\in G_c}\gamma P\qquad(h\in G_c).
\]
 With the stabilizer mass in the class sum this gives
the single series, up to the level character of the change of
representative. This is precisely the trace comparison for
the source value lattice and kernel with its prescribed level
action. Writing \(f\) for this integral class form and \(\chi_U\)
for the selected character, take the unnormalized trace
\[
       T_{\chi_U} f=\sum_{u\in U}\chi_U(u)^{-1}u f .
\]
The coefficients are integral and the translates preserve
the value lattices, so \(T_{\chi_U}f\in M_t\).

Equivariance gives \(\theta(T_{\chi_U}f)=T_{\chi_U}\theta(f)\).
This is the level trace of the single series, with the same
root of unity as in the class comparison. This equality can be tested on
all lowest-component vectors in the faithful complex realization.
For a nonzero image, take a nonzero coordinate of the two
rational vectors in a rational basis. Their ratio belongs
to the coefficient field and is that root of unity, hence
an integral unit.
For a coordinate \(y_i\ne0\) and its phase multiple \(x_i\), this reads
\[
  \zeta=x_i/y_i\in K,\qquad v_w(\zeta)=0 .
\]
 It can be incorporated in the original
lattice. A zero generator already belongs to the image.
Consequently projecting the tuple generators gives the
integral inclusion \(\mathcal X_3\subset e\theta(M_t)\).
\label{cor:paired-twist-index-4}

For the reverse span, let \(f_c\) be the
stabilizer-averaged numerator for a class \(c\) with stabilizer
\(G_c\). Thus 
\[
  f=\sum_c |G_c|^{-1}f_c
\]
 over the coefficient
field. The selected character condition and Lemma~\ref{lem:unnormalized-row-character} give \(T_{\chi_U}f=|U|f\), and hence
\[
 e\theta(f)
   =\frac1{|U|}\sum_c\frac{e\theta(T_{\chi_U}f_c)}{|G_c|}.
\]
Each trace in the numerator is an integral level form whose
theta value is in the tuple span. Thus \(e\theta(M_t)\) lies in the rational span of
\(\mathcal X_3\). Its rational span is the good space on the
base lift-spanning condition.

In detail, at the faithful
complex embedding the theta lifts of the chosen level
sources span the base good space. Their class sums give the same vectors up to the common nonzero phase.
A rational matrix acquires the same rank after embedding,
so the projected class-sum map surjects rationally onto
\(W\).
For the rational class-sum matrix \(A\), this descent is the equality
\[
  \operatorname{rank}_K A=\operatorname{rank}_{\mathbb C}\sigma(A).
\]
 This uses the simultaneous level basis, the class-sum
identity, and the base lift span for this tuple. Both modules
are now lattices in \(W\), the larger one finitely generated
by the image of \(M_t\) and the smaller one by Noetherianity.

\proofstep{Equivariance after good projection.}
Finally the matched kernel intertwines the good projectors,
\(\theta e_X=e\theta\). On \(\operatorname{im}e_X\) the
difference between a vector and its chosen lift projection
\(\operatorname{pr}\) belongs to \(\ker\theta\). Therefore

\[
  \theta\operatorname{pr} e_X(M_t)=e\theta(M_t),
\]
 an
identity on the image of the integral source itself. This
last step uses the two projector identities on the matched
level data. The same source tuple serves at the conjugate
place by its specified conjugation rule.
\label{cor:paired-twist-index-5}
\end{proof}


\begin{lemma}[Signed index after the twist]
With the sign convention of Section~\ref{subsec:division-length},
the transported index satisfies
\[
 \frac{\mathcal I_v}{e_v}
 \ \le\ (d_v+u_v)D_0+
              \frac{r_v/e_v+R_j u_v}{I}+\epsilon_{vj},
 \qquad
 d_v=\frac{E-v(\sigma_v)}{e_v},\quad
 u_v=\frac{v(|U|)}{e_v}.
\]
Moreover \(R_j/I\le3D_0\).
\end{lemma}
\begin{proof}
Divide the preceding length estimate by \(Ie_v\). The scalar in
the theta reading subtracts \(v(\sigma_v)D_0/e_v\), and the trace and
twist comparison adds \(u_vD_0+\epsilon_{vj}\). This gives the
displayed sign for an index positive on a smaller lattice.
Counting the strictly upper triangular entries gives
\(u_v=fr(r-1)/2\) at the two level primes and zero elsewhere.
If \(d_0\) is the character-space dimension, the rank estimates give
\[
 (1-y)^2(1-\delta)R_j
 \le (1-y)(1-\delta)|U|d_0
 \le I(1-\delta)d_0\le I D_0.
\]
Since \(y\le1/8\) and \(\delta\le1/4\), the coefficient on the left
exceeds \(1/3\). \label{cor:paired-row-indices-5}
\end{proof}

We explain the character and group-order bounds used here.
Over the residue field \(\mathbb F_Q\) of cardinality \(Q\), the mirabolic subgroup is
\(\operatorname{GL}_{r-1}(\mathbb F_Q)\ltimes\mathbb F_Q^{r-1}\).
Counting successive independent columns and using

\[
  |U_Q|=Q^{r(r-1)/2}
\]
 shows that its index divided by
\(|U_Q|\) is \(\prod_{a=1}^{r-1}(1-Q^{-a})\). This formula
also covers the empty product. For the product over the
two level primes the ratio \(I/|U|\) is at least
\(1-\sum_Q 1/(Q-1)\), since the sum of all the powers
\(Q^{-a}\) at one place is at most \(1/(Q-1)\). There
are at most \(2f\) places, and each \(Q\) is at least
\(\min(g_0,g_1)\). The common size record has \(f\le k\) and
\(\min(g_0,g_1)\ge(kb)^8\), so
\(2f/(\min(g_0,g_1)-1)<1/8\).

\begin{lemma}[The character dimension at the definite lift]
\label{lem:paired-character-share}
For the definite two-level cusp model at the selected weight, write \(V\) for the full cusp space of dimension \(R_j\). Its character multiplicity at the given level is
\[
 \mu=(\chi|_U)\psi_0^{-1},\qquad \dim V[\mu]=\frac{R_j}{|U|}.
\]

\end{lemma}

\begin{proof}
\proofstep{Free orbits at the definite lift.}
For the dimension bound use the two-level model of cusp
forms for this same lift. Scalar extension through the
selected complex embedding preserves the ambient dimension
and the dimension of the character space. In fact the
images of the coordinate basis are independent and span,
and the joint eigenspace is the kernel of the stacked
matrices \(\rho(u)-\mu(u)\) for \(u\in U\).
Equivalently, if \(V\) is the model space, the selected subspace is
\[
  V[\mu]=\bigcap_{u\in U}\ker\big(\rho(u)-\mu(u)I\big).
\]
 All their
nonzero minors stay nonzero. Use the model character
\(\mu=(\chi|_U)\psi_0^{-1}\), including the scalar
twist on the model action.

At the definite end the
subgroup acts freely on the components at this level.
A stabilizer would give a unitary rational matrix at a
balanced frame, unipotent modulo the specified split prime
and bounded in the positive definite norms. The
small-element estimate at that level forces it to be the
identity, and faithfulness gives the identity subgroup
element. The direct sum of the component spaces is therefore
a sum of regular representations of \(U\), by identifying
a representative fiber with each of its translates. Its
\(\mu\)-space has dimension \(R_j/|U|\).
\label{cor:paired-character-share-1}

\end{proof}

\subsubsection*{Equivariant realization of the selected row as cusp forms}
\label{subsec:selected-cusp-reading}

An index computed in geometric coordinates controls the selected row only through an equivariant identification. We exhibit the cusp-form realization on which the local cuts and the level actions agree.

Fix a lift $j$ and the run used at a coefficient place $v$.
Write $K_0$ for the construction field of the run, $K$ for
the common row field and $C_j$ for the rational cusp space
at its prescribed weight $\kappa$. The form coordinates
identify $C_j$ with a $K_0$-subspace of $K^{R_j}$ and
identify $K\otimes_{K_0} C_j$ with $K^{R_j}$. If \(b_1,\ldots,b_d\) is the common basis in the row data over \(K_0\), scalar extension preserves every coordinate independently:
\[
  K\otimes_{K_0}\bigoplus_{i=1}^dK_0 b_i=\bigoplus_{i=1}^dK(1\otimes b_i).
\]
 The chosen complex place determines an embedding
$\sigma_0:K_0\hookrightarrow\mathbb C$ and an extension
$\sigma:K\hookrightarrow\mathbb C$. Use the bijection
between the two-level components and the run's components.
At corresponding components the Hermitian data, frame
domains, embeddings and weights agree. Restrict the holomorphic forms to these frame domains.

The ordinary stable-range theta calculation in
Section~\ref{subsec:theta-row-rank} supplies the span before
projection. Here the target has rank $r$ over $E$ of
degree $2f$, the positive definite source has even rank
$n_V\ge 8fr+4$, and the level is odd with chosen split half
coprime to its conjugate.
The degree and rank conditions in the stable range combine to give
\[
  n_V\ge 8fr+4=4[E:\mathbb Q]r+4 .
\]
 Use the run's prescribed first
character Gram comparison and its admissible weight,
including positivity of the last weight. The chart
independence and the short theta transformation specialize
the stable unfolding to readings of raw short theta
sections. Every such short section is cuspidal at this
weight.

On each component extend a bounded holomorphic
section by zero off its frame domain, as a function on
all frame coordinates. The extension belongs to the
space on which the theta span is taken. The stable unfolding
expresses its family of components in a finite complex
span of the read short sections. Restricting the resulting
identity to the domains and transporting by the component
bijection gives the desired expansion of the original
family on the two-level model. The construction uses
the positive bad modulus and the Gram comparison of the
actual run.\label{cor:paired-row-reading-1}

\begin{proposition}[The equivariant cusp-form realization and the rank at the smaller level]
\label{prop:equivariant-selected-row}
For a lift \(j\) of either run in
Proposition~\ref{prop:linked-comparison-rows}, use its rational
cusp space \(C_j\) at weight \(\kappa\) over \(K_0\), with the
form coordinates of Definition~\ref{def:common-row-lattices}.
Extend along the complex embedding \(\sigma_0\) of \(K_0\). The chosen basis extends to a basis of the complex space of holomorphic cusp forms and identifies the corresponding character actions. The retained arithmetic projector is the same
one in the two realizations. For each lift the partition of the retained base labels gives the small-level image dimension
\[
 dZ_0(j)=D_0.
\]
 The full rank of the integral image uses the separate
off-prime trace inclusion.
\end{proposition}

\begin{proof}
\proofstep{Independence in the two coefficient realizations.}
For independence take a finite $K_0$-independent family of
cusp rules. Evaluate it on the framed test algebras through
all maps to $\overline{K_0}$. These evaluations detect a
relation of rules. Indeed the test algebras are reduced
and of finite type, so a nonzero test value survives at
a closed point over a finite extension of $K_0$. The span
of the evaluation rows is invariant under
$\operatorname{Gal}(\overline{K_0}/K_0)$, because composing
each test map with an automorphism retains it in the family.

If the row rank were deficient, the unique reduced row
echelon basis of the annihilator would be Galois invariant.
A nonzero basis row would give a $K_0$-relation among the rules. Hence some set \(I\) of rows, as large as the family of rules, gives a square matrix \(A_I\) with nonzero determinant. Under the chosen complex embedding \(\sigma\),
\[
  \det\big(\sigma(A_I)\big)=\sigma(\det A_I)\ne0 .
\]
 Framed base change realizes these tests
as complex frames. Uniformization and the common Borel
weight multiply all entries of an evaluation row by the
same nonzero scalar when passing to its frame-domain
representative. The resulting functions on frames are therefore independent. Restriction to the model domains preserves independence, since the functions vanish outside these domains and every component occurs.\label{cor:paired-row-reading-2}

\proofstep{The same rational character space on the two levels.}
Read the common $K_0$-basis in the two-level model and
extend its coefficients by $\sigma$. Independence and the
stable span give the isomorphism
\[
 \mathbb C\otimes_{K,\sigma} K^{R_j}\ \simeq\ V_\kappa .
\]
On a rational cusp form the lift of a level change acts on
the right in its moduli coordinate. In frame coordinates
this is the two-level pullback, with the same scalar and
the same change of component as in the run. The equality
on the basis therefore intertwines the full level action
after extension of scalars. In particular the dimensions
of the character kernel and its intersection with the kernel
of the good projector can be computed in $V_\kappa$.
Let \(A\) be the stacked character equations and \(P\) the good projector in these coordinates. Their common kernel is computed by the single matrix
\[
  \ker A\cap\ker P=\ker\begin{pmatrix}A\\P\end{pmatrix}.
\]
Changing the complex embedding preserves their algebraic
ranks.\label{cor:paired-row-reading-3}

\proofstep{Comparing with the retained arithmetic projector.}
The identification with the retained arithmetic rank uses
one more comparison at the selected point. In its row
coordinates the tie is the composite of the map to
$K\otimes_{K_0}C_j$ with the form-coordinate isomorphism
above. It identifies the subgroup of the arithmetic row
with the preimage of the base unipotent subgroup under the
level injection and identifies their characters there.
The determinant character in the coordinate change is
one on that subgroup. Consequently the arithmetic character
space maps bijectively to $V_\lambda$ after extension of
scalars.
Consequently the complex identification preserves the character rank, since scalar extension commutes with the common kernel
\[
  \bigcap_{u}\ker\big(\rho(u)-\mu(u)I\big).
\]
 The second row link uses the base-tuple theta
lift and its rational reading. In the split part of the
chosen realization these give the same reduced level
group and the same target projector, because the scalar
phase at the base is one.

The base-level lift identifies
the retained source space with
$V_\lambda\cap\operatorname{im}e_{\rm good}$ in the target
row. Its source dimension is the common $D_0$ of the
point, for every lift $j$. In the rank comparison the
base-tuple maps identify the dimensions of the summands
indexed by each retained irreducible type. Those summands
exhaust the respective spaces, so the sums of their
dimensions agree. Thus the target character rank
is the arithmetic $d_0$, and the retained rank is
\[
 D_0=d_0-\dim(V_\lambda\cap\ker e_{\rm good}).
\]
The equality uses invariance of the character space under
$e_{\rm good}$ and its decomposition into the image and
kernel of the restricted idempotent. In the integral
comparison all row lattices are transported by
this same tie. The scalar extension here computes ranks
and the lattice comparison uses localization in the
number field.\label{cor:paired-row-reading-4}

\proofstep{The rank in the determinant denominator.}
The small-level rank in the division denominator can be
read from the projection at the base tuple. The level
projection at the point's base index \(t_0\) is surjective
for each lift. The tuple readout sends \(t_0\) to \(\mu_0\),
since the base Hecke character is the identity and the
tuple characters distinguish the used indices.
On the matched row models the level projection induces the
theta adjoint \(\Theta_{\mu_0,j}^{*}\) by the class-sum identity
and the two level actions. Thus this adjoint surjects onto
the reference row space. Since

\[
  \ker\Theta_{\mu_0,j}=(\operatorname{im}\Theta_{\mu_0,j}^*)^\perp,
\]
\(\Theta_{\mu_0,j}\) is injective.

Decompose the reference row into the occurring class types.
The tied eigenspace multiplicity formula for each type gives
the sum of the row counts over that type's record labels.
These label sets partition the \(D_0\) base row labels,
and the class type summands form an internal direct sum.
Summing gives \(D_0\) at every lift. By injectivity the
small-level range \(\operatorname{im}\Theta_{\mu_0,j}\)
has that same dimension. This proves the identity
\(dZ_0(j)=D_0\) used in the cohomological division term.
The full rank of the projected \(J_v\) at a twisted index
follows separately from the local trace inclusion on the
assigned tuple.
\label{cor:paired-row-reading-5}
\end{proof}

\paragraph{The selected levels and the separating section.}

Equivariance identifies the row as a cusp space. Separating its level components also requires a section that distinguishes those components at the chosen embedding.

The components are those of the complex uniformization of the chosen construction at its fixed complex place. A component $c$ has lower
arithmetic group $\Gamma_c$ and its two prescribed lattices.
The change of level sends $c$ to $hc$ and has a hermitian
transporter $t(h,c)$ carrying both lattices to their counterparts.
The transporter multiplication is well defined modulo the lower group of the destination. For \(h,k\) in the level group it gives
\[
  t(h,kc)t(k,c)t(hk,c)^{-1}\in\Gamma_{hkc}.
\]
 A transporter fixing its
component lies in the lower group only for the identity level element.
The lower character on the arithmetic group is trivial. The
components have the common signature, CM embeddings and weight.
Use the equivariant cusp-form realization of
Section~\ref{subsec:selected-cusp-reading} at this weight. It identifies
the scalar extension of the whole row with the space of bounded
holomorphic sections on these components.
\label{cor:paired-cusp-ranks-1}

Let \(H\) reduce to upper unitriangulars at one chosen level prime \(g\), and let \(H_c\) stabilize the component \(c\). Take the upper group \(G\) generated by \(\Gamma_c\) and the transporters of \(H_c\). The transporter product law above and the trivial kernel on level elements give
\[
  1\longrightarrow\Gamma_c\longrightarrow G\longrightarrow H_c\longrightarrow1.
\]
 It preserves both lattices and is unipotent modulo every prime
of $E$ above $g$. The congruence on the chosen half of the primes
follows from the level change and the kernel congruence for $\Gamma_c$.
On the conjugate half, the reduction of a unitary matrix is the
contragredient with respect to the perfect paired form, so it too
is unipotent.
If \(A\) is unipotent, so is its inverse transpose on the conjugate half, since
\[
  A^{-\mathsf t}-I=-A^{-\mathsf t}(A^{\mathsf t}-I),\qquad (A^{\mathsf t}-I)^r=0 .
\]
 The prime $g$ splits completely in $E$ and exceeds $3r$.
These conditions apply to the character subgroup, the two half-level
subgroups used below, and the trivial subgroup.
\label{cor:paired-cusp-ranks-2}

\begin{lemma}[Separation at the actual two-level translates]
\label{lem:separating-level-section}
Let \(x_b\) denote the coset centers in the chosen cut and \(c\) the selected center. There is a cusp section \(s_c\) of the prescribed determinant weight whose values in the chosen frames satisfy
\[
 s_c(x_b)=\begin{cases}1,&b=c,\\0,&b\ne c.\end{cases}
\]
 The comparison uses the marked subgroups at the two
actual levels and the same coefficient realization as the
character calculation.
\end{lemma}

\begin{proof}
\proofstep{A cusp section distinguishing the level translates.}
Use a balanced majorant of the self-dual component lattice, in a
frame $x_c$ for the Hermitian form and that majorant. The property
needed of the majorant is that an endomorphism of the lattice
strictly contracting at every complex embedding is zero. The
semistable-majorant statement gives this property by applying the
degree inequality to the kernel and image of a contraction. The lattice has degree zero at the semistable majorant, and its
sublattices have degree at most zero. A nonzero strict contraction
would increase the degree of the quotient by its saturated kernel
when mapped isomorphically onto its image lattice. This is
incompatible with the degree bound for that image in the lattice.
Indeed for a nonzero strict contraction \(T\) the saturated kernel gives the incompatible inequalities
\[
  0\le\widehat{\deg}(L/\ker T)<\widehat{\deg}(T L)\le0 .
\]
 The property is invariant under the
right triangular change to an orthonormal frame.

Here is the small-element estimate in this frame. If the operator
norm of $\gamma$ is at most $e^{1/(3r)}$ at every embedding, then
for $1\le a\le r$,
\[
 \operatorname{tr}(\gamma^a)-r\in g\mathcal O_E,
 \qquad |\tau(\operatorname{tr}(\gamma^a)-r)|
       \le r(e^{1/3}+1)<g .
\]
The trace congruence uses nilpotence of the reduction of $\gamma-1$
at each prime over $g$ and squarefreeness of $g\mathcal O_E$. The
field norm now forces all the displayed trace differences to vanish.
Newton's identities give the characteristic polynomial $(X-1)^r$.

In an orthonormal basis the squared Hilbert--Schmidt norm of
$\gamma-1$ is at most $r(e^{2/(3r)}-1)\le e^{2/3}-1<1$, using
$\operatorname{tr}\gamma=r$.
The Hilbert--Schmidt computation uses the trace to cancel both linear terms, giving
\[
  \|\gamma-I\|_{\rm HS}^2=\operatorname{tr}(\gamma^*\gamma)-r\le r(e^{2/(3r)}-1)<1 .
\]
 Its operator norm is strictly below
one, so balance gives $\gamma=1$. The radius potential $R$ on
the positive-plane domain measures the change of majorant from the
orthonormal center. If two points of the ball $R<1/(49r^2)$ are
related by $\gamma$, their two Gram comparisons bound the norm
of $\gamma$ at the center by $e^{1/(3r)}$. This gives a ball whose
nonidentity upper arithmetic translates are disjoint. In the
definite case all upper arithmetic matrices are isometries in the
majorant, so the upper group itself is trivial.
\label{cor:paired-cusp-ranks-3}

For $0<p_0<r$ put $n=f r(r-1)/2$, $d_X=f p_0(r-p_0)$ and
\[
  s_0=2r+50r^2d_X+2,\qquad d=\tfrac12s_0(-1^{[p_0]},+1^{[r-p_0]}).
\]
Here the two signed blocks are repeated at all CM embeddings.
We use the separating section of weight $d$ at the free center.
For clarity, the relevant orbit interpolation estimate concerns an
upper arithmetic group $G$ containing $\Gamma_c$, a free ball
$R<\rho_0=1/(49r^2)$, and the weight potential $\varphi_d$.
The signed two-block calculation used in \hyperref[par:slice-interpolation]{the cut interpolation}
gives $\mathcal L(\varphi_d-\log D)\ge (s_0-2r+2)\mathcal L(R)$.
This margin pays for a radial weight $bR$ and for logarithmic
poles of order $d_X$ at the orbit, since
\[
  s_0-2r+2 > b+d_X/\rho_0.
\]

When \(p_0(r-p_0)\ge4\), use the exponent
\[
  b=4fr^2+1.
\]
 For the smaller signatures
one uses the sum estimate on the positive-plane domain with
$b=4(r-1)$, since the compact flag fibers contribute no factor in
this block-constant weight. In either range the corresponding choice of \(b\) gives a constant \(C\) for the individual model with
\[
  \sup_x\sum_{\gamma\in\Gamma_c}\exp\!\left(-\tfrac b2 R(\gamma x)\right)\le C.
\]
 These bounds use separation in
the orthogonal coordinates and integration of the invariant
volume density over the disjoint balls. All these constants are
for the individual model and weight. Only the stated weight shift
will be used across the family.
\label{cor:paired-cusp-ranks-4}

\proofstep{Interpolation with simultaneous vanishing on the other rows.}
On the flag chart take a smooth cutoff equal to one near the center
and supported inside the free ball. Place smoothed logarithmic
poles at the $G$-orbit of the underlying positive plane. The poles
are truncated to zero outside the disjoint balls. Their Levi loss
is at most $(1+\eta)d_X/\rho_0$ times $\mathcal L(R)$ for a fixed
sufficiently small $\eta>0$.
For the smoothed pole potential \(P\), the curvature cost on these balls is
\[
  {\cal L}(P)\ \ge\ -\frac{(1+\eta)d_X}{\rho_0}\,{\cal L}(R).
\]
 The smoothing parameter tends to zero, with the potential
uniformly bounded on the cutoff's derivative annulus.
Solve $\bar\partial$ of the cutoff with total weight
$\varphi_d-\log D-bR$ plus this logarithmic potential. The same
weighted estimate and frame extension as for the cut interpolation
apply. The solution has uniformly bounded weighted norm as the
smoothings tend to their poles.

In a small ball in the $d_X$ transverse
directions around any orbit point, the weight in the squared-norm integral
is bounded below by a fixed multiple of
$(|z|^2+\varepsilon)^{-d_X}$ times the inverse of the radial
factor. The radial mass diverges logarithmically. For \(d_X>0\) and
a fixed small \(\delta_0>0\), it satisfies
\[
 \int_0^{\delta_0}\frac{r^{2d_X-1}}{(r^2+\varepsilon)^{d_X}}\,dr
    \ \ge\ 2^{-d_X}\int_{\sqrt\varepsilon}^{\delta_0}\frac{dr}{r}
    \ \longrightarrow\ \infty .
\]
Near each orbit centre the cutoff is locally constant, so the
correction there is holomorphic. Its squared sphere mean bounds
its squared value at the centre. The uniform weighted norm bound,
divided by this diverging mass and with the radial weight evaluated
at that centre up to the common comparison constant, makes the
correction at the centre tend to zero. The
bound at each other $G$-orbit point has the additional factor
$e^{-bR/2}$ coming from the radial weight at that point.

Thus after
Poincar\'e summation, the cusp section has readings arbitrarily close
to one at the center and zero at all other coset centers.
For a sequence of such sums \(F_n\), with the representatives indexed by \(b\), the limiting readings have the form
\[
  \lim_{n\to\infty}F_n(x_b)=\begin{cases}1&b=c,\\0&b\ne c.\end{cases}
\]
 Summation
uses the orbit bounds just stated. The bounded section is cuspidal
in this block-constant weight with positive cross-block gap, by the
cusp translation and decay criterion. The map from cusp sections
to the finite tuple of readings has a closed image, being a linear
subspace of a finite-dimensional space. It contains the exact delta
tuple.
Writing \(E\) for the finite tuple evaluation and \(\delta_c\) for the desired delta tuple, closedness yields
\[
  \delta_c\in\overline{\operatorname{im}E}=\operatorname{im}E .
\]
 This gives the separating section on a component, extended by
zero on the other components. At a definite component the upper
group is trivial and the determinant power itself is nonzero there.
\cite{hormander-slice}\label{cor:paired-cusp-ranks-5}

\proofstep{Returning the separating section to the chosen components.}
For each $H$-orbit of components choose a representative $c$ and
its separating section. Sum these sections over the orbits. For the
sum $u_0$, the reading of $g^*u_0$ at $x_c$ vanishes unless $g=1$.
Inside the stabilizer this follows from the separation modulo
$\Gamma_c$, and outside it the component supports are disjoint.
At the transported frame of $kc$ the entry $(h,g)$ of the translate
matrix reads $(hgk)^*u_0$ at $x_c$. Solving the group equation gives the unique nonzero entry of row \(h\):
\[
  hgk=1\quad\Longleftrightarrow\quad g=h^{-1}k^{-1}.
\]
 The determinant therefore has a
nonzero value on every component, including when the group is
noncommutative. This construction is available independently for
all the indicated subgroups with the same weight $d$.
\label{cor:paired-cusp-ranks-6}

\end{proof}

\begin{established}[The flat cut and the determinant-weight gaps]
\label{hyp:three-determinant-weight-ranks}\label{cor:paired-cusp-ranks-7}
At an indefinite lift write \(n=f\binom r2\), \(T_0=T\),
\(L=w_L\), \(F=w_F\), and \(\kappa=F+T_0 L\). On the common
flat cut of \hyperref[par:slice-interpolation]{the cut interpolation},
\(n\ge2\), \(11n^2\le T_0\), and \(s_0=2r+50r^2d_X+2\le T_0\).
Number the flag coordinates \(0,\ldots,r-1\), negative block first,
and put \(k_{\rm vol}(i)=2i-(r-1)\). For
\[
                   U=F-k_{\rm vol}-(n+2)L
\]
the adjacent decreasing gaps inside each block are at least zero,
and \(U_b-U_a\ge1\) for every negative coordinate \(a\) and
positive coordinate \(b\). For \(L\) the corresponding lower
bounds are \(1\) and \(19T_0/10\), respectively.
\end{established}
\begin{proof}
The same light columns and auxiliary section give the flat
intersection for the three weights by
\hyperref[par:slice-interpolation]{the cut interpolation}.
Lemma~\ref{hyp:common-slice-scales} gives \(T_0\ge11n^2\);
\(d_X=fp_0(r-p_0)\le n\), and the polynomial bounds in its proof
give \(2r+50r^2n+2\le T_0\) at large scales.

Set \(d=n+2,\ T'=T_0+d\). For an internal gap \(a\) of the
unsigned row of \(\kappa\), the light gap \(l\) is defined by
\(a+2=T'l+\rho,\ 0\le\rho<T'\).
Lemma~\ref{hyp:height-section-weights} gives \(l\ge1\).
Since \(F=\kappa-T_0L\), after reversing and negating the negative
block the increasing shift in the volume weight contributes two:
\[
 a+2=T'l+\rho,\qquad k_{\rm vol}(i+1)-k_{\rm vol}(i)=2,\qquad
 U_i-U_{i+1}=a-(T_0+d)l+2=\rho\ge0 .
\]

At an unsigned terminal entry write
\(a-d_0=T'l+\rho\), where the two offsets \(d_0\) are \(dn_W+2r\)
and zero. Lemma~\ref{hyp:height-section-weights} gives \(a\ge2T_0^2\)
and \(T'\le41T_0/40\). Consequently
\[
 l>\frac{2T_0^2-T_0}{41T_0/40}-1\ \ge\ \frac{19}{10}T_0
 \quad(T_0\ge 100).
\]
The light rows have nonnegative central shifts of opposite signs,
so every positive coordinate minus every negative coordinate has
at least this gap. The unsigned terminal entries of \(F-dL\) are
\(d_0+\rho\). In the difference of the last positive and the first
negative flag coordinate the central shifts of \(F-dL\) contribute
\(n_P-dn_W\); subtracting \(k_{\rm vol}\) costs \(2(r-1)\).
The resulting difference for \(U\) is at least
\[
 (dn_W+2r)+0+n_P-dn_W-2(r-1)=n_P+2\ge1 .
\]
Internal monotonicity makes this the smallest cross-block difference.
\end{proof}

\begin{proposition}[Dimensions at the three determinant weights]
\label{prop:three-selected-weight-ranks}
Use \(n\ge2\), \(T_0\ge11n^2\), and one common flat cut
of Lemma~\ref{hyp:three-determinant-weight-ranks}, with
\(\delta_c\) points on component \(c\). Put
\(\delta=\sum_c\delta_c\). Then the cusp dimensions satisfy
\[
 \dim V_{\kappa-d}\ge\delta\binom{T_0+n}{n},\qquad
 \dim V_{\kappa}\le\delta\binom{T_0+2n+6}{n},\qquad
 \dim V_{\kappa+2d}\le\delta\binom{T_0+3n+10}{n}.
\]
In particular \(\dim V_{\kappa+2d}\le 2\dim V_{\kappa-d}\). The dimensions used in the comparison
are those of the prescribed cusp spaces.
\end{proposition}

\begin{proof}
\proofstep{The lower rank by interpolation.}
For the lower count, subtract $d$ from $F$. In the interpolation
weight $F-d-k_{\rm vol}-nL$ the cross-block coefficient is at least
$1+3.8T_0-s_0$. This exceeds $4fr^2$, since $fr^2\le4n^2$.
Subtracting one further $L$ leaves the cross-block coefficient
positive, with all within-block coefficients nonnegative. These
are the curvature and additional margin required by the cut
interpolation with zero additional potential. It gives bounded
sections $Q_j$ of weight $F-d$ whose readings at the cut are the
coordinate unit vectors.

For each reading $j$ and $|\alpha|\le T_0$,
form
\[
   Q_j\,w_0^{T_0-|\alpha|}\prod_{a=1}^n S_a^{\alpha_a}.
\]
At least one light factor is cuspidal, including when $|\alpha|=T_0$.
In a linear relation, take the least surviving total degree at the
points and use the local coordinates $S_a/w_0$. Their leading
coefficients read the corresponding $Q_j$ alone. They vanish only
when all the coefficients at that degree vanish, so the family is
independent. For $\delta_c$ cut points the lower dimension at
$\kappa-d$ is $\delta_c\binom{T_0+n}{n}$.
\label{cor:paired-cusp-ranks-8}

\proofstep{The upper rank by injecting into jets.}
For the upper count, the jet map at these same points is injective
through degree $T_0+n+6$. For \(f_0\) of weight \(F+vL\), put
\[
 Q=\frac{\|f_0\|^2\|\theta_2\|^2\|\theta_1\|^2\|\eta_0\|^{2/a}}
        {\big(\sum_{i=1}^n\|S_i\|^2\big)^{v+n+6+1/2}}.
\]
The two theta weights are $(n+4)L-F$ and $2L-d_{\rm sign}$,
and the auxiliary weight is $a d_{\rm sign}$ for the determinant
sign weight $d_{\rm sign}=(-1^{[p_0]},+1^{[r-p_0]})$.
The effective numerator weight is
\[
 (F+vL)+((n+4)L-F)+(2L-d_{\rm sign})+d_{\rm sign}
               =(v+n+6)L .
\]
At a cut point write \(\varrho\) for the size of the coordinate
tuple \(S_i/w_0\).

If the jets of \(f_0\) through degree \(v+n+6\) vanish, its
squared norm vanishes to order at least \(2(v+n+7)\). The cut rows are
local coordinates, so their sum of squared norms is comparable to
\(\varrho^2\). In this frame,
\[
 Q=O\!\left(\varrho^{\,2(v+n+7)-2(v+n+6+1/2)}\right)
   =O(\varrho)\qquad(\varrho\to0).
\]
At a noncompact end the lower bound for the light sum and the
numerator decay give
\[
  Q\ll 
       \exp\bigl((-4/g_0+2c(v+n+6+1/2)) P\bigr).
\]
For \(v=T_0\) the coefficient of \(P\) is negative by
\(T_0cg_0\le1\) and \(n+6+1/2<T_0\). A positive maximum could only occur
where the factors in the numerator are nonzero and the light sum
is positive.

At such a point the logarithm of the quotient is a
sum of pluriharmonic terms, a negative multiple of
$\log\sum_a|S_a|^2$, and $\varphi_L/2$. The numerator contributes zero Levi form at this point.
The map \(\lambda:z\mapsto[S_1(z):\cdots:S_n(z)]\) has target
dimension \(n-1\) and therefore has a nonzero tangent vector
\(\xi\) in the kernel of its differential. Along \(\xi\) the
pullback of the Fubini--Study form vanishes, so
\[
   \mathcal L(\log\text{quotient})(\xi,\bar\xi)
        =\tfrac12\mathcal L(\varphi_L)(\xi,\bar\xi)>0 .
\]
The last inequality uses the positive curvature of the selected
light weight on the flag chart. A smooth function with a local
maximum has nonpositive Levi form there. Thus all-zero jets force
\(f_0=0\). The dimension at $\kappa$ is
at most $\delta_c\binom{T_0+2n+6}{n}$.

\proofstep{The second upper weight shift.}
The cut contains a point on each component under consideration.
The bounded base theta section of weight $F$ has a nonzero value
at some frame. At that frame the nonvanishing light family gives
\[
  |w_0|^2+\sum_{i=1}^n|S_i|^2>0.
\]
 Multiply the base section by its
$T_0$th power to obtain a nonzero bounded section of weight
$F+T_0L$. The jet injection just proved applies to that bounded
section and would force it to vanish if the cut were empty.

For the upper count at $\kappa+2d$, the weight $(n+4)L-2d$
has positive within-block and cross-block differences. In the
curvature test subtract both $nL$ and the density weight. After both subtractions the remaining cross-block margin satisfies
\[
  7.6T_0-2s_0-2(r-1)>4fr^2.
\]
 Subtracting a further $L$
leaves nonnegative differences. These bounds give bounded
section interpolation on the same cut. Multiplication by one such nonzero
bounded section is injective by holomorphic continuation, and
puts the result at $F+(T_0+n+4)L$.

The maximum argument applies
there with jet degree $T_0+2n+10$. The extra $n+4$ in the
end estimate uses $2n+10.5<T_0$ and $T_0cg_0\le1$, so
$2c(v+n+6.5)<4/g_0$ for $v=T_0+n+4$. Hence the upper count at $\kappa+2d$ is
$\delta_c\binom{T_0+3n+10}{n}$. The component upper bounds
sum by restriction, and the lower bound sums by extension by zero.

For the upper shift we have
\[
 \frac{\binom{T_0+3n+10}{n}}{\binom{T_0+n}{n}}
 \le \exp\left(\frac{n(2n+10)}{T_0+n}\right)<2
 \quad(n\ge2,\ T_0\ge 11n^2).
\]
The loss in the lower shift relative to the unshifted count is
used below. The numbers of cut points and the interpolation
constants disappear from the ratios.
\label{cor:paired-cusp-ranks-9}
\end{proof}


\begin{proposition}[Uniform proportion of the selected character space]
\label{prop:paired-character-share}
In the same two-level model, let \(V[\mu]\) be the character space for \(\mu=(\chi|_U)\psi_0^{-1}\). On the selected family,
\[
 \dim V[\mu]\ge(1-y)\,\frac{R_j}{|U|},\qquad y\le\tfrac18.
\]

Lemma~\ref{lem:paired-character-share} gives the definite case.
For the other lifts use the separating section and the rank bounds
of Lemma~\ref{lem:separating-level-section} and
Proposition~\ref{prop:three-selected-weight-ranks}.
\end{proposition}

\begin{proof}
Write \(P_\mu=\sum_g\mu(g)\rho(g^{-1})\).
Multiplication by a translate of \(u_0\) carries a shifted
cusp section \(X\) into the cusp space of weight \(\kappa\).
This uses boundedness of the translated section.
The map
\[
 X\longmapsto\big(P_\mu(X\cdot h^*u_0)\big)_{h\in U}
\]
from the shifted cusp space to \(|U|\) copies of the
character space is injective. Indeed, at a frame \(z\)
where the translate matrix is invertible, its rows express
the vanishing of those traces as a linear system on the
values \((g^*X)(z)\) times nonzero character and action
scalars. The column for the identity gives \(X(z)=0\).

On every component the determinant is holomorphic and
nonzero somewhere. The invertible locus is dense in its
connected frame domain, so continuity gives vanishing
throughout. Finally,
\[
 \frac{\binom{T_0+n}{n}}{\binom{T_0+2n+6}{n}}
   =\prod_{a=1}^{n}
       \left(1-\frac{n+6}{T_0+n+6+a}\right)
   \ge 1-\frac{n(n+6)}{T_0}.
\]
The family takes \(n(n+6)/T_0\le1/(4(fk)^3)\le1/8\).
This gives the asserted lower share for the character on
both signatures on the indicated two-level model. Taking \(y\) to be the
larger of the two losses gives \(y\le1/8\).
\end{proof}

\paragraph{The rank removed by the arithmetic projector.}

A separating section identifies the levels. The projector can still remove part of their rank, which is counted in the finite representation at the two level primes.

Use the finite group of the two level primes in the same model.
For this calculation twist its action by the inverse of the
central form character $\chi$, so that the chosen character
becomes the product Whittaker character $\lambda$ of $U$. At
each residue factor its value is a nontrivial additive character
of the sum on the first superdiagonal. This twist only multiplies
the scalar in the pullback action. Write $V$ for the cusp space
of weight $\kappa$ with that action and $V_\lambda$ for its
character space. Its dimension is $d_0$.

\begin{lemma}[Trace of the discarded central idempotent]
\label{lem:paired-projector-rank-loss}
For the character space of dimension \(d_0\), write \(r_{\rm bad}\) for the rank discarded by the bad-type projector. The ratios \(\beta_k\) from \eqref{eq:row-bad-share} give
\[
 r_{\rm bad}\le2d_0(\beta_0+\beta_1).
\]
 This comparison uses the integral
arithmetic projector through its scalar extension.
\end{lemma}

\begin{proof}
\proofstep{The finite type decomposition at the two level primes.}
Let $H_k$ be the whole finite linear group at the $k$th level
prime and the upper unitriangulars at the other. There is an
$H_k$-equivariant embedding of $V$ into $a_k$ copies of the
regular representation, with $a_k |H_k|\le 2|U|d_0$.
For a definite form space use the free action on components,
which gives equality with $|U|d_0$ in this bound. In the
indefinite case, let $w$ be a separating section for $H=H_k$.
Write $w_h$ for its ordinary pullbacks, indexed so that the
ring action of $g$ sends $w_h$ to $w_{gh}$. The scalar twist belongs to the action on the form \(v\), while pullback permutes the separating sections. Consequently
\[
  g(w_h v)=w_{gh}\,g(v).
\]

Let $W$ be the invariant subspace of the span of all $w_h v$
for $v\in V$. Consider the map
\[
 V\longrightarrow \operatorname{Map}(H,W),\qquad
 v\longmapsto
    \left(\sum_{g\in H}g(w_h v)\right)_{h\in H}.
\]
For $kv$ the tuple at $h$ becomes the tuple for $v$ at
$k^{-1}h$, so the map intertwines the left regular action.
If the tuple vanishes, the group determinant of the $w_h$
gives $g(v)=0$ for each $g$ on a dense open set of every
component. This gives injectivity. Conversely the map
$(a_h)\mapsto\sum_h w_h a_h$ from $W^{\oplus H}$ has target weight $\kappa+2d$.

If its value vanishes, applying each $g\in H$ gives the same
invertible matrix of translates on the $a_h$ because $W$
consists of invariants. Writing \(a_k=\dim W\), apply first the injection just constructed,
then the three-weight comparison and finally the trace injection for \(U\):
\[
 a_k|H_k|\le\dim V_{\kappa+2d}
 \le2\dim V_{\kappa-d}\le2|U|d_0 .
\]
The bounds on the irreducible multiplicities for $H_k$ are
unchanged by extending the algebraic coefficient field through
a different complex embedding. Indeed, over an algebraic
closure they are the ranks of the central primitive idempotents
divided by the corresponding irreducible degrees. A change of
embedding permutes these idempotents and preserves their degrees.
The regular representation contains each irreducible with
multiplicity its degree. Thus the embedding bound also holds
at the embedding used for the finite type projectors.
\label{cor:paired-good-share-1}

\proofstep{Tracing the central bad projector.}
In the group algebra let $b_k$ be the central idempotent for
bad tuples of semisimple labels at the $k$th prime, extended
by the identity over the other prime. The two commuting central idempotents give the retained projector
\[
  e_{\rm good}=(1-b_0)(1-b_1)=1-b_0-b_1+b_0b_1.
\]
 Let $e_\lambda$
be the normalized character idempotent of $U$.
The operator $b_ke_\lambda$ is idempotent and belongs to
$\mathbb C[H_k]$. The dimension discarded inside $V_\lambda$
is at most the sum of the two ranks of these idempotents.

For the regular representation of $H_k$ the trace of group
algebra multiplication by $B$ is $|H_k|$ times the coefficient
of $1$ in $B$. Indeed each diagonal entry on the group basis
is that coefficient. Consequently, writing $\beta_k$ for that
coefficient in $b_ke_\lambda$ divided by the coefficient in
$e_\lambda$, the rank discarded is at most
\begin{equation}\label{eq:row-bad-share}
 \sum_k a_k\frac{|H_k|}{|U|}\,\beta_k
 \ \le\ 2d_0(\beta_0+\beta_1).
\end{equation}
Here the projector image has its arithmetic meaning even if the
cusp reading uses a different embedding from the type weights.
Both the character kernel and its intersection with the kernel of
the good idempotent are ranks of stacked algebraic linear maps.
This is the use of the two embeddings in the selected comparison.
\label{cor:paired-good-share-2}

\end{proof}

Let \(G=\operatorname{GL}_r(\mathbb F_g)\), \(g=g_k\), with upper
unitriangular subgroup \(U\). For a monic degree-\(r\) polynomial
\(P\) with unit constant coefficient let \(e_P\) be its central
series projector. Write \(e_\psi\) for the idempotent of a standard
nondegenerate character of \(U\) and \([1]\) for the coefficient
of the identity in a group-algebra element.

\begin{lemma}[Local weights of squarefree type labels]
The numbers \(\omega_g(P)=[1](e_Pe_\psi)/[1]e_\psi\) are
nonnegative, sum to one and are independent of the standard
nontrivial additive character used to define \(\psi\). If \(P\)
is squarefree, then
\[
        \omega_g(P)=\prod_{Q\mid P}(g^{\deg Q}-1)^{-1},
\]
where \(Q\) ranges over monic irreducible factors.
\end{lemma}
\begin{proof}
In the regular representation \(\operatorname{tr}_{\rm reg}(a)
=|G|[1]a\); the projectors \(e_Pe_\psi\) split the range of
\(e_\psi\). Conjugation by a diagonal matrix with the prescribed
successive ratios carries one standard character to the other,
fixing each \(e_P\). This proves the first assertions.

For squarefree \(P\) the dual element has distinct eigenvalues;
its centralizer is the torus \(T=\prod_{Q\mid P}\mathbb F_{g^{\deg Q}}^\times\).
The torus character is in general position. The Deligne--Lusztig
degree and Gelfand--Graev multiplicity formulas for \(G\), whose
center is connected, give one irreducible in this series of degree
\(|G|_{g'}/|T|\), with Whittaker multiplicity one. Thus
\[
 \omega_g(P)=\frac{(|G|_{g'}/|T|)\cdot1}{|G|/|U|}=|T|^{-1}.
\]
\cite[\S\S5.24, 7.1, 7.4 and 10.7]{deligne-lusztig-finite}
\end{proof}
At the level prime \(g_k\), sum the products of these local
weights over bad tuples, one factor per split place, to obtain
\(\beta_k\). \label{cor:paired-good-share-3}

\paragraph{Bounds for the four finite label exclusions.}
Lemma~\ref{lem:finite-type-event-bounds} gives the four mass
bounds used here.
The finite type bound used on the selected family is
\begin{align}
 \beta_k\le{}&\frac{3f}{g}+\frac{23 f\log g}{\sqrt g}
     +2fm\left(\frac1{4\min(q,d')}+\frac{1+\log r}{g-1}\right)
       \notag\\
     &+\frac r2(7/8)^f+\exp\left(-\frac{f}{5\log r}\right).
 \label{eq:row-bad-types}
\end{align}
In this formula $q,d'$ are the fixed auxiliary primes at the
point, and $d'$ is distinct from the determinant weight $d$.
The hypotheses are $r\ge100$, $g\ge4096$, distinct odd $q,d'$
with $\max(q,d')\ge16$ and $16qd'\mid g-1$, and the same
character labels as in the projector. The squarefree mass is used
as a subprobability. In degree $r\ge2$ there are $g^r-g^{r-1}$
monic squarefree polynomials by the Euler product, and removing
zero constant terms costs at most $g^{r-1}$.

Since the torus
order of each remaining label is at most $g^r$, the missing
weight is at most $2/g$.
The retained mass is consequently bounded below by
\[
  \frac{g^r-g^{r-1}-g^{r-1}}{g^r}=1-\frac2g .
\]
 A union over the $f$ places fits the
first displayed term. The other terms use four events on the
squarefree labels, namely a pair of roots with equal prime-to-$\ell$
parts at a nonexcluded prime $\ell$, a matching part against a
prescribed scalar or its inverse, a common proper divisor degree
at all places, and the absence of a prime-degree factor of the
specified large size at every place. These are the four tests of
the finite type. We give the counting argument in the next
paragraph.\label{cor:paired-good-share-4}

For the two root tests, let \(g\) be the chosen level prime, let \(r\) be the rank, and put \(N=r\). For a degree-\(j\) label, the order of its multiplicative group and its \(\ell\)-primary part are
\[
  A_j=g^j-1,\qquad A_{j,\ell}=\ell^{v_\ell(g^j-1)}.
\]
 The
squarefree mass of labels that contain specified distinct
irreducibles of degrees $j_1,\ldots,j_u$ is at most
$\prod_i A_{j_i}^{-1}$.

Strip off the specified factors and sum
the remaining weights, which have total at most one in the
remaining degree. This last assertion also follows without the
character formula: if $I_j$ is the number of irreducibles of
degree $j$ with nonzero constant term, then $jI_j\le A_j$ by
counting their distinct roots in $\mathbb F_{g^j}^\times$.
Thus the mass of polynomials with factor counts $m_j$ is at
most $\prod_j(j^{m_j}m_j!)^{-1}$, the probability of that cycle
count under a uniform permutation. This proves domination for
any event specified by the factor degrees.

\begin{lemma}[Mass of the four excluded finite types]
\label{lem:finite-type-event-bounds}
Use the squarefree subprobability measure and the hypotheses of
\eqref{eq:row-bad-types}. Write \(m_{\rm pair},m_{\rm scal},m_{\rm div},m_{\rm cyc}\)
for the probabilities of a repeated prime-to-\(\ell\) root part, a match
with a prescribed scalar or its inverse, a common proper divisor degree,
and absence of the prescribed large prime degree, respectively. Then
\[
 \begin{aligned}
 m_{\rm pair}&\le 23f\log(g)/\sqrt g,&
 m_{\rm scal}&\le2fm\left(\frac1{4\min(q,d')}+\frac{1+\log r}{g-1}\right),\\
 m_{\rm div}&\le\tfrac r2(7/8)^f,&
 m_{\rm cyc}&\le\exp\!\left(-\frac f{5\log r}\right).
 \end{aligned}
\]
The pair test ranges over the nonexcluded \(\ell\), and the last two
bounds use the product over the \(f\) residue factors. The separate
summand \(3f/g\) in \eqref{eq:row-bad-types} allows for the missing
squarefree mass.
\end{lemma}

\begin{proof}
\proofstep{Repeated roots in the squarefree mass.}
When two different roots of a single irreducible of degree $j$
have the same prime-to-$\ell$ part, they are $z,z^{g^a}$ with
$1\le a<j$. Their equality requires $z^{g^a-1}$ to lie in the
$\ell$-primary group. In the cyclic group of order $A_j$, this
allows at most $(g^{\gcd(a,j)}-1)A_{j,\ell}$ choices of $z$.
For a nonexcluded prime, $A_{j,\ell}\le g^{j/4}$ by its
exclusion test.

Since \(\gcd(a,j)\le j/2\), the root count is bounded by
\[
  (g^{\gcd(a,j)}-1)A_{j,\ell}\le g^{j/2}g^{j/4}=g^{3j/4}.
\]
 An equality between distinct roots can occur only
for $\ell\mid A_j$. There are at most $j\log_2 g$ such primes.
Every represented irreducible has at least $j$ winning roots,
because the test is invariant under Frobenius. Dividing the root
count by $j$ and then by $A_j$ bounds the mass for one factor by
\[
 \frac{4096}{4095}\log_2 g
       \sum_{j\ge2}(j-1)u^j,
 \qquad u=g^{-1/4}\le1/8.
\]

For a pair of distinct irreducibles of degrees $i,j$ use the
cyclic groups on the two roots. The common prime-to-$\ell$
part lies in the group of order $g^{\gcd(i,j)}-1$, and the two
primary parts give at most $A_{i,\ell}A_{j,\ell}$ possibilities.
Thus the number of ordered root pairs is at most $g^{3(i+j)/4}$.
The union over the prime factors of $A_i A_j$, followed by
stripping the two factors from the label, bounds the mass by
\[
 \left(\frac{4096}{4095}\right)^2\log_2 g
                \sum_{i,j\ge1}(i+j)u^{i+j}.
\]
The sums are $u^2/(1-u)^2$ and $2u^2/(1-u)^3$. Together they
fit $23\log g/\sqrt g$, with a constant independent of $N$.
This also bounds the whole-place failure by the union over places.
\label{cor:paired-type-mass-1}

\proofstep{A fixed scalar and its root count.}
For a fixed scalar in $\mathbb F_g^\times$, the number of
roots in $\mathbb F_{g^j}^\times$ with matching prime-to-$\ell$
part is at most $A_{j,\ell}$. It suffices to take the union over
primes dividing $A_j$. If a match occurs at another prime the
root is the scalar itself, which is already counted. Divide by
$j$ for the number of roots of each irreducible, and by $A_j$
for the mass of labels containing it.
For fixed \(j,\ell\), these two divisions bound the mass from a fixed scalar by
\[
  \frac{1}{A_j}\,\frac{A_{j,\ell}}{j} .
\]
 At $j=1$, the sum of the
primary parts of $g-1$, divided by $g-1$, is at most
$3/(16\min(q,d'))$. Indeed the three parts for $2,q,d'$ have
size at least $16,q,d'$, and the product of any two is at least
$16\min(q,d')$.

For previous sum \(S\ge2\) and product \(P\), adjoining a primary part
\(c\ge2\) gives
\[
 \frac{S+c}{cP}\le\frac S P,
 \qquad cS-(S+c)=(c-1)(S-1)-1\ge0 .
\] At $j\ge2$, the product of the parts of
$A_j$ supported on $g-1$ is at most $j(g-1)$ by the lifting
formula for powers, also at two since $g\equiv1\pmod4$.
The product of the other parts is at most $A_j/(g-1)$.
In each group the sum of its primary parts is at most the
product. Thus the degree $j$ contribution is at most
$(g-1)/A_j+1/(j(g-1))$. Summing the resulting estimate uses both
\[
  \frac{g-1}{g^j-1}=\frac{1}{1+g+\cdots+g^{j-1}}\le g^{1-j},\qquad \sum_{j=2}^N\frac1j\le\log N.
\]
 This proves the scalar term
in \eqref{eq:row-bad-types} after summing over the $2fm$
scalars and inverse scalars.
\label{cor:paired-type-mass-2}

\proofstep{The divisor-degree and cycle tests.}
A divisor of degree $a$ for a squarefree label corresponds to
an invariant $a$-subset of a permutation in the degree-count
bound. For $0<a<N$, count stable $a$-subsets with weight eight,
except give weight four when the permutation fixes a point on
each side. For a fixed subset \(A\), let \(E_A\) be the event that
it is stable and \(F_A\) the event of a fixed point on
each side. Given \(E_A\), the two restrictions are
independent uniform permutations. The derangement formula
gives probability at least \(1/2\) of a fixed point on
each side, including a side of size one. Thus the mean
contribution from this subset is
\[
 \mathbb P(E_A)(8-4\mathbb P(F_A\mid E_A))\le7\mathbb P(E_A).
\]

When some subset is
stable the total weight is at least eight. In the weight-four
case, swap a fixed point on each side to obtain a different
stable subset of the same size. Summing the stabilization
probability over subsets gives one. Hence the chance of a
stable $a$-subset is at most $7/8$. Independence of the places
and a union over $1\le a\le N/2$ give the fourth term of
\eqref{eq:row-bad-types}.

A permutation can have at most one cycle of prime length in
$(N/2,N-3]$. For each such \(p\), the choice of its support, its cyclic ordering and the permutation on the complement gives
\[
  \Pr(\text{a cycle of length }p)=\binom Np\frac{(p-1)!(N-p)!}{N!}=\frac1p.
\]
 The sum of these reciprocals is at least
$1/(5\log N)$ for $N\ge100$. For $100\le N<16900$ use the sieve to count primes through
$16899$, and sum the lower bounds $10^{-12}\lfloor10^{12}/p\rfloor$.
Multiplying by $5\lfloor\log_2 N\rfloor\cdot0.6931471803$ gives
at least one for every such $N$, using only integer arithmetic.

For larger
$N$, write $n_1=\lfloor N/2\rfloor$, $s_1=\lfloor\sqrt N\rfloor$
and $t=\#\{p:N/2<p\le N-3\}$. The central binomial lower
bound $4^{n_1}/n_1$ and its prime factorization give
\[
 \frac{N-1}{6}\log 4 \le
                 \big(\pi(\sqrt{2n_1})+t+4\big)\log N.
\]
Indeed the primes at most $\sqrt{2n_1}$ cost at most $2n_1$
each, those between $\sqrt{2n_1}$ and $2n_1/3$ cost at most
the primorial up to $2n_1/3$, bounded by $4^{\lfloor2n_1/3\rfloor}$,
and those between $2n_1/3$ and $n_1$ do not divide the binomial.
At most three primes near the upper endpoint were omitted from
the list for $t$.

Now $\pi(\sqrt{2n_1})\le(s_1+9)/3$ by counting
classes coprime to $6$, and $\log N\le(s_1+1)/64+7.70407$ by
$\log u\le u-1$ applied to $u=(s_1+1)/128$. With $s_1\ge130$,
expanding the two sides gives
\[
 (s_1+21)((s_1+1)/64+7.70407)+0.693148\le0.093147s_1^2.
\]
Since $\log 2>0.69314718$, substitution gives $5t\log N\ge N$.
Each reciprocal in the sum is at least $1/N$. The cycle-count domination gives
$1-1/(5\log N)$ as an upper bound for the single-place mass
without one of the prescribed factors. Raising to the $f$th
power bounds the final event in \eqref{eq:row-bad-types}.
\label{cor:paired-type-mass-3}
\end{proof}


\begin{proposition}[A uniform bound on the removed rank]
\label{prop:paired-good-rank}
Use either run of Proposition~\ref{prop:linked-comparison-rows}
and its retained projector in
Proposition~\ref{prop:equivariant-selected-row}. Assume the
finite labels of Lemma~\ref{lem:finite-type-event-bounds} and
size bounds \eqref{eq:paired-bad-type-sizes}. Writing \(d_{\rm rem}\) for the removed rank at the lift, one has
\[
 \delta=\exp(-x^{e_z/7}),\qquad d_{\rm rem}\le\delta d_0,\qquad (1-\delta)d_0\le D_0.
\]
 The finite type bounds are applied
with the same parameters as the central projector trace.
\end{proposition}

\begin{proof}
\proofstep{Summing the finite tests over the selected labels.}

The size bounds for the type mass are
\begin{equation}\label{eq:paired-bad-type-sizes}
 f\ge x^2,\qquad f,m,r\le x^5,\qquad
 g\ge \exp(8x^\epsilon),\qquad
 \min(q,d')\ge\exp(x^{e_z/6}),\qquad
 0<e_z\le\epsilon<1 .
\end{equation}
The exponents \(e_z,\epsilon\) are fixed for the family. On these bounds every term in \eqref{eq:row-bad-types},
including those with $\log g$, is $o(\exp(-x^{e_z/7}))$, uniformly
in the place, lift and selected point. For the logarithm term use
$\log g/\sqrt g\le O(g^{-1/4})$. For the last two terms use
$\log r\le5\log x$ and $f\ge x^2$. At each level prime the resulting family bound is
\[
  4\beta_k\le\delta,\qquad \delta=\exp(-x^{e_z/7})\le\tfrac14.
\]
This threshold is common to the two runs and all the readings.
By \eqref{eq:row-bad-share} the removed rank is at most
$\delta d_0$. Using the tied rank equality from
Section~\ref{subsec:selected-cusp-reading} at each lift gives
$(1-\delta)d_0\le D_0$ on the retained row. Its numerator
comparisons use the corresponding eigenspaces in the same row.
\label{cor:paired-good-share-5}
\end{proof}



\begin{lemma}[Majorant on a conjugate pair]
Put \(\widehat W=2fb,\ s_B=\widehat W+fS_e,\ s_C=Cm(2+v_d)s_B\).
For the underlying rational prime \(p\), one can choose \(A_j(p)\ge0\) so that
\[
 d_w+d_{\bar w}-2F_j(w)+8u_w\le A_j(p).
\]

One can take \(4T+4fr^2\) at a level prime,
\(2(s_B+T+1)+4{\bf1}_{p=d}fkb'\) at \(q,d\), and
\(2(s_C+T+1)\) at the designated small primes. At the other
primes \(p\le Y_0\) one can take twice the excess over \(F_j(w)\)
in \eqref{eq:paired-ordinary-depth}, and zero above \(Y_0\).
\end{lemma}
\begin{proof}
Use the normalized depth \eqref{eq:slice-finite-baseline}.
Write \(e_v=e_1 e_2\) where \(e_2\) is
absolute ramification in the construction field.
The tie for \(E\) bounds \(E/e_v\) above by
\(T(\nu^*+{\bf1}_{\rm bad})+1/e_2\). This follows by
rounding first in the construction field, then multiplying
by \(e_1\), with \(\nu_\lambda=e_2\nu^*\).

At a level prime \(p\), the unscaled degree bound gives
\(\nu^*\le D_L/(p-1)\), \(TD_L\le D_\kappa<(p-2)/2\).
The field contains the required \(p\)-th roots of unity, so
\(e_2\ge p-1\ge2\). Thus \(T\nu^*<1/2\) and 
\[
  E/e_v< T+1.
\]
The scaling satisfies \(v(\sigma_v)/e_v=n_v+s'_v\)
with \(s'_v\ge0\). The two embeddings of the level
normalization give 
\[
  n_w+n_{\bar w}=-2F_j(w).
\]
 In particular
\(d_v\le2T-n_v\) as \(T\ge1\). Adding at the two places
retains the sign of \(F_j\); omitting the normalization would
discard this cancellation.

At the ramified primes \(q,d\) write
\(\nu^*=\nu_W+b_p\), and 
\[
  v(\sigma_v)/e_v=T b_p+s'_v
\]
with \(s'_v\ge0\). Using \(\nu_W\le D_L/(p-1)\) and
\(TD_L\le D_\kappa\le C_0 k (W_j+fS_e)\) with \(C_0k\le p-1\)
gives \(d_v\le s_B+T+1\), \(s_B=\widehat W+fS_e\). The forced
term vanishes at \(q\); at \(d\) it is nonpositive, of magnitude
at most \(2fkb'\) by the bound on the smaller paired weight.

At the designated small primes the
\hyperref[par:unscaled-light-depth]{Hodge and colength estimate} gives
\[
 \begin{gathered}
 d_v\le s_C+T+1, \\ s_C=Cm(2+v_d)s_B, \\ F_j=0.
\end{gathered}
\] Here \(v_d\) denotes the small-prime
discriminant parameter. At the other primes \(p\le Y_0\)
write the local companion-chart depth bound as
\begin{equation}\label{eq:paired-ordinary-depth}
 d_v\le F_j(w)+1+
 \frac{C(m\widehat W\log k+fkS_e)}p+Ck\widehat W\vartheta_p+
 {\bf1}_{p>X_2}\frac{Ck\widehat W}{p^2}.
\end{equation}
It holds for \(v=w,\bar w\) by the split, inert and
conjugating chart comparisons. The selected numerical record
sets the additional chart error identically zero.

Here \(\vartheta_p=p^{-i_0(k,p)}\), where \(i_0\) is the least
odd \(i\ge3\) with \(p^i>k^{3/2}\). Hence
\(\vartheta_p\le\min(k^{-3/2},p^{-3})\). Apply
\(\sum_{p\le y}\log p\le y\log4\) on \([2,\sqrt k]\) and then
on the intervals \((2^a\sqrt k,2^{a+1}\sqrt k]\):
\[
 \sum_p\vartheta_p\log p
 \le\frac{3}{2k}+\frac3k\sum_{a\ge0}4^{-a}
 =\frac{11}{2k}\le\frac{20}{k}.
\]
At a good
prime \(p>Y_0\ge D_L\), the chosen monomial gives depth zero, the normalization is a unit and the rounded
exponent is zero. Here \(d_v=0\) and \(F_j(w)\ge0\).
\label{cor:paired-row-indices-6}

Adding the two local upper bounds in each case gives the stated
\(A_j(p)\); at a level prime use \(n_w+n_{\bar w}=-2F_j(w)\).
\end{proof}

\begin{corollary}[Weighted contribution of the run remainders]
In the paired index estimate take
\[
  \mathfrak e_j(w)=A_j(p)+
       \frac{r_w/e_w+r_{\bar w}/e_{\bar w}}{I D_0}.
\]
The weighted contribution of its remainder fraction is at most
\(12fmb\delta_2\), uniformly in \(j,K\) and the finite set of places.
\end{corollary}
\begin{proof}
In the coefficients normalized by \(e_w\), the place weight is
\(e_w\omega_w\). Conjugation preserves both \(e_w\) and \(\omega_w\).
Apply \eqref{eq:run-cut} to each of the two sums of remainders:
their combined bound is \(4R_jfmb\delta_2/(I D_0)\), and
\(R_j/I\le3D_0\).
\end{proof}

\paragraph{Application to the restricted sections of the parent.}
Fix a point at index \(i(u_0)\) of the translate family of
Section~\ref{sec:family-window}. The character and level
identifications on its compatible geometric model give two
short-section covers of the very same core, at the prescribed
primes \(p_d,p'_d\). The finite-group action on the core forms
pulls back to both covers. Its homomorphism at lift \(j\) is
chosen on the core once. At places above \(p_d\) take the cover
at \(p'_d\), and at all other places take the cover at \(p_d\).

\begin{proposition}[Uniform choice of remainders on the parent family]\label{prop:uniform-parent-remainders}
There is a common family threshold for both orientations and all
allowed lifts and translates. Beyond it each nonempty place class
\(\mathcal C\) has one remainder function \(r_{\mathcal C}\). If
\(r_w(V)\) denotes the required remainder for a subspace tested in
the parent row at \(w\), the same choice on that class satisfies
\[
 r_w(V)\le r_{\mathcal C}(w)
 \qquad(w\in\mathcal C,\ V\text{ any such tested subspace}).
\]
The empty class may be assigned the zero function.
\end{proposition}
\begin{proof}
The distinguished prime of the cover used at \(v\) is a unit there.
Thus the character trace and
the comparison of its twisted generators apply on both sides of
a conjugate pair. The field for each of the two auxiliary primes
has the level and character comparison of that prime, while the
small coefficient field for the form frame is the common core
field. The half-level ideals in the scalar \(\sigma_v\) come from
that core. The level-coordinate maps, frame weights, and rational
parts required here were identified in
paragraph~\ref{par:selected-form-coordinates}.
\label{cor:paired-row-indices-7}

Fix the section and height constants in \eqref{eq:geometric-run-index}.
Increase the family parameter \(C_P\) to at least \(1000\) and use it for
both orientations. The scale parameter in the model equals
\(x\). Its fixed light coefficient is \(C_{n_W}=10\), and its
level interval gives 
\[
  \log g_i\le 9L+1.
\]
 Since \(x\le frL\),
the fixed power \(C_P^{20}\le frL\) holds once
\(x\ge C_P^{20}\). The numerical conditions in
paragraph~\ref{par:slice-scales} now give the order and room
inequalities at every lift, with the common density of the indices with larger weights
\(\delta_2\).

Denote by \(X_{\rm sec},X_{\rm rank},X_{\rm depth},X_{\rm size}\) the thresholds from the section, rank, depth and family-size estimates, taking the largest size threshold if several occur. The simultaneous range is
\[
  x>\max\{X_{\rm sec},X_{\rm rank},X_{\rm depth},X_{\rm size}\}.
\]
 The constants
were chosen before the point or lift. With the weights of
the form frame and the scaled light and short lattices, the
generic restrictions and their Hasse lifts used to compute
\eqref{eq:geometric-run-index} are precisely the ones of these
two covers. Each nonempty place class therefore uses its one
function of remainders for every subspace, with zero sufficient
on an empty class.
In particular, if \(V_{\rm good}\) denotes the marked good space, the subspace to which this one remainder estimate is next applied is
\[
  \sum_{g\ {\rm mirabolic}} g V_{\rm good} .
\]

In particular it can be evaluated on the
whole sum of the mirabolic translates of the marked good space.
The common integral base rank \(D_0\), its trace action and
its projected lattices at \(\mu(w)\) are those of
Section~\ref{subsec:division-level-pins}. Thus
\eqref{eq:run-cut}, with its rounded depth and the signed
twist comparison above, bounds the two transported indices
of the parent. The forced terms agree with \(F_j(w)\) there
because both read the signed weight sum at the selected chart
of lift \(j\) in the same depth datum. This realizes the
geometric term of the first row error even as \(u_0\) varies.
\label{cor:paired-row-indices-8}
\end{proof}
